% generated by GAPDoc2LaTeX from XML source (Frank Luebeck)
\documentclass[a4paper,11pt]{report}

\usepackage[top=37mm,bottom=37mm,left=27mm,right=27mm]{geometry}
\sloppy
\pagestyle{myheadings}
\usepackage{amssymb}
\usepackage[latin1]{inputenc}
\usepackage{makeidx}
\makeindex
\usepackage{color}
\definecolor{FireBrick}{rgb}{0.5812,0.0074,0.0083}
\definecolor{RoyalBlue}{rgb}{0.0236,0.0894,0.6179}
\definecolor{RoyalGreen}{rgb}{0.0236,0.6179,0.0894}
\definecolor{RoyalRed}{rgb}{0.6179,0.0236,0.0894}
\definecolor{LightBlue}{rgb}{0.8544,0.9511,1.0000}
\definecolor{Black}{rgb}{0.0,0.0,0.0}

\definecolor{linkColor}{rgb}{0.0,0.0,0.554}
\definecolor{citeColor}{rgb}{0.0,0.0,0.554}
\definecolor{fileColor}{rgb}{0.0,0.0,0.554}
\definecolor{urlColor}{rgb}{0.0,0.0,0.554}
\definecolor{promptColor}{rgb}{0.0,0.0,0.589}
\definecolor{brkpromptColor}{rgb}{0.589,0.0,0.0}
\definecolor{gapinputColor}{rgb}{0.589,0.0,0.0}
\definecolor{gapoutputColor}{rgb}{0.0,0.0,0.0}

%%  for a long time these were red and blue by default,
%%  now black, but keep variables to overwrite
\definecolor{FuncColor}{rgb}{0.0,0.0,0.0}
%% strange name because of pdflatex bug:
\definecolor{Chapter }{rgb}{0.0,0.0,0.0}
\definecolor{DarkOlive}{rgb}{0.1047,0.2412,0.0064}


\usepackage{fancyvrb}

\usepackage{mathptmx,helvet}
\usepackage[T1]{fontenc}
\usepackage{textcomp}


\usepackage[
            pdftex=true,
            bookmarks=true,        
            a4paper=true,
            pdftitle={Written with GAPDoc},
            pdfcreator={LaTeX with hyperref package / GAPDoc},
            colorlinks=true,
            backref=page,
            breaklinks=true,
            linkcolor=linkColor,
            citecolor=citeColor,
            filecolor=fileColor,
            urlcolor=urlColor,
            pdfpagemode={UseNone}, 
           ]{hyperref}

\newcommand{\maintitlesize}{\fontsize{50}{55}\selectfont}

% write page numbers to a .pnr log file for online help
\newwrite\pagenrlog
\immediate\openout\pagenrlog =\jobname.pnr
\immediate\write\pagenrlog{PAGENRS := [}
\newcommand{\logpage}[1]{\protect\write\pagenrlog{#1, \thepage,}}
%% were never documented, give conflicts with some additional packages

\newcommand{\GAP}{\textsf{GAP}}

%% nicer description environments, allows long labels
\usepackage{enumitem}
\setdescription{style=nextline}

%% depth of toc
\setcounter{tocdepth}{1}




\usepackage{amsmath}
\usepackage[all]{xy}
\newcommand{\tensor}{\otimes}
\newcommand{\quiverproduct}{\times}
\newcommand{\add}{\operatorname{add}\nolimits}
\newcommand{\Hom}{\operatorname{Hom}\nolimits}
\newcommand{\End}{\operatorname{End}\nolimits}
\newcommand{\Ext}{\operatorname{Ext}\nolimits}
\newcommand{\rad}{\operatorname{rad}\nolimits}
\newcommand{\op}{{\operatorname{op}\nolimits}}
\renewcommand{\Im}{\operatorname{Im}\nolimits}
\DeclareMathOperator{\modc}{mod}
\DeclareMathOperator{\proj}{proj}


%% command for ColorPrompt style examples
\newcommand{\gapprompt}[1]{\color{promptColor}{\bfseries #1}}
\newcommand{\gapbrkprompt}[1]{\color{brkpromptColor}{\bfseries #1}}
\newcommand{\gapinput}[1]{\color{gapinputColor}{#1}}


\begin{document}

\logpage{[ 0, 0, 0 ]}
\begin{titlepage}
\mbox{}\vfill

\begin{center}{\maintitlesize \textbf{QPA\\
\mbox{}}}\\
\vfill

\hypersetup{pdftitle={QPA}}
\markright{\scriptsize \mbox{}\hfill QPA \hfill\mbox{}}
{\Huge \textbf{Quivers and Path Algebras\\
\mbox{}}}\\
\vfill

{\Huge Version 1.37\mbox{}}\\[1cm]
{April 2025\mbox{}}\\[1cm]
\mbox{}\\[2cm]
{\Large \textbf{\strut The QPA\texttt{\symbol{45}}team     \strut\mbox{}}}\\
\hypersetup{pdfauthor={The QPA\texttt{\symbol{45}}team    }}
\mbox{}\\[2cm]
\begin{minipage}{12cm}\noindent
 \end{minipage}

\end{center}\vfill

\mbox{}\\
{\mbox{}\\
\small \noindent \textbf{The QPA\texttt{\symbol{45}}team    }  Email: \href{mailto://oyvind.solberg@ntnu.no} {\texttt{oyvind.solberg@ntnu.no}}\\
  Homepage: \href{https://gap-packages.github.io/qpa/} {\texttt{https://gap\texttt{\symbol{45}}packages.github.io/qpa/}}\\
  Address: \begin{minipage}[t]{8cm}\noindent
 Department of Mathematical Sciences\\
 NTNU\\
 N\texttt{\symbol{45}}7491 Trondheim\\
 Norway \end{minipage}
}\\
\end{titlepage}

\newpage\setcounter{page}{2}
{\small 
\section*{Abstract}
\logpage{[ 0, 0, 1 ]}
 The GAP4 deposited package \textsf{QPA} extends the GAP functionality for computations with finite dimensional
quotients of path algebras. \textsf{QPA} has data structures for quivers, quotients of path algebras, representations
of quivers with relations and complexes of modules. Basic operations on
representations of quivers are implemented as well as constructing minimal
projective resolutions of modules (using using linear algebra). A not
necessarily minimal projective resolution constructed by using Groebner basis
theory and a paper by
Green\texttt{\symbol{45}}Solberg\texttt{\symbol{45}}Zacharia, "Minimal
projective resolutions", has been implemented. A goal is to have a test for
finite representation type. This work has started, but there is a long way
left. Part of this work is to implement/port the functionality and data
structures that was available in CREP. \mbox{}}\\[1cm]
{\small 
\section*{Copyright}
\logpage{[ 0, 0, 2 ]}
{\copyright} 2014\texttt{\symbol{45}}2025 by The QPA\texttt{\symbol{45}}team.

 \textsf{QPA} is free software; you can redistribute it and/or modify it under the terms of
the GNU General Public License as published by the Free Software Foundation;
either version 2 of the License, or (at your option) any later version. For
details, see the FSF{\textquoteright}s own site (\href{https://www.gnu.org/licenses/gpl.html} {\texttt{https://www.gnu.org/licenses/gpl.html}}).

 If you obtained \textsf{QPA}, we would be grateful for a short notification sent to one of members of the
QPA\texttt{\symbol{45}}team. If you publish a result which was partially
obtained with the usage of \textsf{QPA}, please cite it in the following form:

 The QPA\texttt{\symbol{45}}team, \textsf{QPA} \texttt{\symbol{45}} Quivers, path algebras and representations, Version 1.37;
2025 (\href{https://folk.ntnu.no/oyvinso/QPA/} {\texttt{https://folk.ntnu.no/oyvinso/QPA/}}) \mbox{}}\\[1cm]
{\small 
\section*{Acknowledgements}
\logpage{[ 0, 0, 3 ]}
The system design of \textsf{QPA} was initiated by Edward L. Green, Lenwood S. Heath, and Craig A. Struble. It
was continued and completed by Randall Cone and Edward Green. We would like to
thank the following people for their contributions: \begin{center}
\begin{tabular}{ll}Chain complexes&
Kristin Krogh Arnesen and {\O}ystein Skarts{\ae}terhagen\\
Degeneration order for modules in finite type&
Andrzej Mroz\\
GBNP interface (for Groebner bases)&
Randall Cone\\
Homomorphisms of modules&
{\O}yvind Solberg and Anette Wraalsen\\
Koszul duals&
Stephen Corwin\\
Matrix representations of path algebras&
{\O}yvind Solberg and George Yuhasz\\
Opposite algebra and tensor products of algebras&
{\O}ystein Skarts{\ae}terhagen\\
Predefined classes of algebras&
Andrzej Mroz and {\O}yvind Solberg\\
Projective resolutions (using Groebnar basis)&
Randall Cone and {\O}yvind Solberg\\
Projective resolutions (using linear algebra)&
{\O}yvind Solberg\\
Quickstart&
Kristin Krogh Arnesen\\
Quivers, path algebras&
Gerard Brunick\\
The bounded derived category&
Kristin Krogh Arnesen and {\O}ystein Skarts{\ae}terhagen\\
Unitforms&
{\O}yvind Solberg\\
&
\\
&
\\
&
\\
&
\\
&
\\
&
\\
&
\\
&
\\
&
\\
&
\\
&
\\
&
\\
&
\\
&
\\
&
\\
&
\\
\end{tabular}\\[2mm]
\end{center}

 \mbox{}}\\[1cm]
\newpage

\def\contentsname{Contents\logpage{[ 0, 0, 4 ]}}

\tableofcontents
\newpage

 
\chapter{\textcolor{Chapter }{Introduction}}\label{Introduction}
\logpage{[ 1, 0, 0 ]}
\hyperdef{L}{X7DFB63A97E67C0A1}{}
{
 
\section{\textcolor{Chapter }{General aims}}\label{Aims}
\logpage{[ 1, 1, 0 ]}
\hyperdef{L}{X8557083378F2A3B2}{}
{
 The overall aim of \textsf{QPA} is to provide computational tools in basic research in mathematics (algebra).
It seeks to furnish users within and outside the area of representation theory
of finite dimensional algebras with a computational software package that will
help in exploring a problem and testing conjectures. As with all software
development, new ideas and question will arise. The ability to study and
compute examples will help in proving or disproving
well\texttt{\symbol{45}}established questions/conjectures. Furthermore, it
will enable us to consider new examples which were not accessible by hand or
other means before. In this way, we hope that \textsf{QPA} will aid in the development, not only of the area of representation theory of
finite dimensional algebras, but also of a broad variety of other areas of
mathematics, where such structures occur. In addition we aspire to create a
research environment for international cooperation on computational
representation theory of finite dimensional algebras. }

 
\section{\textcolor{Chapter }{Installation and system requirements}}\label{Installation}
\logpage{[ 1, 2, 0 ]}
\hyperdef{L}{X7DB566D5785B7DBC}{}
{
 Since the release of GAP version 4.7.8, the package \textsf{QPA} is distributed with GAP, so if you download this version or any later version
of GAP, the package \textsf{QPA} will be installed at the same time as GAP is installed. Follow the
instructions for installing GAP on the web page \href{https://www.gap-system.org/Download/index.html} {\texttt{https://www.gap\texttt{\symbol{45}}system.org/Download/index.html}} After having successfully started GAP, give the command 
\begin{Verbatim}[commandchars=!@|,fontsize=\small,frame=single,label=Example]
  !gapprompt@gap>| !gapinput@LoadPackage("qpa");|
\end{Verbatim}
 which loads QPA and you can start using QPA.

 One can also clone the QPA git repository from 
\begin{itemize}
\item \href{https://github.com/gap-packages/qpa} {\texttt{https://github.com/gap\texttt{\symbol{45}}packages/qpa}}
\end{itemize}
 to follow the QPA development. Follow the instructions on 
\begin{itemize}
\item \href{https://folk.ntnu.no/oyvinso/QPA/} {\texttt{https://folk.ntnu.no/oyvinso/QPA/}}
\end{itemize}
 to install QPA by cloning the the git repository. }

 }

 
\chapter{\textcolor{Chapter }{Quickstart}}\label{Quickstart}
\logpage{[ 2, 0, 0 ]}
\hyperdef{L}{X7F83DF528480AEA3}{}
{
 This chapter is intended for those who would like to get started with \textsf{QPA} right away by playing with a few examples. We assume that the user is familiar
with \textsf{GAP} syntax, for instance the different ways to display various \textsf{GAP} objects: \texttt{View}, \texttt{Print} and \texttt{Display}. These features are all implemented for the objects defined in \textsf{QPA}, and by using \texttt{Display} on an object, you will get a complete description of it.

 The following examples show how to create the most fundamental algebraic
structures featured in \textsf{QPA}, namely quivers, path algebras and quotients of path algebras, modules and
module homomorphisms. Sometimes, there is more than one way of constructing
such objects. See their respective chapter in the documentation for more on
this. The code from the examples can be found in the \texttt{examples/} directory of the distribution of \textsf{QPA}.

  
\section{\textcolor{Chapter }{Example 1 \texttt{\symbol{45}}\texttt{\symbol{45}} quivers, path algebras and
quotients of path algebras}}\label{Example 1}
\logpage{[ 2, 1, 0 ]}
\hyperdef{L}{X7D3488D984288697}{}
{
 We construct a quiver $Q$, i.e. a finite directed graph, with one vertex and two loops: 
\begin{Verbatim}[commandchars=!@|,fontsize=\small,frame=single,label=Example]
  !gapprompt@gap>| !gapinput@Q := Quiver( 1, [ [1,1,"a"], [1,1,"b"] ] );|
  <quiver with 1 vertices and 2 arrows>
  !gapprompt@gap>| !gapinput@Display(Q);|
  Quiver( ["v1"], [["v1","v1","a"],["v1","v1","b"]] )
\end{Verbatim}
 When displaying $Q$, we observe that the vertex has been named \texttt{v1}, and that this name is used when describing the arrows. (The "Display" style
of viewing a quiver can also be used in construction, i.e., we could have
written \texttt{Q := Quiver( ["v1"], [["v1","v1","a"],["v1","v1","b"]] )} to get the same object.)

 If we want to know the number and names of the vertices and arrows, without
getting the structure of $Q$, we can request this information as shown below. We can also access the
vertices and arrows directly. 
\begin{Verbatim}[commandchars=!@|,fontsize=\small,frame=single,label=Example]
  !gapprompt@gap>| !gapinput@VerticesOfQuiver(Q);|
  [ v1 ]
  !gapprompt@gap>| !gapinput@ArrowsOfQuiver(Q);|
  [ a, b ]
  !gapprompt@gap>| !gapinput@Q.a;|
  a
\end{Verbatim}
 The next step is to create the path algebra $kQ$ from $Q$, where $k$ is the rational numbers (in general, one can chose any field implemented in \textsf{GAP}). 
\begin{Verbatim}[commandchars=!@|,fontsize=\small,frame=single,label=Example]
  !gapprompt@gap>| !gapinput@kQ := PathAlgebra(Rationals, Q);|
  <Rationals[<quiver with 1 vertices and 2 arrows>]>
  !gapprompt@gap>| !gapinput@Display(kQ);|
  <Path algebra of the quiver <quiver with 1 vertices and 2 arrows> 
  over the field Rationals>
\end{Verbatim}
 We know that this algebra has three generators, with the vertex \texttt{v{\textunderscore}1} as the identity. This can be verified by \textsf{QPA}. For convenience, we introduce new variables \texttt{v1}, \texttt{a} and \texttt{b} to get easier access to the generators. 
\begin{Verbatim}[commandchars=!@|,fontsize=\small,frame=single,label=Example]
  !gapprompt@gap>| !gapinput@AssignGeneratorVariables(kQ);|
  #I  Assigned the global variables [ v1, a, b ]
  !gapprompt@gap>| !gapinput@v1; a; b;|
  (1)*v1
  (1)*a
  (1)*b
  !gapprompt@gap>| !gapinput@id := One(kQ);|
  (1)*v1
  !gapprompt@gap>| !gapinput@v1 = id;|
  true
\end{Verbatim}
 Now, we want to construct a finite dimensional algebra, by dividing out some
ideal. The generators of the ideal (the relations) are given in terms of
paths, and it is important to know the convention of writing paths used in \textsf{QPA}. If we first go the arrow $a$ and then the arrow $b$, the path is written as $a*b$. 

 Say that we want our ideal to be generated by the relations \texttt{\texttt{\symbol{92}}\texttt{\symbol{123}}a\texttt{\symbol{94}}2, a*b
\texttt{\symbol{45}} b*a,
b\texttt{\symbol{94}}2\texttt{\symbol{92}}\texttt{\symbol{125}}}. Then we make a list \texttt{relations} consisting of these relations and to construct the quotient we say: \texttt{A := kQ/relations;} on the command line in \textsf{GAP}. 
\begin{Verbatim}[commandchars=!@|,fontsize=\small,frame=single,label=Example]
  !gapprompt@gap>| !gapinput@relations := [a^2,a*b-b*a, b*b];|
  [ (1)*a^2, (1)*a*b+(-1)*b*a, (1)*b^2 ]
  !gapprompt@gap>| !gapinput@A := kQ/relations;|
  <Rationals[<quiver with 1 vertices and 2 arrows>]/<two-sided ideal in 
  <Rationals[<quiver with 1 vertices and 2 arrows>]>, (3 generators)>>
\end{Verbatim}
 See \ref{qpa:Quotients_of_path_algebras} for further remarks on constructing quotients of path algebras. }

 
\section{\textcolor{Chapter }{Example 2 \texttt{\symbol{45}}\texttt{\symbol{45}} Introducing modules}}\logpage{[ 2, 2, 0 ]}
\hyperdef{L}{X7D0E555F79FFD1EE}{}
{
 In representation theory, there are several conventions for expressing modules
of path algebras, and again it is useful to comment on the convention used in \textsf{QPA}. A module (or representation) of an algebra $A = kQ/I$ is, briefly explained, a picture of $Q$ where the vertices are finite dimensional $k$\texttt{\symbol{45}}vectorspaces, and the arrows are linear transformations
between the vector spaces respecting the relations of $I$. The modules are \emph{right} modules, and a linear transformation from $k^n$ to $k^m$ is represented by a $n \times m$\texttt{\symbol{45}}matrix.

 There are several ways of constructing modules in $QPA$. First, we will explore some modules which $QPA$ gives us for free, namely the indecomposable projectives. We start by
constructing a new algebra. The underlying quiver has three vertices and three
arrows and looks like an $A_3$ quiver with both arrows pointing to the right, and one additional loop in the
final vertex. The only relation is to go this loop twice. 
\begin{Verbatim}[commandchars=!@|,fontsize=\small,frame=single,label=Example]
  !gapprompt@gap>| !gapinput@Q := Quiver( 3, [ [1,2,"a"], [2,3,"b"], [3,3,"c"] ]);|
  <quiver with 3 vertices and 3 arrows>
  !gapprompt@gap>| !gapinput@kQ := PathAlgebra(Rationals, Q);|
  <Rationals[<quiver with 3 vertices and 3 arrows>]>
  !gapprompt@gap>| !gapinput@relations := [kQ.c*kQ.c];|
  [ (1)*c^2 ]
  !gapprompt@gap>| !gapinput@A := kQ/relations;|
  <Rationals[<quiver with 3 vertices and 3 arrows>]/
  <two-sided ideal in <Rationals[<quiver with 3 vertices and 3 arrows>]>, 
    (1 generators)>>
\end{Verbatim}
 The indecomposable projectives are easily created with one command. We use \texttt{Display} to explore the modules. 
\begin{Verbatim}[commandchars=!@|,fontsize=\small,frame=single,label=Example]
  !gapprompt@gap>| !gapinput@projectives := IndecProjectiveModules(A);|
  [ <[ 1, 1, 2 ]>, <[ 0, 1, 2 ]>, <[ 0, 0, 2 ]> ]
  !gapprompt@gap>| !gapinput@proj1 := projectives[1];|
  <[ 1, 1, 2 ]>
  !gapprompt@gap>| !gapinput@Display(proj1);|
  <Module over <Rationals[<quiver with 3 vertices and 3 arrows>]/
  <two-sided ideal in <Rationals[<quiver with 3 vertices and 3 arrows>]>, 
    (1 generators)>> with dimension vector 
  [ 1, 1, 2 ]> and linear maps given by
  for arrow a:
  [ [  1 ] ]
  for arrow b:
  [ [  1,  0 ] ]
  for arrow c:
  [ [  0,  1 ],
    [  0,  0 ] ]
\end{Verbatim}
 If we, for some reason, want to use the maps of this module, we can get the
matrices directly by using the command \texttt{MatricesOfPathAlgebraModule(proj1)}: 
\begin{Verbatim}[commandchars=!@|,fontsize=\small,frame=single,label=Example]
  !gapprompt@gap>| !gapinput@M := MatricesOfPathAlgebraModule(proj1);|
  [ [ [ 1 ] ], [ [ 1, 0 ] ], [ [ 0, 1 ], [ 0, 0 ] ] ]
  !gapprompt@gap>| !gapinput@M[1];|
  [ [ 1 ] ]
\end{Verbatim}
 Naturally, the indecomposable injective modules are just as easily
constructed, and so are the simple modules. 
\begin{Verbatim}[commandchars=!@|,fontsize=\small,frame=single,label=Example]
  !gapprompt@gap>| !gapinput@injectives := IndecInjectiveModules(A);|
  [ <[ 1, 0, 0 ]>, <[ 1, 1, 0 ]>, <[ 2, 2, 2 ]> ]
  !gapprompt@gap>| !gapinput@simples := SimpleModules(A);|
  [ <[ 1, 0, 0 ]>, <[ 0, 1, 0 ]>, <[ 0, 0, 1 ]> ]
\end{Verbatim}
 We know for a fact that the simple module in vertex 1 and the indecomposable
injective module in vertex 1 coincide. Let us look at this relationship in \textsf{QPA}: 
\begin{Verbatim}[commandchars=!@|,fontsize=\small,frame=single,label=Example]
  !gapprompt@gap>| !gapinput@s1 := simples[1];|
  <[ 1, 0, 0 ]>
  !gapprompt@gap>| !gapinput@inj1 := injectives[1];|
  <[ 1, 0, 0 ]>
  !gapprompt@gap>| !gapinput@IsIdenticalObj(s1,inj1);|
  false
  !gapprompt@gap>| !gapinput@s1 = inj1;|
  true
  !gapprompt@gap>| !gapinput@IsomorphicModules(s1,inj1);|
  true
\end{Verbatim}
 We observe that \textsf{QPA} recognizes the modules as "the same" (that is, isomorphic); however, they are \emph{not} the same instance and hence the simplest test for equality fails. This is
important to bear in mind \texttt{\symbol{45}}\texttt{\symbol{45}} objects
which are isomorphic and regarded as the same in the "real world", are not
necessarily the same in \textsf{GAP}. }

 
\section{\textcolor{Chapter }{Example 3 \texttt{\symbol{45}}\texttt{\symbol{45}} Constructing modules and
module homomorphisms}}\logpage{[ 2, 3, 0 ]}
\hyperdef{L}{X790BB3A1815A9B4D}{}
{
 Assume we want to construct the following $A$\texttt{\symbol{45}}module $M$, where $A$ is the same algebra as in the previous example $ \xymatrix{0\ar[r]^0 & {\ensuremath{\mathbb Q}}\ar[r]^1 & {\ensuremath{\mathbb
Q}} \ar@(ur,dr)^0} $. This module is neither indecomposable projective or injective, nor simple,
so we need to do the dirty work ourselves. Usually, the easiest way to
construct a module is to state the dimension vector and the
non\texttt{\symbol{45}}zero maps. Here, there is only one
non\texttt{\symbol{45}}zero map, and we write 
\begin{Verbatim}[commandchars=!@|,fontsize=\small,frame=single,label=Example]
  !gapprompt@gap>| !gapinput@M := RightModuleOverPathAlgebra( A, [0,1,1], [ ["b", [[1]] ] ] );|
  <[ 0, 1, 1 ]>
\end{Verbatim}
 To make sure we got everything right, we can use \texttt{Display(M)} to view the maps. The most tricky thing is usually to get the correct numbers
of brackets. Here is a slightly bigger example: $ \xymatrix{{\ensuremath{\mathbb Q}}\ar[r]^{\left(\begin{smallmatrix}0 &
0\end{smallmatrix}\right)} & {\ensuremath{\mathbb
Q}}^2\ar[r]^{\left(\begin{smallmatrix}1 & 0\\ -1 & 0\end{smallmatrix}\right)}
& {\ensuremath{\mathbb Q}}^2\ar@(ur,dr)^{\left(\begin{smallmatrix}0 & 0\\ 1 &
0\end{smallmatrix}\right)}} $. 
\begin{Verbatim}[commandchars=!@|,fontsize=\small,frame=single,label=Example]
  !gapprompt@gap>| !gapinput@N := RightModuleOverPathAlgebra( A, [1,2,2], [ ["a",[[1,1]] ], |
     ["b", [[1,0], [-1,0]] ], ["c", [[0,0],[1,0]] ] ] );
  <[ 1, 2, 2 ]>
\end{Verbatim}
 Now we want to construct a map between the two modules, say $f: M \rightarrow N$, which is non\texttt{\symbol{45}}zero only in vertex 2. This is done by 
\begin{Verbatim}[commandchars=!@|,fontsize=\small,frame=single,label=Example]
  !gapprompt@gap>| !gapinput@f := RightModuleHomOverAlgebra(M,N, [ [[0]], [[1,1]], NullMat(1,2,Rationals)]);|
  <<[ 0, 1, 1 ]> ---> <[ 1, 2, 2 ]>>
  !gapprompt@gap>| !gapinput@Display(f);|
  <<Module over <Rationals[<quiver with 3 vertices and 3 arrows>]/
  <two-sided ideal in <Rationals[<quiver with 3 vertices and 3 arrows>]>, 
    (1 generators)>> with dimension vector 
  [ 0, 1, 1 ]> ---> <Module over <Rationals[<quiver with 3 vertices and 3 arrows>]/
  <two-sided ideal in <Rationals[<quiver with 3 vertices and 3 arrows>]>, 
    (1 generators)>> with dimension vector [ 1, 2, 2 ]>>
  with linear map for vertex number 1:
  [ [  0 ] ]
  linear map for vertex number 2:
  [ [  1,  1 ] ]
  linear map for vertex number 3:
  [ [  0,  0 ] ]
\end{Verbatim}
 Note the two different ways of writing zero maps. Again, we can retrieve the
matrices describing $f$: 
\begin{Verbatim}[commandchars=!@|,fontsize=\small,frame=single,label=Example]
  !gapprompt@gap>| !gapinput@MatricesOfPathAlgebraMatModuleHomomorphism(f);|
  [ [ [ 0 ] ], [ [ 1, 1 ] ], [ [ 0, 0 ] ] ]
\end{Verbatim}
 }

 }

 
\chapter{\textcolor{Chapter }{Quivers}}\label{Quivers}
\logpage{[ 3, 0, 0 ]}
\hyperdef{L}{X7FA7E6B581D41A94}{}
{
 
\section{\textcolor{Chapter }{Information class, Quivers}}\logpage{[ 3, 1, 0 ]}
\hyperdef{L}{X7C14F4617F7E9F09}{}
{
 A quiver $Q$ is a set derived from a labeled directed multigraph with loops $\Gamma$. An element of $Q$ is called a *path*, and falls into one of three classes. The first class is
the set of *vertices* of $\Gamma$. The second class is the set of *walks* in $\Gamma$ of length at least one, each of which is represented by the corresponding
sequence of *arrows* in $\Gamma$. The third class is the singleton set containing the distinguished *zero
path*, usually denoted $0$. An associative multiplication is defined on $Q$.

 This chapter describes the functions in \textsf{QPA} that deal with paths and quivers. The functions for constructing paths in
Section \ref{Constructing Paths} are normally not useful in isolation; typically, they are invoked by the
functions for constructing quivers in Section \ref{Constructing Quivers}. 

\subsection{\textcolor{Chapter }{InfoQuiver}}
\logpage{[ 3, 1, 1 ]}\nobreak
\hyperdef{L}{X786A65F07FA1BB78}{}
{\noindent\textcolor{FuncColor}{$\triangleright$\enspace\texttt{InfoQuiver\index{InfoQuiver@\texttt{InfoQuiver}}
\label{InfoQuiver}
}\hfill{\scriptsize (info class)}}\\


 is the info class for functions dealing with quivers. }

 }

 
\section{\textcolor{Chapter }{Constructing Quivers}}\label{Constructing Quivers}
\logpage{[ 3, 2, 0 ]}
\hyperdef{L}{X860B15D57EAB46D7}{}
{
 

\subsection{\textcolor{Chapter }{Quiver (no. of vertices, list of arrows)}}
\logpage{[ 3, 2, 1 ]}\nobreak
\hyperdef{L}{X7BD8455A7F2C5CA3}{}
{\noindent\textcolor{FuncColor}{$\triangleright$\enspace\texttt{Quiver({\mdseries\slshape N, arrows})\index{Quiver@\texttt{Quiver}!no. of vertices, list of arrows}
\label{Quiver:no. of vertices, list of arrows}
}\hfill{\scriptsize (function)}}\\
\noindent\textcolor{FuncColor}{$\triangleright$\enspace\texttt{Quiver({\mdseries\slshape vertices, arrows})\index{Quiver@\texttt{Quiver}!lists of vertices and arrows}
\label{Quiver:lists of vertices and arrows}
}\hfill{\scriptsize (function)}}\\
\noindent\textcolor{FuncColor}{$\triangleright$\enspace\texttt{Quiver({\mdseries\slshape adjacencymatrix})\index{Quiver@\texttt{Quiver}!adjacenymatrix}
\label{Quiver:adjacenymatrix}
}\hfill{\scriptsize (function)}}\\


Arguments: First construction: \mbox{\texttt{\mdseries\slshape N}} \texttt{\symbol{45}}\texttt{\symbol{45}} number of vertices, \mbox{\texttt{\mdseries\slshape arrows}} \texttt{\symbol{45}}\texttt{\symbol{45}} a list of arrows to specify the graph $\Gamma$. Second construction: \mbox{\texttt{\mdseries\slshape vertices}} \texttt{\symbol{45}}\texttt{\symbol{45}} a list of vertex names, \mbox{\texttt{\mdseries\slshape arrows}} \texttt{\symbol{45}}\texttt{\symbol{45}} a list of arrows. Third construction:
takes an adjacency matrix for the graph $\Gamma$. \\
\textbf{\indent Returns:\ }
a quiver, which is an object from the category \texttt{IsQuiver} (\ref{IsQuiver}).



 In the first and third constructions, the vertices are named `v1, v2, ...'. In
the second construction, unique vertex names are given as strings in the list
that is the first parameter. Each arrow is a list consisting of a source
vertex and a target vertex, followed optionally by an arrow name as a string.

 Vertices and arrows are referenced as record components using the dot (`.')
operator. }

 
\begin{Verbatim}[commandchars=!@|,fontsize=\small,frame=single,label=Example]
  !gapprompt@gap>| !gapinput@q1 := Quiver(["u","v"],[["u","u","a"],["u","v","b"], |
  !gapprompt@>| !gapinput@              ["v","u","c"],["v","v","d"]]);|
  <quiver with 2 vertices and 4 arrows>
  !gapprompt@gap>| !gapinput@VerticesOfQuiver(q1);|
  [ u, v ]
  !gapprompt@gap>| !gapinput@ArrowsOfQuiver(q1);|
  [ a, b, c, d ]
  !gapprompt@gap>| !gapinput@q2 := Quiver(2,[[1,1],[2,1],[1,2]]);|
  <quiver with 2 vertices and 3 arrows>
  !gapprompt@gap>| !gapinput@ArrowsOfQuiver(q2);|
  [ a1, a2, a3 ]
  !gapprompt@gap>| !gapinput@VerticesOfQuiver(q2);|
  [ v1, v2 ]
  !gapprompt@gap>| !gapinput@q3 := Quiver(2,[[1,1,"a"],[2,1,"b"],[1,2,"c"]]);|
  <quiver with 2 vertices and 3 arrows>
  !gapprompt@gap>| !gapinput@ArrowsOfQuiver(q3);|
  [ a, b, c ]
  !gapprompt@gap>| !gapinput@q4 := Quiver([[1,1],[2,1]]);|
  <quiver with 2 vertices and 5 arrows>
  !gapprompt@gap>| !gapinput@VerticesOfQuiver(q4);|
  [ v1, v2 ]
  !gapprompt@gap>| !gapinput@ArrowsOfQuiver(q4);|
  [ a1, a2, a3, a4, a5 ]
  !gapprompt@gap>| !gapinput@SourceOfPath(q4.a2);|
  v1
  !gapprompt@gap>| !gapinput@TargetOfPath(q4.a2);|
  v2
\end{Verbatim}
 

\subsection{\textcolor{Chapter }{DynkinQuiver (DynkinQuiver)}}
\logpage{[ 3, 2, 2 ]}\nobreak
\hyperdef{L}{X7EE08F058702A717}{}
{\noindent\textcolor{FuncColor}{$\triangleright$\enspace\texttt{DynkinQuiver({\mdseries\slshape Delta, n, orientation})\index{DynkinQuiver@\texttt{DynkinQuiver}!DynkinQuiver}
\label{DynkinQuiver:DynkinQuiver}
}\hfill{\scriptsize (operation)}}\\


 Arguments: \mbox{\texttt{\mdseries\slshape Delta}}, \mbox{\texttt{\mdseries\slshape n}}, \mbox{\texttt{\mdseries\slshape orientation}} \texttt{\symbol{45}}\texttt{\symbol{45}} a character (A,D,E), a positive
integer, and a list giving the orientation. \\
\textbf{\indent Returns:\ }
a Dynkin quiver of type \mbox{\texttt{\mdseries\slshape Delta}} ("A", "D", or "E") with index \mbox{\texttt{\mdseries\slshape n}} and orientation of the arrows given by the list \mbox{\texttt{\mdseries\slshape orientation}}. 



 If \mbox{\texttt{\mdseries\slshape Delta}} is equal to "A" with index \mbox{\texttt{\mdseries\slshape n}}, then the list \mbox{\texttt{\mdseries\slshape orientation}} is of the form \texttt{["r", "l", "l", ...,"r", "l"]} of length \mbox{\texttt{\mdseries\slshape n\texttt{\symbol{45}}1}}, where "l" or "r" in coordinate $i$ means that the arrow $a_i$ is oriented to the left or to the right, respectively. The vertices and the
arrows are named as in the following diagram $\xymatrix{1\ar@{-}[r]^-{a_1} & 2\ar@{-}[r]^-{a_2} & \ar@{--}[r] &
\ar@{-}[r]^-{a_{n-2}} & n - 1 \ar@{-}[r]^-{a_{n-1}} & n }$\\
 If \mbox{\texttt{\mdseries\slshape Delta}} is equal to "D" with index \mbox{\texttt{\mdseries\slshape n}} and \mbox{\texttt{\mdseries\slshape n}} greater or equal to $4$, then the list \mbox{\texttt{\mdseries\slshape orientation}} is of the form \texttt{["r", "l", "l", ...,"r", "l"]} of length \mbox{\texttt{\mdseries\slshape n\texttt{\symbol{45}}1}}, where "l" or "r" in coordinate $i$ means that the arrow $a_i$ is oriented to the left or to the right, respectively. The vertices and the
arrows are named as in the following diagram $\xymatrix{1\ar@{-}[dr]^-{a_1} & & & & & \\ & 3\ar@{-}[r]^{a_3} & \ar@{--}[r] &
\ar@{-}[r]^-{a_{n-2}} & n - 1 \ar@{-}[r]^-{a_{n-1}} & n \\ 2\ar@{-}[ur]_-{a_2}
& & & & & }$ \\
 If \mbox{\texttt{\mdseries\slshape Delta}} is equal to "E" with index \mbox{\texttt{\mdseries\slshape n}} and \mbox{\texttt{\mdseries\slshape n}} in $[6,7,8]$, then the list \mbox{\texttt{\mdseries\slshape orientation}} is of the form \texttt{["r", "l", "l", ...,"r", "l","d"]} of length \mbox{\texttt{\mdseries\slshape n\texttt{\symbol{45}}1}}, where "l" or "r" in the \mbox{\texttt{\mdseries\slshape n \texttt{\symbol{45}} 2}} first coordinates and at coordinate $i$ means that the arrow $a_i$ is oriented to the left or to the right, respectively, and the last
orientation parameter is "d" or "u" indicating if the arrow $a_{n-1}$ is oriented down or up. The vertices and the arrows are named as in the
following diagram $\xymatrix{ & & n\ar@{-}[d]^{a_{n-1}} & & & & & \\ 1\ar@{-}[r]^{a_1} &
2\ar@{-}[r]^{a_2} & 3\ar@{-}[r]^-{a_3} & \ar@{--}[r] &\ar@{-}[r] & n - 2
\ar@{-}[r]^-{a_{n-2}} & n - 1} $ \\
 }

 

\subsection{\textcolor{Chapter }{OrderedBy}}
\logpage{[ 3, 2, 3 ]}\nobreak
\hyperdef{L}{X83E498008233E2DD}{}
{\noindent\textcolor{FuncColor}{$\triangleright$\enspace\texttt{OrderedBy({\mdseries\slshape quiver, ordering})\index{OrderedBy@\texttt{OrderedBy}}
\label{OrderedBy}
}\hfill{\scriptsize (function)}}\\
\textbf{\indent Returns:\ }
a copy of \mbox{\texttt{\mdseries\slshape quiver}} whose elements are ordered by \mbox{\texttt{\mdseries\slshape ordering}}. The default ordering of a quiver is length left lexicographic. See Section \ref{qpa:Orderings} for more information.



 }

 }

 
\section{\textcolor{Chapter }{Categories and Properties of Quivers}}\logpage{[ 3, 3, 0 ]}
\hyperdef{L}{X80BB6B6183134D88}{}
{
 

\subsection{\textcolor{Chapter }{IsQuiver}}
\logpage{[ 3, 3, 1 ]}\nobreak
\hyperdef{L}{X7E9C03497FD7778B}{}
{\noindent\textcolor{FuncColor}{$\triangleright$\enspace\texttt{IsQuiver({\mdseries\slshape object})\index{IsQuiver@\texttt{IsQuiver}}
\label{IsQuiver}
}\hfill{\scriptsize (category)}}\\
\textbf{\indent Returns:\ }
 true when \mbox{\texttt{\mdseries\slshape object}} is a quiver. 

}

 

\subsection{\textcolor{Chapter }{IsAcyclicQuiver}}
\logpage{[ 3, 3, 2 ]}\nobreak
\hyperdef{L}{X82E99C5F8624BD86}{}
{\noindent\textcolor{FuncColor}{$\triangleright$\enspace\texttt{IsAcyclicQuiver({\mdseries\slshape quiver})\index{IsAcyclicQuiver@\texttt{IsAcyclicQuiver}}
\label{IsAcyclicQuiver}
}\hfill{\scriptsize (property)}}\\
\textbf{\indent Returns:\ }
 true when \mbox{\texttt{\mdseries\slshape quiver}} is a quiver with no oriented cycles. 

}

 

\subsection{\textcolor{Chapter }{IsUAcyclicQuiver}}
\logpage{[ 3, 3, 3 ]}\nobreak
\hyperdef{L}{X846C44937B3AF09A}{}
{\noindent\textcolor{FuncColor}{$\triangleright$\enspace\texttt{IsUAcyclicQuiver({\mdseries\slshape quiver})\index{IsUAcyclicQuiver@\texttt{IsUAcyclicQuiver}}
\label{IsUAcyclicQuiver}
}\hfill{\scriptsize (property)}}\\
\textbf{\indent Returns:\ }
 true when \mbox{\texttt{\mdseries\slshape quiver}} is a quiver with no unoriented cycles. Note: an oriented cycle is also an
unoriented cycle! 

}

 

\subsection{\textcolor{Chapter }{IsConnectedQuiver}}
\logpage{[ 3, 3, 4 ]}\nobreak
\hyperdef{L}{X7909BF627C5D0D4A}{}
{\noindent\textcolor{FuncColor}{$\triangleright$\enspace\texttt{IsConnectedQuiver({\mdseries\slshape quiver})\index{IsConnectedQuiver@\texttt{IsConnectedQuiver}}
\label{IsConnectedQuiver}
}\hfill{\scriptsize (property)}}\\
\textbf{\indent Returns:\ }
 true when \mbox{\texttt{\mdseries\slshape quiver}} is a connected quiver (i.e. each pair of vertices is connected by an
unoriented path in \mbox{\texttt{\mdseries\slshape quiver}}). 

}

 

\subsection{\textcolor{Chapter }{IsTreeQuiver}}
\logpage{[ 3, 3, 5 ]}\nobreak
\hyperdef{L}{X7A2B55BD7B6F0360}{}
{\noindent\textcolor{FuncColor}{$\triangleright$\enspace\texttt{IsTreeQuiver({\mdseries\slshape quiver})\index{IsTreeQuiver@\texttt{IsTreeQuiver}}
\label{IsTreeQuiver}
}\hfill{\scriptsize (property)}}\\
\textbf{\indent Returns:\ }
 true when \mbox{\texttt{\mdseries\slshape quiver}} is a tree as a graph (i.e. it is connected and contains no unoriented cycles). 

}

 
\begin{Verbatim}[commandchars=!@|,fontsize=\small,frame=single,label=Example]
  !gapprompt@gap>| !gapinput@q1 := Quiver(2,[[1,2]]);|
  <quiver with 2 vertices and 1 arrows>
  !gapprompt@gap>| !gapinput@IsQuiver("v1");|
  false
  !gapprompt@gap>| !gapinput@IsQuiver(q1);|
  true
  !gapprompt@gap>| !gapinput@IsAcyclicQuiver(q1); IsUAcyclicQuiver(q1); |
  true
  true
  !gapprompt@gap>| !gapinput@IsConnectedQuiver(q1); IsTreeQuiver(q1);|
  true
  true
  !gapprompt@gap>| !gapinput@q2 := Quiver(["u","v"],[["u","v"],["v","u"]]);|
  <quiver with 2 vertices and 2 arrows>
  !gapprompt@gap>| !gapinput@IsAcyclicQuiver(q2); IsUAcyclicQuiver(q2); |
  false
  false
  !gapprompt@gap>| !gapinput@IsConnectedQuiver(q2); IsTreeQuiver(q2);|
  true
  false
  !gapprompt@gap>| !gapinput@q3 := Quiver(["u","v"],[["u","v"],["u","v"]]);|
  <quiver with 2 vertices and 2 arrows>
  !gapprompt@gap>| !gapinput@IsAcyclicQuiver(q3); IsUAcyclicQuiver(q3); |
  true
  false
  !gapprompt@gap>| !gapinput@IsConnectedQuiver(q3); IsTreeQuiver(q3); |
  true
  false
  !gapprompt@gap>| !gapinput@q4 := Quiver(2, []);|
  <quiver with 2 vertices and 0 arrows>
  !gapprompt@gap>| !gapinput@IsAcyclicQuiver(q4); IsUAcyclicQuiver(q4); |
  true
  true
  !gapprompt@gap>| !gapinput@IsConnectedQuiver(q4); IsTreeQuiver(q4);|
  false
  false
\end{Verbatim}
 

\subsection{\textcolor{Chapter }{IsDynkinQuiver}}
\logpage{[ 3, 3, 6 ]}\nobreak
\hyperdef{L}{X85EC85B58688CCCC}{}
{\noindent\textcolor{FuncColor}{$\triangleright$\enspace\texttt{IsDynkinQuiver({\mdseries\slshape quiver})\index{IsDynkinQuiver@\texttt{IsDynkinQuiver}}
\label{IsDynkinQuiver}
}\hfill{\scriptsize (property)}}\\
\textbf{\indent Returns:\ }
 true when \mbox{\texttt{\mdseries\slshape quiver}} is a Dynkin quiver (more precisely, when underlying undirected graph of \mbox{\texttt{\mdseries\slshape quiver}} is a Dynkin diagram). 



 This function prints an additional information. If it returns true, it prints
the Dynkin type of \mbox{\texttt{\mdseries\slshape quiver}}, i.e. A{\textunderscore}n, D{\textunderscore}m, E{\textunderscore}6,
E{\textunderscore}7 or E{\textunderscore}8. Moreover, in case \mbox{\texttt{\mdseries\slshape quiver}} is not connected or contains an unoriented cycle, the function also prints a
respective info. }

 
\begin{Verbatim}[commandchars=!@|,fontsize=\small,frame=single,label=Example]
  !gapprompt@gap>| !gapinput@q1 := Quiver(4,[[1,4],[4,2],[3,4]]);|
  <quiver with 4 vertices and 3 arrows>
  !gapprompt@gap>| !gapinput@IsDynkinQuiver(q1);|
  D_4
  true
  !gapprompt@gap>| !gapinput@q2 := Quiver(2,[[1,2],[1,2]]);|
  <quiver with 2 vertices and 2 arrows>
  !gapprompt@gap>| !gapinput@IsDynkinQuiver(q2);|
  Quiver contains an (un)oriented cycle.
  false
  !gapprompt@gap>| !gapinput@q3 := Quiver(5,[[1,5],[2,5],[3,5],[4,5]]);|
  <quiver with 5 vertices and 4 arrows>
\end{Verbatim}
 }

 
\section{\textcolor{Chapter }{Orderings of paths in a quiver}}\label{qpa:Orderings}
\logpage{[ 3, 4, 0 ]}
\hyperdef{L}{X78BBB63B828EB9FB}{}
{
 The only supported ordering on the paths in a quiver is length left
lexicographic ordering. The reason for this is that \textsf{QPA} does not have its own functions for computing Groebner basis. Instead they are
computed using the \textsf{GAP}\texttt{\symbol{45}}package \textsf{GBNP}. The interface with this package, which is provided by the \textsf{QPA}, only supports the length left lexicographic ordering, even though \textsf{GBNP} supports more orderings.

 For constructing a quiver, there are three different methods. TODO: Explain
how the vertices and arrows are ordered.

 }

 
\section{\textcolor{Chapter }{Attributes and Operations for Quivers}}\logpage{[ 3, 5, 0 ]}
\hyperdef{L}{X80CF69E37B54F3C1}{}
{
 

\subsection{\textcolor{Chapter }{. (for quiver)}}
\logpage{[ 3, 5, 1 ]}\nobreak
\hyperdef{L}{X8198B2897FF5AC4B}{}
{\noindent\textcolor{FuncColor}{$\triangleright$\enspace\texttt{.({\mdseries\slshape Q, element})\index{.@\texttt{.}!for quiver}
\label{.:for quiver}
}\hfill{\scriptsize (operation)}}\\


 Arguments: \mbox{\texttt{\mdseries\slshape Q}} \texttt{\symbol{45}}\texttt{\symbol{45}} a quiver, and \mbox{\texttt{\mdseries\slshape element}} \texttt{\symbol{45}}\texttt{\symbol{45}} a vertex or an arrow. 

 The operation \texttt{.} allows access to generators of the quiver. If you have named your vertices and
arrows then the access looks like `\mbox{\texttt{\mdseries\slshape Q}}.\mbox{\texttt{\mdseries\slshape name of element}}'. If you have not named the elements of the quiver, then the default names
are v1, v2, ... and a1, a2, ... in the order they are created. }

 

\subsection{\textcolor{Chapter }{VerticesOfQuiver}}
\logpage{[ 3, 5, 2 ]}\nobreak
\hyperdef{L}{X7C82A4BC7FB329D8}{}
{\noindent\textcolor{FuncColor}{$\triangleright$\enspace\texttt{VerticesOfQuiver({\mdseries\slshape quiver})\index{VerticesOfQuiver@\texttt{VerticesOfQuiver}}
\label{VerticesOfQuiver}
}\hfill{\scriptsize (attribute)}}\\
\textbf{\indent Returns:\ }
 a list of paths that are vertices in \mbox{\texttt{\mdseries\slshape quiver}}. 

}

 

\subsection{\textcolor{Chapter }{ArrowsOfQuiver}}
\logpage{[ 3, 5, 3 ]}\nobreak
\hyperdef{L}{X82C42D7D820D5F9B}{}
{\noindent\textcolor{FuncColor}{$\triangleright$\enspace\texttt{ArrowsOfQuiver({\mdseries\slshape quiver})\index{ArrowsOfQuiver@\texttt{ArrowsOfQuiver}}
\label{ArrowsOfQuiver}
}\hfill{\scriptsize (attribute)}}\\
\textbf{\indent Returns:\ }
 a list of paths that are arrows in \mbox{\texttt{\mdseries\slshape quiver}}. 

}

 

\subsection{\textcolor{Chapter }{AdjacencyMatrixOfQuiver}}
\logpage{[ 3, 5, 4 ]}\nobreak
\hyperdef{L}{X7AF572F081AEFE98}{}
{\noindent\textcolor{FuncColor}{$\triangleright$\enspace\texttt{AdjacencyMatrixOfQuiver({\mdseries\slshape quiver})\index{AdjacencyMatrixOfQuiver@\texttt{AdjacencyMatrixOfQuiver}}
\label{AdjacencyMatrixOfQuiver}
}\hfill{\scriptsize (attribute)}}\\
\textbf{\indent Returns:\ }
 the adjacency matrix of \mbox{\texttt{\mdseries\slshape quiver}}. 

}

 

\subsection{\textcolor{Chapter }{GeneratorsOfQuiver}}
\logpage{[ 3, 5, 5 ]}\nobreak
\hyperdef{L}{X7B4A7F0F813E63FC}{}
{\noindent\textcolor{FuncColor}{$\triangleright$\enspace\texttt{GeneratorsOfQuiver({\mdseries\slshape quiver})\index{GeneratorsOfQuiver@\texttt{GeneratorsOfQuiver}}
\label{GeneratorsOfQuiver}
}\hfill{\scriptsize (attribute)}}\\
\textbf{\indent Returns:\ }
 a list of the vertices and the arrows in \mbox{\texttt{\mdseries\slshape quiver}}. 

}

 

\subsection{\textcolor{Chapter }{NumberOfVertices}}
\logpage{[ 3, 5, 6 ]}\nobreak
\hyperdef{L}{X822BD7F37F8AF016}{}
{\noindent\textcolor{FuncColor}{$\triangleright$\enspace\texttt{NumberOfVertices({\mdseries\slshape quiver})\index{NumberOfVertices@\texttt{NumberOfVertices}}
\label{NumberOfVertices}
}\hfill{\scriptsize (attribute)}}\\
\textbf{\indent Returns:\ }
 the number of vertices in \mbox{\texttt{\mdseries\slshape quiver}}. 

}

 

\subsection{\textcolor{Chapter }{NumberOfArrows}}
\logpage{[ 3, 5, 7 ]}\nobreak
\hyperdef{L}{X7AC77C9C7D069663}{}
{\noindent\textcolor{FuncColor}{$\triangleright$\enspace\texttt{NumberOfArrows({\mdseries\slshape quiver})\index{NumberOfArrows@\texttt{NumberOfArrows}}
\label{NumberOfArrows}
}\hfill{\scriptsize (attribute)}}\\
\textbf{\indent Returns:\ }
 the number of arrows in \mbox{\texttt{\mdseries\slshape quiver}}. 

}

 

\subsection{\textcolor{Chapter }{OrderingOfQuiver}}
\logpage{[ 3, 5, 8 ]}\nobreak
\hyperdef{L}{X84D1D1AA82689B03}{}
{\noindent\textcolor{FuncColor}{$\triangleright$\enspace\texttt{OrderingOfQuiver({\mdseries\slshape quiver})\index{OrderingOfQuiver@\texttt{OrderingOfQuiver}}
\label{OrderingOfQuiver}
}\hfill{\scriptsize (attribute)}}\\
\textbf{\indent Returns:\ }
 the ordering used to order elements in \mbox{\texttt{\mdseries\slshape quiver}}. See Section \ref{qpa:Orderings} for more information. 

}

 

\subsection{\textcolor{Chapter }{OppositeQuiver}}
\logpage{[ 3, 5, 9 ]}\nobreak
\hyperdef{L}{X84B82F6F84A442AB}{}
{\noindent\textcolor{FuncColor}{$\triangleright$\enspace\texttt{OppositeQuiver({\mdseries\slshape quiver})\index{OppositeQuiver@\texttt{OppositeQuiver}}
\label{OppositeQuiver}
}\hfill{\scriptsize (attribute)}}\\
\textbf{\indent Returns:\ }
 the opposite quiver of \mbox{\texttt{\mdseries\slshape quiver}}, where the vertices are labelled "name in original quiver" +
"{\textunderscore}op" and the arrows are labelled "name in original quiver" +
"{\textunderscore}op". 



 This attribute contains the opposite quiver of a quiver, that is, a quiver
which is the same except that every arrow goes in the opposite direction.

 The operation \texttt{OppositePath} (\ref{OppositePath}) takes a path in a quiver to the corresponding path in the opposite quiver.

 The opposite of the opposite of a quiver $Q$ is isomorphic to $Q$. In QPA, we regard these two quivers to be the same, so the call \texttt{OppositeQuiver(OppositeQuiver(Q))} returns the object \texttt{Q}. }

 
\begin{Verbatim}[commandchars=!@|,fontsize=\small,frame=single,label=Example]
  !gapprompt@gap>| !gapinput@q1 := Quiver(["u","v"],[["u","u","a"],["u","v","b"],|
  !gapprompt@>| !gapinput@              ["v","u","c"],["v","v","d"]]);|
  <quiver with 2 vertices and 4 arrows>
  !gapprompt@gap>| !gapinput@q1.a;|
  a
  !gapprompt@gap>| !gapinput@q1.v;|
  v
  !gapprompt@gap>| !gapinput@VerticesOfQuiver(q1);|
  [ u, v ]
  !gapprompt@gap>| !gapinput@ArrowsOfQuiver(q1);|
  [ a, b, c, d ]
  !gapprompt@gap>| !gapinput@AdjacencyMatrixOfQuiver(q1);|
  [ [ 1, 1 ], [ 1, 1 ] ]
  !gapprompt@gap>| !gapinput@GeneratorsOfQuiver(q1);|
  [ u, v, a, b, c, d ]
  !gapprompt@gap>| !gapinput@NumberOfVertices(q1);|
  2
  !gapprompt@gap>| !gapinput@NumberOfArrows(q1);|
  4
  !gapprompt@gap>| !gapinput@OrderingOfQuiver(q1);|
  <length left lexicographic ordering>
  !gapprompt@gap>| !gapinput@q1_op := OppositeQuiver(q1);|
  <quiver with 2 vertices and 4 arrows>
  !gapprompt@gap>| !gapinput@VerticesOfQuiver(q1_op);|
  [ u_op, v_op ]
  !gapprompt@gap>| !gapinput@ArrowsOfQuiver(q1_op);|
  [ a_op, b_op, c_op, d_op ]
\end{Verbatim}
 

\subsection{\textcolor{Chapter }{FullSubquiver}}
\logpage{[ 3, 5, 10 ]}\nobreak
\hyperdef{L}{X87B9B85483CBB238}{}
{\noindent\textcolor{FuncColor}{$\triangleright$\enspace\texttt{FullSubquiver({\mdseries\slshape quiver, list})\index{FullSubquiver@\texttt{FullSubquiver}}
\label{FullSubquiver}
}\hfill{\scriptsize (operation)}}\\
\textbf{\indent Returns:\ }
 This function returns a quiver which is a full subquiver of a \mbox{\texttt{\mdseries\slshape quiver}} induced by the \mbox{\texttt{\mdseries\slshape list}} of its vertices. 



 The names of vertices and arrows in resulting (sub)quiver remain the same as
in original one. The function checks if \mbox{\texttt{\mdseries\slshape list}} consists of vertices of \mbox{\texttt{\mdseries\slshape quiver}}. }

 

\subsection{\textcolor{Chapter }{ConnectedComponentsOfQuiver}}
\logpage{[ 3, 5, 11 ]}\nobreak
\hyperdef{L}{X824D7C768708F006}{}
{\noindent\textcolor{FuncColor}{$\triangleright$\enspace\texttt{ConnectedComponentsOfQuiver({\mdseries\slshape quiver})\index{ConnectedComponentsOfQuiver@\texttt{ConnectedComponentsOfQuiver}}
\label{ConnectedComponentsOfQuiver}
}\hfill{\scriptsize (operation)}}\\
\textbf{\indent Returns:\ }
 This function returns a list of quivers which are all connected components of
a \mbox{\texttt{\mdseries\slshape quiver}}. 



 The names of vertices and arrows in resulting (sub)quiver remain the same as
in original one. The function sets the property \texttt{IsConnectedQuiver} (\ref{IsConnectedQuiver}) to true for all the components. }

 
\begin{Verbatim}[commandchars=!@|,fontsize=\small,frame=single,label=Example]
  !gapprompt@gap>| !gapinput@Q := Quiver(6, [ [1,2],[1,1],[3,2],[4,5],[4,5] ]);|
  <quiver with 6 vertices and 5 arrows>
  !gapprompt@gap>| !gapinput@VerticesOfQuiver(Q);|
  [ v1, v2, v3, v4, v5, v6 ]
  !gapprompt@gap>| !gapinput@FullSubquiver(Q, [Q.v1, Q.v2]);|
  <quiver with 2 vertices and 2 arrows>
  !gapprompt@gap>| !gapinput@ConnectedComponentsOfQuiver(Q);|
  [ <quiver with 3 vertices and 3 arrows>, 
    <quiver with 2 vertices and 2 arrows>, 
    <quiver with 1 vertices and 0 arrows> ]
\end{Verbatim}
 

\subsection{\textcolor{Chapter }{DoubleQuiver}}
\logpage{[ 3, 5, 12 ]}\nobreak
\hyperdef{L}{X7C4BB3E5872FD483}{}
{\noindent\textcolor{FuncColor}{$\triangleright$\enspace\texttt{DoubleQuiver({\mdseries\slshape quiver})\index{DoubleQuiver@\texttt{DoubleQuiver}}
\label{DoubleQuiver}
}\hfill{\scriptsize (attribute)}}\\


Arguments: \mbox{\texttt{\mdseries\slshape quiver}} \texttt{\symbol{45}}\texttt{\symbol{45}} a quiver.\\
\textbf{\indent Returns:\ }
 the double quiver of \mbox{\texttt{\mdseries\slshape quiver}}. 



 The vertices in the double quiver are labelled by the same labels as in the
original quiver. The "old" arrows in the double quiver are labelled by the
same labels as in the original, and the "new" arrows in the double quiver are
labelled by the same labels as in the original quiver with the string "bar"
added. 

 }

 

\subsection{\textcolor{Chapter }{SeparatedQuiver}}
\logpage{[ 3, 5, 13 ]}\nobreak
\hyperdef{L}{X7B4EA2D6869BBAF3}{}
{\noindent\textcolor{FuncColor}{$\triangleright$\enspace\texttt{SeparatedQuiver({\mdseries\slshape quiver})\index{SeparatedQuiver@\texttt{SeparatedQuiver}}
\label{SeparatedQuiver}
}\hfill{\scriptsize (attribute)}}\\


Arguments: \mbox{\texttt{\mdseries\slshape quiver}} \texttt{\symbol{45}}\texttt{\symbol{45}} a quiver.\\
\textbf{\indent Returns:\ }
 the separated quiver of \mbox{\texttt{\mdseries\slshape quiver}}. 



 The vertices in the separated quiver are labelled $v$ and $v'$ for each vertex $v$ in \mbox{\texttt{\mdseries\slshape quiver}}, and for each arrow $a\colon v \to w$ in \mbox{\texttt{\mdseries\slshape quiver}} the arrow $v\to w'$ is labelled $a$.

 }

 }

 
\section{\textcolor{Chapter }{Categories and Properties of Paths}}\logpage{[ 3, 6, 0 ]}
\hyperdef{L}{X862A804B80C5A47D}{}
{
 

\subsection{\textcolor{Chapter }{IsPath}}
\logpage{[ 3, 6, 1 ]}\nobreak
\hyperdef{L}{X78B503DB83F2B6DE}{}
{\noindent\textcolor{FuncColor}{$\triangleright$\enspace\texttt{IsPath({\mdseries\slshape object})\index{IsPath@\texttt{IsPath}}
\label{IsPath}
}\hfill{\scriptsize (category)}}\\


 All path objects are in this category. }

 

\subsection{\textcolor{Chapter }{IsQuiverVertex}}
\logpage{[ 3, 6, 2 ]}\nobreak
\hyperdef{L}{X791474408297F7A0}{}
{\noindent\textcolor{FuncColor}{$\triangleright$\enspace\texttt{IsQuiverVertex({\mdseries\slshape object})\index{IsQuiverVertex@\texttt{IsQuiverVertex}}
\label{IsQuiverVertex}
}\hfill{\scriptsize (category)}}\\


 All vertices are in this category. }

 

\subsection{\textcolor{Chapter }{IsArrow}}
\logpage{[ 3, 6, 3 ]}\nobreak
\hyperdef{L}{X8266C9B8840C12EB}{}
{\noindent\textcolor{FuncColor}{$\triangleright$\enspace\texttt{IsArrow({\mdseries\slshape object})\index{IsArrow@\texttt{IsArrow}}
\label{IsArrow}
}\hfill{\scriptsize (category)}}\\


 All arrows are in this category. }

 

\subsection{\textcolor{Chapter }{IsZeroPath}}
\logpage{[ 3, 6, 4 ]}\nobreak
\hyperdef{L}{X8624A0B0795149CF}{}
{\noindent\textcolor{FuncColor}{$\triangleright$\enspace\texttt{IsZeroPath({\mdseries\slshape object})\index{IsZeroPath@\texttt{IsZeroPath}}
\label{IsZeroPath}
}\hfill{\scriptsize (property)}}\\


 is true when \mbox{\texttt{\mdseries\slshape object}} is the zero path. }

 
\begin{Verbatim}[commandchars=!@|,fontsize=\small,frame=single,label=Example]
  !gapprompt@gap>| !gapinput@q1 := Quiver(["u","v"],[["u","u","a"],["u","v","b"],|
  !gapprompt@>| !gapinput@              ["v","u","c"],["v","v","d"]]);|
  <quiver with 2 vertices and 4 arrows>
  !gapprompt@gap>| !gapinput@IsPath(q1.b);|
  true
  !gapprompt@gap>| !gapinput@IsPath(q1.u);|
  true
  !gapprompt@gap>| !gapinput@IsQuiverVertex(q1.c);|
  false
  !gapprompt@gap>| !gapinput@IsZeroPath(q1.d);|
  false
\end{Verbatim}
 }

 
\section{\textcolor{Chapter }{Attributes and Operations of Paths}}\logpage{[ 3, 7, 0 ]}
\hyperdef{L}{X7C8294338676C80E}{}
{
 

\subsection{\textcolor{Chapter }{SourceOfPath}}
\logpage{[ 3, 7, 1 ]}\nobreak
\hyperdef{L}{X84D8493C7AAF4ACC}{}
{\noindent\textcolor{FuncColor}{$\triangleright$\enspace\texttt{SourceOfPath({\mdseries\slshape path})\index{SourceOfPath@\texttt{SourceOfPath}}
\label{SourceOfPath}
}\hfill{\scriptsize (attribute)}}\\
\textbf{\indent Returns:\ }
 the source (first) vertex of \mbox{\texttt{\mdseries\slshape path}}. 

}

 

\subsection{\textcolor{Chapter }{TargetOfPath}}
\logpage{[ 3, 7, 2 ]}\nobreak
\hyperdef{L}{X86827D3F78F51815}{}
{\noindent\textcolor{FuncColor}{$\triangleright$\enspace\texttt{TargetOfPath({\mdseries\slshape path})\index{TargetOfPath@\texttt{TargetOfPath}}
\label{TargetOfPath}
}\hfill{\scriptsize (attribute)}}\\
\textbf{\indent Returns:\ }
 the target (last) vertex of \mbox{\texttt{\mdseries\slshape path}}. 

}

 

\subsection{\textcolor{Chapter }{LengthOfPath}}
\logpage{[ 3, 7, 3 ]}\nobreak
\hyperdef{L}{X7FF179D17C4F9FAC}{}
{\noindent\textcolor{FuncColor}{$\triangleright$\enspace\texttt{LengthOfPath({\mdseries\slshape path})\index{LengthOfPath@\texttt{LengthOfPath}}
\label{LengthOfPath}
}\hfill{\scriptsize (attribute)}}\\
\textbf{\indent Returns:\ }
 the length of \mbox{\texttt{\mdseries\slshape path}}. 

}

 

\subsection{\textcolor{Chapter }{WalkOfPath}}
\logpage{[ 3, 7, 4 ]}\nobreak
\hyperdef{L}{X781A8E06850E47B4}{}
{\noindent\textcolor{FuncColor}{$\triangleright$\enspace\texttt{WalkOfPath({\mdseries\slshape path})\index{WalkOfPath@\texttt{WalkOfPath}}
\label{WalkOfPath}
}\hfill{\scriptsize (attribute)}}\\
\textbf{\indent Returns:\ }
 a list of the arrows that constitute \mbox{\texttt{\mdseries\slshape path}} in order. 

}

 

\subsection{\textcolor{Chapter }{*}}
\logpage{[ 3, 7, 5 ]}\nobreak
\hyperdef{L}{X7857704878577048}{}
{\noindent\textcolor{FuncColor}{$\triangleright$\enspace\texttt{*({\mdseries\slshape p, q})\index{*@\texttt{*}}
\label{*}
}\hfill{\scriptsize (operation)}}\\


 Arguments: \mbox{\texttt{\mdseries\slshape p}} and \mbox{\texttt{\mdseries\slshape q}} \texttt{\symbol{45}}\texttt{\symbol{45}} two paths in the same quiver. \\
\textbf{\indent Returns:\ }
 the multiplication of the paths. If the paths are not in the same quiver an
error is returned. If the target of \mbox{\texttt{\mdseries\slshape p}} differs from the source of \mbox{\texttt{\mdseries\slshape q}}, then the result is the zero path. Otherwise, if either path is a vertex,
then the result is the other path. Finally, if both are paths of length at
least 1, then the result is the concatenation of the walks of the two paths. 

}

 
\begin{Verbatim}[commandchars=!@|,fontsize=\small,frame=single,label=Example]
  !gapprompt@gap>| !gapinput@q1 := Quiver(["u","v"],[["u","u","a"],["u","v","b"],|
  !gapprompt@>| !gapinput@              ["v","u","c"],["v","v","d"]]);|
  <quiver with 2 vertices and 4 arrows>
  !gapprompt@gap>| !gapinput@SourceOfPath(q1.v);|
  v
  !gapprompt@gap>| !gapinput@p1:=q1.a*q1.b*q1.d*q1.d;|
  a*b*d^2
  !gapprompt@gap>| !gapinput@TargetOfPath(p1);|
  v
  !gapprompt@gap>| !gapinput@p2:=q1.b*q1.b;|
  0
  !gapprompt@gap>| !gapinput@WalkOfPath(p1);|
  [ a, b, d, d ]
  !gapprompt@gap>| !gapinput@WalkOfPath(q1.a);|
  [ a ]
  !gapprompt@gap>| !gapinput@LengthOfPath(p1);|
  4
  !gapprompt@gap>| !gapinput@LengthOfPath(q1.v);|
  0
\end{Verbatim}
 

\subsection{\textcolor{Chapter }{=}}
\logpage{[ 3, 7, 6 ]}\nobreak
\hyperdef{L}{X806A4814806A4814}{}
{\noindent\textcolor{FuncColor}{$\triangleright$\enspace\texttt{=({\mdseries\slshape p, q})\index{=@\texttt{=}}
\label{=}
}\hfill{\scriptsize (operation)}}\\


 Arguments: \mbox{\texttt{\mdseries\slshape p}} and \mbox{\texttt{\mdseries\slshape q}} \texttt{\symbol{45}}\texttt{\symbol{45}} two paths in the same quiver. \\
\textbf{\indent Returns:\ }
 true if the two paths are equal. Two paths are equal if they have the same
source and the same target and if they have the same walks. 

}

 

\subsection{\textcolor{Chapter }{{\textless} (for two paths in a quiver)}}
\logpage{[ 3, 7, 7 ]}\nobreak
\hyperdef{L}{X7DAD2700853E8C21}{}
{\noindent\textcolor{FuncColor}{$\triangleright$\enspace\texttt{{\textless}({\mdseries\slshape p, q})\index{{\textless}@\texttt{{\textless}}!for two paths in a quiver}
\label{<:for two paths in a quiver}
}\hfill{\scriptsize (operation)}}\\


 Arguments: \mbox{\texttt{\mdseries\slshape p}} and \mbox{\texttt{\mdseries\slshape q}} \texttt{\symbol{45}}\texttt{\symbol{45}} two paths in the same quiver. \\
\textbf{\indent Returns:\ }
 a comparison of the two paths with respect to the ordering of the quiver. 

}

 }

 
\section{\textcolor{Chapter }{Attributes of Vertices}}\logpage{[ 3, 8, 0 ]}
\hyperdef{L}{X8158F0D27C4628FB}{}
{
 

\subsection{\textcolor{Chapter }{IncomingArrowsOfVertex}}
\logpage{[ 3, 8, 1 ]}\nobreak
\hyperdef{L}{X85E53C177F80E77E}{}
{\noindent\textcolor{FuncColor}{$\triangleright$\enspace\texttt{IncomingArrowsOfVertex({\mdseries\slshape vertex})\index{IncomingArrowsOfVertex@\texttt{IncomingArrowsOfVertex}}
\label{IncomingArrowsOfVertex}
}\hfill{\scriptsize (attribute)}}\\
\textbf{\indent Returns:\ }
 a list of arrows having \mbox{\texttt{\mdseries\slshape vertex}} as target. Only meaningful if \mbox{\texttt{\mdseries\slshape vertex}} is in a quiver. 

}

 

\subsection{\textcolor{Chapter }{OutgoingArrowsOfVertex}}
\logpage{[ 3, 8, 2 ]}\nobreak
\hyperdef{L}{X8345D79381E46D73}{}
{\noindent\textcolor{FuncColor}{$\triangleright$\enspace\texttt{OutgoingArrowsOfVertex({\mdseries\slshape vertex})\index{OutgoingArrowsOfVertex@\texttt{OutgoingArrowsOfVertex}}
\label{OutgoingArrowsOfVertex}
}\hfill{\scriptsize (attribute)}}\\
\textbf{\indent Returns:\ }
 a list of arrows having \mbox{\texttt{\mdseries\slshape vertex}} as source. 

}

 

\subsection{\textcolor{Chapter }{InDegreeOfVertex}}
\logpage{[ 3, 8, 3 ]}\nobreak
\hyperdef{L}{X7C9CD0527CB9E6EF}{}
{\noindent\textcolor{FuncColor}{$\triangleright$\enspace\texttt{InDegreeOfVertex({\mdseries\slshape vertex})\index{InDegreeOfVertex@\texttt{InDegreeOfVertex}}
\label{InDegreeOfVertex}
}\hfill{\scriptsize (attribute)}}\\
\textbf{\indent Returns:\ }
 the number of arrows having \mbox{\texttt{\mdseries\slshape vertex}} as target. Only meaningful if \mbox{\texttt{\mdseries\slshape vertex}} is in a quiver. 

}

 

\subsection{\textcolor{Chapter }{OutDegreeOfVertex}}
\logpage{[ 3, 8, 4 ]}\nobreak
\hyperdef{L}{X7A09EB648070276D}{}
{\noindent\textcolor{FuncColor}{$\triangleright$\enspace\texttt{OutDegreeOfVertex({\mdseries\slshape vertex})\index{OutDegreeOfVertex@\texttt{OutDegreeOfVertex}}
\label{OutDegreeOfVertex}
}\hfill{\scriptsize (attribute)}}\\
\textbf{\indent Returns:\ }
 the number of arrows having \mbox{\texttt{\mdseries\slshape vertex}} as source. 

}

 

\subsection{\textcolor{Chapter }{NeighborsOfVertex}}
\logpage{[ 3, 8, 5 ]}\nobreak
\hyperdef{L}{X7A557B4C83B7C601}{}
{\noindent\textcolor{FuncColor}{$\triangleright$\enspace\texttt{NeighborsOfVertex({\mdseries\slshape vertex})\index{NeighborsOfVertex@\texttt{NeighborsOfVertex}}
\label{NeighborsOfVertex}
}\hfill{\scriptsize (attribute)}}\\
\textbf{\indent Returns:\ }
 a list of neighbors of \mbox{\texttt{\mdseries\slshape vertex}}, that is, vertices that are targets of arrows having \mbox{\texttt{\mdseries\slshape vertex}} as source. 

}

 
\begin{Verbatim}[commandchars=!@|,fontsize=\small,frame=single,label=Example]
  !gapprompt@gap>| !gapinput@q1 := Quiver(["u","v"],[["u","u","a"],["u","v","b"],|
  !gapprompt@>| !gapinput@              ["v","u","c"],["v","v","d"]]);|
  <quiver with 2 vertices and 4 arrows>
  !gapprompt@gap>| !gapinput@OutgoingArrowsOfVertex(q1.u);|
  [ a, b ]
  !gapprompt@gap>| !gapinput@InDegreeOfVertex(q1.u);|
  2
  !gapprompt@gap>| !gapinput@NeighborsOfVertex(q1.v);|
  [ u, v ]
\end{Verbatim}
 }

 
\section{\textcolor{Chapter }{Posets}}\logpage{[ 3, 9, 0 ]}
\hyperdef{L}{X79540DAB85902432}{}
{
 This implementation of posets was done by the
HomAlg\texttt{\symbol{45}}project. 

\subsection{\textcolor{Chapter }{Poset (for a list P and a set of relations rel)}}
\logpage{[ 3, 9, 1 ]}\nobreak
\hyperdef{L}{X80854BB778E3833E}{}
{\noindent\textcolor{FuncColor}{$\triangleright$\enspace\texttt{Poset({\mdseries\slshape P, rel})\index{Poset@\texttt{Poset}!for a list P and a set of relations rel}
\label{Poset:for a list P and a set of relations rel}
}\hfill{\scriptsize (operation)}}\\


 Arguments: \mbox{\texttt{\mdseries\slshape P}} a list and \mbox{\texttt{\mdseries\slshape rel}} \texttt{\symbol{45}}\texttt{\symbol{45}} a list of pairs from \mbox{\texttt{\mdseries\slshape P}}. \\
\textbf{\indent Returns:\ }
 the poset defined on the points \mbox{\texttt{\mdseries\slshape P}} and the relations generated by \mbox{\texttt{\mdseries\slshape rel}}. 



 The elements in \mbox{\texttt{\mdseries\slshape P}} is given as a list, and for example \texttt{["a", "b", "c", "d"]} and the relations are given as a list of lists, for instance in the above
case: \texttt{[ ["a", "b", "c"], ["b", "d"], ["c", "d"]].} The first list means that $a < b$ and $a < c$, and the second one means $b < d$ and finally the last one means $c < d$. }

 

\subsection{\textcolor{Chapter }{Size}}
\logpage{[ 3, 9, 2 ]}\nobreak
\hyperdef{L}{X858ADA3B7A684421}{}
{\noindent\textcolor{FuncColor}{$\triangleright$\enspace\texttt{Size({\mdseries\slshape P})\index{Size@\texttt{Size}}
\label{Size}
}\hfill{\scriptsize (attribute)}}\\
\textbf{\indent Returns:\ }
 the number of elements of the poset \mbox{\texttt{\mdseries\slshape P}}. 

}

 

\subsection{\textcolor{Chapter }{UnderlyingSet}}
\logpage{[ 3, 9, 3 ]}\nobreak
\hyperdef{L}{X7AF9E0CD850F8B03}{}
{\noindent\textcolor{FuncColor}{$\triangleright$\enspace\texttt{UnderlyingSet({\mdseries\slshape P})\index{UnderlyingSet@\texttt{UnderlyingSet}}
\label{UnderlyingSet}
}\hfill{\scriptsize (operation)}}\\


 Arguments: \mbox{\texttt{\mdseries\slshape P}} \texttt{\symbol{45}}\texttt{\symbol{45}} poset. \\
\textbf{\indent Returns:\ }
 the underlying set of the poset \mbox{\texttt{\mdseries\slshape P}}. 

}

 

\subsection{\textcolor{Chapter }{PartialOrderOfPoset}}
\logpage{[ 3, 9, 4 ]}\nobreak
\hyperdef{L}{X8468B5F77EBA547A}{}
{\noindent\textcolor{FuncColor}{$\triangleright$\enspace\texttt{PartialOrderOfPoset({\mdseries\slshape P})\index{PartialOrderOfPoset@\texttt{PartialOrderOfPoset}}
\label{PartialOrderOfPoset}
}\hfill{\scriptsize (operation)}}\\


 Arguments: \mbox{\texttt{\mdseries\slshape P}} \texttt{\symbol{45}}\texttt{\symbol{45}} poset. \\
\textbf{\indent Returns:\ }
 the partial order of the poset \mbox{\texttt{\mdseries\slshape P}} as a function. 

}

 }

 }

 
\chapter{\textcolor{Chapter }{Path Algebras}}\label{PathAlgebras}
\logpage{[ 4, 0, 0 ]}
\hyperdef{L}{X7E8A43A484CE0BA8}{}
{
  
\section{\textcolor{Chapter }{Introduction}}\logpage{[ 4, 1, 0 ]}
\hyperdef{L}{X7DFB63A97E67C0A1}{}
{
  A path algebra is an algebra constructed from a field $F$ (see Chapter 56 and 57 in the \textsf{GAP} manual for information about fields) and a quiver $Q$. The path algebra $FQ$ contains all finite linear combinations of paths of $Q$. This chapter describes the functions in \textsf{QPA} that deal with path algebras and quotients of path algebras. Path algebras are
algebras, so see Chapter 60: Algebras in the \textsf{GAP} manual for functionality such as generators, basis functions, and mappings.

 The only supported ordering of elements in a path algebra is length left
lexicographic ordering. See \ref{qpa:Orderings} for more information.

        }

 
\section{\textcolor{Chapter }{Constructing Path Algebras}}\label{Constructing Paths}
\logpage{[ 4, 2, 0 ]}
\hyperdef{L}{X848A225A84A15B1E}{}
{
  

\subsection{\textcolor{Chapter }{PathAlgebra}}
\logpage{[ 4, 2, 1 ]}\nobreak
\hyperdef{L}{X7CA1C87B8202C2E9}{}
{\noindent\textcolor{FuncColor}{$\triangleright$\enspace\texttt{PathAlgebra({\mdseries\slshape F, Q})\index{PathAlgebra@\texttt{PathAlgebra}}
\label{PathAlgebra}
}\hfill{\scriptsize (function)}}\\


Arguments: \mbox{\texttt{\mdseries\slshape F}} \texttt{\symbol{45}}\texttt{\symbol{45}} a field, \mbox{\texttt{\mdseries\slshape Q}} \texttt{\symbol{45}}\texttt{\symbol{45}} a quiver.\\
\textbf{\indent Returns:\ }
the path algebra \mbox{\texttt{\mdseries\slshape FQ}} of \mbox{\texttt{\mdseries\slshape Q}} over the field \mbox{\texttt{\mdseries\slshape F}}.



For construction of fields, see the \textsf{GAP} documentation. The elements of the path algebra \mbox{\texttt{\mdseries\slshape FQ}} will be ordered by left length\texttt{\symbol{45}}lexicographic ordering. }

 
\begin{Verbatim}[commandchars=!@|,fontsize=\small,frame=single,label=Example]
  !gapprompt@gap>| !gapinput@Q := Quiver( ["u","v"] , [ ["u","u","a"], ["u","v","b"], |
  !gapprompt@>| !gapinput@["v","u","c"], ["v","v","d"] ] );|
  <quiver with 2 vertices and 4 arrows>
  !gapprompt@gap>| !gapinput@F := Rationals;|
  Rationals
  !gapprompt@gap>| !gapinput@FQ := PathAlgebra(F,Q);|
  <Rationals[<quiver with 2 vertices and 4 arrows>]>
\end{Verbatim}
            }

 
\section{\textcolor{Chapter }{Categories and Properties of Path Algebras}}\logpage{[ 4, 3, 0 ]}
\hyperdef{L}{X85A3A8767E7C11AD}{}
{
  

\subsection{\textcolor{Chapter }{IsPathAlgebra}}
\logpage{[ 4, 3, 1 ]}\nobreak
\hyperdef{L}{X8255FDF78315E1B3}{}
{\noindent\textcolor{FuncColor}{$\triangleright$\enspace\texttt{IsPathAlgebra({\mdseries\slshape object})\index{IsPathAlgebra@\texttt{IsPathAlgebra}}
\label{IsPathAlgebra}
}\hfill{\scriptsize (property)}}\\


 Arguments: \mbox{\texttt{\mdseries\slshape object}} \texttt{\symbol{45}}\texttt{\symbol{45}} any object in \textsf{GAP}. \\
 \textbf{\indent Returns:\ }
true whenever \mbox{\texttt{\mdseries\slshape object}} is a path algebra. 

}

 
\begin{Verbatim}[commandchars=!@|,fontsize=\small,frame=single,label=Example]
  !gapprompt@gap>| !gapinput@IsPathAlgebra(FQ);|
  true
\end{Verbatim}
 }

 
\section{\textcolor{Chapter }{Attributes and Operations for Path Algebras}}\logpage{[ 4, 4, 0 ]}
\hyperdef{L}{X7DE2F2A48492041A}{}
{
 

\subsection{\textcolor{Chapter }{AssociatedGradedAlgebra}}
\logpage{[ 4, 4, 1 ]}\nobreak
\hyperdef{L}{X7D8F5BDF85DA7B14}{}
{\noindent\textcolor{FuncColor}{$\triangleright$\enspace\texttt{AssociatedGradedAlgebra({\mdseries\slshape A})\index{AssociatedGradedAlgebra@\texttt{AssociatedGradedAlgebra}}
\label{AssociatedGradedAlgebra}
}\hfill{\scriptsize (attribute)}}\\


 Arguments: \mbox{\texttt{\mdseries\slshape A}} \texttt{\symbol{45}}\texttt{\symbol{45}} a quiver algebra.\\
 \textbf{\indent Returns:\ }
 the associated graded algebra of \mbox{\texttt{\mdseries\slshape A}} with to the radical filtration of the algebra. 

}

 

\subsection{\textcolor{Chapter }{AssociatedMonomialAlgebra}}
\logpage{[ 4, 4, 2 ]}\nobreak
\hyperdef{L}{X7DF51D297E0E6A8B}{}
{\noindent\textcolor{FuncColor}{$\triangleright$\enspace\texttt{AssociatedMonomialAlgebra({\mdseries\slshape A})\index{AssociatedMonomialAlgebra@\texttt{AssociatedMonomialAlgebra}}
\label{AssociatedMonomialAlgebra}
}\hfill{\scriptsize (attribute)}}\\


 Arguments: \mbox{\texttt{\mdseries\slshape A}} \texttt{\symbol{45}}\texttt{\symbol{45}} a quiver algebra.\\
 \textbf{\indent Returns:\ }
 the associated monomial algebra of \mbox{\texttt{\mdseries\slshape A}} with respect to the Groebner basis the path algebra is endoved with. 

}

 

\subsection{\textcolor{Chapter }{QuiverOfPathAlgebra}}
\logpage{[ 4, 4, 3 ]}\nobreak
\hyperdef{L}{X83FBA499856580B0}{}
{\noindent\textcolor{FuncColor}{$\triangleright$\enspace\texttt{QuiverOfPathAlgebra({\mdseries\slshape FQ})\index{QuiverOfPathAlgebra@\texttt{QuiverOfPathAlgebra}}
\label{QuiverOfPathAlgebra}
}\hfill{\scriptsize (attribute)}}\\


 Arguments: \mbox{\texttt{\mdseries\slshape FQ}} \texttt{\symbol{45}}\texttt{\symbol{45}} a path algebra. \\
 \textbf{\indent Returns:\ }
 the quiver from which \mbox{\texttt{\mdseries\slshape FQ}} was constructed. 

}

 
\begin{Verbatim}[commandchars=!@|,fontsize=\small,frame=single,label=Example]
  !gapprompt@gap>| !gapinput@QuiverOfPathAlgebra(FQ);|
  <quiver with 2 vertices and 4 arrows> 
\end{Verbatim}
 

\subsection{\textcolor{Chapter }{OrderingOfAlgebra}}
\logpage{[ 4, 4, 4 ]}\nobreak
\hyperdef{L}{X8279084B828E5FD7}{}
{\noindent\textcolor{FuncColor}{$\triangleright$\enspace\texttt{OrderingOfAlgebra({\mdseries\slshape FQ})\index{OrderingOfAlgebra@\texttt{OrderingOfAlgebra}}
\label{OrderingOfAlgebra}
}\hfill{\scriptsize (attribute)}}\\


 Arguments: \mbox{\texttt{\mdseries\slshape FQ}} \texttt{\symbol{45}}\texttt{\symbol{45}} a path algebra.\\
 \textbf{\indent Returns:\ }
 the ordering of the quiver of the path algebra. 



\emph{Note:} As of the current version of \textsf{QPA}, only left length lexicographic ordering is supported. }

 

\subsection{\textcolor{Chapter }{. (for a path algebra)}}
\logpage{[ 4, 4, 5 ]}\nobreak
\hyperdef{L}{X86CDD46F7F05ADE9}{}
{\noindent\textcolor{FuncColor}{$\triangleright$\enspace\texttt{.({\mdseries\slshape FQ, generator})\index{.@\texttt{.}!for a path algebra}
\label{.:for a path algebra}
}\hfill{\scriptsize (operation)}}\\


 Arguments: \mbox{\texttt{\mdseries\slshape FQ}} \texttt{\symbol{45}}\texttt{\symbol{45}} a path algebra, \mbox{\texttt{\mdseries\slshape generator}} \texttt{\symbol{45}}\texttt{\symbol{45}} a vertex or an arrow in the quiver \mbox{\texttt{\mdseries\slshape Q}}. \\
 \textbf{\indent Returns:\ }
 the \mbox{\texttt{\mdseries\slshape generator}} as an element of the path algebra. 



Other elements of the path algebra can be constructed as linear combinations
of the generators. For further operations on elements, see below. }

 
\begin{Verbatim}[commandchars=!@|,fontsize=\small,frame=single,label=Example]
  !gapprompt@gap>| !gapinput@FQ.a;|
  (1)*a
  !gapprompt@gap>| !gapinput@FQ.v;|
  (1)*v
  !gapprompt@gap>| !gapinput@elem := 2*FQ.a - 3*FQ.v;|
  (-3)*v+(2)*a
\end{Verbatim}
 }

 
\section{\textcolor{Chapter }{Operations on Path Algebra Elements}}\logpage{[ 4, 5, 0 ]}
\hyperdef{L}{X7CEF60107CE4616B}{}
{
 

\subsection{\textcolor{Chapter }{ElementOfPathAlgebra}}
\logpage{[ 4, 5, 1 ]}\nobreak
\hyperdef{L}{X841C00E87E19528E}{}
{\noindent\textcolor{FuncColor}{$\triangleright$\enspace\texttt{ElementOfPathAlgebra({\mdseries\slshape PA, path})\index{ElementOfPathAlgebra@\texttt{ElementOfPathAlgebra}}
\label{ElementOfPathAlgebra}
}\hfill{\scriptsize (operation)}}\\


 Arguments: \mbox{\texttt{\mdseries\slshape PA}} \texttt{\symbol{45}}\texttt{\symbol{45}} a path algebra, \mbox{\texttt{\mdseries\slshape path}} \texttt{\symbol{45}}\texttt{\symbol{45}} a path in the quiver from which \mbox{\texttt{\mdseries\slshape PA}} was constructed.\\
 \textbf{\indent Returns:\ }
The embedding of \mbox{\texttt{\mdseries\slshape path}} into the path algebra \mbox{\texttt{\mdseries\slshape PA}}, or it returns false if \mbox{\texttt{\mdseries\slshape path}} is not an element of the quiver from which \mbox{\texttt{\mdseries\slshape PA}} was constructed. 

}

 

\subsection{\textcolor{Chapter }{{\textless} (for two elements in a path algebra)}}
\logpage{[ 4, 5, 2 ]}\nobreak
\hyperdef{L}{X7E3FAB1F803E26FF}{}
{\noindent\textcolor{FuncColor}{$\triangleright$\enspace\texttt{{\textless}({\mdseries\slshape a, b})\index{{\textless}@\texttt{{\textless}}!for two elements in a path algebra}
\label{<:for two elements in a path algebra}
}\hfill{\scriptsize (operation)}}\\


 Arguments: \mbox{\texttt{\mdseries\slshape a}} and \mbox{\texttt{\mdseries\slshape b}} \texttt{\symbol{45}}\texttt{\symbol{45}} two elements of the same path
algebra.\\
 \textbf{\indent Returns:\ }
 True whenever \mbox{\texttt{\mdseries\slshape a}} is smaller than \mbox{\texttt{\mdseries\slshape b}}, according to the ordering of the path algebra. 

}

 

\subsection{\textcolor{Chapter }{IsLeftUniform}}
\logpage{[ 4, 5, 3 ]}\nobreak
\hyperdef{L}{X853C8B0B8665BFBB}{}
{\noindent\textcolor{FuncColor}{$\triangleright$\enspace\texttt{IsLeftUniform({\mdseries\slshape element})\index{IsLeftUniform@\texttt{IsLeftUniform}}
\label{IsLeftUniform}
}\hfill{\scriptsize (operation)}}\\


 Arguments: \mbox{\texttt{\mdseries\slshape element}} \texttt{\symbol{45}}\texttt{\symbol{45}} an element of the path algebra.\\
 \textbf{\indent Returns:\ }
 true if each monomial in \mbox{\texttt{\mdseries\slshape element}} has the same source vertex, false otherwise. 

}

 

\subsection{\textcolor{Chapter }{IsRightUniform}}
\logpage{[ 4, 5, 4 ]}\nobreak
\hyperdef{L}{X7C06BE7483992634}{}
{\noindent\textcolor{FuncColor}{$\triangleright$\enspace\texttt{IsRightUniform({\mdseries\slshape element})\index{IsRightUniform@\texttt{IsRightUniform}}
\label{IsRightUniform}
}\hfill{\scriptsize (operation)}}\\


 Arguments: \mbox{\texttt{\mdseries\slshape element}} \texttt{\symbol{45}}\texttt{\symbol{45}} an element of the path algebra.\\
 \textbf{\indent Returns:\ }
 true if each monomial in \mbox{\texttt{\mdseries\slshape element}} has the same target vertex, false otherwise. 

}

 

\subsection{\textcolor{Chapter }{IsUniform}}
\logpage{[ 4, 5, 5 ]}\nobreak
\hyperdef{L}{X8735FBE180797557}{}
{\noindent\textcolor{FuncColor}{$\triangleright$\enspace\texttt{IsUniform({\mdseries\slshape element})\index{IsUniform@\texttt{IsUniform}}
\label{IsUniform}
}\hfill{\scriptsize (operation)}}\\


 Arguments: \mbox{\texttt{\mdseries\slshape element}} \texttt{\symbol{45}}\texttt{\symbol{45}} an element of the path algebra. \\
 \textbf{\indent Returns:\ }
 true whenever \mbox{\texttt{\mdseries\slshape element}} is both left and right uniform. 

}

 
\begin{Verbatim}[commandchars=!@|,fontsize=\small,frame=single,label=Example]
  !gapprompt@gap>| !gapinput@IsLeftUniform(elem);|
  false
  !gapprompt@gap>| !gapinput@IsRightUniform(elem);|
  false
  !gapprompt@gap>| !gapinput@IsUniform(elem);|
  false
  !gapprompt@gap>| !gapinput@another := FQ.a*FQ.b + FQ.b*FQ.d*FQ.c*FQ.b*FQ.d;|
  (1)*a*b+(1)*b*d*c*b*d
  !gapprompt@gap>| !gapinput@IsLeftUniform(another);|
  true
  !gapprompt@gap>| !gapinput@IsRightUniform(another);|
  true
  !gapprompt@gap>| !gapinput@IsUniform(another);|
  true 
\end{Verbatim}
 

\subsection{\textcolor{Chapter }{LeadingTerm}}
\logpage{[ 4, 5, 6 ]}\nobreak
\hyperdef{L}{X84C98E687A3A84D8}{}
{\noindent\textcolor{FuncColor}{$\triangleright$\enspace\texttt{LeadingTerm({\mdseries\slshape element})\index{LeadingTerm@\texttt{LeadingTerm}}
\label{LeadingTerm}
}\hfill{\scriptsize (operation)}}\\
\noindent\textcolor{FuncColor}{$\triangleright$\enspace\texttt{Tip({\mdseries\slshape element})\index{Tip@\texttt{Tip}}
\label{Tip}
}\hfill{\scriptsize (operation)}}\\


 Arguments: \mbox{\texttt{\mdseries\slshape element}} \texttt{\symbol{45}}\texttt{\symbol{45}} an element of the path algebra. \\
 \textbf{\indent Returns:\ }
 the term in \mbox{\texttt{\mdseries\slshape element}} whose monomial is largest among those monomials that have nonzero coefficients
(known as the "tip" of \mbox{\texttt{\mdseries\slshape element}}). 



 \emph{Note: } The two operations are equivalent. }

 

\subsection{\textcolor{Chapter }{LeadingCoefficient}}
\logpage{[ 4, 5, 7 ]}\nobreak
\hyperdef{L}{X80710E9B7D8340BD}{}
{\noindent\textcolor{FuncColor}{$\triangleright$\enspace\texttt{LeadingCoefficient({\mdseries\slshape element})\index{LeadingCoefficient@\texttt{LeadingCoefficient}}
\label{LeadingCoefficient}
}\hfill{\scriptsize (operation)}}\\
\noindent\textcolor{FuncColor}{$\triangleright$\enspace\texttt{TipCoefficient({\mdseries\slshape element})\index{TipCoefficient@\texttt{TipCoefficient}}
\label{TipCoefficient}
}\hfill{\scriptsize (operation)}}\\


 Arguments: \mbox{\texttt{\mdseries\slshape element}} \texttt{\symbol{45}}\texttt{\symbol{45}} an element of the path algebra. \\
 \textbf{\indent Returns:\ }
 the coefficient of the tip of \mbox{\texttt{\mdseries\slshape element}} (which is an element of the field). 



 \emph{Note: } The two operations are equivalent. }

 

\subsection{\textcolor{Chapter }{LeadingMonomial}}
\logpage{[ 4, 5, 8 ]}\nobreak
\hyperdef{L}{X7B3EAE41795598A5}{}
{\noindent\textcolor{FuncColor}{$\triangleright$\enspace\texttt{LeadingMonomial({\mdseries\slshape element})\index{LeadingMonomial@\texttt{LeadingMonomial}}
\label{LeadingMonomial}
}\hfill{\scriptsize (operation)}}\\
\noindent\textcolor{FuncColor}{$\triangleright$\enspace\texttt{TipMonomial({\mdseries\slshape element})\index{TipMonomial@\texttt{TipMonomial}}
\label{TipMonomial}
}\hfill{\scriptsize (operation)}}\\


 Arguments: \mbox{\texttt{\mdseries\slshape element}} \texttt{\symbol{45}}\texttt{\symbol{45}} an element of the path algebra. \\
 \textbf{\indent Returns:\ }
 the monomial of the tip of \mbox{\texttt{\mdseries\slshape element}} (which is an element of the underlying quiver, not of the path algebra). 



 \emph{Note: } The two operations are equivalent. }

 
\begin{Verbatim}[commandchars=!@|,fontsize=\small,frame=single,label=Example]
  !gapprompt@gap>| !gapinput@elem := FQ.a*FQ.b*FQ.c + FQ.b*FQ.d*FQ.c+FQ.d*FQ.d;|
  (1)*d^2+(1)*a*b*c+(1)*b*d*c
  !gapprompt@gap>| !gapinput@LeadingTerm(elem);|
  (1)*b*d*c
  !gapprompt@gap>| !gapinput@LeadingCoefficient(elem);|
  1
  !gapprompt@gap>| !gapinput@mon := LeadingMonomial(elem);|
  b*d*c
  !gapprompt@gap>| !gapinput@mon in FQ;|
  false
  !gapprompt@gap>| !gapinput@mon in Q;|
  true 
\end{Verbatim}
 

\subsection{\textcolor{Chapter }{MakeUniformOnRight}}
\logpage{[ 4, 5, 9 ]}\nobreak
\hyperdef{L}{X8172B40181E1B7D2}{}
{\noindent\textcolor{FuncColor}{$\triangleright$\enspace\texttt{MakeUniformOnRight({\mdseries\slshape elems})\index{MakeUniformOnRight@\texttt{MakeUniformOnRight}}
\label{MakeUniformOnRight}
}\hfill{\scriptsize (operation)}}\\


 Arguments: \mbox{\texttt{\mdseries\slshape elems}} \texttt{\symbol{45}}\texttt{\symbol{45}} a list of elements in a path algebra.\\
 \textbf{\indent Returns:\ }
 a list of right uniform elements generated by each element of \mbox{\texttt{\mdseries\slshape elems}}. 

}

 

\subsection{\textcolor{Chapter }{MappedExpression}}
\logpage{[ 4, 5, 10 ]}\nobreak
\hyperdef{L}{X796249A682818750}{}
{\noindent\textcolor{FuncColor}{$\triangleright$\enspace\texttt{MappedExpression({\mdseries\slshape expr, gens1, gens2})\index{MappedExpression@\texttt{MappedExpression}}
\label{MappedExpression}
}\hfill{\scriptsize (operation)}}\\


 Arguments: \mbox{\texttt{\mdseries\slshape expr}} \texttt{\symbol{45}}\texttt{\symbol{45}} element of a path algebra, \mbox{\texttt{\mdseries\slshape gens1}} and \mbox{\texttt{\mdseries\slshape gens2}} \texttt{\symbol{45}}\texttt{\symbol{45}} equal\texttt{\symbol{45}}length lists
of generators for subalgebras.\\
 \textbf{\indent Returns:\ }
 \mbox{\texttt{\mdseries\slshape expr}} as an element of the subalgebra generated by \mbox{\texttt{\mdseries\slshape gens2}}. 



 The element \mbox{\texttt{\mdseries\slshape expr}} must be in the subalgebra generated by \mbox{\texttt{\mdseries\slshape gens1}}. The lists define a mapping of each generator in \mbox{\texttt{\mdseries\slshape gens1}} to the corresponding generator in \mbox{\texttt{\mdseries\slshape gens2}}. The value returned is the evaluation of the mapping at \mbox{\texttt{\mdseries\slshape expr}}. }

 

\subsection{\textcolor{Chapter }{SupportOfQuiverAlgebraElement}}
\logpage{[ 4, 5, 11 ]}\nobreak
\hyperdef{L}{X7D6DDDF178B0F2D9}{}
{\noindent\textcolor{FuncColor}{$\triangleright$\enspace\texttt{SupportOfQuiverAlgebraElement({\mdseries\slshape algebra, element})\index{SupportOfQuiverAlgebraElement@\texttt{SupportOfQuiverAlgebraElement}}
\label{SupportOfQuiverAlgebraElement}
}\hfill{\scriptsize (operation)}}\\
\noindent\textcolor{FuncColor}{$\triangleright$\enspace\texttt{LeftSupportOfQuiverAlgebraElement({\mdseries\slshape algebra, element})\index{LeftSupportOfQuiverAlgebraElement@\texttt{LeftSupportOfQuiverAlgebraElement}}
\label{LeftSupportOfQuiverAlgebraElement}
}\hfill{\scriptsize (operation)}}\\
\noindent\textcolor{FuncColor}{$\triangleright$\enspace\texttt{RIghtSupportOfQuiverAlgebraElement({\mdseries\slshape algebra, element})\index{RIghtSupportOfQuiverAlgebraElement@\texttt{RIghtSupportOfQuiverAlgebraElement}}
\label{RIghtSupportOfQuiverAlgebraElement}
}\hfill{\scriptsize (operation)}}\\


 Arguments: \mbox{\texttt{\mdseries\slshape algebra}} \texttt{\symbol{45}}\texttt{\symbol{45}} a QuiverAlgebra, \mbox{\texttt{\mdseries\slshape element}} \texttt{\symbol{45}}\texttt{\symbol{45}} an element of the QuiverAlgebra\\
 \textbf{\indent Returns:\ }
the second version returns a list of the index of the vertices such that the
product from the left is non\texttt{\symbol{45}}zero, the third version
returns a list of the index of the vertices such that the product from the
right is non\texttt{\symbol{45}}zero, and the first version returns the both
of the previous lists of indices of vertices. 

}

 

\subsection{\textcolor{Chapter }{VertexPosition}}
\logpage{[ 4, 5, 12 ]}\nobreak
\hyperdef{L}{X849AC0F67A131929}{}
{\noindent\textcolor{FuncColor}{$\triangleright$\enspace\texttt{VertexPosition({\mdseries\slshape element})\index{VertexPosition@\texttt{VertexPosition}}
\label{VertexPosition}
}\hfill{\scriptsize (operation)}}\\


 Arguments: \mbox{\texttt{\mdseries\slshape element}} \texttt{\symbol{45}}\texttt{\symbol{45}} an element of the path algebra on the
form $k*v$, where $v$ is a vertex of the underlying quiver and $k$ is an element of the field.\\
 \textbf{\indent Returns:\ }
 the position of the vertex \mbox{\texttt{\mdseries\slshape v}} in the list of vertices of the quiver. 

}

 

\subsection{\textcolor{Chapter }{RelationsOfAlgebra}}
\logpage{[ 4, 5, 13 ]}\nobreak
\hyperdef{L}{X86795D8E7ED73048}{}
{\noindent\textcolor{FuncColor}{$\triangleright$\enspace\texttt{RelationsOfAlgebra({\mdseries\slshape A})\index{RelationsOfAlgebra@\texttt{RelationsOfAlgebra}}
\label{RelationsOfAlgebra}
}\hfill{\scriptsize (attribute)}}\\


 Arguments: \mbox{\texttt{\mdseries\slshape A}} \texttt{\symbol{45}}\texttt{\symbol{45}} a quiver algebra. \\
 \textbf{\indent Returns:\ }
a set of generators for the ideal in the path algebra $kQ$ from which the algebra \mbox{\texttt{\mdseries\slshape FQ}} was constructed. If \mbox{\texttt{\mdseries\slshape A}} is a path algebra, then an empty list is returned. 

}

 }

 
\section{\textcolor{Chapter }{Constructing Quotients of Path Algebras}}\label{qpa:Quotients_of_path_algebras}
\logpage{[ 4, 6, 0 ]}
\hyperdef{L}{X7F0D555379C97A6E}{}
{
  In the introduction we saw already one way of constructing a quotient of a
path algebra. In addition to this there are at least two other ways of
constructing a quotient of a path algebra; one with factoring out an ideal and
one where a Groebner basis is attached to the quotient. We discuss these two
next.

 For several functions in \textsf{QPA} to function properly one needs to have a Groebner basis attached to the
quotient one wants to construct, or equivalently a Groebner basis for the
ideal one is factoring out. For this to work the ideal must admit a finite
Groebner basis. However, to our knowledge there is no algorithm for
determining if an ideal has a finite Groebner basis. On the other hand, it is
known that if the factor algebra is finite dimensional, then the ideal has a
finite Groebner basis (independent of the ordering of the elements, see \cite{Green} ). In addition to having a finite Groebner basis, several functions also need
that the factoring ideal is admissible. A quotient of a path algebra by an
admissible ideal belongs to the category \texttt{IsAdmissibleQuotientOfPathAlgebra} (\ref{IsAdmissibleQuotientOfPathAlgebra}). The method used in the introduction constructs a quotient in this category.
However, there are situations where it is interesting to analyze quotients of
path algebras by a non\texttt{\symbol{45}}admissible ideal, so we provide also
additional methods.

 In the example below, we construct a factor of a path algebra purely with
commands in \textsf{GAP} (cf. also Chapter 60: Algebras in the \textsf{GAP} manual on how to construct an ideal and a quotient of an algebra). Functions
which use Groebner bases like \texttt{IsFiniteDimensional} (\ref{IsFiniteDimensional}), \texttt{Dimension} (\ref{Dimension}), \texttt{IsSpecialBiserialAlgebra} (\ref{IsSpecialBiserialAlgebra}) or a membership test \texttt{\texttt{\symbol{92}}in} (\ref{bSlashin:elt. in path alg. and ideal}) will work properly (they simply compute the Groebner basis if it is
necessary). But some "older" functions (like \texttt{IndecProjectiveModules} (\ref{IndecProjectiveModules})) can fail or give an incorrect answer! This way of constructing a quotient of
a path algebra can be useful e.g. if we know that computing a Groebner basis
will take a long time and we do not need this because we want to deal only
with modules.

 
\begin{Verbatim}[commandchars=!@|,fontsize=\small,frame=single,label=Example]
  !gapprompt@gap>| !gapinput@Q := Quiver( 1, [ [1,1,"a"], [1,1,"b"] ] );|
  <quiver with 1 vertices and 2 arrows>
  !gapprompt@gap>| !gapinput@kQ := PathAlgebra(Rationals, Q);|
  <Rationals[<quiver with 1 vertices and 2 arrows>]>
  !gapprompt@gap>| !gapinput@gens := GeneratorsOfAlgebra(kQ);|
  [ (1)*v1, (1)*a, (1)*b ]
  !gapprompt@gap>| !gapinput@a := gens[2];|
  (1)*a
  !gapprompt@gap>| !gapinput@b := gens[3];|
  (1)*b
  !gapprompt@gap>| !gapinput@relations := [a^2,a*b-b*a, b*b];|
  [ (1)*a^2, (1)*a*b+(-1)*b*a, (1)*b^2 ]
  !gapprompt@gap>| !gapinput@I := Ideal(kQ,relations);|
  <two-sided ideal in <Rationals[<quiver with 1 vertices and 2 arrows>]>
      , (3 generators)>
  !gapprompt@gap>| !gapinput@A := kQ/I;|
  <Rationals[<quiver with 1 vertices and 2 arrows>]/
  <two-sided ideal in <Rationals[<quiver with 1 vertices and 2 arrows>]>
      , (3 generators)>>
  !gapprompt@gap>| !gapinput@IndecProjectiveModules(A);|
  Compute a Groebner basis of the ideal you are factoring out with before \
  you form the quotient algebra, or you have entered an algebra which \
  is not finite dimensional.
  fail
\end{Verbatim}
 To resolve this matter, we need to compute the Gr{\"o}bner basis of the ideal
generated by the relations in $kQ$ (yes, it seems like we are going in circles here. Remember, then, that an
ideal in the "mathematical sense" may exist independently of the a
corresponding \texttt{Ideal} object in \textsf{GAP}. Also, Gr{\"o}bner bases in \textsf{QPA} are handled by the \textsf{GBNP} package, with constructor methods not dependent on \texttt{Ideal} objects. After creating the ideal $I$, we need to perform yet another Gr{\"o}bner basis operation which just set a
respective attribute for $I$, see \texttt{GroebnerBasis} (\ref{GroebnerBasis}). 
\begin{Verbatim}[commandchars=!@|,fontsize=\small,frame=single,label=Example]
  !gapprompt@gap>| !gapinput@gb := GBNPGroebnerBasis(relations,kQ);|
  [ (1)*a^2, (-1)*a*b+(1)*b*a, (1)*b^2 ]
  !gapprompt@gap>| !gapinput@I := Ideal(kQ,gb);                    |
  <two-sided ideal in <Rationals[<quiver with 1 vertices and 2 arrows>]>, 
    (3 generators)>
  !gapprompt@gap>| !gapinput@GroebnerBasis(I,gb);                  |
  <complete two-sided Groebner basis containing 3 elements>
  !gapprompt@gap>| !gapinput@IndecProjectiveModules(A);                                                |
  fail
  !gapprompt@gap>| !gapinput@A := kQ/I;                                                                |
  <Rationals[<quiver with 1 vertices and 2 arrows>]/
  <two-sided ideal in <Rationals[<quiver with 1 vertices and 2 arrows>]>, 
    (3 generators)>>
  !gapprompt@gap>| !gapinput@IndecProjectiveModules(A);|
  [ <[ 4 ]> ]
\end{Verbatim}
 Note that the instruction \texttt{A := kQ/relations;} used in Introduction is exactly an abbreviation for a sequence of instructions
with Groebner basis as in above example.

 Most \textsf{QPA} operations working on algebras handle path algebras and quotients of path
algebras in the same way (when this makes sense). However, there are still a
few operations which does not work properly when given a quotient of a path
algebra. When constructing a quotient of a path algebra one needs define the
ideal one is factoring out. Above this has been done with the commands 
\begin{Verbatim}[commandchars=!@|,fontsize=\small,frame=single,label=Example]
  !gapprompt@gap>| !gapinput@gens := GeneratorsOfAlgebra(kQ);|
  [ (1)*v1, (1)*a, (1)*b ]
  !gapprompt@gap>| !gapinput@a := gens[2];|
  (1)*a
  !gapprompt@gap>| !gapinput@b := gens[3];|
  (1)*b
\end{Verbatim}
 The following command makes this process easier. 

\subsection{\textcolor{Chapter }{AssignGeneratorVariables}}
\logpage{[ 4, 6, 1 ]}\nobreak
\hyperdef{L}{X814203E281F3272E}{}
{\noindent\textcolor{FuncColor}{$\triangleright$\enspace\texttt{AssignGeneratorVariables({\mdseries\slshape A})\index{AssignGeneratorVariables@\texttt{AssignGeneratorVariables}}
\label{AssignGeneratorVariables}
}\hfill{\scriptsize (operation)}}\\


 Arguments: \mbox{\texttt{\mdseries\slshape A}} \texttt{\symbol{45}}\texttt{\symbol{45}} a quiver algebra.\\
 \textbf{\indent Returns:\ }
 Takes a quiver algebra \mbox{\texttt{\mdseries\slshape A}} as an argument and creates variables, say $v_1,...,v_n$ for the vertices, and $a_1,...,a_t$ for the arrows for the corresponding elements in \mbox{\texttt{\mdseries\slshape A}}, whenever the quiver for the quiver algebra \mbox{\texttt{\mdseries\slshape A}} is was constructed with the vertices being named $v_1,...,v_n$ and the arrows being named $a_1,...,a_t$. 

}

 Here is an example of its use. 
\begin{Verbatim}[commandchars=!@|,fontsize=\small,frame=single,label=Example]
  !gapprompt@gap>| !gapinput@AssignGeneratorVariables(kQ);|
  #I  Assigned the global variables [ v1, a, b ]
  !gapprompt@gap>| !gapinput@v1; a; b;|
  (1)*v1
  (1)*a
  (1)*b
\end{Verbatim}
 }

 
\section{\textcolor{Chapter }{Ideals and operations on ideals}}\label{IdealsAndOperations}
\logpage{[ 4, 7, 0 ]}
\hyperdef{L}{X812C0F8D7E4B1134}{}
{
 

\subsection{\textcolor{Chapter }{Ideal}}
\logpage{[ 4, 7, 1 ]}\nobreak
\hyperdef{L}{X82ACACDD7D8E9B25}{}
{\noindent\textcolor{FuncColor}{$\triangleright$\enspace\texttt{Ideal({\mdseries\slshape FQ, elems})\index{Ideal@\texttt{Ideal}}
\label{Ideal}
}\hfill{\scriptsize (function)}}\\


Arguments: \mbox{\texttt{\mdseries\slshape FQ}} \texttt{\symbol{45}}\texttt{\symbol{45}} a path algebra, \mbox{\texttt{\mdseries\slshape elems}} \texttt{\symbol{45}}\texttt{\symbol{45}} a list of elements in \mbox{\texttt{\mdseries\slshape FQ}}.\\
 \textbf{\indent Returns:\ }
the ideal of \mbox{\texttt{\mdseries\slshape FQ}} generated by \mbox{\texttt{\mdseries\slshape elems}} with the property \texttt{IsIdealInPathAlgebra} (\ref{IsIdealInPathAlgebra}). 



For more on ideals, see the \textsf{GAP} reference manual (Chapter 60.6). \\
 \emph{Technical info:} Ideal is a synonym for a global GAP function TwoSidedIdeal which calls an
operation TwoSidedIdealByGenerators (synonym IdealByGenerators) for an algebra
(FLMLOR). }

 
\begin{Verbatim}[commandchars=!@|,fontsize=\small,frame=single,label=Example]
  !gapprompt@gap>| !gapinput@gb := GBNPGroebnerBasis(relations,kQ);|
  [ (1)*a^2, (-1)*a*b+(1)*b*a, (1)*b^2 ]
  !gapprompt@gap>| !gapinput@I := Ideal(kQ,gb);|
  <two-sided ideal in <Rationals[<quiver with 1 vertices and 2 arrows>]>
      , (3 generators)>
  !gapprompt@gap>| !gapinput@GroebnerBasis(I,gb);|
  <complete two-sided Groebner basis containing 3 elements>
  !gapprompt@gap>| !gapinput@IndecProjectiveModules(A);|
  [ <[ 4 ]> ]
  !gapprompt@gap>| !gapinput@A := kQ/I;|
  <Rationals[<quiver with 1 vertices and 2 arrows>]/
  <two-sided ideal in <Rationals[<quiver with 1 vertices and 2 arrows>]>
      , (3 generators)>>
  !gapprompt@gap>| !gapinput@IndecProjectiveModules(A);|
  [ <[ 4 ]> ]
  true
\end{Verbatim}
 

\subsection{\textcolor{Chapter }{IdealOfQuotient}}
\logpage{[ 4, 7, 2 ]}\nobreak
\hyperdef{L}{X7EACC0D285D18E19}{}
{\noindent\textcolor{FuncColor}{$\triangleright$\enspace\texttt{IdealOfQuotient({\mdseries\slshape A})\index{IdealOfQuotient@\texttt{IdealOfQuotient}}
\label{IdealOfQuotient}
}\hfill{\scriptsize (attribute)}}\\


 Arguments: \mbox{\texttt{\mdseries\slshape A}} \texttt{\symbol{45}}\texttt{\symbol{45}} a quiver algebra. \\
 \textbf{\indent Returns:\ }
the ideal in the path algebra $kQ$ from which \mbox{\texttt{\mdseries\slshape A}} was constructed. 

}

 

\subsection{\textcolor{Chapter }{PathsOfLengthTwo}}
\logpage{[ 4, 7, 3 ]}\nobreak
\hyperdef{L}{X859C987B7C5F0D8D}{}
{\noindent\textcolor{FuncColor}{$\triangleright$\enspace\texttt{PathsOfLengthTwo({\mdseries\slshape Q})\index{PathsOfLengthTwo@\texttt{PathsOfLengthTwo}}
\label{PathsOfLengthTwo}
}\hfill{\scriptsize (operation)}}\\


 Arguments: \mbox{\texttt{\mdseries\slshape Q}} \texttt{\symbol{45}}\texttt{\symbol{45}} a quiver.\\
 \textbf{\indent Returns:\ }
 a list of all paths of length two in \mbox{\texttt{\mdseries\slshape Q}}, sorted by \texttt{{\textless}}. Fails with error message if \mbox{\texttt{\mdseries\slshape Q}} is not a Quiver object. 

}

 

\subsection{\textcolor{Chapter }{NthPowerOfArrowIdeal}}
\logpage{[ 4, 7, 4 ]}\nobreak
\hyperdef{L}{X830187497E0BD4F0}{}
{\noindent\textcolor{FuncColor}{$\triangleright$\enspace\texttt{NthPowerOfArrowIdeal({\mdseries\slshape FQ, n})\index{NthPowerOfArrowIdeal@\texttt{NthPowerOfArrowIdeal}}
\label{NthPowerOfArrowIdeal}
}\hfill{\scriptsize (operation)}}\\


 Arguments: \mbox{\texttt{\mdseries\slshape FQ}} \texttt{\symbol{45}}\texttt{\symbol{45}} a path algebra, \mbox{\texttt{\mdseries\slshape n}} \texttt{\symbol{45}}\texttt{\symbol{45}} a positive integer.\\
 \textbf{\indent Returns:\ }
 the ideal generated all the paths of length \mbox{\texttt{\mdseries\slshape n}} in \mbox{\texttt{\mdseries\slshape FQ}}. 

}

 

\subsection{\textcolor{Chapter }{AddNthPowerToRelations}}
\logpage{[ 4, 7, 5 ]}\nobreak
\hyperdef{L}{X824D550E8371098C}{}
{\noindent\textcolor{FuncColor}{$\triangleright$\enspace\texttt{AddNthPowerToRelations({\mdseries\slshape FQ, rels, n})\index{AddNthPowerToRelations@\texttt{AddNthPowerToRelations}}
\label{AddNthPowerToRelations}
}\hfill{\scriptsize (operation)}}\\


 Arguments: \mbox{\texttt{\mdseries\slshape FQ}} \texttt{\symbol{45}}\texttt{\symbol{45}} a path algebra, \mbox{\texttt{\mdseries\slshape rels}} \texttt{\symbol{45}}\texttt{\symbol{45}} a (possibly empty) list of elements
in \mbox{\texttt{\mdseries\slshape FQ}}, \mbox{\texttt{\mdseries\slshape n}} \texttt{\symbol{45}}\texttt{\symbol{45}} a positive integer.\\
 \textbf{\indent Returns:\ }
 the list \mbox{\texttt{\mdseries\slshape rels}} with the paths of length \mbox{\texttt{\mdseries\slshape n}} of \mbox{\texttt{\mdseries\slshape FQ}} appended (will change the list \mbox{\texttt{\mdseries\slshape rels}}). 

}

 

\subsection{\textcolor{Chapter }{\texttt{\symbol{92}}in (elt. in path alg. and ideal)}}
\logpage{[ 4, 7, 6 ]}\nobreak
\hyperdef{L}{X83E8D45B82356D8E}{}
{\noindent\textcolor{FuncColor}{$\triangleright$\enspace\texttt{\texttt{\symbol{92}}in({\mdseries\slshape elt, I})\index{\texttt{\symbol{92}}in@\texttt{\texttt{\symbol{92}}in}!elt. in path alg. and ideal}
\label{bSlashin:elt. in path alg. and ideal}
}\hfill{\scriptsize (operation)}}\\


 Arguments: \mbox{\texttt{\mdseries\slshape elt}} \texttt{\symbol{45}} an element in a path algebra, \mbox{\texttt{\mdseries\slshape I}} \texttt{\symbol{45}} an ideal in the same path algebra (i.e. with \texttt{IsIdealInPathAlgebra} (\ref{IsIdealInPathAlgebra}) property).\\
 \textbf{\indent Returns:\ }
true, if \mbox{\texttt{\mdseries\slshape elt}} belongs to \mbox{\texttt{\mdseries\slshape I}}.



It performs the membership test for an ideal in path algebra using completely
reduced Groebner bases machinery. \\
 \emph{Technical info:} For the efficiency reasons, it computes Groebner basis for \mbox{\texttt{\mdseries\slshape I}} only if it has not been computed yet. Similarly, it performs
CompletelyReduceGroebnerBasis only if it has not been reduced yet. The method
can change the existing Groebner basis.\\
 \emph{Remark:} It works only in case \mbox{\texttt{\mdseries\slshape I}} is in the arrow ideal. }

 }

 
\section{\textcolor{Chapter }{Categories and properties of ideals}}\label{PropertiesOfIdeals}
\logpage{[ 4, 8, 0 ]}
\hyperdef{L}{X82A1683E7A402E73}{}
{
 

\subsection{\textcolor{Chapter }{IsAdmissibleIdeal}}
\logpage{[ 4, 8, 1 ]}\nobreak
\hyperdef{L}{X7F40193B877D76BC}{}
{\noindent\textcolor{FuncColor}{$\triangleright$\enspace\texttt{IsAdmissibleIdeal({\mdseries\slshape I})\index{IsAdmissibleIdeal@\texttt{IsAdmissibleIdeal}}
\label{IsAdmissibleIdeal}
}\hfill{\scriptsize (property)}}\\


 Arguments: \mbox{\texttt{\mdseries\slshape I}} \texttt{\symbol{45}}\texttt{\symbol{45}} an IsIdealInPathAlgebra object. \\
 \textbf{\indent Returns:\ }
true whenever \mbox{\texttt{\mdseries\slshape I}} is an \emph{admissible} ideal in a path algebra, i.e. \mbox{\texttt{\mdseries\slshape I}} is a subset of $J^2$ and \mbox{\texttt{\mdseries\slshape I}} contains $J^n$ for some $n$, where $J$ is the arrow ideal. 



 \emph{Technical note:} The second condition is checked by the nilpotency index of the radical and
checking if the ideal generated by the arrows to one plus this index is in the
ideal of the relations (this uses Groebner bases machinery). }

 

\subsection{\textcolor{Chapter }{IsIdealInPathAlgebra}}
\logpage{[ 4, 8, 2 ]}\nobreak
\hyperdef{L}{X818EE2B9789BB175}{}
{\noindent\textcolor{FuncColor}{$\triangleright$\enspace\texttt{IsIdealInPathAlgebra({\mdseries\slshape I})\index{IsIdealInPathAlgebra@\texttt{IsIdealInPathAlgebra}}
\label{IsIdealInPathAlgebra}
}\hfill{\scriptsize (property)}}\\


 Arguments: \mbox{\texttt{\mdseries\slshape I}} \texttt{\symbol{45}}\texttt{\symbol{45}} an IsFLMLOR object. \\
 \textbf{\indent Returns:\ }
true whenever \mbox{\texttt{\mdseries\slshape I}} is an ideal in a path algebra. 

}

 

\subsection{\textcolor{Chapter }{IsHomogeneousListOfElements}}
\logpage{[ 4, 8, 3 ]}\nobreak
\hyperdef{L}{X830F0CFC7F30B49D}{}
{\noindent\textcolor{FuncColor}{$\triangleright$\enspace\texttt{IsHomogeneousListOfElements({\mdseries\slshape gens, n})\index{IsHomogeneousListOfElements@\texttt{IsHomogeneousListOfElements}}
\label{IsHomogeneousListOfElements}
}\hfill{\scriptsize (operation)}}\\


 Arguments: \mbox{\texttt{\mdseries\slshape gens}}, \mbox{\texttt{\mdseries\slshape n}} \texttt{\symbol{45}}\texttt{\symbol{45}} a list of elements in a path algebra
and a non\texttt{\symbol{45}}negative integer.\\
 \textbf{\indent Returns:\ }
true whenever \mbox{\texttt{\mdseries\slshape gens}} is a list of elements in the linear span of degree \mbox{\texttt{\mdseries\slshape n}} elements or zero of a path algebra. It returns false whenever \mbox{\texttt{\mdseries\slshape gens}} is a list of elements in a path algebra, but not in the linear span of degree \mbox{\texttt{\mdseries\slshape n}} of a path algebra and zero. Otherwise it returns fail. 

}

 

\subsection{\textcolor{Chapter }{IsMonomialIdeal}}
\logpage{[ 4, 8, 4 ]}\nobreak
\hyperdef{L}{X7C40D53785D67A9E}{}
{\noindent\textcolor{FuncColor}{$\triangleright$\enspace\texttt{IsMonomialIdeal({\mdseries\slshape I})\index{IsMonomialIdeal@\texttt{IsMonomialIdeal}}
\label{IsMonomialIdeal}
}\hfill{\scriptsize (property)}}\\


 Arguments: \mbox{\texttt{\mdseries\slshape I}} \texttt{\symbol{45}}\texttt{\symbol{45}} an IsIdealInPathAlgebra object. \\
 \textbf{\indent Returns:\ }
true whenever \mbox{\texttt{\mdseries\slshape I}} is a \emph{monomial} ideal in a path algebra, i.e. \mbox{\texttt{\mdseries\slshape I}} is generated by a set of monomials (= "zero\texttt{\symbol{45}}relations"). 



 \emph{Technical note:} It uses the observation: \mbox{\texttt{\mdseries\slshape I}} is a monomial ideal iff Groebner basis of \mbox{\texttt{\mdseries\slshape I}} is a set of monomials. It computes Groebner basis for \mbox{\texttt{\mdseries\slshape I}} only in case it has not been computed yet and a usual set of generators
(GeneratorsOfIdeal) is not a set of monomials. }

 

\subsection{\textcolor{Chapter }{IsQuadraticIdeal}}
\logpage{[ 4, 8, 5 ]}\nobreak
\hyperdef{L}{X841058F8850FA9D3}{}
{\noindent\textcolor{FuncColor}{$\triangleright$\enspace\texttt{IsQuadraticIdeal({\mdseries\slshape rels})\index{IsQuadraticIdeal@\texttt{IsQuadraticIdeal}}
\label{IsQuadraticIdeal}
}\hfill{\scriptsize (operation)}}\\


 Arguments: \mbox{\texttt{\mdseries\slshape rels}} \texttt{\symbol{45}}\texttt{\symbol{45}} a list of elements in a path algebra.\\
 \textbf{\indent Returns:\ }
true whenever \mbox{\texttt{\mdseries\slshape rels}} is a list of elements in the linear span of degree two elements of a path
algebra. It returns false whenever \mbox{\texttt{\mdseries\slshape rels}} is a list of elements in a path algebra, but not in the linear span of degree
two of a path algebra. Otherwise it returns fail. 

}

 }

 
\section{\textcolor{Chapter }{Operations on ideals}}\label{OperationsOnIdeals}
\logpage{[ 4, 9, 0 ]}
\hyperdef{L}{X7C45A01B7A587D9E}{}
{
 

\subsection{\textcolor{Chapter }{ProductOfIdeals}}
\logpage{[ 4, 9, 1 ]}\nobreak
\hyperdef{L}{X7C1D2A2481599348}{}
{\noindent\textcolor{FuncColor}{$\triangleright$\enspace\texttt{ProductOfIdeals({\mdseries\slshape I, J})\index{ProductOfIdeals@\texttt{ProductOfIdeals}}
\label{ProductOfIdeals}
}\hfill{\scriptsize (operation)}}\\


 Arguments: \mbox{\texttt{\mdseries\slshape I, J}} \texttt{\symbol{45}}\texttt{\symbol{45}} two ideals in a path algebra \mbox{\texttt{\mdseries\slshape KQ}}.\\
 \textbf{\indent Returns:\ }
the ideal formed by the product of the ideals \mbox{\texttt{\mdseries\slshape I}} and \mbox{\texttt{\mdseries\slshape J}}, whenever the ideal \mbox{\texttt{\mdseries\slshape J}} admits finitely many nontips in \mbox{\texttt{\mdseries\slshape KQ}}. 



 The function checks if the two ideals are ideals in the same path algebra and
that \mbox{\texttt{\mdseries\slshape J}} admits finitely many nontips in \mbox{\texttt{\mdseries\slshape KQ}}. }

 

\subsection{\textcolor{Chapter }{QuadraticPerpOfPathAlgebraIdeal}}
\logpage{[ 4, 9, 2 ]}\nobreak
\hyperdef{L}{X7B0035AF7B030BDF}{}
{\noindent\textcolor{FuncColor}{$\triangleright$\enspace\texttt{QuadraticPerpOfPathAlgebraIdeal({\mdseries\slshape rels})\index{QuadraticPerpOfPathAlgebraIdeal@\texttt{QuadraticPerpOfPathAlgebraIdeal}}
\label{QuadraticPerpOfPathAlgebraIdeal}
}\hfill{\scriptsize (operation)}}\\


 Arguments: \mbox{\texttt{\mdseries\slshape rels}} \texttt{\symbol{45}}\texttt{\symbol{45}} a list of elements in a path algebra.\\
 \textbf{\indent Returns:\ }
fail if \mbox{\texttt{\mdseries\slshape rels}} is not a list of elements in the linear span of degree two elements of a path
algebra \mbox{\texttt{\mdseries\slshape KQ}}. Otherwise it returns a list of length two, where the first element is a set
of generators for the ideal $\mbox{\texttt{\mdseries\slshape rels}}^\perp$ in opposite algebra of \mbox{\texttt{\mdseries\slshape KQ}} and the second element is the opposite algebra of \mbox{\texttt{\mdseries\slshape KQ}}. 

}

 }

 
\section{\textcolor{Chapter }{Attributes of ideals}}\label{AttributesOfIdeals}
\logpage{[ 4, 10, 0 ]}
\hyperdef{L}{X85D4E72B787B1C49}{}
{
 For many of the functions related to quotients, you will need to compute a
Groebner basis of the ideal. This is done with the GBNP package. The following
example shows how to set a Groebner basis for an ideal (note that this must be
done before the quotient is constructed). See the next two chapters for more
on Groebner bases. 
\begin{Verbatim}[commandchars=!@|,fontsize=\small,frame=single,label=Example]
  !gapprompt@gap>| !gapinput@rels := [FQ.a - FQ.b*FQ.c, FQ.d*FQ.d];|
  [ (1)*a+(-1)*b*c, (1)*d^2 ]
  !gapprompt@gap>| !gapinput@gb := GBNPGroebnerBasis(rels, FQ); |
  [ (-1)*a+(1)*b*c, (1)*d^2 ]
  !gapprompt@gap>| !gapinput@I := Ideal(FQ, gb);|
  <two-sided ideal in <Rationals[<quiver with 2 vertices and 4 arrows>]>
      , (2 generators)>
  !gapprompt@gap>| !gapinput@GroebnerBasis(I, gb);|
  <complete two-sided Groebner basis containing 2 elements>
  !gapprompt@gap>| !gapinput@quot := FQ/I;|
  <Rationals[<quiver with 2 vertices and 4 arrows>]/
  <two-sided ideal in <Rationals[<quiver with 2 vertices and 4 arrows>]>
      , (2 generators)>>
\end{Verbatim}
 

\subsection{\textcolor{Chapter }{GroebnerBasisOfIdeal}}
\logpage{[ 4, 10, 1 ]}\nobreak
\hyperdef{L}{X7D2896F27C976231}{}
{\noindent\textcolor{FuncColor}{$\triangleright$\enspace\texttt{GroebnerBasisOfIdeal({\mdseries\slshape I})\index{GroebnerBasisOfIdeal@\texttt{GroebnerBasisOfIdeal}}
\label{GroebnerBasisOfIdeal}
}\hfill{\scriptsize (attribute)}}\\


 Arguments: \mbox{\texttt{\mdseries\slshape I}} \texttt{\symbol{45}}\texttt{\symbol{45}} an ideal in path algebra.\\
 \textbf{\indent Returns:\ }
a Groebner basis of ideal \mbox{\texttt{\mdseries\slshape I}} (if it has been already computed!). 



 This attribute is set only by an operation \texttt{GroebnerBasis} (\ref{GroebnerBasis}). }

 }

 
\section{\textcolor{Chapter }{Categories and Properties of Quotients of Path Algebras}}\logpage{[ 4, 11, 0 ]}
\hyperdef{L}{X7A3CA333873389AD}{}
{
  

\subsection{\textcolor{Chapter }{IsAdmissibleQuotientOfPathAlgebra}}
\logpage{[ 4, 11, 1 ]}\nobreak
\hyperdef{L}{X7ACDD33087F98B88}{}
{\noindent\textcolor{FuncColor}{$\triangleright$\enspace\texttt{IsAdmissibleQuotientOfPathAlgebra({\mdseries\slshape A})\index{IsAdmissibleQuotientOfPathAlgebra@\texttt{IsAdmissibleQuotientOfPathAlgebra}}
\label{IsAdmissibleQuotientOfPathAlgebra}
}\hfill{\scriptsize (filter)}}\\


 Arguments: \mbox{\texttt{\mdseries\slshape A}} \texttt{\symbol{45}}\texttt{\symbol{45}} any object.\\
 \textbf{\indent Returns:\ }
true whenever \mbox{\texttt{\mdseries\slshape A}} is a quotient of a path algebra by an admissible ideal constructed by the
command \texttt{\texttt{\symbol{92}}/} with arguments a path algebra and a list of relations, \texttt{KQ/rels}, where \texttt{rels} is a list of relations. Otherwise it returns an error message. 

}

 

\subsection{\textcolor{Chapter }{IsQuotientOfPathAlgebra}}
\logpage{[ 4, 11, 2 ]}\nobreak
\hyperdef{L}{X790DB9BF831B577D}{}
{\noindent\textcolor{FuncColor}{$\triangleright$\enspace\texttt{IsQuotientOfPathAlgebra({\mdseries\slshape object})\index{IsQuotientOfPathAlgebra@\texttt{IsQuotientOfPathAlgebra}}
\label{IsQuotientOfPathAlgebra}
}\hfill{\scriptsize (property)}}\\


 Argument: \mbox{\texttt{\mdseries\slshape object}} \texttt{\symbol{45}}\texttt{\symbol{45}} any object in \textsf{GAP}.\\
 \textbf{\indent Returns:\ }
 true whenever \mbox{\texttt{\mdseries\slshape object}} is a quotient of a path algebra. 

}

 
\begin{Verbatim}[commandchars=!@|,fontsize=\small,frame=single,label=Example]
  !gapprompt@gap>| !gapinput@quot := FQ/I;|
  <Rationals[<quiver with 2 vertices and 4 arrows>]/
  <two-sided ideal in <Rationals[<quiver with 2 vertices and 4 arrows>]>
      , (2 generators)>>
  !gapprompt@gap>| !gapinput@IsQuotientOfPathAlgebra(quot);|
  true
  !gapprompt@gap>| !gapinput@IsQuotientOfPathAlgebra(FQ);|
  false
\end{Verbatim}
 

\subsection{\textcolor{Chapter }{IsFiniteDimensional}}
\logpage{[ 4, 11, 3 ]}\nobreak
\hyperdef{L}{X802DB9FB824B0167}{}
{\noindent\textcolor{FuncColor}{$\triangleright$\enspace\texttt{IsFiniteDimensional({\mdseries\slshape A})\index{IsFiniteDimensional@\texttt{IsFiniteDimensional}}
\label{IsFiniteDimensional}
}\hfill{\scriptsize (property)}}\\


 Arguments: \mbox{\texttt{\mdseries\slshape A}} \texttt{\symbol{45}}\texttt{\symbol{45}} a path algebra or a quotient of a
path algebra. \\
 \textbf{\indent Returns:\ }
true whenever \mbox{\texttt{\mdseries\slshape A}} is a finite dimensional algebra. 



\emph{Technical note:} For a path algebra it uses a standard \textsf{GAP} method. For a quotient of a path algebra it uses Groebner bases machinery (it
computes Groebner basis for the ideal only in case it has not been computed
yet). }

 

\subsection{\textcolor{Chapter }{IsCanonicalAlgebra}}
\logpage{[ 4, 11, 4 ]}\nobreak
\hyperdef{L}{X8523C98A870CF7B5}{}
{\noindent\textcolor{FuncColor}{$\triangleright$\enspace\texttt{IsCanonicalAlgebra({\mdseries\slshape A})\index{IsCanonicalAlgebra@\texttt{IsCanonicalAlgebra}}
\label{IsCanonicalAlgebra}
}\hfill{\scriptsize (property)}}\\


 Arguments: \mbox{\texttt{\mdseries\slshape A}} \texttt{\symbol{45}}\texttt{\symbol{45}} a path algebra or a quotient of a
path algebra.\\
 \textbf{\indent Returns:\ }
 true if \mbox{\texttt{\mdseries\slshape A}} has been constructed by the operation \texttt{CanonicalAlgebra} (\ref{CanonicalAlgebra}), otherwise "Error, no method found". 

}

 

\subsection{\textcolor{Chapter }{IsDistributiveAlgebra}}
\logpage{[ 4, 11, 5 ]}\nobreak
\hyperdef{L}{X78C9AA2085058DFA}{}
{\noindent\textcolor{FuncColor}{$\triangleright$\enspace\texttt{IsDistributiveAlgebra({\mdseries\slshape A})\index{IsDistributiveAlgebra@\texttt{IsDistributiveAlgebra}}
\label{IsDistributiveAlgebra}
}\hfill{\scriptsize (property)}}\\


 Arguments: \mbox{\texttt{\mdseries\slshape A}} \texttt{\symbol{45}}\texttt{\symbol{45}} a path algebra or a quotient of a
path algebra.\\
 \textbf{\indent Returns:\ }
 true if \mbox{\texttt{\mdseries\slshape A}} is an admissible quotient of a path algebra and distributive. Otherwise it
returns false if \mbox{\texttt{\mdseries\slshape A}} is an admissible quotient of a path algebra and distributive. If \mbox{\texttt{\mdseries\slshape A}} is a quotient of a path algebra, but not an admissible quotient, then it looks
for other methods. There are not further methods implemented in QPA as of now. 

}

 

\subsection{\textcolor{Chapter }{IsFiniteGlobalDimensionAlgebra}}
\logpage{[ 4, 11, 6 ]}\nobreak
\hyperdef{L}{X7EF4868B84CC749E}{}
{\noindent\textcolor{FuncColor}{$\triangleright$\enspace\texttt{IsFiniteGlobalDimensionAlgebra({\mdseries\slshape A})\index{IsFiniteGlobalDimensionAlgebra@\texttt{IsFiniteGlobalDimensionAlgebra}}
\label{IsFiniteGlobalDimensionAlgebra}
}\hfill{\scriptsize (property)}}\\


 Arguments: \mbox{\texttt{\mdseries\slshape A}} \texttt{\symbol{45}} an algebra over a field.\\
 \textbf{\indent Returns:\ }
 true if it is known that the entered algebra \mbox{\texttt{\mdseries\slshape A}} has finite global dimension. 



 There is no method associated to this, so if it is not known that the algebra
has finite global dimension, then an error message saying "no method found!"
is return. }

 

\subsection{\textcolor{Chapter }{IsGentleAlgebra}}
\logpage{[ 4, 11, 7 ]}\nobreak
\hyperdef{L}{X796625487F5F92A7}{}
{\noindent\textcolor{FuncColor}{$\triangleright$\enspace\texttt{IsGentleAlgebra({\mdseries\slshape A})\index{IsGentleAlgebra@\texttt{IsGentleAlgebra}}
\label{IsGentleAlgebra}
}\hfill{\scriptsize (property)}}\\


 Arguments: \mbox{\texttt{\mdseries\slshape A}} \texttt{\symbol{45}}\texttt{\symbol{45}} a path algebra or a quotient of a
path algebra.\\
 \textbf{\indent Returns:\ }
 true if the algebra \mbox{\texttt{\mdseries\slshape A}} is a gentle algebra. Otherwise false. 

}

 

\subsection{\textcolor{Chapter }{IsGorensteinAlgebra}}
\logpage{[ 4, 11, 8 ]}\nobreak
\hyperdef{L}{X7E6BE2187B48691D}{}
{\noindent\textcolor{FuncColor}{$\triangleright$\enspace\texttt{IsGorensteinAlgebra({\mdseries\slshape A})\index{IsGorensteinAlgebra@\texttt{IsGorensteinAlgebra}}
\label{IsGorensteinAlgebra}
}\hfill{\scriptsize (property)}}\\


 Arguments: \mbox{\texttt{\mdseries\slshape A}} \texttt{\symbol{45}}\texttt{\symbol{45}} an algebra.\\
 \textbf{\indent Returns:\ }
 true if it is known that \mbox{\texttt{\mdseries\slshape A}} is a Gorenstein algebra. If unknown it returns an error message saying "no
method found!". 



 There is no method installed for this yet. }

 

\subsection{\textcolor{Chapter }{IsHereditaryAlgebra}}
\logpage{[ 4, 11, 9 ]}\nobreak
\hyperdef{L}{X859FF8F2865D0A3A}{}
{\noindent\textcolor{FuncColor}{$\triangleright$\enspace\texttt{IsHereditaryAlgebra({\mdseries\slshape A})\index{IsHereditaryAlgebra@\texttt{IsHereditaryAlgebra}}
\label{IsHereditaryAlgebra}
}\hfill{\scriptsize (property)}}\\


 Arguments: \mbox{\texttt{\mdseries\slshape A}} \texttt{\symbol{45}}\texttt{\symbol{45}} an admissible quotient of a path
algebra.\\
 \textbf{\indent Returns:\ }
 true if \mbox{\texttt{\mdseries\slshape A}} is a hereditary algebra and false otherwise. 

}

 

\subsection{\textcolor{Chapter }{IsKroneckerAlgebra}}
\logpage{[ 4, 11, 10 ]}\nobreak
\hyperdef{L}{X7809E044817388D1}{}
{\noindent\textcolor{FuncColor}{$\triangleright$\enspace\texttt{IsKroneckerAlgebra({\mdseries\slshape A})\index{IsKroneckerAlgebra@\texttt{IsKroneckerAlgebra}}
\label{IsKroneckerAlgebra}
}\hfill{\scriptsize (property)}}\\


 Arguments: \mbox{\texttt{\mdseries\slshape A}} \texttt{\symbol{45}}\texttt{\symbol{45}} a path algebra or a quotient of a
path algebra.\\
 \textbf{\indent Returns:\ }
 true if \mbox{\texttt{\mdseries\slshape A}} has been constructed by the operation \texttt{KroneckerAlgebra} (\ref{KroneckerAlgebra}), otherwise "Error, no method found". 

}

 

\subsection{\textcolor{Chapter }{IsMonomialAlgebra}}
\logpage{[ 4, 11, 11 ]}\nobreak
\hyperdef{L}{X7952044F8303A688}{}
{\noindent\textcolor{FuncColor}{$\triangleright$\enspace\texttt{IsMonomialAlgebra({\mdseries\slshape A})\index{IsMonomialAlgebra@\texttt{IsMonomialAlgebra}}
\label{IsMonomialAlgebra}
}\hfill{\scriptsize (property)}}\\


 Arguments: \mbox{\texttt{\mdseries\slshape A}} \texttt{\symbol{45}}\texttt{\symbol{45}} a quiver algebra.\\
 \textbf{\indent Returns:\ }
 true when \mbox{\texttt{\mdseries\slshape A}} is given as \texttt{kQ/I} and \texttt{I} is a monomial ideal in \texttt{kQ}, otherwise it returns false. 

}

 

\subsection{\textcolor{Chapter }{IsNakayamaAlgebra}}
\logpage{[ 4, 11, 12 ]}\nobreak
\hyperdef{L}{X7CB110A7873F7942}{}
{\noindent\textcolor{FuncColor}{$\triangleright$\enspace\texttt{IsNakayamaAlgebra({\mdseries\slshape A})\index{IsNakayamaAlgebra@\texttt{IsNakayamaAlgebra}}
\label{IsNakayamaAlgebra}
}\hfill{\scriptsize (property)}}\\


 Arguments: \mbox{\texttt{\mdseries\slshape A}} \texttt{\symbol{45}}\texttt{\symbol{45}} a path algebra or a quotient of a
path algebra.\\
 \textbf{\indent Returns:\ }
 true if \mbox{\texttt{\mdseries\slshape A}} has been constructed by the operation \texttt{NakayamaAlgebra} (\ref{NakayamaAlgebra}), otherwise "Error, no method found". 

}

 

\subsection{\textcolor{Chapter }{IsQuiverAlgebra}}
\logpage{[ 4, 11, 13 ]}\nobreak
\hyperdef{L}{X8798B8BA7A145A2D}{}
{\noindent\textcolor{FuncColor}{$\triangleright$\enspace\texttt{IsQuiverAlgebra({\mdseries\slshape A})\index{IsQuiverAlgebra@\texttt{IsQuiverAlgebra}}
\label{IsQuiverAlgebra}
}\hfill{\scriptsize (filter)}}\\


 Arguments: \mbox{\texttt{\mdseries\slshape A}} \texttt{\symbol{45}}\texttt{\symbol{45}} an algebra.\\
 \textbf{\indent Returns:\ }
 true if \mbox{\texttt{\mdseries\slshape A}} is a path algebra or a quotient of a path algebra algebra, otherwise false. 

}

 

\subsection{\textcolor{Chapter }{IsRadicalSquareZeroAlgebra}}
\logpage{[ 4, 11, 14 ]}\nobreak
\hyperdef{L}{X7A8D13FE8379776E}{}
{\noindent\textcolor{FuncColor}{$\triangleright$\enspace\texttt{IsRadicalSquareZeroAlgebra({\mdseries\slshape A})\index{IsRadicalSquareZeroAlgebra@\texttt{IsRadicalSquareZeroAlgebra}}
\label{IsRadicalSquareZeroAlgebra}
}\hfill{\scriptsize (property)}}\\


 Arguments: \mbox{\texttt{\mdseries\slshape A}} \texttt{\symbol{45}}\texttt{\symbol{45}} an algebra.\\
 \textbf{\indent Returns:\ }
 true if \mbox{\texttt{\mdseries\slshape A}} is a radical square zero algebra, otherwise false. 

}

 

\subsection{\textcolor{Chapter }{IsSchurianAlgebra}}
\logpage{[ 4, 11, 15 ]}\nobreak
\hyperdef{L}{X82F25BFD7D43AB10}{}
{\noindent\textcolor{FuncColor}{$\triangleright$\enspace\texttt{IsSchurianAlgebra({\mdseries\slshape A})\index{IsSchurianAlgebra@\texttt{IsSchurianAlgebra}}
\label{IsSchurianAlgebra}
}\hfill{\scriptsize (property)}}\\


 Arguments: \mbox{\texttt{\mdseries\slshape A}} \texttt{\symbol{45}}\texttt{\symbol{45}} a path algebra or a quotient of a
path algebra.\\
 \textbf{\indent Returns:\ }
true if \mbox{\texttt{\mdseries\slshape A}} is a schurian algebra. By definition it means that: for all $x,y\in Q_0$ we have $\dim A(x,y)\leq 1 $. 



 \emph{Note:} This method fail when a Groebner basis for ideal has not been computed before
creating a quotient! }

 

\subsection{\textcolor{Chapter }{IsSelfinjectiveAlgebra}}
\logpage{[ 4, 11, 16 ]}\nobreak
\hyperdef{L}{X8555FC6B85FE9C6D}{}
{\noindent\textcolor{FuncColor}{$\triangleright$\enspace\texttt{IsSelfinjectiveAlgebra({\mdseries\slshape A})\index{IsSelfinjectiveAlgebra@\texttt{IsSelfinjectiveAlgebra}}
\label{IsSelfinjectiveAlgebra}
}\hfill{\scriptsize (property)}}\\


 Arguments: \mbox{\texttt{\mdseries\slshape A}} \texttt{\symbol{45}}\texttt{\symbol{45}} a path algebra or a quotient of a
path algebra.\\
 \textbf{\indent Returns:\ }
fail if \mbox{\texttt{\mdseries\slshape A}} is not finite dimensional. Otherwise it returns true or false according to
whether \mbox{\texttt{\mdseries\slshape A}} is selfinjective or not. 

}

 

\subsection{\textcolor{Chapter }{IsSemicommutativeAlgebra}}
\logpage{[ 4, 11, 17 ]}\nobreak
\hyperdef{L}{X8558D44A79AA16CD}{}
{\noindent\textcolor{FuncColor}{$\triangleright$\enspace\texttt{IsSemicommutativeAlgebra({\mdseries\slshape A})\index{IsSemicommutativeAlgebra@\texttt{IsSemicommutativeAlgebra}}
\label{IsSemicommutativeAlgebra}
}\hfill{\scriptsize (property)}}\\


 Arguments: \mbox{\texttt{\mdseries\slshape A}} \texttt{\symbol{45}}\texttt{\symbol{45}} a path algebra or a quotient of a
path algebra.\\
 \textbf{\indent Returns:\ }
true if \mbox{\texttt{\mdseries\slshape A}} is a semicommutative algebra. By definition it means that:\\
 1. \mbox{\texttt{\mdseries\slshape A}} is schurian (cf. \texttt{IsSchurianAlgebra} (\ref{IsSchurianAlgebra})).\\
 2. Quiver $Q$ of \mbox{\texttt{\mdseries\slshape A}} is acyclic (cf. \texttt{IsAcyclicQuiver} (\ref{IsAcyclicQuiver})).\\
 3. For all pairs of vertices $(x,y)$ the following condition is satisfied: for every two paths $P,P'$ from $x$ to $y$: $P\in I \Leftrightarrow P'\in I$. 



 \emph{Note:} This method fail when a Groebner basis for ideal has not been computed before
creating a quotient! }

 

\subsection{\textcolor{Chapter }{IsSemisimpleAlgebra}}
\logpage{[ 4, 11, 18 ]}\nobreak
\hyperdef{L}{X85A8EC2287F35DC1}{}
{\noindent\textcolor{FuncColor}{$\triangleright$\enspace\texttt{IsSemisimpleAlgebra({\mdseries\slshape A})\index{IsSemisimpleAlgebra@\texttt{IsSemisimpleAlgebra}}
\label{IsSemisimpleAlgebra}
}\hfill{\scriptsize (property)}}\\


 Arguments: \mbox{\texttt{\mdseries\slshape A}} \texttt{\symbol{45}} an algebra over a field.\\
 \textbf{\indent Returns:\ }
 true if the entered algebra \mbox{\texttt{\mdseries\slshape A}} is semisimple, false otherwise. 



 Checks if the algebra is finite dimensional. If it is an admissible quotients
of a path algebra, it only checks if the underlying quiver has any arrows or
not. Otherwise, it computes the radical of the algebra and checks if it is
zero. }

 

\subsection{\textcolor{Chapter }{IsSpecialBiserialAlgebra}}
\logpage{[ 4, 11, 19 ]}\nobreak
\hyperdef{L}{X7D7AC1D07A9607DF}{}
{\noindent\textcolor{FuncColor}{$\triangleright$\enspace\texttt{IsSpecialBiserialAlgebra({\mdseries\slshape A})\index{IsSpecialBiserialAlgebra@\texttt{IsSpecialBiserialAlgebra}}
\label{IsSpecialBiserialAlgebra}
}\hfill{\scriptsize (property)}}\\


 Arguments: \mbox{\texttt{\mdseries\slshape A}} \texttt{\symbol{45}}\texttt{\symbol{45}} a path algebra or a quotient of a
path algebra. \\
 \textbf{\indent Returns:\ }
true whenever \mbox{\texttt{\mdseries\slshape A}} is a \emph{special biserial algebra}, i.e. \mbox{\texttt{\mdseries\slshape A=KQ/I}}, where \mbox{\texttt{\mdseries\slshape Q}} is \texttt{IsSpecialBiserialQuiver} (\ref{IsSpecialBiserialQuiver}), \mbox{\texttt{\mdseries\slshape I}} is an admissible ideal (\texttt{IsAdmissibleIdeal} (\ref{IsAdmissibleIdeal})) and \mbox{\texttt{\mdseries\slshape I}} satisfies the "special biserial" conditions, i.e.: 
\begin{itemize}
\item for any arrow $a$ there exists at most one arrow $b$ such that $ab$ does not belong to \mbox{\texttt{\mdseries\slshape I}}
\item there exists at most one arrow $c$ such that $ca$ does not belong to $I$.
\end{itemize}
 



 \emph{Note:} e.g. a path algebra of one loop IS NOT special biserial, but one loop IS
special biserial quiver (see \texttt{IsSpecialBiserialQuiver} (\ref{IsSpecialBiserialQuiver}) for examples). }

 

\subsection{\textcolor{Chapter }{IsStringAlgebra}}
\logpage{[ 4, 11, 20 ]}\nobreak
\hyperdef{L}{X86F0E4AF7C9916CB}{}
{\noindent\textcolor{FuncColor}{$\triangleright$\enspace\texttt{IsStringAlgebra({\mdseries\slshape A})\index{IsStringAlgebra@\texttt{IsStringAlgebra}}
\label{IsStringAlgebra}
}\hfill{\scriptsize (property)}}\\


 Arguments: \mbox{\texttt{\mdseries\slshape A}} \texttt{\symbol{45}}\texttt{\symbol{45}} a path algebra or a quotient of a
path algebra. \\
 \textbf{\indent Returns:\ }
true whenever \mbox{\texttt{\mdseries\slshape A}} is a \emph{string} (special biserial) algebra, i.e. \mbox{\texttt{\mdseries\slshape A=KQ/I}} is a special biserial algebra (\texttt{IsSpecialBiserialAlgebra} (\ref{IsSpecialBiserialAlgebra}) and \mbox{\texttt{\mdseries\slshape I}} is generated by monomials (= "zero\texttt{\symbol{45}}relations") (cf. \texttt{IsMonomialIdeal} (\ref{IsMonomialIdeal})). See \texttt{IsSpecialBiserialQuiver} (\ref{IsSpecialBiserialQuiver}) for examples. 

}

 

\subsection{\textcolor{Chapter }{IsSymmetricAlgebra}}
\logpage{[ 4, 11, 21 ]}\nobreak
\hyperdef{L}{X79DA912C82D01EE8}{}
{\noindent\textcolor{FuncColor}{$\triangleright$\enspace\texttt{IsSymmetricAlgebra({\mdseries\slshape A})\index{IsSymmetricAlgebra@\texttt{IsSymmetricAlgebra}}
\label{IsSymmetricAlgebra}
}\hfill{\scriptsize (property)}}\\


 Arguments: \mbox{\texttt{\mdseries\slshape A}} \texttt{\symbol{45}}\texttt{\symbol{45}} a path algebra or a quotient of a
path algebra.\\
 \textbf{\indent Returns:\ }
fail if \mbox{\texttt{\mdseries\slshape A}} is not finite dimensional or does not have a Groebner basis. Otherwise it
returns true or false according to whether \mbox{\texttt{\mdseries\slshape A}} is symmetric or not. 

}

 

\subsection{\textcolor{Chapter }{IsTriangularReduced}}
\logpage{[ 4, 11, 22 ]}\nobreak
\hyperdef{L}{X87EC06D18021AD76}{}
{\noindent\textcolor{FuncColor}{$\triangleright$\enspace\texttt{IsTriangularReduced({\mdseries\slshape A})\index{IsTriangularReduced@\texttt{IsTriangularReduced}}
\label{IsTriangularReduced}
}\hfill{\scriptsize (property)}}\\


 Arguments: \mbox{\texttt{\mdseries\slshape A}} \texttt{\symbol{45}} a finite dimensional QuiverAlgebra.\\
 \textbf{\indent Returns:\ }
\texttt{false} if the algebra \mbox{\texttt{\mdseries\slshape A}} is triangular reducable, that is, there is a sum over vertices $e$ such that $eA(1 - e) = (0)$ for $e\neq 0,1$. Otherwise, it returns \texttt{true}. 



 The function checks if the algebra \mbox{\texttt{\mdseries\slshape A}} is finite dimensional and gives an error message otherwise. }

 

\subsection{\textcolor{Chapter }{IsWeaklySymmetricAlgebra}}
\logpage{[ 4, 11, 23 ]}\nobreak
\hyperdef{L}{X7EF281F980319375}{}
{\noindent\textcolor{FuncColor}{$\triangleright$\enspace\texttt{IsWeaklySymmetricAlgebra({\mdseries\slshape A})\index{IsWeaklySymmetricAlgebra@\texttt{IsWeaklySymmetricAlgebra}}
\label{IsWeaklySymmetricAlgebra}
}\hfill{\scriptsize (property)}}\\


 Arguments: \mbox{\texttt{\mdseries\slshape A}} \texttt{\symbol{45}}\texttt{\symbol{45}} a path algebra or a quotient of a
path algebra.\\
 \textbf{\indent Returns:\ }
fail if \mbox{\texttt{\mdseries\slshape A}} is not finite dimensional or does not have a Groebner basis. Otherwise it
returns true or false according to whether \mbox{\texttt{\mdseries\slshape A}} is weakly symmetric or not. 

}

 

\subsection{\textcolor{Chapter }{BongartzTest}}
\logpage{[ 4, 11, 24 ]}\nobreak
\hyperdef{L}{X81C569797E900AE9}{}
{\noindent\textcolor{FuncColor}{$\triangleright$\enspace\texttt{BongartzTest({\mdseries\slshape A, bound})\index{BongartzTest@\texttt{BongartzTest}}
\label{BongartzTest}
}\hfill{\scriptsize (property)}}\\


 Arguments: \mbox{\texttt{\mdseries\slshape A, bound}} \texttt{\symbol{45}}\texttt{\symbol{45}} a path algebra or a quotient of a
path algebra and an integer.\\
 \textbf{\indent Returns:\ }
false if there exists a $\tau$\texttt{\symbol{45}}translate of a simple, injective or projective \mbox{\texttt{\mdseries\slshape A}}\texttt{\symbol{45}}module up to the power \mbox{\texttt{\mdseries\slshape bound}} has dimension greater or equal to $\max\{2\dim A, 30\}$. Then \mbox{\texttt{\mdseries\slshape A}} is of infinite representation type. Returns fail otherwise. 

}

 

\subsection{\textcolor{Chapter }{IsFiniteTypeAlgebra}}
\logpage{[ 4, 11, 25 ]}\nobreak
\hyperdef{L}{X8140D2557A23CDAC}{}
{\noindent\textcolor{FuncColor}{$\triangleright$\enspace\texttt{IsFiniteTypeAlgebra({\mdseries\slshape A})\index{IsFiniteTypeAlgebra@\texttt{IsFiniteTypeAlgebra}}
\label{IsFiniteTypeAlgebra}
}\hfill{\scriptsize (property)}}\\


 Arguments: \mbox{\texttt{\mdseries\slshape A}} \texttt{\symbol{45}}\texttt{\symbol{45}} a path algebra or a quotient of a
path algebra.\\
 \textbf{\indent Returns:\ }
Returns true if \mbox{\texttt{\mdseries\slshape A}} is of finite representation type. Returns false if \mbox{\texttt{\mdseries\slshape A}} is of infinite representation type. Returns fail if we can not determine the
representation type (i.e. it impossible from theoretical/algorithmic point of
view or a suitable criterion has not been implemented yet; the implementation
is in progress). Note: in case \mbox{\texttt{\mdseries\slshape A}} is a path algebra the function is completely implemented. 

}

 
\begin{Verbatim}[commandchars=@|B,fontsize=\small,frame=single,label=Example]
  @gapprompt|gap>B @gapinput|Q := Quiver(5, [ [1,2,"a"], [2,4,"b"], [3,2,"c"], [2,5,"d"] ]);B
  <quiver with 5 vertices and 4 arrows>
  @gapprompt|gap>B @gapinput|A := PathAlgebra(Rationals, Q);B
  <Rationals[<quiver with 5 vertices and 4 arrows>]>
  @gapprompt|gap>B @gapinput|IsFiniteTypeAlgebra(A);B
  Infinite type!
  Quiver is not a (union of) Dynkin quiver(s).
  false
  @gapprompt|gap>B @gapinput|quo := A/[A.a*A.b, A.c*A.d];;B
  @gapprompt|gap>B @gapinput|IsFiniteTypeAlgebra(quo);B
  Finite type!
  Special biserial algebra with no unoriented cycles in Q.
  true 
\end{Verbatim}
 }

 
\section{\textcolor{Chapter }{ Operations on String Algebras}}\logpage{[ 4, 12, 0 ]}
\hyperdef{L}{X861E6670814290D0}{}
{
 A quiver is a quadruple $Q = (Q_0, Q_1, s, t)$ where $Q_0$ is the set of vertices, $Q_1$ is the set of arrows and $s, t : Q_1 \rightarrow Q_0$ are source and target functions where for each $ \alpha \in Q_1$, $s(\alpha)$ and $t(\alpha)$ denotes starting and ending vertex of $\alpha$ respectively. For the functions defined in this section the quiver has to be
constructed using first construction of \ref{Constructing Quivers}. The arrows in $Q_1$ should be small Roman letters defined in alphabetical order as in the example
below. These operations can handle string algebras with number of arrows less
than or equal to 26. Say that a finite sequence $\alpha_n\alpha_{n-1}...\alpha_1$ is a path in $Q$ if $s(\alpha_{i+1}) = t(\alpha_i)$ for all $1 \leq i \leq n-1$ The source and target of the path are $s(\alpha_1)$ and $t(\alpha_n)$ respectively. We denote by $\rho$ a set of monomial relations on $Q$. Every relation in $\rho$ is a finite sequence $\alpha_1 \alpha_{2}...\alpha_n$ of elements of $Q_1$ with $n \geq 2$ such that $\alpha_n \alpha_{n-1}...\alpha_1$ is a path. $\newline$ For a quiver $Q$, we define two functions $\sigma, \epsilon : Q_1 \rightarrow \{-1, 1\}$ such that 
\begin{itemize}
\item  if $b_1 \neq b_2$ such that $s(b_1) = s(b_2)$, then $\sigma(b_1) = -\sigma(b_2)$;
\item  if $b_1 \neq b_2$ such that $t(b_1) = t(b_2)$, then $\epsilon(b_1) = -\epsilon(b_2)$;
\item  if $s(b_1) = t(b_2)$ and $b_1b_2 \notin \rho$, then $\sigma(b_1) = -\epsilon(b_2)$.
\end{itemize}
 Let us denote the inverse of the set of arrows in $Q_1$ by the corresponding capital letters. For example if $b$ is an arrow in $Q_1$ the its inverse is denoted by $B$. The set of the inverse arrows is denoted by $Q_1^{-1}$. We extend the definition of $\sigma$ and $\epsilon$ to $Q_1 \cup Q_1^{-1}$ so that $\sigma(B) = \epsilon(b)$ and $\epsilon(B) = \sigma(b)$ for all $b \in Q_1$. Note that a \emph{string} of positive length is a walk along the arrows of $Q$ and their inverses such that it does not contain paths obtained by reversing
relations in $\rho$ and their inverses as substrings. We also extend the definition of $\sigma$ and $\epsilon$ to all strings; if $\alpha = \alpha_n ...\alpha_2 \alpha_1$, then $\sigma(\alpha) = \sigma(\alpha_1)$ and $\epsilon(\alpha) = \epsilon(\alpha_n)$. $\newline$ For each vertex $v\in Q_0$, we have a zero length lazy string starting and ending at vertex $v$, denoted by $1_v$. The inverse of this string is $1_v$ itself. The main purpose of defining $\sigma$ and $\epsilon$ is to distinguish $1_v$ with its inverse. We denote this zero length string as $(v,1)$ and $(v,-1)$ and set $\sigma((v,i)) := -i$ and $\epsilon((v,i)) := i$. The concatenation $(v,i)u$ of strings is defined if and only if $\epsilon(u) = i$ and the concatenation $u(v,i)$ is defined if and only if $\sigma(u) = -i$. Moreover if $\alpha = \alpha_n \alpha_{n-1} ... \alpha_2 \alpha_1$ and $\beta = \beta_m \beta_{m-1} ... \beta_2 \beta_1$ are two strings such that $s(\alpha) = t(\beta)$, their concatenation is the string $\alpha\beta = \alpha_n ... \alpha_1 \beta_m ... \beta_1$. Note that, for a given quiver $Q$, the functions $\sigma, \epsilon$ are not unique. A choice of $\sigma$ and $\epsilon$ functions is computed within our codes and utilized in other functions. 

\subsection{\textcolor{Chapter }{IsValidString}}
\logpage{[ 4, 12, 1 ]}\nobreak
\hyperdef{L}{X85F2FFFD78355788}{}
{\noindent\textcolor{FuncColor}{$\triangleright$\enspace\texttt{IsValidString({\mdseries\slshape Q, rho, input{\textunderscore}str})\index{IsValidString@\texttt{IsValidString}}
\label{IsValidString}
}\hfill{\scriptsize (operation)}}\\


 Arguments: \mbox{\texttt{\mdseries\slshape Q}} \texttt{\symbol{45}}\texttt{\symbol{45}} a quiver, \mbox{\texttt{\mdseries\slshape rho}} \texttt{\symbol{45}}\texttt{\symbol{45}} a list of monomial relations, \mbox{\texttt{\mdseries\slshape input{\textunderscore}str}} \texttt{\symbol{45}}\texttt{\symbol{45}} an input collection of characters.\\
 \textbf{\indent Returns:\ }
 "Not a String Algebra" if the algebra presented by \mbox{\texttt{\mdseries\slshape (Q,rho)}} is not a string algebra, "Valid Positive Length String" if \mbox{\texttt{\mdseries\slshape input{\textunderscore}str}} is a string of positive length in the corresponding string algebra, "Valid
Zero Length String", if \mbox{\texttt{\mdseries\slshape input{\textunderscore}str}} is a zero length string in the corresponding string algebra, "Invalid String"
if \mbox{\texttt{\mdseries\slshape input{\textunderscore}str}} is not a string in the corresponding string algebra 

}

 

\subsection{\textcolor{Chapter }{StringsLessThan}}
\logpage{[ 4, 12, 2 ]}\nobreak
\hyperdef{L}{X85B82A4086AA53D6}{}
{\noindent\textcolor{FuncColor}{$\triangleright$\enspace\texttt{StringsLessThan({\mdseries\slshape Q, rho, level})\index{StringsLessThan@\texttt{StringsLessThan}}
\label{StringsLessThan}
}\hfill{\scriptsize (operation)}}\\


 Arguments: \mbox{\texttt{\mdseries\slshape Q}} \texttt{\symbol{45}}\texttt{\symbol{45}} a quiver, \mbox{\texttt{\mdseries\slshape rho}} \texttt{\symbol{45}}\texttt{\symbol{45}} a list of monomial relations, \mbox{\texttt{\mdseries\slshape level}} \texttt{\symbol{45}}\texttt{\symbol{45}} an integer.\\
 \textbf{\indent Returns:\ }
 "Not a String Algebra" if the algebra presented by \mbox{\texttt{\mdseries\slshape (Q,rho)}} is not a string algebra, an empty list if \mbox{\texttt{\mdseries\slshape level}} is a negative integer, a list of all strings of length less than or equal to \mbox{\texttt{\mdseries\slshape level}} if \mbox{\texttt{\mdseries\slshape level}} is a non\texttt{\symbol{45}}negative integer. 

}

 A string $u$ is called $\textbf{cyclic}$ if $s(u) = t(u)$. A cyclic string is called $\textbf{primitive}$ if it can not be written as an n\texttt{\symbol{45}}fold concatenation $\mathfrak{v}^n$ for any string $\mathfrak{v}$ and any $n \in \mathbb{N}$. If $\mathfrak{b}$ is a primitive cyclic string such that $\mathfrak{b}^m$ exists, for all $m \in \mathbb{N}$, then $\mathfrak{b} $ is called a $\textbf{band}$. 

\subsection{\textcolor{Chapter }{IsABand}}
\logpage{[ 4, 12, 3 ]}\nobreak
\hyperdef{L}{X7F5B0A1A7AAF2C18}{}
{\noindent\textcolor{FuncColor}{$\triangleright$\enspace\texttt{IsABand({\mdseries\slshape Q, rho, input{\textunderscore}str})\index{IsABand@\texttt{IsABand}}
\label{IsABand}
}\hfill{\scriptsize (operation)}}\\


 Arguments: \mbox{\texttt{\mdseries\slshape Q}} \texttt{\symbol{45}}\texttt{\symbol{45}} a quiver, \mbox{\texttt{\mdseries\slshape rho}} \texttt{\symbol{45}}\texttt{\symbol{45}} a list of monomial relations, \mbox{\texttt{\mdseries\slshape input{\textunderscore}str}} \texttt{\symbol{45}}\texttt{\symbol{45}} an input collection of characters.\\
 \textbf{\indent Returns:\ }
 "Not a String Algebra" if the algebra presented by \mbox{\texttt{\mdseries\slshape (Q,rho)}} is not a string algebra, "Invalid String" if \mbox{\texttt{\mdseries\slshape input{\textunderscore}str}} is not a string in the corresponding string algebra, $0$ if \mbox{\texttt{\mdseries\slshape input{\textunderscore}str}} is a valid string but not a band, $1$ if \mbox{\texttt{\mdseries\slshape input{\textunderscore}str}} is a band in the corresponding string algebra. 

}

 

\subsection{\textcolor{Chapter }{BandsLessThan}}
\logpage{[ 4, 12, 4 ]}\nobreak
\hyperdef{L}{X85FFA183800621EA}{}
{\noindent\textcolor{FuncColor}{$\triangleright$\enspace\texttt{BandsLessThan({\mdseries\slshape Q, rho, level})\index{BandsLessThan@\texttt{BandsLessThan}}
\label{BandsLessThan}
}\hfill{\scriptsize (operation)}}\\


 Arguments: \mbox{\texttt{\mdseries\slshape Q}} \texttt{\symbol{45}}\texttt{\symbol{45}} a quiver, \mbox{\texttt{\mdseries\slshape rho}} \texttt{\symbol{45}}\texttt{\symbol{45}} a list of monomial relations, \mbox{\texttt{\mdseries\slshape level}} \texttt{\symbol{45}}\texttt{\symbol{45}} an integer.\\
 \textbf{\indent Returns:\ }
 "Not a String Algebra" if the algebra presented by \mbox{\texttt{\mdseries\slshape (Q,rho)}} is not a string algebra, an empty list if \mbox{\texttt{\mdseries\slshape level}} is a negative integer, a list of all bands of length less than or equal to \mbox{\texttt{\mdseries\slshape level}} if \mbox{\texttt{\mdseries\slshape level}} is a non\texttt{\symbol{45}}negative integer. 

}

 Two bands $\mathfrak{b}_1$ and $\mathfrak{b}_2$ of a string algebra are said to be equivalent if they are cyclic permutations
of each other. For any band $\mathfrak{b}$, output of $\texttt{BandsLessThan}$ function contains all bands equivalent to $\mathfrak{b}$. A representative of each equivalence class is chosen such that the first
syllable from right is an inverse syllable and the last syllable from right is
a direct syllable. This new list is returned as the output of the following
function. 

\subsection{\textcolor{Chapter }{BandRepresentativesLessThan}}
\logpage{[ 4, 12, 5 ]}\nobreak
\hyperdef{L}{X870DB8577F0ABF0E}{}
{\noindent\textcolor{FuncColor}{$\triangleright$\enspace\texttt{BandRepresentativesLessThan({\mdseries\slshape Q, rho, level})\index{BandRepresentativesLessThan@\texttt{BandRepresentativesLessThan}}
\label{BandRepresentativesLessThan}
}\hfill{\scriptsize (operation)}}\\


 Arguments: \mbox{\texttt{\mdseries\slshape Q}} \texttt{\symbol{45}}\texttt{\symbol{45}} a quiver, \mbox{\texttt{\mdseries\slshape rho}} \texttt{\symbol{45}}\texttt{\symbol{45}} a list of monomial relations, \mbox{\texttt{\mdseries\slshape level}} \texttt{\symbol{45}}\texttt{\symbol{45}} an integer.\\
 \textbf{\indent Returns:\ }
 "Not a String Algebra" if the algebra presented by \mbox{\texttt{\mdseries\slshape (Q,rho)}} is not a string algebra, an empty list if \mbox{\texttt{\mdseries\slshape level}} is negative integer, a list of all band representatives of length less than or
equal to \mbox{\texttt{\mdseries\slshape level}} if \mbox{\texttt{\mdseries\slshape level}} is a non\texttt{\symbol{45}}negative integer. 

}

 A string algebra with finite number of bands is called a domestic string
algebra. 

\subsection{\textcolor{Chapter }{IsDomesticStringAlgebra}}
\logpage{[ 4, 12, 6 ]}\nobreak
\hyperdef{L}{X7C42DD687CE572DF}{}
{\noindent\textcolor{FuncColor}{$\triangleright$\enspace\texttt{IsDomesticStringAlgebra({\mdseries\slshape Q, rho})\index{IsDomesticStringAlgebra@\texttt{IsDomesticStringAlgebra}}
\label{IsDomesticStringAlgebra}
}\hfill{\scriptsize (operation)}}\\


 Arguments: \mbox{\texttt{\mdseries\slshape Q}} \texttt{\symbol{45}}\texttt{\symbol{45}} a quiver, \mbox{\texttt{\mdseries\slshape rho}} \texttt{\symbol{45}}\texttt{\symbol{45}} a list of monomial relations. \\
 \textbf{\indent Returns:\ }
 "Not a String Algebra" if the algebra presented by \mbox{\texttt{\mdseries\slshape (Q,rho)}} is not a string algebra, $0$ if the corresponding string algebra is non\texttt{\symbol{45}}domestic, $1$ if the corresponding string algebra is domestic. 

}

 A string is $\textbf{band-free}$ if it does not contain any cyclic permutation of a band as a substring. Let $Ba(\Lambda)$ denote the collection of all bands up to cyclic permutation and $Q_0^{Ba}$ denote a fixed set of representatives of bands in $Ba(\Lambda)$. For $\mathfrak{b}_1, \mathfrak{b}_2 \in Q_0^{Ba}$, we say that a finite string $u$ is a weak bridge $\mathfrak{b}_1 \xrightarrow{u} \mathfrak{b}_2$ if it is band\texttt{\symbol{45}}free and if $\mathfrak{b}_2 u \mathfrak{b}_1$ is a string.$\newline$ We say that a weak bridge $\mathfrak{b}_1 \xrightarrow{u} \mathfrak{b}_2$ is a bridge if there is $\textbf{no}$ $\mathfrak{b}\in Q_0^{Ba}$ and weak bridges $\mathfrak{b}_1 \xrightarrow{u_1} \mathfrak{b}$ and $\mathfrak{b} \xrightarrow{u_2} \mathfrak{b}_2$ such that one of the following holds : 
\begin{itemize}
\item  $\mathfrak{u} = \mathfrak{u_2}\mathfrak{u_1}$, where $\mathfrak{u_1}$ and $\mathfrak{u_2}$ are positive length strings, 
\item  $\mathfrak{u} = \mathfrak{u_2'}\mathfrak{u_1'}$, for some positive length strings $\mathfrak{u_1'}, \mathfrak{u_2'}$ such that $\mathfrak{u_2} = \mathfrak{u_2'} \mathfrak{u_2''}$, $\mathfrak{u_1} = \mathfrak{u_1''} \mathfrak{u_1'}$ and $\mathfrak{b} = \mathfrak{u_1''} \mathfrak{u_2''}$. 
\end{itemize}
 

\subsection{\textcolor{Chapter }{BridgeQuiver}}
\logpage{[ 4, 12, 7 ]}\nobreak
\hyperdef{L}{X85D0616C82375B5C}{}
{\noindent\textcolor{FuncColor}{$\triangleright$\enspace\texttt{BridgeQuiver({\mdseries\slshape Q, rho})\index{BridgeQuiver@\texttt{BridgeQuiver}}
\label{BridgeQuiver}
}\hfill{\scriptsize (operation)}}\\


 Arguments: \mbox{\texttt{\mdseries\slshape Q}} a quiver, \mbox{\texttt{\mdseries\slshape rho}} a list of monomial relations.\\
 \textbf{\indent Returns:\ }
 "Not a String Algebra" if the algebra presented by \mbox{\texttt{\mdseries\slshape (Q,rho)}} is not a string algebra, $0$ if the corresponding string algebra is non\texttt{\symbol{45}}domestic, the
bridge quiver of the corresponding string algebra if it is domestic. 

}

 

\subsection{\textcolor{Chapter }{LocalARQuiver}}
\logpage{[ 4, 12, 8 ]}\nobreak
\hyperdef{L}{X82DBBB737885C73B}{}
{\noindent\textcolor{FuncColor}{$\triangleright$\enspace\texttt{LocalARQuiver({\mdseries\slshape Q, rho, input{\textunderscore}str})\index{LocalARQuiver@\texttt{LocalARQuiver}}
\label{LocalARQuiver}
}\hfill{\scriptsize (operation)}}\\


 Arguments: \mbox{\texttt{\mdseries\slshape Q}} \texttt{\symbol{45}}\texttt{\symbol{45}} a quiver, \mbox{\texttt{\mdseries\slshape rho}} \texttt{\symbol{45}}\texttt{\symbol{45}} a list of monomial relations, \mbox{\texttt{\mdseries\slshape input{\textunderscore}str}} \texttt{\symbol{45}}\texttt{\symbol{45}} an input collection of characters.\\
 \textbf{\indent Returns:\ }
 "Not a String Algebra" if the algebra presented by \mbox{\texttt{\mdseries\slshape (Q,rho)}} is not a string algebra, the list of (equivalence classes of) irreducible
maps, i.e., the arrows of the Auslander\texttt{\symbol{45}}Reiten quiver,
whose source or target is the string module corresponding to \mbox{\texttt{\mdseries\slshape input{\textunderscore}str}}. 

}

 
\begin{Verbatim}[commandchars=!@|,fontsize=\small,frame=single,label=Example]
  !gapprompt@gap>| !gapinput@Q := Quiver(4, [[1,2,"a"], [1,2,"b"], [2,3,"c"], [3,4,"d"], [3,4,"e"]]);|
  <quiver with 4 vertices and 5 arrows>
  !gapprompt@gap>| !gapinput@rho := ["bc", "cd"];|
  [ "bc", "cd" ]
  !gapprompt@gap>| !gapinput@IsValidString(Q,rho,"eca");|
  "Valid Positive Length String"
  !gapprompt@gap>| !gapinput@StringsLessThan(Q,rho,2);|
  [ "(1,1)", "b", "Ab", "(1,-1)", "a", "ca", "Ba", "(2,1)", "c", "B", "ec", "aB",
   "(2,-1)", "A", "bA", "(3,1)", "e", "De", "(3,-1)", "d", "C", "Ed", "AC",
   "(4,1)", "E", "dE", "CE", "(4,-1)", "D", "eD" ]
  !gapprompt@gap>| !gapinput@IsABand(Q,rho,"eca");|
  0
  !gapprompt@gap>| !gapinput@IsABand(Q,rho,"Ab");|
  1
  !gapprompt@gap>| !gapinput@BandsLessThan(Q,rho,3);|
  [ "Ab", "Ba", "aB", "bA", "De", "Ed", "dE", "eD" ]
  !gapprompt@gap>| !gapinput@BandRepresentativesLessThan(Q,rho,3);|
  [ "aB", "bA", "dE", "eD" ]
  !gapprompt@gap>| !gapinput@IsDomesticStringAlgebra(Q,rho);|
  true
  !gapprompt@gap>| !gapinput@Q1 := BridgeQuiver(Q,rho);|
  <quiver with 4 vertices and 2 arrows>
  !gapprompt@gap>| !gapinput@Display(Q1);|
  Quiver( ["v1","v2","v3","v4"], [["v3","v2","CE"],["v1","v4","ec"]] )
  !gapprompt@gap>| !gapinput@LocalARQuiver(Q,rho,"eca");|
  [ "The canonical embedding StringModule(eca) to StringModule(eDeca)",
    "The canonical projection StringModule(ecaBa) to StringModule(eca)" ] 
\end{Verbatim}
 }

 
\section{\textcolor{Chapter }{Attributes and Operations (for Quotients) of Path Algebras}}\logpage{[ 4, 13, 0 ]}
\hyperdef{L}{X86647D317A961513}{}
{
          

\subsection{\textcolor{Chapter }{CartanMatrix}}
\logpage{[ 4, 13, 1 ]}\nobreak
\hyperdef{L}{X84E3FEF587CB66C3}{}
{\noindent\textcolor{FuncColor}{$\triangleright$\enspace\texttt{CartanMatrix({\mdseries\slshape A})\index{CartanMatrix@\texttt{CartanMatrix}}
\label{CartanMatrix}
}\hfill{\scriptsize (operation)}}\\


 Arguments: \mbox{\texttt{\mdseries\slshape A}} \texttt{\symbol{45}}\texttt{\symbol{45}} a path algebra or a quotient of a
path algebra.\\
 \textbf{\indent Returns:\ }
the Cartan matrix of the algebra \mbox{\texttt{\mdseries\slshape A}}, after having checked that \mbox{\texttt{\mdseries\slshape A}} is a finite dimensional quotient of a path algebra. 

}

 

\subsection{\textcolor{Chapter }{Centre/Center}}
\logpage{[ 4, 13, 2 ]}\nobreak
\hyperdef{L}{X7F8084A67A3BE874}{}
{\noindent\textcolor{FuncColor}{$\triangleright$\enspace\texttt{Centre/Center({\mdseries\slshape A})\index{Centre/Center@\texttt{Centre/Center}}
\label{Centre/Center}
}\hfill{\scriptsize (operation)}}\\


 Arguments: \mbox{\texttt{\mdseries\slshape A}} \texttt{\symbol{45}}\texttt{\symbol{45}} a path algebra or a quotient of a
path algebra.\\
 \textbf{\indent Returns:\ }
 the centre of the algebra \mbox{\texttt{\mdseries\slshape A}} as a subalgebra. 



 The function computes the center of \mbox{\texttt{\mdseries\slshape A}} if \mbox{\texttt{\mdseries\slshape A}} is a finite dimensional quotient of a path algebra or \mbox{\texttt{\mdseries\slshape A}} is a path algebra with on restriction on the underlying quiver. }

 

\subsection{\textcolor{Chapter }{ComplexityOfAlgebra}}
\logpage{[ 4, 13, 3 ]}\nobreak
\hyperdef{L}{X835A161E8524797A}{}
{\noindent\textcolor{FuncColor}{$\triangleright$\enspace\texttt{ComplexityOfAlgebra({\mdseries\slshape A, n})\index{ComplexityOfAlgebra@\texttt{ComplexityOfAlgebra}}
\label{ComplexityOfAlgebra}
}\hfill{\scriptsize (operation)}}\\


 Arguments: \mbox{\texttt{\mdseries\slshape A}} \texttt{\symbol{45}}\texttt{\symbol{45}} a path algebra or a quotient of a
path algebra, \mbox{\texttt{\mdseries\slshape n}} \texttt{\symbol{45}}\texttt{\symbol{45}} a positive integer.\\
 \textbf{\indent Returns:\ }
an estimate of the complexity of the algebra \mbox{\texttt{\mdseries\slshape A}}. 



 The function checks if the algebra \mbox{\texttt{\mdseries\slshape A}} is known to have finite global dimension. If so, it returns complexity zero.
Otherwise it tries to estimate the complexity in the following way. Recall
that if a function $f(x)$ is a polynomial in $x$, the degree of $f(x)$ is given by $\lim_{n\to\infty} \frac{\log |f(n)|}{\log n}$. So then this function computes an estimate of the maximal complexity of the
simple modules over \mbox{\texttt{\mdseries\slshape A}} by approximating the complexity of each simple module $S$ by considering the limit $\lim_{m\to \infty} \log \frac{\dim(P(S)(m))}{\log m}$ where $P(S)(m)$ is the $m$\texttt{\symbol{45}}th projective in a minimal projective resolution of \mbox{\texttt{\mdseries\slshape S}} at stage $m$. This limit is estimated by $\frac{\log \dim(P(S)(n))}{\log n}$. }

 

\subsection{\textcolor{Chapter }{CoxeterMatrix}}
\logpage{[ 4, 13, 4 ]}\nobreak
\hyperdef{L}{X815CB1D47CB174ED}{}
{\noindent\textcolor{FuncColor}{$\triangleright$\enspace\texttt{CoxeterMatrix({\mdseries\slshape A})\index{CoxeterMatrix@\texttt{CoxeterMatrix}}
\label{CoxeterMatrix}
}\hfill{\scriptsize (attribute)}}\\


 Arguments: \mbox{\texttt{\mdseries\slshape A}} \texttt{\symbol{45}}\texttt{\symbol{45}} a path algebra or a quotient of a
path algebra.\\
 \textbf{\indent Returns:\ }
 the Coxeter matrix of the algebra \mbox{\texttt{\mdseries\slshape A}}, after having checked that \mbox{\texttt{\mdseries\slshape A}} is a finite dimensional quotient of a path algebra. 

}

 

\subsection{\textcolor{Chapter }{CoxeterPolynomial}}
\logpage{[ 4, 13, 5 ]}\nobreak
\hyperdef{L}{X7F6F526C86052150}{}
{\noindent\textcolor{FuncColor}{$\triangleright$\enspace\texttt{CoxeterPolynomial({\mdseries\slshape A})\index{CoxeterPolynomial@\texttt{CoxeterPolynomial}}
\label{CoxeterPolynomial}
}\hfill{\scriptsize (attribute)}}\\


 Arguments: \mbox{\texttt{\mdseries\slshape A}} \texttt{\symbol{45}}\texttt{\symbol{45}} a path algebra or a quotient of a
path algebra.\\
 \textbf{\indent Returns:\ }
 the Coxeter polynomial of the algebra \mbox{\texttt{\mdseries\slshape A}}, after having checked that \mbox{\texttt{\mdseries\slshape A}} is a finite dimensional quotient of a path algebra. 

}

 

\subsection{\textcolor{Chapter }{Dimension}}
\logpage{[ 4, 13, 6 ]}\nobreak
\hyperdef{L}{X7E6926C6850E7C4E}{}
{\noindent\textcolor{FuncColor}{$\triangleright$\enspace\texttt{Dimension({\mdseries\slshape A})\index{Dimension@\texttt{Dimension}}
\label{Dimension}
}\hfill{\scriptsize (attribute)}}\\


 Arguments: \mbox{\texttt{\mdseries\slshape A}} \texttt{\symbol{45}}\texttt{\symbol{45}} a path algebra or a quotient of a
path algebra.\\
 \textbf{\indent Returns:\ }
 the dimension of the algebra \mbox{\texttt{\mdseries\slshape A}} or \emph{infinity} in case \mbox{\texttt{\mdseries\slshape A}} is an infinite dimensional algebra. 



 For a quotient of a path algebra it uses Groebner bases machinery (it computes
Groebner basis for the ideal only in case it has not been computed yet). }

 

\subsection{\textcolor{Chapter }{FrobeniusForm}}
\logpage{[ 4, 13, 7 ]}\nobreak
\hyperdef{L}{X7A3EACE782DC2198}{}
{\noindent\textcolor{FuncColor}{$\triangleright$\enspace\texttt{FrobeniusForm({\mdseries\slshape A})\index{FrobeniusForm@\texttt{FrobeniusForm}}
\label{FrobeniusForm}
}\hfill{\scriptsize (attribute)}}\\


 Arguments: \mbox{\texttt{\mdseries\slshape A}} \texttt{\symbol{45}}\texttt{\symbol{45}} a quotient of a path algebra.\\
 \textbf{\indent Returns:\ }
 false if \mbox{\texttt{\mdseries\slshape A}} is not selfinjective algebra. Otherwise it returns the Frobenius form of \mbox{\texttt{\mdseries\slshape A}}. 

}

 

\subsection{\textcolor{Chapter }{FrobeniusLinearFunctional}}
\logpage{[ 4, 13, 8 ]}\nobreak
\hyperdef{L}{X789D1DB97C1B9A0D}{}
{\noindent\textcolor{FuncColor}{$\triangleright$\enspace\texttt{FrobeniusLinearFunctional({\mdseries\slshape A})\index{FrobeniusLinearFunctional@\texttt{FrobeniusLinearFunctional}}
\label{FrobeniusLinearFunctional}
}\hfill{\scriptsize (attribute)}}\\


 Arguments: \mbox{\texttt{\mdseries\slshape A}} \texttt{\symbol{45}}\texttt{\symbol{45}} a quotient of a path algebra.\\
 \textbf{\indent Returns:\ }
 false if \mbox{\texttt{\mdseries\slshape A}} is not selfinjective algebra. Otherwise it returns the Frobenius linear
functional of \mbox{\texttt{\mdseries\slshape A}}, which is used to construct the Frobenius form. 

}

 

\subsection{\textcolor{Chapter }{GlobalDimension}}
\logpage{[ 4, 13, 9 ]}\nobreak
\hyperdef{L}{X7D511B3E7A50AB2A}{}
{\noindent\textcolor{FuncColor}{$\triangleright$\enspace\texttt{GlobalDimension({\mdseries\slshape A})\index{GlobalDimension@\texttt{GlobalDimension}}
\label{GlobalDimension}
}\hfill{\scriptsize (attribute)}}\\


 Arguments: \mbox{\texttt{\mdseries\slshape A}} \texttt{\symbol{45}}\texttt{\symbol{45}} an algebra.\\
 \textbf{\indent Returns:\ }
 the global dimension of the algebra \mbox{\texttt{\mdseries\slshape A}} if it is known. Otherwise it returns an error message saying "no method
found!". 



 There is no method installed for this yet. }

 

\subsection{\textcolor{Chapter }{LoewyLength}}
\logpage{[ 4, 13, 10 ]}\nobreak
\hyperdef{L}{X7BEA44FB819910B6}{}
{\noindent\textcolor{FuncColor}{$\triangleright$\enspace\texttt{LoewyLength({\mdseries\slshape A})\index{LoewyLength@\texttt{LoewyLength}}
\label{LoewyLength}
}\hfill{\scriptsize (attribute)}}\\


 Arguments: \mbox{\texttt{\mdseries\slshape A}} \texttt{\symbol{45}}\texttt{\symbol{45}} a path algebra or a quotient of a
path algebra.\\
 \textbf{\indent Returns:\ }
 fail if \mbox{\texttt{\mdseries\slshape A}} is not finite dimensional. Otherwise it returns the Loewy length of the
algebra \mbox{\texttt{\mdseries\slshape A}}. 

}

 

\subsection{\textcolor{Chapter }{NakayamaAutomorphism}}
\logpage{[ 4, 13, 11 ]}\nobreak
\hyperdef{L}{X818AF5A979F8E539}{}
{\noindent\textcolor{FuncColor}{$\triangleright$\enspace\texttt{NakayamaAutomorphism({\mdseries\slshape A})\index{NakayamaAutomorphism@\texttt{NakayamaAutomorphism}}
\label{NakayamaAutomorphism}
}\hfill{\scriptsize (attribute)}}\\


 Arguments: \mbox{\texttt{\mdseries\slshape A}} \texttt{\symbol{45}}\texttt{\symbol{45}} a path algebra or a quotient of a
path algebra.\\
 \textbf{\indent Returns:\ }
 false if \mbox{\texttt{\mdseries\slshape A}} is not selfinjective algebra. Otherwise it returns the Nakayama automorphism
of \mbox{\texttt{\mdseries\slshape A}}. 

}

 

\subsection{\textcolor{Chapter }{NakayamaPermutation}}
\logpage{[ 4, 13, 12 ]}\nobreak
\hyperdef{L}{X7C068FE379FBCE18}{}
{\noindent\textcolor{FuncColor}{$\triangleright$\enspace\texttt{NakayamaPermutation({\mdseries\slshape A})\index{NakayamaPermutation@\texttt{NakayamaPermutation}}
\label{NakayamaPermutation}
}\hfill{\scriptsize (attribute)}}\\


 Arguments: \mbox{\texttt{\mdseries\slshape A}} \texttt{\symbol{45}}\texttt{\symbol{45}} a path algebra or a quotient of a
path algebra.\\
 \textbf{\indent Returns:\ }
 false if \mbox{\texttt{\mdseries\slshape A}} is not selfinjective algebra. Otherwise it returns a list of two elements
where the first is the Nakayama permutation on the simple modules and the
second is the Nakayama permutation on the index set of the simple modules of \mbox{\texttt{\mdseries\slshape A}}. 

}

 

\subsection{\textcolor{Chapter }{OrderOfNakayamaAutomorphism}}
\logpage{[ 4, 13, 13 ]}\nobreak
\hyperdef{L}{X78E378EB83BA3D8A}{}
{\noindent\textcolor{FuncColor}{$\triangleright$\enspace\texttt{OrderOfNakayamaAutomorphism({\mdseries\slshape A})\index{OrderOfNakayamaAutomorphism@\texttt{OrderOfNakayamaAutomorphism}}
\label{OrderOfNakayamaAutomorphism}
}\hfill{\scriptsize (attribute)}}\\


 Arguments: \mbox{\texttt{\mdseries\slshape A}} \texttt{\symbol{45}}\texttt{\symbol{45}} a path algebra or a quotient of a
path algebra.\\
 \textbf{\indent Returns:\ }
 false if \mbox{\texttt{\mdseries\slshape A}} is not selfinjective algebra. Otherwise it returns the order of the Nakayama
autormorphism of \mbox{\texttt{\mdseries\slshape A}} AS AN ELEMENT OF $\operatorname{Aut}(A)$. 

}

 

\subsection{\textcolor{Chapter }{RadicalSeriesOfAlgebra}}
\logpage{[ 4, 13, 14 ]}\nobreak
\hyperdef{L}{X782EFE477EC0C1C6}{}
{\noindent\textcolor{FuncColor}{$\triangleright$\enspace\texttt{RadicalSeriesOfAlgebra({\mdseries\slshape A})\index{RadicalSeriesOfAlgebra@\texttt{RadicalSeriesOfAlgebra}}
\label{RadicalSeriesOfAlgebra}
}\hfill{\scriptsize (attribute)}}\\


 Arguments: \mbox{\texttt{\mdseries\slshape A}} \texttt{\symbol{45}}\texttt{\symbol{45}} an algebra.\\
 \textbf{\indent Returns:\ }
 the radical series of the algebra \mbox{\texttt{\mdseries\slshape A}} in a list, where the first element is the algebra \mbox{\texttt{\mdseries\slshape A}} itself, then radical of \mbox{\texttt{\mdseries\slshape A}}, radical square of \mbox{\texttt{\mdseries\slshape A}}, and so on. 

}

 }

 
\section{\textcolor{Chapter }{Attributes and Operations on Elements of Quotients of Path Algebra}}\label{Attributes_and_Operations_on_elts_quotients_path_algs}
\logpage{[ 4, 14, 0 ]}
\hyperdef{L}{X7BD7DB497917893C}{}
{
  

\subsection{\textcolor{Chapter }{IsElementOfQuotientOfPathAlgebra}}
\logpage{[ 4, 14, 1 ]}\nobreak
\hyperdef{L}{X87495684791B5742}{}
{\noindent\textcolor{FuncColor}{$\triangleright$\enspace\texttt{IsElementOfQuotientOfPathAlgebra({\mdseries\slshape object})\index{IsElementOfQuotientOfPathAlgebra@\texttt{IsElementOfQuotientOfPathAlgebra}}
\label{IsElementOfQuotientOfPathAlgebra}
}\hfill{\scriptsize (property)}}\\


 Arguments: \mbox{\texttt{\mdseries\slshape object}} \texttt{\symbol{45}}\texttt{\symbol{45}} any object in \textsf{GAP}.\\
 \textbf{\indent Returns:\ }
 true whenever \mbox{\texttt{\mdseries\slshape object}} is an element of some quotient of a path algebra. 

}

 
\begin{Verbatim}[commandchars=!@|,fontsize=\small,frame=single,label=Example]
  !gapprompt@gap>| !gapinput@elem := quot.a*quot.b;|
  [(1)*a*b]
  !gapprompt@gap>| !gapinput@IsElementOfQuotientOfPathAlgebra(elem);|
  true
  !gapprompt@gap>| !gapinput@IsElementOfQuotientOfPathAlgebra(FQ.a*FQ.b);    |
  false 
\end{Verbatim}
 

\subsection{\textcolor{Chapter }{Coefficients}}
\logpage{[ 4, 14, 2 ]}\nobreak
\hyperdef{L}{X80B32F667BF6AFD8}{}
{\noindent\textcolor{FuncColor}{$\triangleright$\enspace\texttt{Coefficients({\mdseries\slshape B, element})\index{Coefficients@\texttt{Coefficients}}
\label{Coefficients}
}\hfill{\scriptsize (operation)}}\\


 Arguments: \mbox{\texttt{\mdseries\slshape B, element}} \texttt{\symbol{45}}\texttt{\symbol{45}} a basis for a quotient of a path
algebra and element thereof.\\
 \textbf{\indent Returns:\ }
 the coefficients of the \mbox{\texttt{\mdseries\slshape element}} in terms of the canonical basis \mbox{\texttt{\mdseries\slshape B}} of the quotient of a path algebra in which \mbox{\texttt{\mdseries\slshape element}} is an element. 

}

 

\subsection{\textcolor{Chapter }{IsNormalForm}}
\logpage{[ 4, 14, 3 ]}\nobreak
\hyperdef{L}{X8271E6F27C2C826E}{}
{\noindent\textcolor{FuncColor}{$\triangleright$\enspace\texttt{IsNormalForm({\mdseries\slshape element})\index{IsNormalForm@\texttt{IsNormalForm}}
\label{IsNormalForm}
}\hfill{\scriptsize (operation)}}\\


 Arguments: \mbox{\texttt{\mdseries\slshape element}} \texttt{\symbol{45}}\texttt{\symbol{45}} an element of a path algebra.\\
 \textbf{\indent Returns:\ }
true if \mbox{\texttt{\mdseries\slshape element}} is known to be in normal form. 

}

 
\begin{Verbatim}[commandchars=!@|,fontsize=\small,frame=single,label=Example]
   gap> IsNormalForm(elem);  
   true  
\end{Verbatim}
 

\subsection{\textcolor{Chapter }{{\textless} (for two elements of a path algebra)}}
\logpage{[ 4, 14, 4 ]}\nobreak
\hyperdef{L}{X837DD99B7A233FB5}{}
{\noindent\textcolor{FuncColor}{$\triangleright$\enspace\texttt{{\textless}({\mdseries\slshape a, b})\index{{\textless}@\texttt{{\textless}}!for two elements of a path algebra}
\label{<:for two elements of a path algebra}
}\hfill{\scriptsize (operation)}}\\


 Arguments: \mbox{\texttt{\mdseries\slshape a}} and \mbox{\texttt{\mdseries\slshape b}} \texttt{\symbol{45}}\texttt{\symbol{45}} elements from a path algebra.\\
 \textbf{\indent Returns:\ }
true whenever \mbox{\texttt{\mdseries\slshape a}} {\textless} \mbox{\texttt{\mdseries\slshape b}}. 

}

  

\subsection{\textcolor{Chapter }{ElementOfQuotientOfPathAlgebra}}
\logpage{[ 4, 14, 5 ]}\nobreak
\hyperdef{L}{X7F2527747A3D0D6D}{}
{\noindent\textcolor{FuncColor}{$\triangleright$\enspace\texttt{ElementOfQuotientOfPathAlgebra({\mdseries\slshape family, element, computenormal})\index{ElementOfQuotientOfPathAlgebra@\texttt{ElementOfQuotientOfPathAlgebra}}
\label{ElementOfQuotientOfPathAlgebra}
}\hfill{\scriptsize (operation)}}\\


 Arguments: \mbox{\texttt{\mdseries\slshape family}} \texttt{\symbol{45}}\texttt{\symbol{45}} a family of elements, \mbox{\texttt{\mdseries\slshape element}} \texttt{\symbol{45}}\texttt{\symbol{45}} an element of a path algebra, \mbox{\texttt{\mdseries\slshape computenormal}} \texttt{\symbol{45}}\texttt{\symbol{45}} true or false.\\
 \textbf{\indent Returns:\ }
The projection of \mbox{\texttt{\mdseries\slshape element}} into the quotient given by \mbox{\texttt{\mdseries\slshape family}}. If \mbox{\texttt{\mdseries\slshape computenormal}} is false, then the normal form of the projection of \mbox{\texttt{\mdseries\slshape element}} is returned. 



\mbox{\texttt{\mdseries\slshape family}} is the ElementsFamily of the family of the algebra \mbox{\texttt{\mdseries\slshape element}} is projected into. }

 

\subsection{\textcolor{Chapter }{ElementsOfUnits}}
\logpage{[ 4, 14, 6 ]}\nobreak
\hyperdef{L}{X7AEB3DA18675A263}{}
{\noindent\textcolor{FuncColor}{$\triangleright$\enspace\texttt{ElementsOfUnits({\mdseries\slshape A})\index{ElementsOfUnits@\texttt{ElementsOfUnits}}
\label{ElementsOfUnits}
}\hfill{\scriptsize (attribute)}}\\


 Arguments: \mbox{\texttt{\mdseries\slshape A}} \texttt{\symbol{45}}\texttt{\symbol{45}} a quiver algebra.\\
 \textbf{\indent Returns:\ }
 all the units in the finite dimensional quiver algebra \mbox{\texttt{\mdseries\slshape A}} over a finite field if it has been computed via the command \texttt{Units(A)}. (cf. \texttt{Units} (\ref{Units})) 

}

 

\subsection{\textcolor{Chapter }{OriginalPathAlgebra}}
\logpage{[ 4, 14, 7 ]}\nobreak
\hyperdef{L}{X8088721187BA8D82}{}
{\noindent\textcolor{FuncColor}{$\triangleright$\enspace\texttt{OriginalPathAlgebra({\mdseries\slshape algebra})\index{OriginalPathAlgebra@\texttt{OriginalPathAlgebra}}
\label{OriginalPathAlgebra}
}\hfill{\scriptsize (attribute)}}\\


 Arguments: \mbox{\texttt{\mdseries\slshape algebra}} \texttt{\symbol{45}}\texttt{\symbol{45}} an algebra.\\
 \textbf{\indent Returns:\ }
a path algebra.



If \mbox{\texttt{\mdseries\slshape algebra}} is a quotient of a path algebra or just a path algebra itself, the returned
algebra is the path algebra it was constructed from. Otherwise it returns an
error saying that the algebra entered was not given as a quotient of a path
algebra. }

 

\subsection{\textcolor{Chapter }{Units}}
\logpage{[ 4, 14, 8 ]}\nobreak
\hyperdef{L}{X853C045B7BA6A580}{}
{\noindent\textcolor{FuncColor}{$\triangleright$\enspace\texttt{Units({\mdseries\slshape A})\index{Units@\texttt{Units}}
\label{Units}
}\hfill{\scriptsize (attribute)}}\\


 Arguments: \mbox{\texttt{\mdseries\slshape A}} \texttt{\symbol{45}}\texttt{\symbol{45}} a quiver algebra.\\
 \textbf{\indent Returns:\ }
a representation of the group of the units in the finite dimensional quiver
algebra \mbox{\texttt{\mdseries\slshape A}} over a finite field. It also compute all the invertible elements in \mbox{\texttt{\mdseries\slshape A}}. It can be accessed via the command \texttt{ElementsOfUnits} (\texttt{ElementsOfUnits} (\ref{ElementsOfUnits})). 

}

 }

 
\section{\textcolor{Chapter }{Predefined classes and classes of (quotients of) path algebras}}\label{qpa:predefinedalgebras}
\logpage{[ 4, 15, 0 ]}
\hyperdef{L}{X7B209E0A7DD93C08}{}
{
  

\subsection{\textcolor{Chapter }{BrauerConfigurationAlgebra}}
\logpage{[ 4, 15, 1 ]}\nobreak
\hyperdef{L}{X820AE0208636B9AA}{}
{\noindent\textcolor{FuncColor}{$\triangleright$\enspace\texttt{BrauerConfigurationAlgebra({\mdseries\slshape field, brauer{\textunderscore}configuration})\index{BrauerConfigurationAlgebra@\texttt{BrauerConfigurationAlgebra}}
\label{BrauerConfigurationAlgebra}
}\hfill{\scriptsize (function)}}\\


 Arguments: \mbox{\texttt{\mdseries\slshape field}} \texttt{\symbol{45}}\texttt{\symbol{45}} a field, \mbox{\texttt{\mdseries\slshape brauer{\textunderscore}configuration}} \texttt{\symbol{45}}\texttt{\symbol{45}} a list of the form [[vertices],
[edges/polygons], [orientations]].\\
 \textbf{\indent Returns:\ }
 the Brauer configuration algebra corresponding to the brauer configuration \mbox{\texttt{\mdseries\slshape brauer{\textunderscore}configuration}} over the field \mbox{\texttt{\mdseries\slshape field}}. If the brauer configuration entered is not valid, fail is returned. 



The \mbox{\texttt{\mdseries\slshape brauer{\textunderscore}configuration}} consists of vertices, polygons/edges, and orientations corresponding to a
Brauer Configuration or Brauer Tree. Each vertex must have the form
["vertexname", multiplicity]. Each edge/polygon must have the form
["edgename", "vertex1name", "vertex2name", ...]. There must be an orientation
corresponding to each vertex. Orientations must have the form
["edge1name/polygon1name", "edge2name/polygon2name", ...].}

 
\begin{Verbatim}[commandchars=!@|,fontsize=\small,frame=single,label=Example]
  !gapprompt@gap>| !gapinput@alg := BrauerConfigurationAlgebra(Rationals, [ [ [ "v1", 1 ], [ "v2", 1 ], [ "v3", 2 ] ], [ [ "e1", "v1", "v2" ], [ "e2", "v2", "v3" ] ], [ [ "e1" ], [ "e1", "e2" ], ["e2" ] ]]);|
  <A quotient of the path algebra <Rationals[<quiver with 2 vertices and 
  3 arrows>]> modulo the ideal 
  <two-sided ideal in <Rationals[<quiver with 2 vertices and 3 arrows>]>, 
    (5 generators)>>
\end{Verbatim}
 

\subsection{\textcolor{Chapter }{CanonicalAlgebra}}
\logpage{[ 4, 15, 2 ]}\nobreak
\hyperdef{L}{X7E63A4F37856A075}{}
{\noindent\textcolor{FuncColor}{$\triangleright$\enspace\texttt{CanonicalAlgebra({\mdseries\slshape field, weights[, relcoeff]})\index{CanonicalAlgebra@\texttt{CanonicalAlgebra}}
\label{CanonicalAlgebra}
}\hfill{\scriptsize (operation)}}\\


 Arguments: \mbox{\texttt{\mdseries\slshape field}} \texttt{\symbol{45}}\texttt{\symbol{45}} a field, \mbox{\texttt{\mdseries\slshape weights}} \texttt{\symbol{45}}\texttt{\symbol{45}} a list of positive integers, [, \mbox{\texttt{\mdseries\slshape relcoeff}} \texttt{\symbol{45}}\texttt{\symbol{45}} a list of non\texttt{\symbol{45}}zero
elements in the field.\\
 \textbf{\indent Returns:\ }
the canonical algebra over the \mbox{\texttt{\mdseries\slshape field}} with the quiver given by the weight sequence \mbox{\texttt{\mdseries\slshape weights}} and the relations given by the coefficients \mbox{\texttt{\mdseries\slshape relcoeff}}. 



It function checks if all the \mbox{\texttt{\mdseries\slshape weights}} are greater or equal to two, the number of weights is at least two, the number
of coefficients is the number of \mbox{\texttt{\mdseries\slshape weights}} \texttt{\symbol{45}} 2, the coefficients for the relations are in field and
non\texttt{\symbol{45}}zero. If only the two first arguments are given, then
the number of weights must be two. }

 

\subsection{\textcolor{Chapter }{FinteGlobalDimGreen}}
\logpage{[ 4, 15, 3 ]}\nobreak
\hyperdef{L}{X822FF74F8025DD5F}{}
{\noindent\textcolor{FuncColor}{$\triangleright$\enspace\texttt{FinteGlobalDimGreen({\mdseries\slshape field, n})\index{FinteGlobalDimGreen@\texttt{FinteGlobalDimGreen}}
\label{FinteGlobalDimGreen}
}\hfill{\scriptsize (operation)}}\\


 Arguments: \mbox{\texttt{\mdseries\slshape field}} \texttt{\symbol{45}}\texttt{\symbol{45}} a field, \mbox{\texttt{\mdseries\slshape n}} \texttt{\symbol{45}}\texttt{\symbol{45}} a positive integer.\\
 \textbf{\indent Returns:\ }
a finite dimensional algebra over the field \mbox{\texttt{\mdseries\slshape field}} with two simple modules of finite global dimension $2n$ for$n\geq 2$ and $rad^{2n}\neq 0$. 



The algebras were defined by Edward L. Green in the paper "Remarks on
projective resolutions" Representation theory, II (Proc. Second Internat.
Conf., Carleton Univ., Ottawa, Ont., 1979), pp. 259\texttt{\symbol{45}}279,
Lecture Notes in Math. 832, Springer Berlin (1980), }

 

\subsection{\textcolor{Chapter }{FinteGlobalDimKirkKuz}}
\logpage{[ 4, 15, 4 ]}\nobreak
\hyperdef{L}{X7C1C6ED178D38C31}{}
{\noindent\textcolor{FuncColor}{$\triangleright$\enspace\texttt{FinteGlobalDimKirkKuz({\mdseries\slshape field, n})\index{FinteGlobalDimKirkKuz@\texttt{FinteGlobalDimKirkKuz}}
\label{FinteGlobalDimKirkKuz}
}\hfill{\scriptsize (operation)}}\\


 Arguments: \mbox{\texttt{\mdseries\slshape field}} \texttt{\symbol{45}}\texttt{\symbol{45}} a field, \mbox{\texttt{\mdseries\slshape n}} \texttt{\symbol{45}}\texttt{\symbol{45}} a positive integer.\\
 \textbf{\indent Returns:\ }
a finite dimensional algebra over the field \mbox{\texttt{\mdseries\slshape field}} with two simple modules of finite global dimension $2n + 1$ for $n\geq 2$ and $2$ for $n = 1$ and $rad^4 = 0$ for $n\geq 2$ and $rad^3 = 0$ for $n = 1$. 



The algebras were defined by Ellen Kirkman and James Kuzmanovich in the paper
"Algebras with large homological dimension", Proc. Amer. Math. Soc. 109, no.
4, 903\texttt{\symbol{45}}906 (1990). }

 Returns the finite dimensional algebra with two simple modules 

\subsection{\textcolor{Chapter }{KroneckerAlgebra}}
\logpage{[ 4, 15, 5 ]}\nobreak
\hyperdef{L}{X83498D3D856CC08A}{}
{\noindent\textcolor{FuncColor}{$\triangleright$\enspace\texttt{KroneckerAlgebra({\mdseries\slshape field, n})\index{KroneckerAlgebra@\texttt{KroneckerAlgebra}}
\label{KroneckerAlgebra}
}\hfill{\scriptsize (operation)}}\\


 Arguments: \mbox{\texttt{\mdseries\slshape field}} \texttt{\symbol{45}}\texttt{\symbol{45}} a field, \mbox{\texttt{\mdseries\slshape n}} \texttt{\symbol{45}}\texttt{\symbol{45}} a positive integer.\\
 \textbf{\indent Returns:\ }
the \mbox{\texttt{\mdseries\slshape n}}\texttt{\symbol{45}}Kronecker algebra over the field \mbox{\texttt{\mdseries\slshape field}}. 



It function checks if the number \mbox{\texttt{\mdseries\slshape n}} of arrows is greater or equal to two and returns an error message if not. }

 

\subsection{\textcolor{Chapter }{NakayamaAlgebra}}
\logpage{[ 4, 15, 6 ]}\nobreak
\hyperdef{L}{X7C678A08836F77CC}{}
{\noindent\textcolor{FuncColor}{$\triangleright$\enspace\texttt{NakayamaAlgebra({\mdseries\slshape admiss\texttt{\symbol{45}}seq, field})\index{NakayamaAlgebra@\texttt{NakayamaAlgebra}}
\label{NakayamaAlgebra}
}\hfill{\scriptsize (function)}}\\


 Arguments: \mbox{\texttt{\mdseries\slshape field}} \texttt{\symbol{45}}\texttt{\symbol{45}} a field, \mbox{\texttt{\mdseries\slshape admiss\texttt{\symbol{45}}seq}} \texttt{\symbol{45}}\texttt{\symbol{45}} a list of positive integers.\\
 \textbf{\indent Returns:\ }
 the Nakayama algebra corresponding to the admissible sequence \mbox{\texttt{\mdseries\slshape admiss\texttt{\symbol{45}}seq}} over the field \mbox{\texttt{\mdseries\slshape field}}. If the entered sequence is not an admissible sequence, the sequence is
returned. 



The \mbox{\texttt{\mdseries\slshape admiss\texttt{\symbol{45}}seq}} consists of the dimensions of the projective representations.}

 

\subsection{\textcolor{Chapter }{AdmissibleSequenceGenerator}}
\logpage{[ 4, 15, 7 ]}\nobreak
\hyperdef{L}{X8369BB398212101C}{}
{\noindent\textcolor{FuncColor}{$\triangleright$\enspace\texttt{AdmissibleSequenceGenerator({\mdseries\slshape num{\textunderscore}vert, height})\index{AdmissibleSequenceGenerator@\texttt{AdmissibleSequenceGenerator}}
\label{AdmissibleSequenceGenerator}
}\hfill{\scriptsize (operation)}}\\


 Arguments: \mbox{\texttt{\mdseries\slshape num{\textunderscore}vert}}, \mbox{\texttt{\mdseries\slshape height}} \texttt{\symbol{45}}\texttt{\symbol{45}} two positive integers.\\
 \textbf{\indent Returns:\ }
all admissible sequences of Nakayama algebras with \mbox{\texttt{\mdseries\slshape num{\textunderscore}vert}} number of vertices, and indecomposable projective modules of length at most \mbox{\texttt{\mdseries\slshape height}}. 

}

 
\begin{Verbatim}[commandchars=!@|,fontsize=\small,frame=single,label=Example]
  !gapprompt@gap>| !gapinput@alg := NakayamaAlgebra([2,1], Rationals);|
  <Rationals[<quiver with 2 vertices and 1 arrows>]>
  !gapprompt@gap>| !gapinput@QuiverOfPathAlgebra(alg);|
  <quiver with 2 vertices and 1 arrows>
\end{Verbatim}
 

\subsection{\textcolor{Chapter }{PosetAlgebra}}
\logpage{[ 4, 15, 8 ]}\nobreak
\hyperdef{L}{X78C675F2836D1B18}{}
{\noindent\textcolor{FuncColor}{$\triangleright$\enspace\texttt{PosetAlgebra({\mdseries\slshape F, P})\index{PosetAlgebra@\texttt{PosetAlgebra}}
\label{PosetAlgebra}
}\hfill{\scriptsize (operation)}}\\


 Arguments: \mbox{\texttt{\mdseries\slshape F}} \texttt{\symbol{45}}\texttt{\symbol{45}} a field, \mbox{\texttt{\mdseries\slshape P}} \texttt{\symbol{45}}\texttt{\symbol{45}} a poset.\\
 \textbf{\indent Returns:\ }
the poset algebra associated to the poset \mbox{\texttt{\mdseries\slshape P}} over the field \mbox{\texttt{\mdseries\slshape K}}. 

}

 

\subsection{\textcolor{Chapter }{PosetOfPosetAlgebra}}
\logpage{[ 4, 15, 9 ]}\nobreak
\hyperdef{L}{X7F877C5F839A3AA9}{}
{\noindent\textcolor{FuncColor}{$\triangleright$\enspace\texttt{PosetOfPosetAlgebra({\mdseries\slshape A})\index{PosetOfPosetAlgebra@\texttt{PosetOfPosetAlgebra}}
\label{PosetOfPosetAlgebra}
}\hfill{\scriptsize (attribute)}}\\


 Arguments: \mbox{\texttt{\mdseries\slshape A}} \texttt{\symbol{45}}\texttt{\symbol{45}} a poset algebra.\\
 \textbf{\indent Returns:\ }
 the poset from which that poset algebra \mbox{\texttt{\mdseries\slshape A}} is constructed. 

}

 

\subsection{\textcolor{Chapter }{TruncatedPathAlgebra}}
\logpage{[ 4, 15, 10 ]}\nobreak
\hyperdef{L}{X7EF1AE62790D7486}{}
{\noindent\textcolor{FuncColor}{$\triangleright$\enspace\texttt{TruncatedPathAlgebra({\mdseries\slshape F, Q, n})\index{TruncatedPathAlgebra@\texttt{TruncatedPathAlgebra}}
\label{TruncatedPathAlgebra}
}\hfill{\scriptsize (operation)}}\\


 Arguments: \mbox{\texttt{\mdseries\slshape F}} \texttt{\symbol{45}}\texttt{\symbol{45}} a field, \mbox{\texttt{\mdseries\slshape Q}} \texttt{\symbol{45}}\texttt{\symbol{45}} a quiver, \mbox{\texttt{\mdseries\slshape n}} \texttt{\symbol{45}}\texttt{\symbol{45}} a positive integer.\\
 \textbf{\indent Returns:\ }
the truncated path algebra \mbox{\texttt{\mdseries\slshape KQ/I}}, where \mbox{\texttt{\mdseries\slshape I}} is the ideal generated by all paths of length \mbox{\texttt{\mdseries\slshape n}} in \mbox{\texttt{\mdseries\slshape KQ}}. 

}

 

\subsection{\textcolor{Chapter }{IsSpecialBiserialQuiver}}
\logpage{[ 4, 15, 11 ]}\nobreak
\hyperdef{L}{X87E17E137E2B0FC4}{}
{\noindent\textcolor{FuncColor}{$\triangleright$\enspace\texttt{IsSpecialBiserialQuiver({\mdseries\slshape Q})\index{IsSpecialBiserialQuiver@\texttt{IsSpecialBiserialQuiver}}
\label{IsSpecialBiserialQuiver}
}\hfill{\scriptsize (property)}}\\


 Arguments: \mbox{\texttt{\mdseries\slshape Q}} \texttt{\symbol{45}}\texttt{\symbol{45}} a quiver. \\
 \textbf{\indent Returns:\ }
true whenever \mbox{\texttt{\mdseries\slshape Q}} is a \emph{"special biserial"} quiver, i.e. every vertex in \mbox{\texttt{\mdseries\slshape Q}} is a source (resp. target) of at most 2 arrows. 



 \emph{Note:} e.g. a path algebra of one loop IS NOT special biserial, but one loop IS
special biserial quiver (cf. \texttt{IsSpecialBiserialAlgebra} (\ref{IsSpecialBiserialAlgebra}) and also an Example below). }

 
\begin{Verbatim}[commandchars=!@|,fontsize=\small,frame=single,label=Example]
  !gapprompt@gap>| !gapinput@Q := Quiver(1, [ [1,1,"a"], [1,1,"b"] ]);;  |
  !gapprompt@gap>| !gapinput@A := PathAlgebra(Rationals, Q);;|
  !gapprompt@gap>| !gapinput@IsSpecialBiserialAlgebra(A); IsStringAlgebra(A);|
  false
  false
  !gapprompt@gap>| !gapinput@rel1 := [A.a*A.b, A.a^2, A.b^2];  |
  [ (1)*a*b, (1)*a^2, (1)*b^2 ]
  !gapprompt@gap>| !gapinput@quo1 := A/rel1;;|
  !gapprompt@gap>| !gapinput@IsSpecialBiserialAlgebra(quo1); IsStringAlgebra(quo1);|
  true
  true
  !gapprompt@gap>| !gapinput@rel2 := [A.a*A.b-A.b*A.a, A.a^2, A.b^2];  |
  [ (1)*a*b+(-1)*b*a, (1)*a^2, (1)*b^2 ]
  !gapprompt@gap>| !gapinput@quo2 := A/rel2;;|
  !gapprompt@gap>| !gapinput@IsSpecialBiserialAlgebra(quo2); IsStringAlgebra(quo2);|
  true
  false
  !gapprompt@gap>| !gapinput@rel3 := [A.a*A.b+A.b*A.a, A.a^2, A.b^2, A.b*A.a];   |
  [ (1)*a*b+(1)*b*a, (1)*a^2, (1)*b^2, (1)*b*a ]
  !gapprompt@gap>| !gapinput@quo3 := A/rel3;;|
  !gapprompt@gap>| !gapinput@IsSpecialBiserialAlgebra(quo3); IsStringAlgebra(quo3);|
  true
  true
  !gapprompt@gap>| !gapinput@rel4 := [A.a*A.b, A.a^2, A.b^3]; |
  [ (1)*a*b, (1)*a^2, (1)*b^3 ]
  !gapprompt@gap>| !gapinput@quo4 := A/rel4;;|
  !gapprompt@gap>| !gapinput@IsSpecialBiserialAlgebra(quo4); IsStringAlgebra(quo4);|
  false
  false 
\end{Verbatim}
 }

 
\section{\textcolor{Chapter }{Opposite algebras}}\label{qpa:opposite}
\logpage{[ 4, 16, 0 ]}
\hyperdef{L}{X840BAB827C62AA4C}{}
{
  

\subsection{\textcolor{Chapter }{OppositePath}}
\logpage{[ 4, 16, 1 ]}\nobreak
\hyperdef{L}{X85794BE082B632B9}{}
{\noindent\textcolor{FuncColor}{$\triangleright$\enspace\texttt{OppositePath({\mdseries\slshape p})\index{OppositePath@\texttt{OppositePath}}
\label{OppositePath}
}\hfill{\scriptsize (operation)}}\\


 Arguments: \mbox{\texttt{\mdseries\slshape p}} \texttt{\symbol{45}}\texttt{\symbol{45}} a path.\\
 \textbf{\indent Returns:\ }
 the path corresponding to \mbox{\texttt{\mdseries\slshape p}} in the opposite quiver. 

}

 The following example illustrates the use of \texttt{OppositeQuiver} (\ref{OppositeQuiver}) and \texttt{OppositePath} (\ref{OppositePath}). 
\begin{Verbatim}[commandchars=!@|,fontsize=\small,frame=single,label=Example]
  !gapprompt@gap>| !gapinput@Q := Quiver( [ "u", "v" ], [ [ "u", "u", "a" ], |
  !gapprompt@>| !gapinput@             [ "u", "v", "b" ] ] );|
  <quiver with 2 vertices and 2 arrows>
  !gapprompt@gap>| !gapinput@Qop := OppositeQuiver(Q);|
  <quiver with 2 vertices and 2 arrows>
  !gapprompt@gap>| !gapinput@VerticesOfQuiver( Qop );|
  [ u_op, v_op ]
  !gapprompt@gap>| !gapinput@ArrowsOfQuiver( Qop );|
  [ a_op, b_op ]
  !gapprompt@gap>| !gapinput@OppositePath( Q.a * Q.b );|
  b_op*a_op
  !gapprompt@gap>| !gapinput@IsIdenticalObj( Q, OppositeQuiver( Qop ) );|
  true
  !gapprompt@gap>| !gapinput@OppositePath( Qop.b_op * Qop.a_op );|
  a*b
\end{Verbatim}
 

\subsection{\textcolor{Chapter }{OppositePathAlgebra}}
\logpage{[ 4, 16, 2 ]}\nobreak
\hyperdef{L}{X87A86AFB782211D6}{}
{\noindent\textcolor{FuncColor}{$\triangleright$\enspace\texttt{OppositePathAlgebra({\mdseries\slshape A})\index{OppositePathAlgebra@\texttt{OppositePathAlgebra}}
\label{OppositePathAlgebra}
}\hfill{\scriptsize (attribute)}}\\


 Arguments: \mbox{\texttt{\mdseries\slshape A}} \texttt{\symbol{45}}\texttt{\symbol{45}} a path algebra or quotient of path
algebra.\\
 \textbf{\indent Returns:\ }
 the opposite algebra $\mbox{\texttt{\mdseries\slshape A}}^\text{op}$. 



 This attribute contains the opposite algebra of an algebra.

 The opposite algebra of a path algebra is the path algebra over the opposite
quiver (as given by \texttt{OppositeQuiver} (\ref{OppositeQuiver})). The opposite algebra of a quotient of a path algebra has the opposite
quiver and the opposite relations of the original algebra.

 The function \texttt{OppositePathAlgebraElement} (\ref{OppositePathAlgebraElement}) takes an algebra element to the corresponding element in the opposite algebra.

 The opposite of the opposite of an algebra $A$ is isomorphic to $A$. In QPA, we regard these two algebras to be the same, so the call \texttt{OppositePathAlgebra(OppositePathAlgebra(A))} returns the object \texttt{A}. }

 

\subsection{\textcolor{Chapter }{OppositePathAlgebraElement}}
\logpage{[ 4, 16, 3 ]}\nobreak
\hyperdef{L}{X82C303CE808D54C1}{}
{\noindent\textcolor{FuncColor}{$\triangleright$\enspace\texttt{OppositePathAlgebraElement({\mdseries\slshape x})\index{OppositePathAlgebraElement@\texttt{OppositePathAlgebraElement}}
\label{OppositePathAlgebraElement}
}\hfill{\scriptsize (function)}}\\


 Arguments: \mbox{\texttt{\mdseries\slshape x}} \texttt{\symbol{45}}\texttt{\symbol{45}} a path.\\
 \textbf{\indent Returns:\ }
 the element corresponding to \mbox{\texttt{\mdseries\slshape x}} in the opposite algebra. 

}

 

\subsection{\textcolor{Chapter }{OppositeAlgebraHomomorphism}}
\logpage{[ 4, 16, 4 ]}\nobreak
\hyperdef{L}{X86B3EDE679B2493E}{}
{\noindent\textcolor{FuncColor}{$\triangleright$\enspace\texttt{OppositeAlgebraHomomorphism({\mdseries\slshape A})\index{OppositeAlgebraHomomorphism@\texttt{OppositeAlgebraHomomorphism}}
\label{OppositeAlgebraHomomorphism}
}\hfill{\scriptsize (attribute)}}\\


 Arguments: \mbox{\texttt{\mdseries\slshape A}} \texttt{\symbol{45}}\texttt{\symbol{45}} an algebra homomorphism of two
(quotients) of path algebras.\\
 \textbf{\indent Returns:\ }
 algebra homomorphism between the corresponding opposite algebras. 

}

 The following example illustrates the use of \texttt{OppositePathAlgebra} (\ref{OppositePathAlgebra}) and \texttt{OppositePathAlgebraElement} (\ref{OppositePathAlgebraElement}). 
\begin{Verbatim}[commandchars=!@|,fontsize=\small,frame=single,label=Example]
  !gapprompt@gap>| !gapinput@Q := Quiver( [ "u", "v" ], [ [ "u", "u", "a" ], |
  !gapprompt@>| !gapinput@             [ "u", "v", "b" ] ] );|
  <quiver with 2 vertices and 2 arrows>
  !gapprompt@gap>| !gapinput@A := PathAlgebra( Rationals, Q );|
  <Rationals[<quiver with 2 vertices and 2 arrows>]>
  !gapprompt@gap>| !gapinput@OppositePathAlgebra( A );|
  <Rationals[<quiver with 2 vertices and 2 arrows>]>
  !gapprompt@gap>| !gapinput@OppositePathAlgebraElement( A.u + 2*A.a + 5*A.a*A.b );|
  (1)*u_op+(2)*a_op+(5)*b_op*a_op
  !gapprompt@gap>| !gapinput@IsIdenticalObj( A, |
  !gapprompt@>| !gapinput@        OppositePathAlgebra( OppositePathAlgebra( A ) ) );|
  true
\end{Verbatim}
 }

 
\section{\textcolor{Chapter }{Tensor products of path algebras}}\label{qpa:patensor}
\logpage{[ 4, 17, 0 ]}
\hyperdef{L}{X842527EC7F90C8C5}{}
{
  If $\Lambda$ and $\Gamma$ are quotients of path algebras over the same field $F$, then their tensor product $\Lambda \tensor_F \Gamma$ is also a quotient of a path algebra over $F$.

 The quiver for the tensor product path algebra is the \texttt{QuiverProduct} (\ref{QuiverProduct}) of the quivers of the original algebras.

 The operation \texttt{TensorProductOfAlgebras} (\ref{TensorProductOfAlgebras}) computes the tensor products of two quotients of path algebras as a quotient
of a path algebra. 

\subsection{\textcolor{Chapter }{QuiverProduct}}
\logpage{[ 4, 17, 1 ]}\nobreak
\hyperdef{L}{X835BBBE18104654A}{}
{\noindent\textcolor{FuncColor}{$\triangleright$\enspace\texttt{QuiverProduct({\mdseries\slshape Q1, Q2})\index{QuiverProduct@\texttt{QuiverProduct}}
\label{QuiverProduct}
}\hfill{\scriptsize (operation)}}\\


 Arguments: \mbox{\texttt{\mdseries\slshape Q1}} and \mbox{\texttt{\mdseries\slshape Q2}} \texttt{\symbol{45}}\texttt{\symbol{45}} quivers.\\
 \textbf{\indent Returns:\ }
 the product quiver $\mbox{\texttt{\mdseries\slshape Q1}} \quiverproduct \mbox{\texttt{\mdseries\slshape Q2}}$. 



 A vertex in $\mbox{\texttt{\mdseries\slshape Q1}} \quiverproduct \mbox{\texttt{\mdseries\slshape Q2}}$ which is made by combining a vertex named \texttt{u} in \mbox{\texttt{\mdseries\slshape Q1}} with a vertex \texttt{v} in \mbox{\texttt{\mdseries\slshape Q2}} is named \texttt{u{\textunderscore}v}. Arrows are named similarly (they are made by combining an arrow from one
quiver with a vertex from the other). }

 

\subsection{\textcolor{Chapter }{QuiverProductDecomposition}}
\logpage{[ 4, 17, 2 ]}\nobreak
\hyperdef{L}{X858517C18242C2F1}{}
{\noindent\textcolor{FuncColor}{$\triangleright$\enspace\texttt{QuiverProductDecomposition({\mdseries\slshape Q})\index{QuiverProductDecomposition@\texttt{QuiverProductDecomposition}}
\label{QuiverProductDecomposition}
}\hfill{\scriptsize (attribute)}}\\


 Arguments: \mbox{\texttt{\mdseries\slshape Q}} \texttt{\symbol{45}}\texttt{\symbol{45}} a quiver.\\
 \textbf{\indent Returns:\ }
 the original quivers \mbox{\texttt{\mdseries\slshape Q}} is a product of, if \mbox{\texttt{\mdseries\slshape Q}} was created by the \texttt{QuiverProduct} (\ref{QuiverProduct}) operation. 



 The value of this attribute is an object in the category \texttt{IsQuiverProductDecomposition} (\ref{IsQuiverProductDecomposition}). }

 

\subsection{\textcolor{Chapter }{IsQuiverProductDecomposition}}
\logpage{[ 4, 17, 3 ]}\nobreak
\hyperdef{L}{X80E3731882B80106}{}
{\noindent\textcolor{FuncColor}{$\triangleright$\enspace\texttt{IsQuiverProductDecomposition({\mdseries\slshape object})\index{IsQuiverProductDecomposition@\texttt{IsQuiverProductDecomposition}}
\label{IsQuiverProductDecomposition}
}\hfill{\scriptsize (category)}}\\


 Arguments: \mbox{\texttt{\mdseries\slshape object}} \texttt{\symbol{45}}\texttt{\symbol{45}} any object in \textsf{GAP}. 

 Category for objects containing information about the relation between a
product quiver and the quivers it is a product of. The quiver factors can be
extracted from the decomposition object by using the [] notation (like
accessing elements of a list). The decomposition object is also used by the
operations \texttt{IncludeInProductQuiver} (\ref{IncludeInProductQuiver}) and \texttt{ProjectFromProductQuiver} (\ref{ProjectFromProductQuiver}). }

 

\subsection{\textcolor{Chapter }{IncludeInProductQuiver}}
\logpage{[ 4, 17, 4 ]}\nobreak
\hyperdef{L}{X856E8B5B7F550647}{}
{\noindent\textcolor{FuncColor}{$\triangleright$\enspace\texttt{IncludeInProductQuiver({\mdseries\slshape L, Q})\index{IncludeInProductQuiver@\texttt{IncludeInProductQuiver}}
\label{IncludeInProductQuiver}
}\hfill{\scriptsize (operation)}}\\


 Arguments: \mbox{\texttt{\mdseries\slshape L}} \texttt{\symbol{45}}\texttt{\symbol{45}} a list containing the paths $q_1$ and $q_2$, $Q$ \texttt{\symbol{45}}\texttt{\symbol{45}} a product quiver.\\
 \textbf{\indent Returns:\ }
a path in \mbox{\texttt{\mdseries\slshape Q}}. 



 Includes paths $q_1$ and $q_2$ from two quivers into the product of these quivers, \mbox{\texttt{\mdseries\slshape Q}}. If at least one of $q_1$ and $q_2$ is a vertex, there is exactly one possible inclusion. If they are both
non\texttt{\symbol{45}}trivial paths, there are several possibilities. This
operation constructs the path which is the inclusion of $q_1$ at the source of $q_2$ multiplied with the inclusion of $q_2$ at the target of $q_1$. }

 

\subsection{\textcolor{Chapter }{ProjectFromProductQuiver}}
\logpage{[ 4, 17, 5 ]}\nobreak
\hyperdef{L}{X8455692C7E282C6C}{}
{\noindent\textcolor{FuncColor}{$\triangleright$\enspace\texttt{ProjectFromProductQuiver({\mdseries\slshape i, p})\index{ProjectFromProductQuiver@\texttt{ProjectFromProductQuiver}}
\label{ProjectFromProductQuiver}
}\hfill{\scriptsize (operation)}}\\


 Arguments: \mbox{\texttt{\mdseries\slshape i}} \texttt{\symbol{45}}\texttt{\symbol{45}} a positive integer, \mbox{\texttt{\mdseries\slshape p}} \texttt{\symbol{45}}\texttt{\symbol{45}} a path in the product quiver.\\
 \textbf{\indent Returns:\ }
the projection of the product quiver path \mbox{\texttt{\mdseries\slshape p}} to one of the factors. Which factor it should be projected to is specified by
the argument \mbox{\texttt{\mdseries\slshape i}}. 

}

 The following example shows how the operations related to quiver products are
used. 
\begin{Verbatim}[commandchars=!@|,fontsize=\small,frame=single,label=Example]
  !gapprompt@gap>| !gapinput@q1 := Quiver( [ "u1", "u2" ], [ [ "u1", "u2", "a" ] ] );|
  <quiver with 2 vertices and 1 arrows>
  !gapprompt@gap>| !gapinput@q2 := Quiver( [ "v1", "v2", "v3" ],|
                        [ [ "v1", "v2", "b" ],
                          [ "v2", "v3", "c" ] ] );
  <quiver with 3 vertices and 2 arrows>
  !gapprompt@gap>| !gapinput@q1_q2 := QuiverProduct( q1, q2 );|
  <quiver with 6 vertices and 7 arrows>
  !gapprompt@gap>| !gapinput@q1_q2.u1_b * q1_q2.a_v2;|
  u1_b*a_v2
  !gapprompt@gap>| !gapinput@IncludeInProductQuiver( [ q1.a, q2.b * q2.c ], q1_q2 );|
  a_v1*u2_b*u2_c
  !gapprompt@gap>| !gapinput@ProjectFromProductQuiver( 2, q1_q2.a_v1 * q1_q2.u2_b * q1_q2.u2_c );|
  b*c
  !gapprompt@gap>| !gapinput@q1_q2_dec := QuiverProductDecomposition( q1_q2 );|
  <object>
  !gapprompt@gap>| !gapinput@q1_q2_dec[ 1 ];|
  <quiver with 2 vertices and 1 arrows>
  !gapprompt@gap>| !gapinput@q1_q2_dec[ 1 ] = q1;|
  true  
\end{Verbatim}
 

\subsection{\textcolor{Chapter }{TensorProductOfAlgebras}}
\logpage{[ 4, 17, 6 ]}\nobreak
\hyperdef{L}{X7A9026937BDDFA6C}{}
{\noindent\textcolor{FuncColor}{$\triangleright$\enspace\texttt{TensorProductOfAlgebras({\mdseries\slshape FQ1, FQ2})\index{TensorProductOfAlgebras@\texttt{TensorProductOfAlgebras}}
\label{TensorProductOfAlgebras}
}\hfill{\scriptsize (operation)}}\\


 Arguments: \mbox{\texttt{\mdseries\slshape FQ1}} and \mbox{\texttt{\mdseries\slshape FQ2}} \texttt{\symbol{45}}\texttt{\symbol{45}} (quotients of) path algebras.\\
 \textbf{\indent Returns:\ }
 The tensor product of \mbox{\texttt{\mdseries\slshape FQ1}} and \mbox{\texttt{\mdseries\slshape FQ2}}. 



 The result is a quotient of a path algebra, whose quiver is the \texttt{QuiverProduct} (\ref{QuiverProduct}) of the quivers of the operands. }

 

\subsection{\textcolor{Chapter }{TensorAlgebrasInclusion}}
\logpage{[ 4, 17, 7 ]}\nobreak
\hyperdef{L}{X7EE8921D787C8377}{}
{\noindent\textcolor{FuncColor}{$\triangleright$\enspace\texttt{TensorAlgebrasInclusion({\mdseries\slshape T, n})\index{TensorAlgebrasInclusion@\texttt{TensorAlgebrasInclusion}}
\label{TensorAlgebrasInclusion}
}\hfill{\scriptsize (operation)}}\\


 Arguments: \mbox{\texttt{\mdseries\slshape T}} \texttt{\symbol{45}}\texttt{\symbol{45}} quiver algebra, \mbox{\texttt{\mdseries\slshape n}} \texttt{\symbol{45}}\texttt{\symbol{45}} 1 or 2.\\
 \textbf{\indent Returns:\ }
 Returns the inclusion $A \hookrightarrow A \otimes B$ or the inclusion $B \hookrightarrow A \otimes B$ if $n = 1$ or $n = 2$, respectively. 

}

 

\subsection{\textcolor{Chapter }{SimpleTensor}}
\logpage{[ 4, 17, 8 ]}\nobreak
\hyperdef{L}{X7B31F4F680135E72}{}
{\noindent\textcolor{FuncColor}{$\triangleright$\enspace\texttt{SimpleTensor({\mdseries\slshape L, T})\index{SimpleTensor@\texttt{SimpleTensor}}
\label{SimpleTensor}
}\hfill{\scriptsize (operation)}}\\


 Arguments: \mbox{\texttt{\mdseries\slshape L}} \texttt{\symbol{45}}\texttt{\symbol{45}} a list containing two elements $x$ and $y$ of two (quotients of) path algebras, \mbox{\texttt{\mdseries\slshape T}} \texttt{\symbol{45}}\texttt{\symbol{45}} the tensor product of these algebras.\\
 \textbf{\indent Returns:\ }
 the simple tensor $x \tensor y$. 



 $x \tensor y$ is in the tensor product \mbox{\texttt{\mdseries\slshape T}} (produced by \texttt{TensorProductOfAlgebras} (\ref{TensorProductOfAlgebras})). }

 

\subsection{\textcolor{Chapter }{TensorProductDecomposition}}
\logpage{[ 4, 17, 9 ]}\nobreak
\hyperdef{L}{X7F0EBF88866A537D}{}
{\noindent\textcolor{FuncColor}{$\triangleright$\enspace\texttt{TensorProductDecomposition({\mdseries\slshape T})\index{TensorProductDecomposition@\texttt{TensorProductDecomposition}}
\label{TensorProductDecomposition}
}\hfill{\scriptsize (attribute)}}\\


 Arguments: \mbox{\texttt{\mdseries\slshape T}} \texttt{\symbol{45}}\texttt{\symbol{45}} a tensor product of path algebras.\\
 \textbf{\indent Returns:\ }
a list of the factors in the tensor product. 



\mbox{\texttt{\mdseries\slshape T}} should be produced by \texttt{TensorProductOfAlgebras} (\ref{TensorProductOfAlgebras})). }

 The following example shows how the operations for tensor products of
quotients of path algebras are used. 
\begin{Verbatim}[commandchars=!@|,fontsize=\small,frame=single,label=Example]
  !gapprompt@gap>| !gapinput@q1 := Quiver( [ "u1", "u2" ], [ [ "u1", "u2", "a" ] ] );|
  <quiver with 2 vertices and 1 arrows>
  !gapprompt@gap>| !gapinput@q2 := Quiver( [ "v1", "v2", "v3", "v4" ],|
  !gapprompt@>| !gapinput@                      [ [ "v1", "v2", "b" ],|
  !gapprompt@>| !gapinput@                        [ "v1", "v3", "c" ],|
  !gapprompt@>| !gapinput@                        [ "v2", "v4", "d" ],|
  !gapprompt@>| !gapinput@                        [ "v3", "v4", "e" ] ] );|
  <quiver with 4 vertices and 4 arrows>
  !gapprompt@gap>| !gapinput@fq1 := PathAlgebra( Rationals, q1 );|
  <Rationals[<quiver with 2 vertices and 1 arrows>]>
  !gapprompt@gap>| !gapinput@fq2 := PathAlgebra( Rationals, q2 );|
  <Rationals[<quiver with 4 vertices and 4 arrows>]>
  !gapprompt@gap>| !gapinput@I := Ideal( fq2, [ fq2.b * fq2.d - fq2.c * fq2.e ] );|
  <two-sided ideal in <Rationals[<quiver with 4 vertices and 4 arrows>]>
      , (1 generators)>
  !gapprompt@gap>| !gapinput@quot := fq2 / I;|
  <Rationals[<quiver with 4 vertices and 4 arrows>]/
  <two-sided ideal in <Rationals[<quiver with 4 vertices and 4 arrows>]>
      , (1 generators)>>
  !gapprompt@gap>| !gapinput@t := TensorProductOfAlgebras( fq1, quot );|
  <Rationals[<quiver with 8 vertices and 12 arrows>]/
  <two-sided ideal in <Rationals[<quiver with 8 vertices and 
      12 arrows>]>, (6 generators)>>
  !gapprompt@gap>| !gapinput@SimpleTensor( [ fq1.a, quot.b ], t );|
  [(1)*u1_b*a_v2]
  !gapprompt@gap>| !gapinput@t_dec := TensorProductDecomposition( t );|
  [ <Rationals[<quiver with 2 vertices and 1 arrows>]>, 
    <Rationals[<quiver with 4 vertices and 4 arrows>]/
      <two-sided ideal in <Rationals[<quiver with 4 vertices and 
          4 arrows>]>, (1 generators)>> ]
  !gapprompt@gap>| !gapinput@t_dec[ 1 ] = fq1;|
  true
\end{Verbatim}
 

\subsection{\textcolor{Chapter }{EnvelopingAlgebra}}
\logpage{[ 4, 17, 10 ]}\nobreak
\hyperdef{L}{X820195C47E2BE7E0}{}
{\noindent\textcolor{FuncColor}{$\triangleright$\enspace\texttt{EnvelopingAlgebra({\mdseries\slshape A})\index{EnvelopingAlgebra@\texttt{EnvelopingAlgebra}}
\label{EnvelopingAlgebra}
}\hfill{\scriptsize (attribute)}}\\


 Arguments: \mbox{\texttt{\mdseries\slshape A}} \texttt{\symbol{45}}\texttt{\symbol{45}} a (quotient of) a path algebra. \\
 \textbf{\indent Returns:\ }
the enveloping algebra $\mbox{\texttt{\mdseries\slshape A}}^\text{e} = \mbox{\texttt{\mdseries\slshape A}}^\text{op} \tensor \mbox{\texttt{\mdseries\slshape A}}$ of \mbox{\texttt{\mdseries\slshape A}}. 

}

 

\subsection{\textcolor{Chapter }{EnvelopingAlgebraHomomorphism}}
\logpage{[ 4, 17, 11 ]}\nobreak
\hyperdef{L}{X8210F6627AB95229}{}
{\noindent\textcolor{FuncColor}{$\triangleright$\enspace\texttt{EnvelopingAlgebraHomomorphism({\mdseries\slshape f})\index{EnvelopingAlgebraHomomorphism@\texttt{EnvelopingAlgebraHomomorphism}}
\label{EnvelopingAlgebraHomomorphism}
}\hfill{\scriptsize (attribute)}}\\


 Arguments: \mbox{\texttt{\mdseries\slshape f}} \texttt{\symbol{45}}\texttt{\symbol{45}} an algebra homomorphism of two
(quotients of) a path algebras. \\
 \textbf{\indent Returns:\ }
the algebra homomorphism between the corresponding enveloping algebras the
source of \mbox{\texttt{\mdseries\slshape f}} and the range of \mbox{\texttt{\mdseries\slshape f}}. 

}

 

\subsection{\textcolor{Chapter }{IsEnvelopingAlgebra}}
\logpage{[ 4, 17, 12 ]}\nobreak
\hyperdef{L}{X7AE13B567B5F72FC}{}
{\noindent\textcolor{FuncColor}{$\triangleright$\enspace\texttt{IsEnvelopingAlgebra({\mdseries\slshape A})\index{IsEnvelopingAlgebra@\texttt{IsEnvelopingAlgebra}}
\label{IsEnvelopingAlgebra}
}\hfill{\scriptsize (property)}}\\


 Arguments: \mbox{\texttt{\mdseries\slshape A}} \texttt{\symbol{45}}\texttt{\symbol{45}} an algebra.\\
 \textbf{\indent Returns:\ }
true if and only if \mbox{\texttt{\mdseries\slshape A}} is the result of a call to \texttt{EnvelopingAlgebra} (\ref{EnvelopingAlgebra}). 

}

 

\subsection{\textcolor{Chapter }{AlgebraAsModuleOverEnvelopingAlgebra}}
\logpage{[ 4, 17, 13 ]}\nobreak
\hyperdef{L}{X80D827747ACD76FA}{}
{\noindent\textcolor{FuncColor}{$\triangleright$\enspace\texttt{AlgebraAsModuleOverEnvelopingAlgebra({\mdseries\slshape A})\index{AlgebraAsModuleOverEnvelopingAlgebra@\texttt{Algebra}\-\texttt{As}\-\texttt{Module}\-\texttt{Over}\-\texttt{Enveloping}\-\texttt{Algebra}}
\label{AlgebraAsModuleOverEnvelopingAlgebra}
}\hfill{\scriptsize (attribute)}}\\


 Arguments: \mbox{\texttt{\mdseries\slshape A}} \texttt{\symbol{45}}\texttt{\symbol{45}} a (quotient of a) path algebra \mbox{\texttt{\mdseries\slshape A}}. \\
 \textbf{\indent Returns:\ }
the algebra \mbox{\texttt{\mdseries\slshape A}} as a right module over the enveloping algebra of \mbox{\texttt{\mdseries\slshape A}}. 

}

 

\subsection{\textcolor{Chapter }{DualOfAlgebraAsModuleOverEnvelopingAlgebra}}
\logpage{[ 4, 17, 14 ]}\nobreak
\hyperdef{L}{X7A4C262287D74AB0}{}
{\noindent\textcolor{FuncColor}{$\triangleright$\enspace\texttt{DualOfAlgebraAsModuleOverEnvelopingAlgebra({\mdseries\slshape A})\index{DualOfAlgebraAsModuleOverEnvelopingAlgebra@\texttt{Dual}\-\texttt{Of}\-\texttt{Algebra}\-\texttt{As}\-\texttt{Module}\-\texttt{Over}\-\texttt{Enveloping}\-\texttt{Algebra}}
\label{DualOfAlgebraAsModuleOverEnvelopingAlgebra}
}\hfill{\scriptsize (attribute)}}\\


 Arguments: \mbox{\texttt{\mdseries\slshape A}} \texttt{\symbol{45}}\texttt{\symbol{45}} a finite dimensional (admissible
quotient of) a path algebra \mbox{\texttt{\mdseries\slshape A}}. \\
 \textbf{\indent Returns:\ }
the algebra \mbox{\texttt{\mdseries\slshape A}} as a right module over the enveloping algebra of \mbox{\texttt{\mdseries\slshape A}}. 

}

 

\subsection{\textcolor{Chapter }{TrivialExtensionOfQuiverAlgebra}}
\logpage{[ 4, 17, 15 ]}\nobreak
\hyperdef{L}{X8101415F7FFB34CF}{}
{\noindent\textcolor{FuncColor}{$\triangleright$\enspace\texttt{TrivialExtensionOfQuiverAlgebra({\mdseries\slshape A})\index{TrivialExtensionOfQuiverAlgebra@\texttt{TrivialExtensionOfQuiverAlgebra}}
\label{TrivialExtensionOfQuiverAlgebra}
}\hfill{\scriptsize (attribute)}}\\


 Arguments: \mbox{\texttt{\mdseries\slshape A}} \texttt{\symbol{45}}\texttt{\symbol{45}} a finite dimensional (admissible
quotient of) a path algebra \mbox{\texttt{\mdseries\slshape A}}. \\
 \textbf{\indent Returns:\ }
the trivial extension algebra $T(A)=A\oplus D(A)$ of the entered algebra \mbox{\texttt{\mdseries\slshape A}}. 

}

 

\subsection{\textcolor{Chapter }{TrivialExtensionOfQuiverAlgebraProjection}}
\logpage{[ 4, 17, 16 ]}\nobreak
\hyperdef{L}{X81D98182822E8911}{}
{\noindent\textcolor{FuncColor}{$\triangleright$\enspace\texttt{TrivialExtensionOfQuiverAlgebraProjection({\mdseries\slshape A})\index{TrivialExtensionOfQuiverAlgebraProjection@\texttt{Trivial}\-\texttt{Extension}\-\texttt{Of}\-\texttt{Quiver}\-\texttt{Algebra}\-\texttt{Projection}}
\label{TrivialExtensionOfQuiverAlgebraProjection}
}\hfill{\scriptsize (attribute)}}\\


 Arguments: \mbox{\texttt{\mdseries\slshape A}} \texttt{\symbol{45}}\texttt{\symbol{45}} a finite dimensional (admissible
quotient of) a path algebra \mbox{\texttt{\mdseries\slshape A}}. \\
 \textbf{\indent Returns:\ }
the natural algebra projection form trivial extension $T(A)$ to \mbox{\texttt{\mdseries\slshape A}}. 

}

 }

 
\section{\textcolor{Chapter }{Operations on quiver algebras}}\label{qpa:operationsonquiveralgebras}
\logpage{[ 4, 18, 0 ]}
\hyperdef{L}{X79B8B91F8097BB80}{}
{
  

\subsection{\textcolor{Chapter }{QuiverAlgebraOfAmodAeA}}
\logpage{[ 4, 18, 1 ]}\nobreak
\hyperdef{L}{X7EE5A11883B86971}{}
{\noindent\textcolor{FuncColor}{$\triangleright$\enspace\texttt{QuiverAlgebraOfAmodAeA({\mdseries\slshape A, elist})\index{QuiverAlgebraOfAmodAeA@\texttt{QuiverAlgebraOfAmodAeA}}
\label{QuiverAlgebraOfAmodAeA}
}\hfill{\scriptsize (operation)}}\\


 Arguments: \mbox{\texttt{\mdseries\slshape A}}, \mbox{\texttt{\mdseries\slshape elist}} \texttt{\symbol{45}} a finite dimensional quiver algebra and a list of
integers.\\
 \textbf{\indent Returns:\ }
a quiver algebra isomorphic \mbox{\texttt{\mdseries\slshape A}} modulo the ideal generated by a sum of vertices in \mbox{\texttt{\mdseries\slshape A}}. 



 Given a quiver algebra \mbox{\texttt{\mdseries\slshape A}} and a sum of vertices $e$, this function computes the quiver algebra $A/AeA$. The list \mbox{\texttt{\mdseries\slshape elist}} is a list of integers, where each integer occurring in the list corresponds to
the position of the vertex in the vertices defining the idempotent $e$. }

 

\subsection{\textcolor{Chapter }{QuiverAlgebraOfeAe}}
\logpage{[ 4, 18, 2 ]}\nobreak
\hyperdef{L}{X7B2D7385829F5EC6}{}
{\noindent\textcolor{FuncColor}{$\triangleright$\enspace\texttt{QuiverAlgebraOfeAe({\mdseries\slshape A, elist})\index{QuiverAlgebraOfeAe@\texttt{QuiverAlgebraOfeAe}}
\label{QuiverAlgebraOfeAe}
}\hfill{\scriptsize (operation)}}\\


 Arguments: \mbox{\texttt{\mdseries\slshape A}}, \mbox{\texttt{\mdseries\slshape e}} \texttt{\symbol{45}} a finite dimensional quiver algebra and an idempotent.\\
 \textbf{\indent Returns:\ }
a quiver algebra isomorphic $eAe$, where \mbox{\texttt{\mdseries\slshape A}} is the entered algebra and \mbox{\texttt{\mdseries\slshape e}} is the entered idempotent. 



 Given a quiver algebra \mbox{\texttt{\mdseries\slshape A}} and an idempotent \mbox{\texttt{\mdseries\slshape e}}, this function computes the quiver algebra isomorphic to $eAe$. The function checks if the entered element \mbox{\texttt{\mdseries\slshape e}} is an idempotent in \mbox{\texttt{\mdseries\slshape A}}. }

 }

 
\section{\textcolor{Chapter }{Finite dimensional algebras over finite fields}}\label{qpa:moregeneralalgebrastuff}
\logpage{[ 4, 19, 0 ]}
\hyperdef{L}{X8561BCB6835D561F}{}
{
  

\subsection{\textcolor{Chapter }{AlgebraAsQuiverAlgebra}}
\logpage{[ 4, 19, 1 ]}\nobreak
\hyperdef{L}{X810A29FB7E6EA24D}{}
{\noindent\textcolor{FuncColor}{$\triangleright$\enspace\texttt{AlgebraAsQuiverAlgebra({\mdseries\slshape A})\index{AlgebraAsQuiverAlgebra@\texttt{AlgebraAsQuiverAlgebra}}
\label{AlgebraAsQuiverAlgebra}
}\hfill{\scriptsize (operation)}}\\


 Arguments: \mbox{\texttt{\mdseries\slshape A}} \texttt{\symbol{45}} a finite dimensional algebra over a finite field.\\
 \textbf{\indent Returns:\ }
a (quotient of a) path algebra isomorphic to the entered algebra \mbox{\texttt{\mdseries\slshape A}} whenever possible and a list of the images of the vertices and the arrows in
this path algebra in \mbox{\texttt{\mdseries\slshape A}}. 



 The operation only applies when \mbox{\texttt{\mdseries\slshape A}} is a finite dimensional indecomposable algebra over a finite field, otherwise
it returns an error message. It checks the algebra \mbox{\texttt{\mdseries\slshape A}} is basic and elementary over some field and otherwise it returns an error
message. In the list of images the images of the vertices are listed first and
then the images of the arrows. }

 

\subsection{\textcolor{Chapter }{BlocksOfAlgebra}}
\logpage{[ 4, 19, 2 ]}\nobreak
\hyperdef{L}{X7E9078077EE8B51B}{}
{\noindent\textcolor{FuncColor}{$\triangleright$\enspace\texttt{BlocksOfAlgebra({\mdseries\slshape A})\index{BlocksOfAlgebra@\texttt{BlocksOfAlgebra}}
\label{BlocksOfAlgebra}
}\hfill{\scriptsize (operation)}}\\


 Arguments: \mbox{\texttt{\mdseries\slshape A}} \texttt{\symbol{45}} a finite dimensional algebra.\\
 \textbf{\indent Returns:\ }
a block decomposition of the entered algebra \mbox{\texttt{\mdseries\slshape A}} as a list of indecomposable algebras. 

}

 

\subsection{\textcolor{Chapter }{IsBasicAlgebra}}
\logpage{[ 4, 19, 3 ]}\nobreak
\hyperdef{L}{X84B423137F933795}{}
{\noindent\textcolor{FuncColor}{$\triangleright$\enspace\texttt{IsBasicAlgebra({\mdseries\slshape A})\index{IsBasicAlgebra@\texttt{IsBasicAlgebra}}
\label{IsBasicAlgebra}
}\hfill{\scriptsize (property)}}\\


 Arguments: \mbox{\texttt{\mdseries\slshape A}} \texttt{\symbol{45}} a finite dimensional algebra over a finite field.\\
 \textbf{\indent Returns:\ }
 true if the entered algebra \mbox{\texttt{\mdseries\slshape A}} is a (finite dimensional) basic algebra and false otherwise. This method only
applies to algebras over finite fields. 

}

 

\subsection{\textcolor{Chapter }{IsElementaryAlgebra}}
\logpage{[ 4, 19, 4 ]}\nobreak
\hyperdef{L}{X7D30E9C878221B42}{}
{\noindent\textcolor{FuncColor}{$\triangleright$\enspace\texttt{IsElementaryAlgebra({\mdseries\slshape A})\index{IsElementaryAlgebra@\texttt{IsElementaryAlgebra}}
\label{IsElementaryAlgebra}
}\hfill{\scriptsize (property)}}\\


 Arguments: \mbox{\texttt{\mdseries\slshape A}} \texttt{\symbol{45}} a finite dimensional algebra over a finite field.\\
 \textbf{\indent Returns:\ }
 true if the entered algebra \mbox{\texttt{\mdseries\slshape A}} is a (finite dimensional) elementary algebra and false otherwise. This method
only applies to algebras over finite fields. 



 The algebra \mbox{\texttt{\mdseries\slshape A}} need not to be an elementary algebra over the field which it is defined, but
be an elementary algebra over a field extension. }

 

\subsection{\textcolor{Chapter }{PreprojectiveAlgebra}}
\logpage{[ 4, 19, 5 ]}\nobreak
\hyperdef{L}{X7B35109B8176FE56}{}
{\noindent\textcolor{FuncColor}{$\triangleright$\enspace\texttt{PreprojectiveAlgebra({\mdseries\slshape M, n, or, A})\index{PreprojectiveAlgebra@\texttt{PreprojectiveAlgebra}}
\label{PreprojectiveAlgebra}
}\hfill{\scriptsize (operation)}}\\


 Arguments (first version): \mbox{\texttt{\mdseries\slshape M}} \texttt{\symbol{45}} a path algebra module over finite dimensional hereditary
algebra over a finite field, \mbox{\texttt{\mdseries\slshape n}} \texttt{\symbol{45}} an integer.\\
 Arguments (second version): \mbox{\texttt{\mdseries\slshape A}} \texttt{\symbol{45}} a path algebra.\\
 \textbf{\indent Returns:\ }
in the first version the preprojective algebra of the module \mbox{\texttt{\mdseries\slshape M}} if it only has support up to degree \mbox{\texttt{\mdseries\slshape n}}. In the second version it returns the preprojective algebra of a hereditary
algebra given by a path algebra over a field. 

}

 

\subsection{\textcolor{Chapter }{PrimitiveIdempotents}}
\logpage{[ 4, 19, 6 ]}\nobreak
\hyperdef{L}{X80C0C6C37C4A2ABD}{}
{\noindent\textcolor{FuncColor}{$\triangleright$\enspace\texttt{PrimitiveIdempotents({\mdseries\slshape A})\index{PrimitiveIdempotents@\texttt{PrimitiveIdempotents}}
\label{PrimitiveIdempotents}
}\hfill{\scriptsize (attribute)}}\\


 Arguments: \mbox{\texttt{\mdseries\slshape A}} \texttt{\symbol{45}} a finite dimensional simple algebra over a finite field.\\
 \textbf{\indent Returns:\ }
 a complete set of primitive idempotents $\{ e_i \}$ such that $A \simeq Ae_1 + ... + Ae_n$. 



 TODO: Understand what this function actually does. }

 }

 
\section{\textcolor{Chapter }{Algebras}}\label{qpa:idempotentalgebrastuff}
\logpage{[ 4, 20, 0 ]}
\hyperdef{L}{X7DDBF6F47A2E021C}{}
{
  

\subsection{\textcolor{Chapter }{LiftingCompleteSetOfOrthogonalIdempotents}}
\logpage{[ 4, 20, 1 ]}\nobreak
\hyperdef{L}{X821B7B047871B42D}{}
{\noindent\textcolor{FuncColor}{$\triangleright$\enspace\texttt{LiftingCompleteSetOfOrthogonalIdempotents({\mdseries\slshape f, e})\index{LiftingCompleteSetOfOrthogonalIdempotents@\texttt{Lifting}\-\texttt{Complete}\-\texttt{Set}\-\texttt{Of}\-\texttt{Orthogonal}\-\texttt{Idempotents}}
\label{LiftingCompleteSetOfOrthogonalIdempotents}
}\hfill{\scriptsize (operation)}}\\


 Arguments: \mbox{\texttt{\mdseries\slshape map}} \texttt{\symbol{45}} an algebra homomorphism, \mbox{\texttt{\mdseries\slshape idempotents}} elements in the range of \mbox{\texttt{\mdseries\slshape map}}.\\
 \textbf{\indent Returns:\ }
a complete set of orthogonal idempotents in \texttt{Source(f)} which are liftings of the entered \mbox{\texttt{\mdseries\slshape idempotents}} whenever possible. 



 The operation only applies when the domain of \mbox{\texttt{\mdseries\slshape f}} is finite dimensional. It checks if the list of elements \mbox{\texttt{\mdseries\slshape idempotents}} is a set complete set of orthogonal idempotents. If some idempotent in \mbox{\texttt{\mdseries\slshape idempotents}} is not in the image of \mbox{\texttt{\mdseries\slshape map}}, then an error message is returned. If all idempotents in \mbox{\texttt{\mdseries\slshape idempotents}} has a preimage, then this operation returns a complete set of orthogonal of
idempotents which is a lifting of the idempotents in \mbox{\texttt{\mdseries\slshape idempotents}} to the source of \mbox{\texttt{\mdseries\slshape f}} whenever possible or it returns fail. }

 

\subsection{\textcolor{Chapter }{LiftingIdempotent}}
\logpage{[ 4, 20, 2 ]}\nobreak
\hyperdef{L}{X83041BDF78BF3CCA}{}
{\noindent\textcolor{FuncColor}{$\triangleright$\enspace\texttt{LiftingIdempotent({\mdseries\slshape f, e})\index{LiftingIdempotent@\texttt{LiftingIdempotent}}
\label{LiftingIdempotent}
}\hfill{\scriptsize (operation)}}\\


 Arguments: \mbox{\texttt{\mdseries\slshape f}} \texttt{\symbol{45}} an algebra homomorphism, \mbox{\texttt{\mdseries\slshape e}} an element in the range of \mbox{\texttt{\mdseries\slshape f}}.\\
 \textbf{\indent Returns:\ }
an idempotent $a$ in \texttt{Source(f)} such that \texttt{ImageElm(f, a) = e} whenever possible. If \mbox{\texttt{\mdseries\slshape e}} is not in the image of \mbox{\texttt{\mdseries\slshape f}}, an error message is given, and if \mbox{\texttt{\mdseries\slshape e}} does not have a preimage by \mbox{\texttt{\mdseries\slshape f}} \texttt{fail} is returned. 



 The operation only applies when the domain of \mbox{\texttt{\mdseries\slshape f}} is finite dimensional. It checks if the element \mbox{\texttt{\mdseries\slshape e}} is an idempotent. If \mbox{\texttt{\mdseries\slshape e}} is not in the image of \mbox{\texttt{\mdseries\slshape f}}, then an error message is returned. If \mbox{\texttt{\mdseries\slshape e}} has a preimage, then this operation returns a lifting of \mbox{\texttt{\mdseries\slshape e}} to the source of \mbox{\texttt{\mdseries\slshape f}} whenever possible or it returns fail. Using the algorithm described in the
proof of Proposition 27.1 in Anderson and Fuller, Rings and categories of
modules, second edition, GMT, Springer\texttt{\symbol{45}}Verlag. }

 }

 
\section{\textcolor{Chapter }{Saving and reading quotients of path algebras to and from a file}}\label{qpa:read-save}
\logpage{[ 4, 21, 0 ]}
\hyperdef{L}{X850B9F12806FF76B}{}
{
  

\subsection{\textcolor{Chapter }{ReadAlgebra}}
\logpage{[ 4, 21, 1 ]}\nobreak
\hyperdef{L}{X82A638C77FA75549}{}
{\noindent\textcolor{FuncColor}{$\triangleright$\enspace\texttt{ReadAlgebra({\mdseries\slshape string})\index{ReadAlgebra@\texttt{ReadAlgebra}}
\label{ReadAlgebra}
}\hfill{\scriptsize (operation)}}\\


 Arguments: \mbox{\texttt{\mdseries\slshape string}} \texttt{\symbol{45}} a name of a file.\\
 \textbf{\indent Returns:\ }
 a finite dimension quotient $A$ of a path algebra saved by command \texttt{SaveAlgebra} to the file \mbox{\texttt{\mdseries\slshape string}}. This function creates the algebra $A$ again, which can be saved to a file again with the function \texttt{SaveAlgebra}. 

}

 

\subsection{\textcolor{Chapter }{SaveAlgebra}}
\logpage{[ 4, 21, 2 ]}\nobreak
\hyperdef{L}{X7E60DDCE848CB739}{}
{\noindent\textcolor{FuncColor}{$\triangleright$\enspace\texttt{SaveAlgebra({\mdseries\slshape A, string, action})\index{SaveAlgebra@\texttt{SaveAlgebra}}
\label{SaveAlgebra}
}\hfill{\scriptsize (operation)}}\\


 Arguments: \mbox{\texttt{\mdseries\slshape A}} \texttt{\symbol{45}} an algebra, \mbox{\texttt{\mdseries\slshape string}} \texttt{\symbol{45}} a name of a file, \mbox{\texttt{\mdseries\slshape action}} \texttt{\symbol{45}} a string.\\
 \textbf{\indent Returns:\ }
 or creates a file with name \mbox{\texttt{\mdseries\slshape string}}, storing the algebra \mbox{\texttt{\mdseries\slshape A}}, which can be opened again with the function \texttt{ReadAlgebra} and reconstructed. The last argument \mbox{\texttt{\mdseries\slshape action}} decides if the file \mbox{\texttt{\mdseries\slshape string}}, if it exists already, should be overwritten, not overwritten or the user
should be prompted for an answer to this question. The corresponding user
inputs are: "delete", "keep" or "query". 

}

 }

 }

 
\chapter{\textcolor{Chapter }{Groebner Basis}}\label{Groebner-Basis}
\logpage{[ 5, 0, 0 ]}
\hyperdef{L}{X8371E66387CB2E49}{}
{
 This chapter contains the declarations and implementations needed for Groebner
basis. Currently, we do not provide algorithms to actually compute Groebner
basis; instead, the declarations and implementations are provided here for \textsf{GAP} objects and the actual elements of Groebner basis are computed by the \textsf{GBNP} package. 
\section{\textcolor{Chapter }{Constructing a Groebner Basis}}\logpage{[ 5, 1, 0 ]}
\hyperdef{L}{X850B47047FD4D709}{}
{
 

\subsection{\textcolor{Chapter }{InfoGroebnerBasis}}
\logpage{[ 5, 1, 1 ]}\nobreak
\hyperdef{L}{X8451936885F68598}{}
{\noindent\textcolor{FuncColor}{$\triangleright$\enspace\texttt{InfoGroebnerBasis\index{InfoGroebnerBasis@\texttt{InfoGroebnerBasis}}
\label{InfoGroebnerBasis}
}\hfill{\scriptsize (info class)}}\\


 is the info class for functions dealing with Groebner basis. }

 

\subsection{\textcolor{Chapter }{GroebnerBasis}}
\logpage{[ 5, 1, 2 ]}\nobreak
\hyperdef{L}{X7A43611E876B7560}{}
{\noindent\textcolor{FuncColor}{$\triangleright$\enspace\texttt{GroebnerBasis({\mdseries\slshape I, rels})\index{GroebnerBasis@\texttt{GroebnerBasis}}
\label{GroebnerBasis}
}\hfill{\scriptsize (operation)}}\\


Arguments: \mbox{\texttt{\mdseries\slshape I}} \texttt{\symbol{45}}\texttt{\symbol{45}} an ideal, \mbox{\texttt{\mdseries\slshape rels}} \texttt{\symbol{45}}\texttt{\symbol{45}} a list of relations generating \mbox{\texttt{\mdseries\slshape I}}.\\
 \textbf{\indent Returns:\ }
 an object \mbox{\texttt{\mdseries\slshape GB}} in the \texttt{IsGroebnerBasis} (\ref{IsGroebnerBasis}) category with \texttt{IsCompleteGroebnerBasis} (\ref{IsCompleteGroebnerBasis}) property set on true.



 Sets also \mbox{\texttt{\mdseries\slshape GB}} as a value of the attribute \texttt{GroebnerBasisOfIdeal} (\ref{GroebnerBasisOfIdeal}) for \mbox{\texttt{\mdseries\slshape I}} (so one has an access to it by calling GroebnerBasisOfIdeal(\mbox{\texttt{\mdseries\slshape I}})).\\
 There are absolutely no computations and no checks for correctness in this
function. Giving a set of relations that does not form a Groebner basis may
result in incorrect answers or unexpected errors. This function is intended to
be used by packages providing access to external Groebner basis programs and
should be invoked before further computations on Groebner basis or ideal I
(cf. also \texttt{IsCompleteGroebnerBasis} (\ref{IsCompleteGroebnerBasis})). }

 }

 
\section{\textcolor{Chapter }{Categories and Properties of Groebner Basis}}\logpage{[ 5, 2, 0 ]}
\hyperdef{L}{X79C7DC5D873A14D0}{}
{
 

\subsection{\textcolor{Chapter }{IsCompletelyReducedGroebnerBasis}}
\logpage{[ 5, 2, 1 ]}\nobreak
\hyperdef{L}{X85C0C1CD87C70AAA}{}
{\noindent\textcolor{FuncColor}{$\triangleright$\enspace\texttt{IsCompletelyReducedGroebnerBasis({\mdseries\slshape gb})\index{IsCompletelyReducedGroebnerBasis@\texttt{IsCompletelyReducedGroebnerBasis}}
\label{IsCompletelyReducedGroebnerBasis}
}\hfill{\scriptsize (property)}}\\


 Arguments: \mbox{\texttt{\mdseries\slshape GB}} \texttt{\symbol{45}}\texttt{\symbol{45}} a Groebner basis.\\
 \textbf{\indent Returns:\ }
 true when \mbox{\texttt{\mdseries\slshape GB}} is a Groebner basis which is completely reduced. 

}

 

\subsection{\textcolor{Chapter }{IsCompleteGroebnerBasis}}
\logpage{[ 5, 2, 2 ]}\nobreak
\hyperdef{L}{X86E7D0AE87CA048D}{}
{\noindent\textcolor{FuncColor}{$\triangleright$\enspace\texttt{IsCompleteGroebnerBasis({\mdseries\slshape gb})\index{IsCompleteGroebnerBasis@\texttt{IsCompleteGroebnerBasis}}
\label{IsCompleteGroebnerBasis}
}\hfill{\scriptsize (property)}}\\


 Arguments: \mbox{\texttt{\mdseries\slshape GB}} \texttt{\symbol{45}}\texttt{\symbol{45}} a Groebner basis.\\
 \textbf{\indent Returns:\ }
true when \mbox{\texttt{\mdseries\slshape GB}} is a complete Groebner basis. 



 While philosophically something that isn't a complete Groebner basis isn't a
Groebner basis at all, this property can be used in conjunction with other
properties to see if the the Groebner basis contains enough information for
computations. An example of a system that creates incomplete Groebner bases is
`Opal'.\\
 \emph{Note:} The current package used for creating Groebner bases is \textsf{GBNP}, and this package does not create incomplete Groebner bases. }

 

\subsection{\textcolor{Chapter }{IsGroebnerBasis}}
\logpage{[ 5, 2, 3 ]}\nobreak
\hyperdef{L}{X7BFD28E687AADFBB}{}
{\noindent\textcolor{FuncColor}{$\triangleright$\enspace\texttt{IsGroebnerBasis({\mdseries\slshape object})\index{IsGroebnerBasis@\texttt{IsGroebnerBasis}}
\label{IsGroebnerBasis}
}\hfill{\scriptsize (category)}}\\


 Arguments: \mbox{\texttt{\mdseries\slshape object}} \texttt{\symbol{45}}\texttt{\symbol{45}} any object in \textsf{GAP}. \\
 \textbf{\indent Returns:\ }
 true when \mbox{\texttt{\mdseries\slshape object}} is a Groebner basis and false otherwise. 



The function only returns true for Groebner bases that has been set as such
using the \texttt{GroebnerBasis} function, as illustrated in the following example. }

 

\subsection{\textcolor{Chapter }{IsHomogeneousGroebnerBasis}}
\logpage{[ 5, 2, 4 ]}\nobreak
\hyperdef{L}{X799FD421784D1FFC}{}
{\noindent\textcolor{FuncColor}{$\triangleright$\enspace\texttt{IsHomogeneousGroebnerBasis({\mdseries\slshape gb})\index{IsHomogeneousGroebnerBasis@\texttt{IsHomogeneousGroebnerBasis}}
\label{IsHomogeneousGroebnerBasis}
}\hfill{\scriptsize (property)}}\\


 Arguments: \mbox{\texttt{\mdseries\slshape GB}} \texttt{\symbol{45}}\texttt{\symbol{45}} a Groebner basis.\\
 \textbf{\indent Returns:\ }
 true when \mbox{\texttt{\mdseries\slshape GB}} is a Groebner basis which is homogeneous. 

}

 
\begin{Verbatim}[commandchars=!@|,fontsize=\small,frame=single,label=Example]
  !gapprompt@gap>| !gapinput@Q := Quiver( 3, [ [1,2,"a"], [2,3,"b"] ] );|
  <quiver with 3 vertices and 2 arrows>
  !gapprompt@gap>| !gapinput@PA := PathAlgebra( Rationals, Q );|
  <Rationals[<quiver with 3 vertices and 2 arrows>]>
  !gapprompt@gap>| !gapinput@rels := [ PA.a*PA.b ];|
  [ (1)*a*b ]
  !gapprompt@gap>| !gapinput@gb := GBNPGroebnerBasis( rels, PA );|
  [ (1)*a*b ]
  !gapprompt@gap>| !gapinput@I := Ideal( PA, gb );|
  <two-sided ideal in <Rationals[<quiver with 3 vertices and 2 arrows>]>
      , (1 generators)>
  !gapprompt@gap>| !gapinput@grb := GroebnerBasis( I, gb );|
  <complete two-sided Groebner basis containing 1 elements>
  !gapprompt@gap>| !gapinput@alg := PA/I;|
  <Rationals[<quiver with 3 vertices and 2 arrows>]/
  <two-sided ideal in <Rationals[<quiver with 3 vertices and 2 arrows>]>
      , (1 generators)>>
  !gapprompt@gap>| !gapinput@IsGroebnerBasis(gb);|
  false
  !gapprompt@gap>| !gapinput@IsGroebnerBasis(grb);|
  true 
\end{Verbatim}
 

\subsection{\textcolor{Chapter }{IsTipReducedGroebnerBasis}}
\logpage{[ 5, 2, 5 ]}\nobreak
\hyperdef{L}{X8592E4C87E41A15A}{}
{\noindent\textcolor{FuncColor}{$\triangleright$\enspace\texttt{IsTipReducedGroebnerBasis({\mdseries\slshape gb})\index{IsTipReducedGroebnerBasis@\texttt{IsTipReducedGroebnerBasis}}
\label{IsTipReducedGroebnerBasis}
}\hfill{\scriptsize (property)}}\\


 Arguments: \mbox{\texttt{\mdseries\slshape GB}} \texttt{\symbol{45}}\texttt{\symbol{45}} a Groebner Basis.\\
 \textbf{\indent Returns:\ }
 true when \mbox{\texttt{\mdseries\slshape GB}} is a Groebner basis which is tip reduced. 

}

 }

 
\section{\textcolor{Chapter }{Attributes and Operations for Groebner Basis}}\logpage{[ 5, 3, 0 ]}
\hyperdef{L}{X84EF455882169920}{}
{
 

\subsection{\textcolor{Chapter }{AdmitsFinitelyManyNontips}}
\logpage{[ 5, 3, 1 ]}\nobreak
\hyperdef{L}{X8722C1577C236116}{}
{\noindent\textcolor{FuncColor}{$\triangleright$\enspace\texttt{AdmitsFinitelyManyNontips({\mdseries\slshape GB})\index{AdmitsFinitelyManyNontips@\texttt{AdmitsFinitelyManyNontips}}
\label{AdmitsFinitelyManyNontips}
}\hfill{\scriptsize (operation)}}\\


 Arguments: \mbox{\texttt{\mdseries\slshape GB}} \texttt{\symbol{45}}\texttt{\symbol{45}} a complete Groebner basis.\\
 \textbf{\indent Returns:\ }
true if the Groebner basis admits only finitely many nontips and false
otherwise. 

}

 

\subsection{\textcolor{Chapter }{CompletelyReduce}}
\logpage{[ 5, 3, 2 ]}\nobreak
\hyperdef{L}{X878AC1107E9671BA}{}
{\noindent\textcolor{FuncColor}{$\triangleright$\enspace\texttt{CompletelyReduce({\mdseries\slshape GB, a})\index{CompletelyReduce@\texttt{CompletelyReduce}}
\label{CompletelyReduce}
}\hfill{\scriptsize (operation)}}\\


 Arguments: \mbox{\texttt{\mdseries\slshape GB}} \texttt{\symbol{45}}\texttt{\symbol{45}} a Groebner basis, \mbox{\texttt{\mdseries\slshape a}} \texttt{\symbol{45}}\texttt{\symbol{45}} an element in a path algebra.\\
 \textbf{\indent Returns:\ }
\mbox{\texttt{\mdseries\slshape a}} reduced by \mbox{\texttt{\mdseries\slshape GB}}. 



 If \mbox{\texttt{\mdseries\slshape a}} is already completely reduced, the original element \mbox{\texttt{\mdseries\slshape a}} is returned. }

 

\subsection{\textcolor{Chapter }{CompletelyReduceGroebnerBasis}}
\logpage{[ 5, 3, 3 ]}\nobreak
\hyperdef{L}{X7FF28B7B80759D24}{}
{\noindent\textcolor{FuncColor}{$\triangleright$\enspace\texttt{CompletelyReduceGroebnerBasis({\mdseries\slshape GB})\index{CompletelyReduceGroebnerBasis@\texttt{CompletelyReduceGroebnerBasis}}
\label{CompletelyReduceGroebnerBasis}
}\hfill{\scriptsize (operation)}}\\


 Arguments: \mbox{\texttt{\mdseries\slshape GB}} \texttt{\symbol{45}}\texttt{\symbol{45}} a Groebner basis.\\
 \textbf{\indent Returns:\ }
 the completely reduced Groebner basis of the ideal generated by \mbox{\texttt{\mdseries\slshape GB}}. 



 The operation modifies a Groebner basis \mbox{\texttt{\mdseries\slshape GB}} such that each relation in \mbox{\texttt{\mdseries\slshape GB}} is completely reduced. The \texttt{IsCompletelyReducedGroebnerBasis} and \texttt{IsTipReducedGroebnerBasis} properties are set as a result of this operation. The resulting relations will
be placed in sorted order according to the ordering of \mbox{\texttt{\mdseries\slshape GB}}. }

 

\subsection{\textcolor{Chapter }{Enumerator}}
\logpage{[ 5, 3, 4 ]}\nobreak
\hyperdef{L}{X7EF8910F82B45EC7}{}
{\noindent\textcolor{FuncColor}{$\triangleright$\enspace\texttt{Enumerator({\mdseries\slshape GB})\index{Enumerator@\texttt{Enumerator}}
\label{Enumerator}
}\hfill{\scriptsize (operation)}}\\


 Arguments: \mbox{\texttt{\mdseries\slshape GB}} \texttt{\symbol{45}}\texttt{\symbol{45}} a Groebner basis. \\
 \textbf{\indent Returns:\ }
an enumerat that enumerates the relations making up the Groebner basis. 



These relations should be enumerated in ascending order with respect to the
ordering for the family the elements are contained in. }

 

\subsection{\textcolor{Chapter }{IsPrefixOfTipInTipIdeal}}
\logpage{[ 5, 3, 5 ]}\nobreak
\hyperdef{L}{X8137C99A7934C1CA}{}
{\noindent\textcolor{FuncColor}{$\triangleright$\enspace\texttt{IsPrefixOfTipInTipIdeal({\mdseries\slshape GB, R})\index{IsPrefixOfTipInTipIdeal@\texttt{IsPrefixOfTipInTipIdeal}}
\label{IsPrefixOfTipInTipIdeal}
}\hfill{\scriptsize (operation)}}\\


 Arguments: \mbox{\texttt{\mdseries\slshape GB}} \texttt{\symbol{45}}\texttt{\symbol{45}} a Groebner basis, \mbox{\texttt{\mdseries\slshape R}} \texttt{\symbol{45}}\texttt{\symbol{45}} a relation.\\
 \textbf{\indent Returns:\ }
true if the tip of the relation \mbox{\texttt{\mdseries\slshape R}} is in the tip ideal generated by the tips of \mbox{\texttt{\mdseries\slshape GB}}. 



 This is used mainly for the construction of right Groebner basis, but is made
available for general use in case there are other unforeseen applications. }

 

\subsection{\textcolor{Chapter }{Iterator}}
\logpage{[ 5, 3, 6 ]}\nobreak
\hyperdef{L}{X83ADF8287ED0668E}{}
{\noindent\textcolor{FuncColor}{$\triangleright$\enspace\texttt{Iterator({\mdseries\slshape GB})\index{Iterator@\texttt{Iterator}}
\label{Iterator}
}\hfill{\scriptsize (operation)}}\\


 Arguments: \mbox{\texttt{\mdseries\slshape GB}} \texttt{\symbol{45}}\texttt{\symbol{45}} a Groebner basis.\\
 \textbf{\indent Returns:\ }
an iterator (in the IsIterator category, see the \textsf{GAP} manual, chapter 28.7). 



 Creates an iterator that iterates over the relations making up the Groebner
basis. These relations are iterated over in ascending order with respect to
the ordering for the family the elements are contained in. }

 

\subsection{\textcolor{Chapter }{Nontips}}
\logpage{[ 5, 3, 7 ]}\nobreak
\hyperdef{L}{X7EAA029F8071ACC6}{}
{\noindent\textcolor{FuncColor}{$\triangleright$\enspace\texttt{Nontips({\mdseries\slshape GB})\index{Nontips@\texttt{Nontips}}
\label{Nontips}
}\hfill{\scriptsize (attribute)}}\\


 Arguments: \mbox{\texttt{\mdseries\slshape GB}} \texttt{\symbol{45}}\texttt{\symbol{45}} a Groebner basis.\\
 \textbf{\indent Returns:\ }
a list of nontip elements for \mbox{\texttt{\mdseries\slshape GB}}. 



 In order to compute the nontip elements, the Groebner basis must be complete
and tip reduced, and there must be a finite number of nontips. If there are an
infinite number of nontips, the operation returns `fail'. }

 

\subsection{\textcolor{Chapter }{NontipSize}}
\logpage{[ 5, 3, 8 ]}\nobreak
\hyperdef{L}{X7840F54D8240C288}{}
{\noindent\textcolor{FuncColor}{$\triangleright$\enspace\texttt{NontipSize({\mdseries\slshape GB})\index{NontipSize@\texttt{NontipSize}}
\label{NontipSize}
}\hfill{\scriptsize (operation)}}\\


 Arguments: \mbox{\texttt{\mdseries\slshape GB}} \texttt{\symbol{45}}\texttt{\symbol{45}} a complete Groebner basis.\\
 \textbf{\indent Returns:\ }
the number of nontips admitted by \mbox{\texttt{\mdseries\slshape GB}}. 

}

 

\subsection{\textcolor{Chapter }{TipReduce}}
\logpage{[ 5, 3, 9 ]}\nobreak
\hyperdef{L}{X7C6CD739788E7F59}{}
{\noindent\textcolor{FuncColor}{$\triangleright$\enspace\texttt{TipReduce({\mdseries\slshape GB, a})\index{TipReduce@\texttt{TipReduce}}
\label{TipReduce}
}\hfill{\scriptsize (operation)}}\\


 Arguments: \mbox{\texttt{\mdseries\slshape GB}} \texttt{\symbol{45}}\texttt{\symbol{45}} a Groebner basis, \mbox{\texttt{\mdseries\slshape a}} \texttt{\symbol{45}} an element in a path algebra. \\
 \textbf{\indent Returns:\ }
the element \mbox{\texttt{\mdseries\slshape a}} tip reduced by the Groebner basis. 



If \mbox{\texttt{\mdseries\slshape a}} is already tip reduced, then the original \mbox{\texttt{\mdseries\slshape a}} is returned. }

 

\subsection{\textcolor{Chapter }{TipReduceGroebnerBasis}}
\logpage{[ 5, 3, 10 ]}\nobreak
\hyperdef{L}{X7B4F38D6852DF8B6}{}
{\noindent\textcolor{FuncColor}{$\triangleright$\enspace\texttt{TipReduceGroebnerBasis({\mdseries\slshape GB})\index{TipReduceGroebnerBasis@\texttt{TipReduceGroebnerBasis}}
\label{TipReduceGroebnerBasis}
}\hfill{\scriptsize (operation)}}\\


 Arguments: \mbox{\texttt{\mdseries\slshape GB}} \texttt{\symbol{45}}\texttt{\symbol{45}} a Groebner basis.\\
 \textbf{\indent Returns:\ }
a tip reduced Groebner basis. 



 The returned Groebner basis is equivalent to \mbox{\texttt{\mdseries\slshape GB}} If \mbox{\texttt{\mdseries\slshape GB}} is already tip reduced, this function returns the original object \mbox{\texttt{\mdseries\slshape GB}}, possibly with the addition of the `IsTipReduced`' property set. }

 }

 
\section{\textcolor{Chapter }{Right Groebner Basis}}\logpage{[ 5, 4, 0 ]}
\hyperdef{L}{X82C1C09486934532}{}
{
 In this section we support right Groebner basis for
two\texttt{\symbol{45}}sided ideals with Groebner basis. More general cases
may be supported in the future. 

\subsection{\textcolor{Chapter }{IsRightGroebnerBasis}}
\logpage{[ 5, 4, 1 ]}\nobreak
\hyperdef{L}{X86EC39527F33EABE}{}
{\noindent\textcolor{FuncColor}{$\triangleright$\enspace\texttt{IsRightGroebnerBasis({\mdseries\slshape object})\index{IsRightGroebnerBasis@\texttt{IsRightGroebnerBasis}}
\label{IsRightGroebnerBasis}
}\hfill{\scriptsize (property)}}\\


 Arguments: \mbox{\texttt{\mdseries\slshape object}} \texttt{\symbol{45}}\texttt{\symbol{45}} any object in \textsf{GAP}.\\
 \textbf{\indent Returns:\ }
true when \mbox{\texttt{\mdseries\slshape object}} is a right Groebner basis. 

}

 

\subsection{\textcolor{Chapter }{RightGroebnerBasis}}
\logpage{[ 5, 4, 2 ]}\nobreak
\hyperdef{L}{X7B29B9207D20EA9E}{}
{\noindent\textcolor{FuncColor}{$\triangleright$\enspace\texttt{RightGroebnerBasis({\mdseries\slshape I})\index{RightGroebnerBasis@\texttt{RightGroebnerBasis}}
\label{RightGroebnerBasis}
}\hfill{\scriptsize (operation)}}\\


 Arguments: \mbox{\texttt{\mdseries\slshape I}} \texttt{\symbol{45}}\texttt{\symbol{45}} a right ideal.\\
 \textbf{\indent Returns:\ }
a right Groebner basis for \mbox{\texttt{\mdseries\slshape I}}, which must support a right Groebner basis theory. Right now, this requires
that \mbox{\texttt{\mdseries\slshape I}} has a complete Groebner basis. 

}

 

\subsection{\textcolor{Chapter }{RightGroebnerBasisOfIdeal}}
\logpage{[ 5, 4, 3 ]}\nobreak
\hyperdef{L}{X812BFF79867FF73A}{}
{\noindent\textcolor{FuncColor}{$\triangleright$\enspace\texttt{RightGroebnerBasisOfIdeal({\mdseries\slshape I})\index{RightGroebnerBasisOfIdeal@\texttt{RightGroebnerBasisOfIdeal}}
\label{RightGroebnerBasisOfIdeal}
}\hfill{\scriptsize (attribute)}}\\


 Arguments: \mbox{\texttt{\mdseries\slshape I}} \texttt{\symbol{45}}\texttt{\symbol{45}} a right ideal.\\
 \textbf{\indent Returns:\ }
a right Groebner basis of a right ideal, \mbox{\texttt{\mdseries\slshape I}}, if one has been computed. 

}

 }

 
\section{\textcolor{Chapter }{Constructing a Groebner Basis for inverse length ordeing}}\logpage{[ 5, 5, 0 ]}
\hyperdef{L}{X79E642237BD3931E}{}
{
 

\subsection{\textcolor{Chapter }{TipInvLenInvLex}}
\logpage{[ 5, 5, 1 ]}\nobreak
\hyperdef{L}{X7D04F1807CA77001}{}
{\noindent\textcolor{FuncColor}{$\triangleright$\enspace\texttt{TipInvLenInvLex({\mdseries\slshape e})\index{TipInvLenInvLex@\texttt{TipInvLenInvLex}}
\label{TipInvLenInvLex}
}\hfill{\scriptsize (operation)}}\\


 Arguments: \mbox{\texttt{\mdseries\slshape e}} \texttt{\symbol{45}}\texttt{\symbol{45}} an element of a path algebra.\\
 \textbf{\indent Returns:\ }
the tip of the element \mbox{\texttt{\mdseries\slshape e}} under the inverse length inverse lexicographic ordering. 

}

 

\subsection{\textcolor{Chapter }{TipCoefficientInvLenInvLex}}
\logpage{[ 5, 5, 2 ]}\nobreak
\hyperdef{L}{X7E8CDBBF8029A9AA}{}
{\noindent\textcolor{FuncColor}{$\triangleright$\enspace\texttt{TipCoefficientInvLenInvLex({\mdseries\slshape e})\index{TipCoefficientInvLenInvLex@\texttt{TipCoefficientInvLenInvLex}}
\label{TipCoefficientInvLenInvLex}
}\hfill{\scriptsize (operation)}}\\


 Arguments: \mbox{\texttt{\mdseries\slshape e}} \texttt{\symbol{45}}\texttt{\symbol{45}} an element of a path algebra.\\
 \textbf{\indent Returns:\ }
the coefficient of the tip of the element \mbox{\texttt{\mdseries\slshape e}} under the inverse length inverse lexicographic ordering. 

}

 

\subsection{\textcolor{Chapter }{TipMonomialInvLenInvLex}}
\logpage{[ 5, 5, 3 ]}\nobreak
\hyperdef{L}{X7BFF3AB87D32B76A}{}
{\noindent\textcolor{FuncColor}{$\triangleright$\enspace\texttt{TipMonomialInvLenInvLex({\mdseries\slshape e})\index{TipMonomialInvLenInvLex@\texttt{TipMonomialInvLenInvLex}}
\label{TipMonomialInvLenInvLex}
}\hfill{\scriptsize (operation)}}\\


 Arguments: \mbox{\texttt{\mdseries\slshape e}} \texttt{\symbol{45}}\texttt{\symbol{45}} element of a path algebra.\\
 \textbf{\indent Returns:\ }
the path/monomial which is the tip of the element \mbox{\texttt{\mdseries\slshape e}} under the inverse length inverse lexicographic ordering. 

}

 

\subsection{\textcolor{Chapter }{TipWalkInvLenInvLex}}
\logpage{[ 5, 5, 4 ]}\nobreak
\hyperdef{L}{X8493D6E87F831A7A}{}
{\noindent\textcolor{FuncColor}{$\triangleright$\enspace\texttt{TipWalkInvLenInvLex({\mdseries\slshape e})\index{TipWalkInvLenInvLex@\texttt{TipWalkInvLenInvLex}}
\label{TipWalkInvLenInvLex}
}\hfill{\scriptsize (operation)}}\\


 Arguments: \mbox{\texttt{\mdseries\slshape e}} \texttt{\symbol{45}}\texttt{\symbol{45}} element of a path algebra.\\
 \textbf{\indent Returns:\ }
the path/monomial which is the tip of the element \mbox{\texttt{\mdseries\slshape e}} under the inverse length inverse lexicographic ordering as a list of arrows in
the underlying quiver. 

}

 

\subsection{\textcolor{Chapter }{GroebnerBasisInvLenInvLex}}
\logpage{[ 5, 5, 5 ]}\nobreak
\hyperdef{L}{X80EA7F088416E2D4}{}
{\noindent\textcolor{FuncColor}{$\triangleright$\enspace\texttt{GroebnerBasisInvLenInvLex({\mdseries\slshape els, A})\index{GroebnerBasisInvLenInvLex@\texttt{GroebnerBasisInvLenInvLex}}
\label{GroebnerBasisInvLenInvLex}
}\hfill{\scriptsize (operation)}}\\


 Arguments: \mbox{\texttt{\mdseries\slshape els}}, \mbox{\texttt{\mdseries\slshape A}} \texttt{\symbol{45}}\texttt{\symbol{45}} a list of elements in an algebra and
an algebra.\\
 \textbf{\indent Returns:\ }
a Groebner basis for an admissible ideal with generators \mbox{\texttt{\mdseries\slshape els}} in a path algebra \mbox{\texttt{\mdseries\slshape A}}. 

}

 

\subsection{\textcolor{Chapter }{RemainderOfDivisionInvLenInvLexGroebnerBasis}}
\logpage{[ 5, 5, 6 ]}\nobreak
\hyperdef{L}{X857EAAD47D52A736}{}
{\noindent\textcolor{FuncColor}{$\triangleright$\enspace\texttt{RemainderOfDivisionInvLenInvLexGroebnerBasis({\mdseries\slshape y, X, A})\index{RemainderOfDivisionInvLenInvLexGroebnerBasis@\texttt{Remainder}\-\texttt{Of}\-\texttt{Division}\-\texttt{Inv}\-\texttt{Len}\-\texttt{Inv}\-\texttt{Lex}\-\texttt{Groebner}\-\texttt{Basis}}
\label{RemainderOfDivisionInvLenInvLexGroebnerBasis}
}\hfill{\scriptsize (operation)}}\\


 Arguments: \mbox{\texttt{\mdseries\slshape y}}, \mbox{\texttt{\mdseries\slshape X}}, \mbox{\texttt{\mdseries\slshape A}}\texttt{\symbol{45}}\texttt{\symbol{45}} an element and a list of elements in
a path algebra, and a path algebra.\\
 \textbf{\indent Returns:\ }
the remainder of reducing \mbox{\texttt{\mdseries\slshape y}} modulo the list of elements \mbox{\texttt{\mdseries\slshape X}} in the path algebra \mbox{\texttt{\mdseries\slshape A}}. 

}

 

\subsection{\textcolor{Chapter }{ReducedListInvLenInvLexQPA}}
\logpage{[ 5, 5, 7 ]}\nobreak
\hyperdef{L}{X8390998285C727E1}{}
{\noindent\textcolor{FuncColor}{$\triangleright$\enspace\texttt{ReducedListInvLenInvLexQPA({\mdseries\slshape })\index{ReducedListInvLenInvLexQPA@\texttt{ReducedListInvLenInvLexQPA}}
\label{ReducedListInvLenInvLexQPA}
}\hfill{\scriptsize (operation)}}\\


 Arguments: \mbox{\texttt{\mdseries\slshape els}}, \mbox{\texttt{\mdseries\slshape A}} \texttt{\symbol{45}}\texttt{\symbol{45}} a list of elements in a path algebra
and a path algebra.\\
 \textbf{\indent Returns:\ }
a reduced list of the elements \mbox{\texttt{\mdseries\slshape els}} by computing the remainder of an element on the list by the previous elements
on the list by first starting with the last element on the list and moving
forward on the list. 

}

 

\subsection{\textcolor{Chapter }{TipReducedListInvLenInvLex}}
\logpage{[ 5, 5, 8 ]}\nobreak
\hyperdef{L}{X7D96FEDA7F69B266}{}
{\noindent\textcolor{FuncColor}{$\triangleright$\enspace\texttt{TipReducedListInvLenInvLex({\mdseries\slshape els, A})\index{TipReducedListInvLenInvLex@\texttt{TipReducedListInvLenInvLex}}
\label{TipReducedListInvLenInvLex}
}\hfill{\scriptsize (operation)}}\\


 Arguments: \mbox{\texttt{\mdseries\slshape els}}, \mbox{\texttt{\mdseries\slshape A}} \texttt{\symbol{45}}\texttt{\symbol{45}} a list of elements in a path algebra
and a path algebra.\\
 \textbf{\indent Returns:\ }
a reduced list of which generate the same ideal of tips. 

}

 

\subsection{\textcolor{Chapter }{LeftmostOccurrenceInvLenInvLex}}
\logpage{[ 5, 5, 9 ]}\nobreak
\hyperdef{L}{X7A6E562C8414E151}{}
{\noindent\textcolor{FuncColor}{$\triangleright$\enspace\texttt{LeftmostOccurrenceInvLenInvLex({\mdseries\slshape b, c})\index{LeftmostOccurrenceInvLenInvLex@\texttt{LeftmostOccurrenceInvLenInvLex}}
\label{LeftmostOccurrenceInvLenInvLex}
}\hfill{\scriptsize (operation)}}\\


 Arguments: \mbox{\texttt{\mdseries\slshape b}}, \mbox{\texttt{\mdseries\slshape c}} \texttt{\symbol{45}}\texttt{\symbol{45}} two list of elements.\\
 \textbf{\indent Returns:\ }
the leftmost occurrence of the list \mbox{\texttt{\mdseries\slshape c}} in the list \mbox{\texttt{\mdseries\slshape b}} if it exists, otherwise it returns fail. 

}

 

\subsection{\textcolor{Chapter }{BiggestHomogeneousPartInvLenInvLex}}
\logpage{[ 5, 5, 10 ]}\nobreak
\hyperdef{L}{X84B52C118749406D}{}
{\noindent\textcolor{FuncColor}{$\triangleright$\enspace\texttt{BiggestHomogeneousPartInvLenInvLex({\mdseries\slshape e, or, elems})\index{BiggestHomogeneousPartInvLenInvLex@\texttt{BiggestHomogeneousPartInvLenInvLex}}
\label{BiggestHomogeneousPartInvLenInvLex}
}\hfill{\scriptsize (operation)}}\\


 Arguments: (first version) \mbox{\texttt{\mdseries\slshape e}} \texttt{\symbol{45}}\texttt{\symbol{45}} an
IsElementOfMagmaRingModuloRelations.\\
 (second version) \mbox{\texttt{\mdseries\slshape elems}} \texttt{\symbol{45}}\texttt{\symbol{45}} a list of elements in a path algebra.\\
 \textbf{\indent Returns:\ }
(first version) the biggest homogeneous part of the element \mbox{\texttt{\mdseries\slshape e}} under the inverse length inverse lexicographic ordering. 

\textbf{\indent Returns:\ }
(second version) the list of the biggest homogeneous part of each element in
the list \mbox{\texttt{\mdseries\slshape elems}} under the inverse length inverse lexicographic ordering. 

}

 

\subsection{\textcolor{Chapter }{TruncateElement}}
\logpage{[ 5, 5, 11 ]}\nobreak
\hyperdef{L}{X7CE1DF847FDA0162}{}
{\noindent\textcolor{FuncColor}{$\triangleright$\enspace\texttt{TruncateElement({\mdseries\slshape elem, s})\index{TruncateElement@\texttt{TruncateElement}}
\label{TruncateElement}
}\hfill{\scriptsize (operation)}}\\


 Arguments: \mbox{\texttt{\mdseries\slshape elem}}, \mbox{\texttt{\mdseries\slshape s}} \texttt{\symbol{45}}\texttt{\symbol{45}} an element in a path algebra and a
positive integer.\\
 \textbf{\indent Returns:\ }
the part of the element \mbox{\texttt{\mdseries\slshape elem}} which consists of the paths of length less than \mbox{\texttt{\mdseries\slshape s}}. 

}

 }

 }

 
\chapter{\textcolor{Chapter }{Right Modules over Path Algebras}}\label{Right-Modules}
\logpage{[ 6, 0, 0 ]}
\hyperdef{L}{X87EFC38F7BC77B27}{}
{
 There are two implementations of right modules over path algebras. The first
type are matrix modules that are defined by vector spaces and linear
transformations. The second type is presentations defined by vertex projective
modules (see \ref{vertexproj}). 
\section{\textcolor{Chapter }{Modules of matrix type}}\logpage{[ 6, 1, 0 ]}
\hyperdef{L}{X86DD15DF877834FD}{}
{
 The first implementation of right modules over path algebras views them as a
collection of vector spaces and linear transformations. Each vertex in the
path algebra is associated with a vector space over the field of the algebra.
For each vertex $v$ of the algebra there is a vector space $V$. Arrows of the algebra are then associated with linear transformations which
map the vector space of the source vertex to the vector space of the target
vertex. For example, if $a$ is an arrow from $v$ to $w$, then there is a transformation from vector space $V$ to $W$. Given the dimension vector of the module we want to construct, the
information we need to provide is the non\texttt{\symbol{45}}zero linear
transformations. The size of the matrices for the zero linear transformation
are given when we know the dimension vector. Alternatively, if we enter all
the transformations, we can create the vector spaces of the correct dimension,
and check to make sure the dimensions all agree. We can create a module in
this way as follows.

 

\subsection{\textcolor{Chapter }{RightModuleOverPathAlgebra (with dimension vector)}}
\logpage{[ 6, 1, 1 ]}\nobreak
\hyperdef{L}{X85E5097D82D9BE62}{}
{\noindent\textcolor{FuncColor}{$\triangleright$\enspace\texttt{RightModuleOverPathAlgebra({\mdseries\slshape A, dim{\textunderscore}vector, gens})\index{RightModuleOverPathAlgebra@\texttt{RightModuleOverPathAlgebra}!with dimension vector}
\label{RightModuleOverPathAlgebra:with dimension vector}
}\hfill{\scriptsize (operation)}}\\
\noindent\textcolor{FuncColor}{$\triangleright$\enspace\texttt{RightModuleOverPathAlgebra({\mdseries\slshape A, mats})\index{RightModuleOverPathAlgebra@\texttt{RightModuleOverPathAlgebra}!no dimension vector}
\label{RightModuleOverPathAlgebra:no dimension vector}
}\hfill{\scriptsize (operation)}}\\
\noindent\textcolor{FuncColor}{$\triangleright$\enspace\texttt{RightModuleOverPathAlgebraNC({\mdseries\slshape A, mats})\index{RightModuleOverPathAlgebraNC@\texttt{RightModuleOverPathAlgebraNC}!no dimension vector}
\label{RightModuleOverPathAlgebraNC:no dimension vector}
}\hfill{\scriptsize (operation)}}\\


 Arguments: \mbox{\texttt{\mdseries\slshape A}} \texttt{\symbol{45}}\texttt{\symbol{45}} a (quotient of a) path algebra and \mbox{\texttt{\mdseries\slshape dim{\textunderscore}vector}} \texttt{\symbol{45}}\texttt{\symbol{45}} the dimension vector of the module, \mbox{\texttt{\mdseries\slshape gens}} or \mbox{\texttt{\mdseries\slshape mats}} \texttt{\symbol{45}}\texttt{\symbol{45}} a list of matrices. For further
explanations, see below. \\
\textbf{\indent Returns:\ }
a module over a path algebra or over a qoutient of a path algebra.



 In the first function call, the second argument \mbox{\texttt{\mdseries\slshape dim{\textunderscore}vector}} is the dimension vector of the module, and the last argument \mbox{\texttt{\mdseries\slshape gens}} (maybe an empty list \texttt{[]}) is a list of elements of the form \texttt{["label",matrix]}. This function constructs a right module over a (quotient of a) path algebra \mbox{\texttt{\mdseries\slshape A}} with dimension vector \mbox{\texttt{\mdseries\slshape dim{\textunderscore}vector}}, and where the generators/arrows with a non\texttt{\symbol{45}}zero action is
given in the list \mbox{\texttt{\mdseries\slshape gens}}. The format of the list \mbox{\texttt{\mdseries\slshape gens}} is [["a",[matrix{\textunderscore}a]],["b",[matrix{\textunderscore}b]],...],
where "a" and "b" are labels of arrows used when the underlying quiver was
created and matrix{\textunderscore}? is the action of the algebra element
corresponding to the arrow with label "?". The action of the arrows can be
entered in any order. The function checks (i) if the algebra \mbox{\texttt{\mdseries\slshape A}} is a (quotient of a) path algebra, (ii) if the matrices of the action of the
arrows have the correct size according to the dimension vector entered, (iii)
also whether or not the relations of the algebra are satisfied and (iv) if all
matrices are over the correct field.

 In the second function call, the list of matrices \mbox{\texttt{\mdseries\slshape mats}} can take on three different forms. The function checks (i), (ii), (iii) and
(iv) as above.

 1) The argument \mbox{\texttt{\mdseries\slshape mats}} can be a list of blocks of matrices where each block is of the form, `["name
of arrow",matrix]'. So if you named your arrows when you created the quiver,
then you can associate a matrix with that arrow explicitly.

 2) The argument \mbox{\texttt{\mdseries\slshape mats}} is just a list of matrices, and the matrices will be associated to the arrows
in the order of arrow creation. If when creating the quiver, the arrow $a$ was created first, then $a$ would be associated with the first matrix.

 3) The method is very much the same as the second method. If \mbox{\texttt{\mdseries\slshape arrows}} is a list of the arrows of the quiver (obtained for instance through \texttt{arrows := ArrowsOfQuiver(Q);}), the argument \mbox{\texttt{\mdseries\slshape mats}} can have the format \texttt{[[arrows[1],matrix{\textunderscore}1],[arrows[2],matrix{\textunderscore}2],....
].}

 If you would like the trivial vector space at any vertex, then for each
incoming arrow "a", associate it with a list of the form \texttt{["a",[n,0]]} where n is the dimension of the vector space at the source vertex of the
arrow. Likewise for all outgoing arrows "b", associate them to a block of form \texttt{["b",[0,n]]} where n is the dimension of the vector space at the target vertex of the
arrow.

 The third function call is the same as the second except that the check (iv)
is not performed.

 A warning though, the function assumes that you do not mix the styles of
inputting the matrices/linear transformations associated to the arrows in the
quiver. Furthermore in the two last versions, each arrow needs to be assigned
a matrix, otherwise an error will be returned.

 }

 

\subsection{\textcolor{Chapter }{RightAlgebraModuleToPathAlgebraMatModule}}
\logpage{[ 6, 1, 2 ]}\nobreak
\hyperdef{L}{X844DD9B386A6AC56}{}
{\noindent\textcolor{FuncColor}{$\triangleright$\enspace\texttt{RightAlgebraModuleToPathAlgebraMatModule({\mdseries\slshape M})\index{RightAlgebraModuleToPathAlgebraMatModule@\texttt{Right}\-\texttt{Algebra}\-\texttt{Module}\-\texttt{To}\-\texttt{Path}\-\texttt{Algebra}\-\texttt{Mat}\-\texttt{Module}}
\label{RightAlgebraModuleToPathAlgebraMatModule}
}\hfill{\scriptsize (operation)}}\\


 Arguments: \mbox{\texttt{\mdseries\slshape M}} \texttt{\symbol{45}}\texttt{\symbol{45}} a right module over an algebra. \\
\textbf{\indent Returns:\ }
a module over a (qoutient of a) path algebra.



 This function constructs a right module over a (quotient of a) path algebra $A$ from a RightAlgebraModule over the same algebra $A$. The function checks if $A$ actually is a quotient of a path algebra and if the module $M$ is finite dimensional and if not, it returns an error message. }

 

\subsection{\textcolor{Chapter }{\texttt{\symbol{92}}= (for two path algebra matrix modules)}}
\logpage{[ 6, 1, 3 ]}\nobreak
\hyperdef{L}{X823D5739809A2D9A}{}
{\noindent\textcolor{FuncColor}{$\triangleright$\enspace\texttt{\texttt{\symbol{92}}=({\mdseries\slshape M, N})\index{\texttt{\symbol{92}}=@\texttt{\texttt{\symbol{92}}=}!for two path algebra matrix modules}
\label{bSlash=:for two path algebra matrix modules}
}\hfill{\scriptsize (operation)}}\\


 Arguments: \mbox{\texttt{\mdseries\slshape M, N}} \texttt{\symbol{45}}\texttt{\symbol{45}} two path algebra matrix modules. \\
\textbf{\indent Returns:\ }
true if \mbox{\texttt{\mdseries\slshape M}} and \mbox{\texttt{\mdseries\slshape N}} has the same dimension vectors and the same matrices defining the module
structure.

}

 
\begin{Verbatim}[commandchars=!@|,fontsize=\small,frame=single,label=Example]
  !gapprompt@gap>| !gapinput@Q := Quiver(2, [[1, 2, "a"], [2, 1, "b"],[1, 1, "c"]]);|
  <quiver with 2 vertices and 3 arrows>
  !gapprompt@gap>| !gapinput@P := PathAlgebra(Rationals, Q);|
  <Rationals[<quiver with 2 vertices and 3 arrows>]>
  !gapprompt@gap>| !gapinput@matrices := [["a", [[1,0,0],[0,1,0]]], |
  !gapprompt@>| !gapinput@ ["b", [[0,1],[1,0],[0,1]]],|
  !gapprompt@>| !gapinput@ ["c", [[0,0],[1,0]]]];|
  [ [ "a", [ [ 1, 0, 0 ], [ 0, 1, 0 ] ] ], 
    [ "b", [ [ 0, 1 ], [ 1, 0 ], [ 0, 1 ] ] ], 
    [ "c", [ [ 0, 0 ], [ 1, 0 ] ] ] ]
  !gapprompt@gap>| !gapinput@M := RightModuleOverPathAlgebra(P,matrices);|
  <[ 2, 3 ]>
  !gapprompt@gap>| !gapinput@mats := [ [[1,0,0], [0,1,0]], [[0,1],[1,0],[0,1]], |
  !gapprompt@>| !gapinput@          [[0,0],[1,0]] ];; |
  !gapprompt@gap>| !gapinput@N := RightModuleOverPathAlgebra(P,mats); |
  <[ 2, 3 ]>
  !gapprompt@gap>| !gapinput@arrows := ArrowsOfQuiver(Q);|
  [ a, b, c ]
  !gapprompt@gap>| !gapinput@mats := [[arrows[1], [[1,0,0],[0,1,0]]],|
  !gapprompt@>| !gapinput@         [arrows[2], [[0,1],[1,0],[0,1]]], |
  !gapprompt@>| !gapinput@         [arrows[3], [[0,0],[1,0]]]];;|
  !gapprompt@gap>| !gapinput@N := RightModuleOverPathAlgebra(P,mats); |
  <[ 2, 3 ]>
  !gapprompt@gap>| !gapinput@# Next we give the vertex simple associate to vertex 1. |
  !gapprompt@gap>| !gapinput@M := RightModuleOverPathAlgebra(P,[["a",[1,0]],["b",[0,1]],|
  !gapprompt@>| !gapinput@             ["c",[[0]]]]);|
  <[ 1, 0 ]>
  !gapprompt@gap>| !gapinput@# The zero module. |
  !gapprompt@gap>| !gapinput@M := RightModuleOverPathAlgebra(P,[["a",[0,0]],["b",[0,0]],|
  !gapprompt@>| !gapinput@             ["c",[0,0]]]);|
  <[ 0, 0 ]>
  !gapprompt@gap>| !gapinput@Dimension(M);|
  0
  !gapprompt@gap>| !gapinput@Basis(M);|
  Basis( <[ 0, 0 ]>, ... )
  !gapprompt@gap>| !gapinput@matrices := [["a", [[1,0,0],[0,1,0]]], ["b",|
  !gapprompt@>| !gapinput@ [[0,1],[1,0],[0,1]]], ["c", [[0,0],[1,0]]]];|
  [ [ "a", [ [ 1, 0, 0 ], [ 0, 1, 0 ] ] ], 
    [ "b", [ [ 0, 1 ], [ 1, 0 ], [ 0, 1 ] ] ], 
    [ "c", [ [ 0, 0 ], [ 1, 0 ] ] ] ]
  !gapprompt@gap>| !gapinput@M := RightModuleOverPathAlgebra(P,[2,3],matrices);|
  <[ 2, 3 ]>
  !gapprompt@gap>| !gapinput@M := RightModuleOverPathAlgebra(P,[2,3],[]);  |
  <[ 2, 3 ]>
  !gapprompt@gap>| !gapinput@A := P/[P.c^2 - P.a*P.b, P.a*P.b*P.c, P.b*P.c];         |
  <Rationals[<quiver with 2 vertices and 3 arrows>]/
  <two-sided ideal in <Rationals[<quiver with 2 vertices and 3 arrows>]>
      , (4 generators)>>
  !gapprompt@gap>| !gapinput@Dimension(A);|
  9
  !gapprompt@gap>| !gapinput@Amod := RightAlgebraModule(A,\*,A);                       |
  <9-dimensional right-module over <Rationals[<quiver with 
  2 vertices and 3 arrows>]/
  <two-sided ideal in <Rationals[<quiver with 2 vertices and 3 arrows>]>
      , (4 generators)>>>
  !gapprompt@gap>| !gapinput@RightAlgebraModuleToPathAlgebraMatModule(Amod);|
  <[ 4, 5 ]> 
\end{Verbatim}
 }

 
\section{\textcolor{Chapter }{Categories Of Matrix Modules}}\logpage{[ 6, 2, 0 ]}
\hyperdef{L}{X869F4DD2877A99BA}{}
{
 

\subsection{\textcolor{Chapter }{IsPathAlgebraMatModule}}
\logpage{[ 6, 2, 1 ]}\nobreak
\hyperdef{L}{X8710CC447F1F7B17}{}
{\noindent\textcolor{FuncColor}{$\triangleright$\enspace\texttt{IsPathAlgebraMatModule({\mdseries\slshape object})\index{IsPathAlgebraMatModule@\texttt{IsPathAlgebraMatModule}}
\label{IsPathAlgebraMatModule}
}\hfill{\scriptsize (filter)}}\\
\textbf{\indent Returns:\ }
 true or false depending on whether \mbox{\texttt{\mdseries\slshape object}} belongs to the category \texttt{IsPathAlgebraMatModule}.



 These matrix modules fall under the category `IsAlgebraModule' with the added
filter of `IsPathAlgebraMatModule'. Operations available for algebra modules
can be applied to path algebra modules. See the main
GAP\texttt{\symbol{45}}manual under "representations of algebras" for more
details. These modules are also vector spaces over the field of the path
algebra. So refer to the main GAP\texttt{\symbol{45}}manual under "vector
spaces" for descriptions of the basis and elementwise operations available. }

 }

 
\section{\textcolor{Chapter }{Acting on Module Elements}}\logpage{[ 6, 3, 0 ]}
\hyperdef{L}{X862F510485ADBC67}{}
{
 

\subsection{\textcolor{Chapter }{\texttt{\symbol{94}} (a PathAlgebraMatModule element and a PathAlgebra element)}}
\logpage{[ 6, 3, 1 ]}\nobreak
\hyperdef{L}{X7BB4066E7D5B15B8}{}
{\noindent\textcolor{FuncColor}{$\triangleright$\enspace\texttt{\texttt{\symbol{94}}({\mdseries\slshape m, p})\index{\texttt{\symbol{94}}@\texttt{\texttt{\symbol{94}}}!a PathAlgebraMatModule element and a PathAlgebra element}
\label{^:a PathAlgebraMatModule element and a PathAlgebra element}
}\hfill{\scriptsize (operation)}}\\


 Arguments: \mbox{\texttt{\mdseries\slshape m}} \texttt{\symbol{45}}\texttt{\symbol{45}} an element in a module, \mbox{\texttt{\mdseries\slshape p}} \texttt{\symbol{45}}\texttt{\symbol{45}} an element in a quiver algebra. \\
\textbf{\indent Returns:\ }
 the element \mbox{\texttt{\mdseries\slshape m}} multiplied with \mbox{\texttt{\mdseries\slshape p}}. 



 When you act on an module element $m$ by an arrow $a$ from $v$ to $w$, the component of $m$ from $V$ is acted on by $L$ the transformation associated to $a$ and placed in the component $W$. All other components are given the value $0$. }

 
\begin{Verbatim}[commandchars=!@|,fontsize=\small,frame=single,label=Example]
  !gapprompt@gap>| !gapinput@# Using the path algebra P from the above example. |
  !gapprompt@gap>| !gapinput@matrices := [["a", [[1,0,0],[0,1,0]]], |
  !gapprompt@>| !gapinput@["b", [[0,1],[1,0],[0,1]]], ["c", [[0,0],[1,0]]]];|
  [ [ "a", [ [ 1, 0, 0 ], [ 0, 1, 0 ] ] ], 
    [ "b", [ [ 0, 1 ], [ 1, 0 ], [ 0, 1 ] ] ], 
    [ "c", [ [ 0, 0 ], [ 1, 0 ] ] ] ]
  !gapprompt@gap>| !gapinput@M := RightModuleOverPathAlgebra(P,matrices);|
  <[ 2, 3 ]>
  !gapprompt@gap>| !gapinput@B:=BasisVectors(Basis(M));|
  [ [ [ 1, 0 ], [ 0, 0, 0 ] ], [ [ 0, 1 ], [ 0, 0, 0 ] ], 
    [ [ 0, 0 ], [ 1, 0, 0 ] ], [ [ 0, 0 ], [ 0, 1, 0 ] ], 
    [ [ 0, 0 ], [ 0, 0, 1 ] ] ]
  !gapprompt@gap>| !gapinput@B[1] + B[3];|
  [ [ 1, 0 ], [ 1, 0, 0 ] ]
  !gapprompt@gap>| !gapinput@4*B[2];|
  [ [ 0, 4 ], [ 0, 0, 0 ] ]
  !gapprompt@gap>| !gapinput@m := 5*B[1] + 2*B[4]+B[5];|
  [ [ 5, 0 ], [ 0, 2, 1 ] ]
  !gapprompt@gap>| !gapinput@m^(P.a*P.b-P.c);|
  [ [ 0, 5 ], [ 0, 0, 0 ] ]
  !gapprompt@gap>| !gapinput@B[1]^P.a;|
  [ [ 0, 0 ], [ 1, 0, 0 ] ]
  !gapprompt@gap>| !gapinput@B[2]^P.b;|
  [ [ 0, 0 ], [ 0, 0, 0 ] ]
  !gapprompt@gap>| !gapinput@B[4]^(P.b*P.c);|
  [ [ 0, 0 ], [ 0, 0, 0 ] ] 
\end{Verbatim}
 }

 
\section{\textcolor{Chapter }{Operations on representations}}\logpage{[ 6, 4, 0 ]}
\hyperdef{L}{X7E5B84B1832D839E}{}
{
 
\begin{Verbatim}[commandchars=!@|,fontsize=\small,frame=single,label=Example]
  !gapprompt@gap>| !gapinput@Q  := Quiver(3,[[1,2,"a"],[1,2,"b"],[2,2,"c"],[2,3,"d"],|
  !gapprompt@>| !gapinput@[3,1,"e"]]);|
  <quiver with 3 vertices and 5 arrows>
  !gapprompt@gap>| !gapinput@KQ := PathAlgebra(Rationals, Q);|
  <Rationals[<quiver with 3 vertices and 5 arrows>]>
  !gapprompt@gap>| !gapinput@gens := GeneratorsOfAlgebra(KQ);|
  [ (1)*v1, (1)*v2, (1)*v3, (1)*a, (1)*b, (1)*c, (1)*d, (1)*e ]
  !gapprompt@gap>| !gapinput@u := gens[1];; v := gens[2];;|
  !gapprompt@gap>| !gapinput@w := gens[3];; a := gens[4];;|
  !gapprompt@gap>| !gapinput@b := gens[5];; c := gens[6];;|
  !gapprompt@gap>| !gapinput@d := gens[7];; e := gens[8];;|
  !gapprompt@gap>| !gapinput@rels := [d*e,c^2,a*c*d-b*d,e*a];;|
  !gapprompt@gap>| !gapinput@A := KQ/rels;|
  <Rationals[<quiver with 3 vertices and 5 arrows>]/
  <two-sided ideal in <Rationals[<quiver with 3 vertices and 5 arrows>]>
      , (5 generators)>>
  !gapprompt@gap>| !gapinput@mat := [["a",[[1,2],[0,3],[1,5]]],["b",[[2,0],[3,0],[5,0]]],|
  !gapprompt@>| !gapinput@["c",[[0,0],[1,0]]],["d",[[1,2],[0,1]]],["e",[[0,0,0],[0,0,0]]]];;|
  !gapprompt@gap>| !gapinput@N := RightModuleOverPathAlgebra(A,mat);   |
  <[ 3, 2, 2 ]> 
\end{Verbatim}
 

\subsection{\textcolor{Chapter }{AnnihilatorOfModule}}
\logpage{[ 6, 4, 1 ]}\nobreak
\hyperdef{L}{X830B834D7F2F0FAD}{}
{\noindent\textcolor{FuncColor}{$\triangleright$\enspace\texttt{AnnihilatorOfModule({\mdseries\slshape M})\index{AnnihilatorOfModule@\texttt{AnnihilatorOfModule}}
\label{AnnihilatorOfModule}
}\hfill{\scriptsize (operation)}}\\


 Arguments: \mbox{\texttt{\mdseries\slshape M}} \texttt{\symbol{45}}\texttt{\symbol{45}} a path algebra module. \\
\textbf{\indent Returns:\ }
a basis of the annihilator of the module \mbox{\texttt{\mdseries\slshape M}} in the finite dimensional algebra over which \mbox{\texttt{\mdseries\slshape M}} is a module. 

}

 

\subsection{\textcolor{Chapter }{BasicVersionOfModule}}
\logpage{[ 6, 4, 2 ]}\nobreak
\hyperdef{L}{X83404C0B7C15E7D0}{}
{\noindent\textcolor{FuncColor}{$\triangleright$\enspace\texttt{BasicVersionOfModule({\mdseries\slshape M})\index{BasicVersionOfModule@\texttt{BasicVersionOfModule}}
\label{BasicVersionOfModule}
}\hfill{\scriptsize (operation)}}\\


 Arguments: \mbox{\texttt{\mdseries\slshape M}} \texttt{\symbol{45}}\texttt{\symbol{45}} a path algebra module. \\
\textbf{\indent Returns:\ }
 a basic version of the entered module \mbox{\texttt{\mdseries\slshape M}}, that is, if $M \simeq M_1^{n_1} \oplus \cdots \oplus M_t^{n_t},$ where $M_i$ is indecomposable, then $M_1\oplus \cdots \oplus M_t$ is returned. At present, this function only work at best for finite
dimensional (quotients of a) path algebra over a finite field. If \mbox{\texttt{\mdseries\slshape M}} is zero, then \mbox{\texttt{\mdseries\slshape M}} is returned. 

}

 

\subsection{\textcolor{Chapter }{BlockDecompositionOfModule}}
\logpage{[ 6, 4, 3 ]}\nobreak
\hyperdef{L}{X7CD05524803C7777}{}
{\noindent\textcolor{FuncColor}{$\triangleright$\enspace\texttt{BlockDecompositionOfModule({\mdseries\slshape M})\index{BlockDecompositionOfModule@\texttt{BlockDecompositionOfModule}}
\label{BlockDecompositionOfModule}
}\hfill{\scriptsize (operation)}}\\


 Arguments: \mbox{\texttt{\mdseries\slshape M}} \texttt{\symbol{45}}\texttt{\symbol{45}} a path algebra module. \\
\textbf{\indent Returns:\ }
a set of modules $\{M_1,..., M_t\}$ such that $M \simeq M_1\oplus \cdots \oplus M_t,$ where each $M_i$ is isomorphic to $X_i^{n_i}$ for some indecomposable module $X_i$ and positive integer $n_i$ for all $i$, where $X_i\not\simeq X_j$ for $i\neq j$. 

}

 

\subsection{\textcolor{Chapter }{BlockSplittingIdempotents}}
\logpage{[ 6, 4, 4 ]}\nobreak
\hyperdef{L}{X85AC0E7C7E3E697D}{}
{\noindent\textcolor{FuncColor}{$\triangleright$\enspace\texttt{BlockSplittingIdempotents({\mdseries\slshape M})\index{BlockSplittingIdempotents@\texttt{BlockSplittingIdempotents}}
\label{BlockSplittingIdempotents}
}\hfill{\scriptsize (operation)}}\\


 Arguments: \mbox{\texttt{\mdseries\slshape M}} \texttt{\symbol{45}}\texttt{\symbol{45}} a path algebra module. \\
\textbf{\indent Returns:\ }
a set $\{e_1,..., e_t\}$ of idempotents in the endomorphism of \mbox{\texttt{\mdseries\slshape M}} such that $M \simeq \Im e_1\oplus \cdots \oplus \Im e_t,$ where each $\Im e_i$ is isomorphic to $X_i^{n_i}$ for some module $X_i$ and positive integer $n_i$ for all $i$. 

}

 

\subsection{\textcolor{Chapter }{CommonDirectSummand}}
\logpage{[ 6, 4, 5 ]}\nobreak
\hyperdef{L}{X8687EC4878E755CC}{}
{\noindent\textcolor{FuncColor}{$\triangleright$\enspace\texttt{CommonDirectSummand({\mdseries\slshape M, N})\index{CommonDirectSummand@\texttt{CommonDirectSummand}}
\label{CommonDirectSummand}
}\hfill{\scriptsize (operation)}}\\


 Arguments: \mbox{\texttt{\mdseries\slshape M}} and \mbox{\texttt{\mdseries\slshape N}} \texttt{\symbol{45}}\texttt{\symbol{45}} two path algebra modules. \\
\textbf{\indent Returns:\ }
a list of four modules [\mbox{\texttt{\mdseries\slshape X}},\mbox{\texttt{\mdseries\slshape U}},\mbox{\texttt{\mdseries\slshape X}}, \mbox{\texttt{\mdseries\slshape V}}], where \mbox{\texttt{\mdseries\slshape X}} is one common non\texttt{\symbol{45}}zero direct summand of \mbox{\texttt{\mdseries\slshape M}} and \mbox{\texttt{\mdseries\slshape N}}, the sum of \mbox{\texttt{\mdseries\slshape X}} and \mbox{\texttt{\mdseries\slshape U}} is \mbox{\texttt{\mdseries\slshape M}} and the sum of \mbox{\texttt{\mdseries\slshape X}} and \mbox{\texttt{\mdseries\slshape V}} is \mbox{\texttt{\mdseries\slshape N}}, if such a non\texttt{\symbol{45}}zero direct summand exists. Otherwise it
returns false.



The function checks if \mbox{\texttt{\mdseries\slshape M}} and \mbox{\texttt{\mdseries\slshape N}} are \texttt{PathAlgebraMatModule}s over the same (quotient of a) path algebra. }

 

\subsection{\textcolor{Chapter }{ComplexityOfModule}}
\logpage{[ 6, 4, 6 ]}\nobreak
\hyperdef{L}{X8376806384B96066}{}
{\noindent\textcolor{FuncColor}{$\triangleright$\enspace\texttt{ComplexityOfModule({\mdseries\slshape M, n})\index{ComplexityOfModule@\texttt{ComplexityOfModule}}
\label{ComplexityOfModule}
}\hfill{\scriptsize (operation)}}\\


 Arguments: \mbox{\texttt{\mdseries\slshape M}} \texttt{\symbol{45}}\texttt{\symbol{45}} path algebdra module, \mbox{\texttt{\mdseries\slshape n}} \texttt{\symbol{45}}\texttt{\symbol{45}} a positive integer. \\
\textbf{\indent Returns:\ }
an estimate of the complexity of the module \mbox{\texttt{\mdseries\slshape M}}. 



 The function checks if the algebra over which the module \mbox{\texttt{\mdseries\slshape M}} lives is known to have finite global dimension. If so, it returns complexity
zero. Otherwise it tries to estimate the complexity in the following way.
Recall that if a function $f(x)$ is a polynomial in $x$, the degree of $f(x)$ is given by $\lim_{n\to\infty} \frac{\log |f(n)|}{\log n}$. So then this function computes an estimate of the complexity of \mbox{\texttt{\mdseries\slshape M}} by approximating the complexity by considering the limit $\lim_{m\to \infty} \log \frac{\dim(P(M)(m))}{\log m}$ where $P(M)(m)$ is the $m$\texttt{\symbol{45}}th projective in a minimal projective resolution of \mbox{\texttt{\mdseries\slshape M}} at stage $m$. This limit is estimated by $\frac{\log \dim(P(M)(n))}{\log n}$. }

 

\subsection{\textcolor{Chapter }{DecomposeModule}}
\logpage{[ 6, 4, 7 ]}\nobreak
\hyperdef{L}{X7B0588497D5008A4}{}
{\noindent\textcolor{FuncColor}{$\triangleright$\enspace\texttt{DecomposeModule({\mdseries\slshape M})\index{DecomposeModule@\texttt{DecomposeModule}}
\label{DecomposeModule}
}\hfill{\scriptsize (operation)}}\\
\noindent\textcolor{FuncColor}{$\triangleright$\enspace\texttt{DecomposeModuleWithInclusions({\mdseries\slshape M})\index{DecomposeModuleWithInclusions@\texttt{DecomposeModuleWithInclusions}}
\label{DecomposeModuleWithInclusions}
}\hfill{\scriptsize (operation)}}\\


 Arguments: \mbox{\texttt{\mdseries\slshape M}} \texttt{\symbol{45}}\texttt{\symbol{45}} a path algebra module. \\
\textbf{\indent Returns:\ }
a list of indecomposable modules whose direct sum is isomorphic to the module \mbox{\texttt{\mdseries\slshape M}} in first variant. The second variant returns a list of inclusions into \mbox{\texttt{\mdseries\slshape M}} with the sum of the images is isomorphic to the module \mbox{\texttt{\mdseries\slshape M}}. 



Warning: the function is not properly tested and it at best only works
properly over finite fields. }

 

\subsection{\textcolor{Chapter }{DecomposeModuleProbabilistic}}
\logpage{[ 6, 4, 8 ]}\nobreak
\hyperdef{L}{X81EF6AE97A4F77FA}{}
{\noindent\textcolor{FuncColor}{$\triangleright$\enspace\texttt{DecomposeModuleProbabilistic({\mdseries\slshape HomMM, M})\index{DecomposeModuleProbabilistic@\texttt{DecomposeModuleProbabilistic}}
\label{DecomposeModuleProbabilistic}
}\hfill{\scriptsize (operation)}}\\


 Arguments: \mbox{\texttt{\mdseries\slshape HomMM}}, \mbox{\texttt{\mdseries\slshape M}} \texttt{\symbol{45}}\texttt{\symbol{45}} a list of basis elements of the
Hom\texttt{\symbol{45}}space of the entered module and a path algebra module. \\
\textbf{\indent Returns:\ }
with (hopefully high probability) a list of indecomposable modules whose
direct sum is isomorphic to the module \mbox{\texttt{\mdseries\slshape M}}. 



Given a module \mbox{\texttt{\mdseries\slshape M}} over a finite dimensional quotient of a path algebra over a finite field, this
function tries to decompose the entered module \mbox{\texttt{\mdseries\slshape M}} by choosing random elements in the endomorphism ring of \mbox{\texttt{\mdseries\slshape M}} which are non\texttt{\symbol{45}}nilpotent and
non\texttt{\symbol{45}}invertible. Such elements splits the module in two
direct summands, and the procedure does this as long as it finds such
elements. The output is not guaranteed to be a list of indecomposable modules,
but their direct sum is isomorphic to the entered module \mbox{\texttt{\mdseries\slshape M}}. This was constructed as joint effort by the participants at the workshop
"Persistence, Representations, and Computation", February 26th
\texttt{\symbol{45}} March 2nd, 2018". This is an experimental function, so
use with caution. }

 

\subsection{\textcolor{Chapter }{DecomposeModuleViaCharPoly}}
\logpage{[ 6, 4, 9 ]}\nobreak
\hyperdef{L}{X7A4A2B0F7BB16D0D}{}
{\noindent\textcolor{FuncColor}{$\triangleright$\enspace\texttt{DecomposeModuleViaCharPoly({\mdseries\slshape M})\index{DecomposeModuleViaCharPoly@\texttt{DecomposeModuleViaCharPoly}}
\label{DecomposeModuleViaCharPoly}
}\hfill{\scriptsize (operation)}}\\


 Arguments: \mbox{\texttt{\mdseries\slshape M}} \texttt{\symbol{45}}\texttt{\symbol{45}} a path algebra module. \\
\textbf{\indent Returns:\ }
a list with high probability of indecomposable modules whose direct sum is
isomorphic to the module \mbox{\texttt{\mdseries\slshape M}}. 



Given a module \mbox{\texttt{\mdseries\slshape M}} over a finite dimensional quotient of a path algebra over a finite field, this
function decomposes the entered module \mbox{\texttt{\mdseries\slshape M}} by computing the endomorphism ring of \mbox{\texttt{\mdseries\slshape M}} and choosing random elements in it. This is an experimental function, so use
with caution. }

 

\subsection{\textcolor{Chapter }{DecomposeModuleViaTop}}
\logpage{[ 6, 4, 10 ]}\nobreak
\hyperdef{L}{X848EAC6C833C95B1}{}
{\noindent\textcolor{FuncColor}{$\triangleright$\enspace\texttt{DecomposeModuleViaTop({\mdseries\slshape M})\index{DecomposeModuleViaTop@\texttt{DecomposeModuleViaTop}}
\label{DecomposeModuleViaTop}
}\hfill{\scriptsize (operation)}}\\


 Arguments: \mbox{\texttt{\mdseries\slshape M}} \texttt{\symbol{45}}\texttt{\symbol{45}} a path algebra module. \\
\textbf{\indent Returns:\ }
a list of indecomposable modules whose direct sum is isomorphic to the module \mbox{\texttt{\mdseries\slshape M}}. 



Given a module \mbox{\texttt{\mdseries\slshape M}} over a finite dimensional quotient of a path algebra over a finite field, this
function decomposes the entered module \mbox{\texttt{\mdseries\slshape M}} by finding the image $\Sigma$ of the endomorphism ring of \mbox{\texttt{\mdseries\slshape M}} in the endomorphism ring of the top of \mbox{\texttt{\mdseries\slshape M}}, in $\Sigma$ finds a complete set of primitive idempotents, lifts them back to the
endomorphism ring of \mbox{\texttt{\mdseries\slshape M}} and decomposes \mbox{\texttt{\mdseries\slshape M}}. This is an experimental function, so use with caution. }

 

\subsection{\textcolor{Chapter }{DecomposeModuleWithMultiplicities}}
\logpage{[ 6, 4, 11 ]}\nobreak
\hyperdef{L}{X7E957BFE7897F504}{}
{\noindent\textcolor{FuncColor}{$\triangleright$\enspace\texttt{DecomposeModuleWithMultiplicities({\mdseries\slshape M})\index{DecomposeModuleWithMultiplicities@\texttt{DecomposeModuleWithMultiplicities}}
\label{DecomposeModuleWithMultiplicities}
}\hfill{\scriptsize (operation)}}\\


 Arguments: \mbox{\texttt{\mdseries\slshape M}} \texttt{\symbol{45}}\texttt{\symbol{45}} a path algebra module. \\
\textbf{\indent Returns:\ }
a list of length two, where the first entry is a list of all indecomposable
non\texttt{\symbol{45}}isomorphic direct summands of \mbox{\texttt{\mdseries\slshape M}} and the second entry is the list of the multiplicities of these direct summand
in the module \mbox{\texttt{\mdseries\slshape M}}. 



Warning: the function is not properly tested and it at best only works
properly over finite fields. }

 

\subsection{\textcolor{Chapter }{Dimension (for a PathAlgebraMatModule)}}
\logpage{[ 6, 4, 12 ]}\nobreak
\hyperdef{L}{X7E38D5A48344C173}{}
{\noindent\textcolor{FuncColor}{$\triangleright$\enspace\texttt{Dimension({\mdseries\slshape M})\index{Dimension@\texttt{Dimension}!for a PathAlgebraMatModule}
\label{Dimension:for a PathAlgebraMatModule}
}\hfill{\scriptsize (operation)}}\\


 Arguments: \mbox{\texttt{\mdseries\slshape M}} \texttt{\symbol{45}}\texttt{\symbol{45}} a path algebra module (\texttt{PathAlgebraMatModule}). \\
\textbf{\indent Returns:\ }
the dimension of the representation \mbox{\texttt{\mdseries\slshape M}}. 

}

 

\subsection{\textcolor{Chapter }{DimensionVector}}
\logpage{[ 6, 4, 13 ]}\nobreak
\hyperdef{L}{X7B5EA4B0820DE28C}{}
{\noindent\textcolor{FuncColor}{$\triangleright$\enspace\texttt{DimensionVector({\mdseries\slshape M})\index{DimensionVector@\texttt{DimensionVector}}
\label{DimensionVector}
}\hfill{\scriptsize (attribute)}}\\


 Arguments: \mbox{\texttt{\mdseries\slshape M}} \texttt{\symbol{45}}\texttt{\symbol{45}} a path algebra module (\texttt{PathAlgebraMatModule}). \\
\textbf{\indent Returns:\ }
the dimension vector of the representation \mbox{\texttt{\mdseries\slshape M}}. 

}

 

\subsection{\textcolor{Chapter }{DirectSumOfQPAModules}}
\logpage{[ 6, 4, 14 ]}\nobreak
\hyperdef{L}{X7BDD77707A013FBE}{}
{\noindent\textcolor{FuncColor}{$\triangleright$\enspace\texttt{DirectSumOfQPAModules({\mdseries\slshape L})\index{DirectSumOfQPAModules@\texttt{DirectSumOfQPAModules}}
\label{DirectSumOfQPAModules}
}\hfill{\scriptsize (operation)}}\\


 Arguments: \mbox{\texttt{\mdseries\slshape L}} \texttt{\symbol{45}}\texttt{\symbol{45}} a list of \texttt{PathAlgebraMatModule}s over the same (quotient of a) path algebra. \\
\textbf{\indent Returns:\ }
the direct sum of the representations contained in the list \mbox{\texttt{\mdseries\slshape L}}.



In addition three attributes are attached to the result, \texttt{IsDirectSumOfModules} (\ref{IsDirectSumOfModules}), \texttt{DirectSumProjections} (\ref{DirectSumProjections}) \texttt{DirectSumInclusions} (\ref{DirectSumInclusions}). }

 

\subsection{\textcolor{Chapter }{DirectSumInclusions}}
\logpage{[ 6, 4, 15 ]}\nobreak
\hyperdef{L}{X857807CF8560B3C4}{}
{\noindent\textcolor{FuncColor}{$\triangleright$\enspace\texttt{DirectSumInclusions({\mdseries\slshape M})\index{DirectSumInclusions@\texttt{DirectSumInclusions}}
\label{DirectSumInclusions}
}\hfill{\scriptsize (attribute)}}\\


 Arguments: \mbox{\texttt{\mdseries\slshape M}} \texttt{\symbol{45}}\texttt{\symbol{45}} a path algebra module (\texttt{PathAlgebraMatModule}). \\
\textbf{\indent Returns:\ }
 the list of inclusions from the individual modules to their direct sum, when a
direct sum has been constructed using \texttt{DirectSumOfQPAModules} (\ref{DirectSumOfQPAModules}). 

}

 

\subsection{\textcolor{Chapter }{DirectSumProjections}}
\logpage{[ 6, 4, 16 ]}\nobreak
\hyperdef{L}{X80CFB7E47A785E12}{}
{\noindent\textcolor{FuncColor}{$\triangleright$\enspace\texttt{DirectSumProjections({\mdseries\slshape M})\index{DirectSumProjections@\texttt{DirectSumProjections}}
\label{DirectSumProjections}
}\hfill{\scriptsize (attribute)}}\\


 Arguments: \mbox{\texttt{\mdseries\slshape M}} \texttt{\symbol{45}}\texttt{\symbol{45}} a path algebra module (\texttt{PathAlgebraMatModule}). \\
\textbf{\indent Returns:\ }
 the list of projections from the direct sum to the individual modules used to
construct direct sum, when a direct sum has been constructed using \texttt{DirectSumOfQPAModules} (\ref{DirectSumOfQPAModules}). 

}

 

\subsection{\textcolor{Chapter }{FromIdentityToDoubleStarHomomorphism}}
\logpage{[ 6, 4, 17 ]}\nobreak
\hyperdef{L}{X837AE1BF7F31AD7C}{}
{\noindent\textcolor{FuncColor}{$\triangleright$\enspace\texttt{FromIdentityToDoubleStarHomomorphism({\mdseries\slshape M})\index{FromIdentityToDoubleStarHomomorphism@\texttt{From}\-\texttt{Identity}\-\texttt{To}\-\texttt{Double}\-\texttt{Star}\-\texttt{Homomorphism}}
\label{FromIdentityToDoubleStarHomomorphism}
}\hfill{\scriptsize (attribute)}}\\


 Arguments: \mbox{\texttt{\mdseries\slshape M}} \texttt{\symbol{45}}\texttt{\symbol{45}} a path algebra module (\texttt{PathAlgebraMatModule}). \\
\textbf{\indent Returns:\ }
the homomorphism from \mbox{\texttt{\mdseries\slshape M}} to the double star of the module \mbox{\texttt{\mdseries\slshape M}}. 

}

 

\subsection{\textcolor{Chapter }{IntersectionOfSubmodules}}
\logpage{[ 6, 4, 18 ]}\nobreak
\hyperdef{L}{X7FE61CFE7A138755}{}
{\noindent\textcolor{FuncColor}{$\triangleright$\enspace\texttt{IntersectionOfSubmodules({\mdseries\slshape list})\index{IntersectionOfSubmodules@\texttt{IntersectionOfSubmodules}}
\label{IntersectionOfSubmodules}
}\hfill{\scriptsize (operation)}}\\


 Arguments: \mbox{\texttt{\mdseries\slshape f, g}} or \mbox{\texttt{\mdseries\slshape list}} \texttt{\symbol{45}}\texttt{\symbol{45}} two homomorphisms of
PathAlgebraMatModules or a list of such. \\
\textbf{\indent Returns:\ }
the subrepresentation given by the intersection of all the submodules given by
the inclusions \mbox{\texttt{\mdseries\slshape f}} and \mbox{\texttt{\mdseries\slshape g}} or \mbox{\texttt{\mdseries\slshape list}}. 



The function checks if \mbox{\texttt{\mdseries\slshape list}} is non\texttt{\symbol{45}}empty and if $\mbox{\texttt{\mdseries\slshape f}}\colon M\to X$ and $\mbox{\texttt{\mdseries\slshape g}}\colon N\to X$ or all the homomorphism in \mbox{\texttt{\mdseries\slshape list}} have the same range and if they all are inclusions. If the function is given
two arguments \mbox{\texttt{\mdseries\slshape f}} and \mbox{\texttt{\mdseries\slshape g}}, then it returns $[f',g',g'*f]$, where $f'\colon E\to N$, $g'\colon E\to M$, and $E$ is the pullback of \mbox{\texttt{\mdseries\slshape f}} and \mbox{\texttt{\mdseries\slshape g}}. For a list of inclusions it returns a monomorphism from a module isomorphic
to the intersection to $X$. }

 

\subsection{\textcolor{Chapter }{IsDirectSummand}}
\logpage{[ 6, 4, 19 ]}\nobreak
\hyperdef{L}{X7E24DCE07E98E50D}{}
{\noindent\textcolor{FuncColor}{$\triangleright$\enspace\texttt{IsDirectSummand({\mdseries\slshape M, N})\index{IsDirectSummand@\texttt{IsDirectSummand}}
\label{IsDirectSummand}
}\hfill{\scriptsize (operation)}}\\


 Arguments: \mbox{\texttt{\mdseries\slshape M, N}} \texttt{\symbol{45}}\texttt{\symbol{45}} two path algebra modules (\texttt{PathAlgebraMatModule}s). \\
\textbf{\indent Returns:\ }
true if \mbox{\texttt{\mdseries\slshape M}} is isomorphic to a direct summand of \mbox{\texttt{\mdseries\slshape N}}, otherwise false.



The function checks if \mbox{\texttt{\mdseries\slshape M}} and \mbox{\texttt{\mdseries\slshape N}} are \texttt{PathAlgebraMatModule}s over the same (quotient of a) path algebra. }

 

\subsection{\textcolor{Chapter }{IsDirectSumOfModules}}
\logpage{[ 6, 4, 20 ]}\nobreak
\hyperdef{L}{X7A50C15B87236111}{}
{\noindent\textcolor{FuncColor}{$\triangleright$\enspace\texttt{IsDirectSumOfModules({\mdseries\slshape M})\index{IsDirectSumOfModules@\texttt{IsDirectSumOfModules}}
\label{IsDirectSumOfModules}
}\hfill{\scriptsize (filter)}}\\


 Arguments: \mbox{\texttt{\mdseries\slshape M}} \texttt{\symbol{45}}\texttt{\symbol{45}} a path algebra module (\texttt{PathAlgebraMatModule}). \\
\textbf{\indent Returns:\ }
 true if \mbox{\texttt{\mdseries\slshape M}} is constructed via the command \texttt{DirectSumOfQPAModules} (\ref{DirectSumOfQPAModules}). 

}

 Using the example above. 
\begin{Verbatim}[commandchars=!@|,fontsize=\small,frame=single,label=Example]
  !gapprompt@gap>| !gapinput@N2 := DirectSumOfQPAModules([N,N]);|
  <[ 6, 4, 4 ]>
  !gapprompt@gap>| !gapinput@proj := DirectSumProjections(N2);|
  [ <<[ 6, 4, 4 ]> ---> <[ 3, 2, 2 ]>>
      , <<[ 6, 4, 4 ]> ---> <[ 3, 2, 2 ]>>
       ]
  !gapprompt@gap>| !gapinput@inc := DirectSumInclusions(N2);|
  [ <<[ 3, 2, 2 ]> ---> <[ 6, 4, 4 ]>>
      , <<[ 3, 2, 2 ]> ---> <[ 6, 4, 4 ]>>
       ] 
\end{Verbatim}
 

\subsection{\textcolor{Chapter }{IsExceptionalModule}}
\logpage{[ 6, 4, 21 ]}\nobreak
\hyperdef{L}{X7C9BFC678598DBF6}{}
{\noindent\textcolor{FuncColor}{$\triangleright$\enspace\texttt{IsExceptionalModule({\mdseries\slshape M})\index{IsExceptionalModule@\texttt{IsExceptionalModule}}
\label{IsExceptionalModule}
}\hfill{\scriptsize (property)}}\\


 Arguments: \mbox{\texttt{\mdseries\slshape M}} \texttt{\symbol{45}}\texttt{\symbol{45}} a path algebra module (\texttt{PathAlgebraMatModule}). \\
\textbf{\indent Returns:\ }
true if \mbox{\texttt{\mdseries\slshape M}} is an exceptional module, otherwise false, if the field, over which the
algebra \mbox{\texttt{\mdseries\slshape M}} is defined over, is finite. 



 The module \mbox{\texttt{\mdseries\slshape M}} is an exceptional module, if it is indecomposable and $\Ext^1(M,M)=(0)$. }

 

\subsection{\textcolor{Chapter }{IsIndecomposableModule}}
\logpage{[ 6, 4, 22 ]}\nobreak
\hyperdef{L}{X82102F847994003E}{}
{\noindent\textcolor{FuncColor}{$\triangleright$\enspace\texttt{IsIndecomposableModule({\mdseries\slshape M})\index{IsIndecomposableModule@\texttt{IsIndecomposableModule}}
\label{IsIndecomposableModule}
}\hfill{\scriptsize (property)}}\\


 Arguments: \mbox{\texttt{\mdseries\slshape M}} \texttt{\symbol{45}}\texttt{\symbol{45}} a path algebra module (\texttt{PathAlgebraMatModule}). \\
\textbf{\indent Returns:\ }
true if \mbox{\texttt{\mdseries\slshape M}} is an indecomposable module, if the field, over which the algebra \mbox{\texttt{\mdseries\slshape M}} is defined over, is finite. If \mbox{\texttt{\mdseries\slshape M}} is the zero module, then \texttt{false} is returned. 

}

 

\subsection{\textcolor{Chapter }{IsInAdditiveClosure}}
\logpage{[ 6, 4, 23 ]}\nobreak
\hyperdef{L}{X7E5246D4831DB250}{}
{\noindent\textcolor{FuncColor}{$\triangleright$\enspace\texttt{IsInAdditiveClosure({\mdseries\slshape M, N})\index{IsInAdditiveClosure@\texttt{IsInAdditiveClosure}}
\label{IsInAdditiveClosure}
}\hfill{\scriptsize (operation)}}\\


 Arguments: \mbox{\texttt{\mdseries\slshape M, N}} \texttt{\symbol{45}}\texttt{\symbol{45}} two path algebra modules (\texttt{PathAlgebraMatModule}s). \\
\textbf{\indent Returns:\ }
true if \mbox{\texttt{\mdseries\slshape M}} is in the additive closure of the module \mbox{\texttt{\mdseries\slshape N}}, otherwise false.



The function checks if \mbox{\texttt{\mdseries\slshape M}} and \mbox{\texttt{\mdseries\slshape N}} are \texttt{PathAlgebraMatModule}s over the same (quotient of a) path algebra. }

 

\subsection{\textcolor{Chapter }{IsInjectiveModule}}
\logpage{[ 6, 4, 24 ]}\nobreak
\hyperdef{L}{X803C2799861FFBC5}{}
{\noindent\textcolor{FuncColor}{$\triangleright$\enspace\texttt{IsInjectiveModule({\mdseries\slshape M})\index{IsInjectiveModule@\texttt{IsInjectiveModule}}
\label{IsInjectiveModule}
}\hfill{\scriptsize (property)}}\\


 Arguments: \mbox{\texttt{\mdseries\slshape M}} \texttt{\symbol{45}}\texttt{\symbol{45}} a path algebra module (\texttt{PathAlgebraMatModule}). \\
\textbf{\indent Returns:\ }
true if the representation \mbox{\texttt{\mdseries\slshape M}} is injective. 

}

 

\subsection{\textcolor{Chapter }{IsomorphicModules}}
\logpage{[ 6, 4, 25 ]}\nobreak
\hyperdef{L}{X7D198BB5808D38F2}{}
{\noindent\textcolor{FuncColor}{$\triangleright$\enspace\texttt{IsomorphicModules({\mdseries\slshape M, N})\index{IsomorphicModules@\texttt{IsomorphicModules}}
\label{IsomorphicModules}
}\hfill{\scriptsize (operation)}}\\


 Arguments: \mbox{\texttt{\mdseries\slshape M, N}} \texttt{\symbol{45}}\texttt{\symbol{45}} two path algebra modules (\texttt{PathAlgebraMatModule}s). \\
\textbf{\indent Returns:\ }
true or false depending on whether \mbox{\texttt{\mdseries\slshape M}} and \mbox{\texttt{\mdseries\slshape N}} are isomorphic or not.



 The function first checks if the modules \mbox{\texttt{\mdseries\slshape M}} and \mbox{\texttt{\mdseries\slshape N}} are modules over the same algebra, and returns fail if not. The function
returns true if the modules are isomorphic, otherwise false. }

 

\subsection{\textcolor{Chapter }{IsProjectiveModule}}
\logpage{[ 6, 4, 26 ]}\nobreak
\hyperdef{L}{X8359AC9585777CA1}{}
{\noindent\textcolor{FuncColor}{$\triangleright$\enspace\texttt{IsProjectiveModule({\mdseries\slshape M})\index{IsProjectiveModule@\texttt{IsProjectiveModule}}
\label{IsProjectiveModule}
}\hfill{\scriptsize (property)}}\\


 Arguments: \mbox{\texttt{\mdseries\slshape M}} \texttt{\symbol{45}}\texttt{\symbol{45}} a path algebra module (\texttt{PathAlgebraMatModule}). \\
\textbf{\indent Returns:\ }
true if the representation \mbox{\texttt{\mdseries\slshape M}} is projective. 

}

 

\subsection{\textcolor{Chapter }{IsRigidModule}}
\logpage{[ 6, 4, 27 ]}\nobreak
\hyperdef{L}{X7ECEEC6F873A7BA6}{}
{\noindent\textcolor{FuncColor}{$\triangleright$\enspace\texttt{IsRigidModule({\mdseries\slshape M})\index{IsRigidModule@\texttt{IsRigidModule}}
\label{IsRigidModule}
}\hfill{\scriptsize (property)}}\\


 Arguments: \mbox{\texttt{\mdseries\slshape M}} \texttt{\symbol{45}}\texttt{\symbol{45}} a path algebra module (\texttt{PathAlgebraMatModule}). \\
\textbf{\indent Returns:\ }
true if \mbox{\texttt{\mdseries\slshape M}} is a rigid module, otherwise false. 



 The module \mbox{\texttt{\mdseries\slshape M}} is a rigid module, if $\Ext^1(M,M)=(0)$. }

 

\subsection{\textcolor{Chapter }{IsSemisimpleModule}}
\logpage{[ 6, 4, 28 ]}\nobreak
\hyperdef{L}{X7A8BC26E866E44DD}{}
{\noindent\textcolor{FuncColor}{$\triangleright$\enspace\texttt{IsSemisimpleModule({\mdseries\slshape M})\index{IsSemisimpleModule@\texttt{IsSemisimpleModule}}
\label{IsSemisimpleModule}
}\hfill{\scriptsize (property)}}\\


 Arguments: \mbox{\texttt{\mdseries\slshape M}} \texttt{\symbol{45}}\texttt{\symbol{45}} a path algebra module (\texttt{PathAlgebraMatModule}). \\
\textbf{\indent Returns:\ }
true if the representation \mbox{\texttt{\mdseries\slshape M}} is semisimple. 

}

 

\subsection{\textcolor{Chapter }{IsSimpleQPAModule}}
\logpage{[ 6, 4, 29 ]}\nobreak
\hyperdef{L}{X7D2067E57E0244F8}{}
{\noindent\textcolor{FuncColor}{$\triangleright$\enspace\texttt{IsSimpleQPAModule({\mdseries\slshape M})\index{IsSimpleQPAModule@\texttt{IsSimpleQPAModule}}
\label{IsSimpleQPAModule}
}\hfill{\scriptsize (property)}}\\


 Arguments: \mbox{\texttt{\mdseries\slshape M}} \texttt{\symbol{45}}\texttt{\symbol{45}} a path algebra module (\texttt{PathAlgebraMatModule}). \\
\textbf{\indent Returns:\ }
true if the representation \mbox{\texttt{\mdseries\slshape M}} is simple. 

}

 

\subsection{\textcolor{Chapter }{IsTauRigidModule}}
\logpage{[ 6, 4, 30 ]}\nobreak
\hyperdef{L}{X7B766424838EE6EA}{}
{\noindent\textcolor{FuncColor}{$\triangleright$\enspace\texttt{IsTauRigidModule({\mdseries\slshape M})\index{IsTauRigidModule@\texttt{IsTauRigidModule}}
\label{IsTauRigidModule}
}\hfill{\scriptsize (property)}}\\


 Arguments: \mbox{\texttt{\mdseries\slshape M}} \texttt{\symbol{45}}\texttt{\symbol{45}} a path algebra module (\texttt{PathAlgebraMatModule}). \\
\textbf{\indent Returns:\ }
true if \mbox{\texttt{\mdseries\slshape M}} is a $\tau$\texttt{\symbol{45}}rigid module, otherwise false. 



 The module \mbox{\texttt{\mdseries\slshape M}} is a $\tau$\texttt{\symbol{45}}rigid module, if $\Hom(M,\tau M)=(0)$. }

 

\subsection{\textcolor{Chapter }{LoewyLength (for a PathAlgebraMatModule)}}
\logpage{[ 6, 4, 31 ]}\nobreak
\hyperdef{L}{X82663AD8832DD57F}{}
{\noindent\textcolor{FuncColor}{$\triangleright$\enspace\texttt{LoewyLength({\mdseries\slshape M})\index{LoewyLength@\texttt{LoewyLength}!for a PathAlgebraMatModule}
\label{LoewyLength:for a PathAlgebraMatModule}
}\hfill{\scriptsize (attribute)}}\\


 Arguments: \mbox{\texttt{\mdseries\slshape M}} \texttt{\symbol{45}}\texttt{\symbol{45}} a path algebra module (\texttt{PathAlgebraMatModule}). \\
\textbf{\indent Returns:\ }
the Loewy length of the module \mbox{\texttt{\mdseries\slshape M}}.



 The function checks that the module \mbox{\texttt{\mdseries\slshape M}} is a module over a finite dimensional quotient of a path algebra, and returns
fail otherwise (This is not implemented yet). }

 

\subsection{\textcolor{Chapter }{IsZero(module)}}
\logpage{[ 6, 4, 32 ]}\nobreak
\hyperdef{L}{X7DC6D98E7F97FFEE}{}
{\noindent\textcolor{FuncColor}{$\triangleright$\enspace\texttt{IsZero(module)({\mdseries\slshape M})\index{IsZero(module)@\texttt{IsZero(module)}}
\label{IsZero(module)}
}\hfill{\scriptsize (property)}}\\


 Arguments: \mbox{\texttt{\mdseries\slshape M}} \texttt{\symbol{45}}\texttt{\symbol{45}} a path algebra module (\texttt{PathAlgebraMatModule}). \\
\textbf{\indent Returns:\ }
true if \mbox{\texttt{\mdseries\slshape M}} is the zero module, otherwise false. 

}

 

\subsection{\textcolor{Chapter }{MakeBasicListOfModules}}
\logpage{[ 6, 4, 33 ]}\nobreak
\hyperdef{L}{X877B3C2480408B5C}{}
{\noindent\textcolor{FuncColor}{$\triangleright$\enspace\texttt{MakeBasicListOfModules({\mdseries\slshape L})\index{MakeBasicListOfModules@\texttt{MakeBasicListOfModules}}
\label{MakeBasicListOfModules}
}\hfill{\scriptsize (operation)}}\\


 Arguments: \mbox{\texttt{\mdseries\slshape L}} \texttt{\symbol{45}}\texttt{\symbol{45}} a list of path algebra modules. \\
\textbf{\indent Returns:\ }
 a basic version of the entered list modules \mbox{\texttt{\mdseries\slshape L}}, that is, if the direct sum $M$ of all the modules in the list \mbox{\texttt{\mdseries\slshape L}} is: $M \simeq M_1^{n_1} \oplus \cdots \oplus M_t^{n_t},$ where $M_i$ is indecomposable, then the list $\{M_1,M_2,\ldots,M_t\}$ is returned. At present, this function only work for finite dimensional
(quotients of a) path algebra over a finite field. 

}

 

\subsection{\textcolor{Chapter }{MatricesOfPathAlgebraModule}}
\logpage{[ 6, 4, 34 ]}\nobreak
\hyperdef{L}{X81BB198380631A9B}{}
{\noindent\textcolor{FuncColor}{$\triangleright$\enspace\texttt{MatricesOfPathAlgebraModule({\mdseries\slshape M})\index{MatricesOfPathAlgebraModule@\texttt{MatricesOfPathAlgebraModule}}
\label{MatricesOfPathAlgebraModule}
}\hfill{\scriptsize (operation)}}\\


 Arguments: \mbox{\texttt{\mdseries\slshape M}} \texttt{\symbol{45}}\texttt{\symbol{45}} a path algebra module (\texttt{PathAlgebraMatModule}). \\
\textbf{\indent Returns:\ }
a list of the matrices that defines the representation \mbox{\texttt{\mdseries\slshape M}} as a right module of the acting path algebra. 



 The list of matrices that are returned are not the same identical to the
matrices entered to define the representation if there is zero vector space in
at least one vertex. Then zero matrices of the appropriate size are returned. }

 

\subsection{\textcolor{Chapter }{MaximalCommonDirectSummand}}
\logpage{[ 6, 4, 35 ]}\nobreak
\hyperdef{L}{X7B480767836D0764}{}
{\noindent\textcolor{FuncColor}{$\triangleright$\enspace\texttt{MaximalCommonDirectSummand({\mdseries\slshape M, N})\index{MaximalCommonDirectSummand@\texttt{MaximalCommonDirectSummand}}
\label{MaximalCommonDirectSummand}
}\hfill{\scriptsize (operation)}}\\


 Arguments: \mbox{\texttt{\mdseries\slshape M, N}} \texttt{\symbol{45}}\texttt{\symbol{45}} two path algebra modules (\texttt{PathAlgebraMatModule}s). \\
\textbf{\indent Returns:\ }
a list of three modules [\mbox{\texttt{\mdseries\slshape X}},\mbox{\texttt{\mdseries\slshape U}},\mbox{\texttt{\mdseries\slshape V}}], where \mbox{\texttt{\mdseries\slshape X}} is a maximal common non\texttt{\symbol{45}}zero direct summand of \mbox{\texttt{\mdseries\slshape M}} and \mbox{\texttt{\mdseries\slshape N}}, the sum of \mbox{\texttt{\mdseries\slshape X}} and \mbox{\texttt{\mdseries\slshape U}} is \mbox{\texttt{\mdseries\slshape M}} and the sum of \mbox{\texttt{\mdseries\slshape X}} and \mbox{\texttt{\mdseries\slshape V}} is \mbox{\texttt{\mdseries\slshape N}}, if such a non\texttt{\symbol{45}}zero maximal direct summand exists.
Otherwise it returns false.



The function checks if \mbox{\texttt{\mdseries\slshape M}} and \mbox{\texttt{\mdseries\slshape N}} are \texttt{PathAlgebraMatModule}s over the same (quotient of a) path algebra. }

 

\subsection{\textcolor{Chapter }{NumberOfNonIsoDirSummands}}
\logpage{[ 6, 4, 36 ]}\nobreak
\hyperdef{L}{X7CD05FF985A1D1B3}{}
{\noindent\textcolor{FuncColor}{$\triangleright$\enspace\texttt{NumberOfNonIsoDirSummands({\mdseries\slshape M})\index{NumberOfNonIsoDirSummands@\texttt{NumberOfNonIsoDirSummands}}
\label{NumberOfNonIsoDirSummands}
}\hfill{\scriptsize (operation)}}\\


 Arguments: \mbox{\texttt{\mdseries\slshape M}} \texttt{\symbol{45}}\texttt{\symbol{45}} a path algebra modules (\texttt{PathAlgebraMatModule}s). \\
\textbf{\indent Returns:\ }
a list with two elements: (1) the number of non\texttt{\symbol{45}}isomorphic
indecomposable direct summands of the module \mbox{\texttt{\mdseries\slshape M}} and (2) the dimensions of the simple blocks of the semisimple ring $\End(M)/\rad \End(M)$.

}

 

\subsection{\textcolor{Chapter }{MinimalGeneratingSetOfModule}}
\logpage{[ 6, 4, 37 ]}\nobreak
\hyperdef{L}{X821FA104861FF19B}{}
{\noindent\textcolor{FuncColor}{$\triangleright$\enspace\texttt{MinimalGeneratingSetOfModule({\mdseries\slshape M})\index{MinimalGeneratingSetOfModule@\texttt{MinimalGeneratingSetOfModule}}
\label{MinimalGeneratingSetOfModule}
}\hfill{\scriptsize (attribute)}}\\


 Arguments: \mbox{\texttt{\mdseries\slshape M}} \texttt{\symbol{45}}\texttt{\symbol{45}} a path algebra module (\texttt{PathAlgebraMatModule}). \\
\textbf{\indent Returns:\ }
a minimal generator set of the module \mbox{\texttt{\mdseries\slshape M}} as a module of the path algebra it is defined over. 

}

 

\subsection{\textcolor{Chapter }{RadicalOfModule}}
\logpage{[ 6, 4, 38 ]}\nobreak
\hyperdef{L}{X7E44920683157DE2}{}
{\noindent\textcolor{FuncColor}{$\triangleright$\enspace\texttt{RadicalOfModule({\mdseries\slshape M})\index{RadicalOfModule@\texttt{RadicalOfModule}}
\label{RadicalOfModule}
}\hfill{\scriptsize (operation)}}\\


 Arguments: \mbox{\texttt{\mdseries\slshape M}} \texttt{\symbol{45}}\texttt{\symbol{45}} a path algebra module (\texttt{PathAlgebraMatModule}). \\
\textbf{\indent Returns:\ }
the radical of the module \mbox{\texttt{\mdseries\slshape M}}.



 This returns only the representation given by the radical of the module \mbox{\texttt{\mdseries\slshape M}}. The operation \texttt{RadicalOfModuleInclusion} (\ref{RadicalOfModuleInclusion}) computes the inclusion of the radical of \mbox{\texttt{\mdseries\slshape M}} into \mbox{\texttt{\mdseries\slshape M}}. This function applies when the algebra over which \mbox{\texttt{\mdseries\slshape M}} is defined is an admissible quotient of a path algebra. }

 

\subsection{\textcolor{Chapter }{RadicalSeries}}
\logpage{[ 6, 4, 39 ]}\nobreak
\hyperdef{L}{X7929281B848A9FBE}{}
{\noindent\textcolor{FuncColor}{$\triangleright$\enspace\texttt{RadicalSeries({\mdseries\slshape M})\index{RadicalSeries@\texttt{RadicalSeries}}
\label{RadicalSeries}
}\hfill{\scriptsize (operation)}}\\


 Arguments: \mbox{\texttt{\mdseries\slshape M}} \texttt{\symbol{45}}\texttt{\symbol{45}} a path algebra module (\texttt{PathAlgebraMatModule}). \\
\textbf{\indent Returns:\ }
the radical series of the module \mbox{\texttt{\mdseries\slshape M}}.



 The function gives the radical series as a list of vectors \texttt{[n{\textunderscore}1,...,n{\textunderscore}s]}, where the algebra has $s$ isomorphism classes of simple modules and the numbers give the multiplicity of
each simple. The first vector listed corresponds to the top layer, and so on. }

 

\subsection{\textcolor{Chapter }{RandomModule}}
\logpage{[ 6, 4, 40 ]}\nobreak
\hyperdef{L}{X798A50297D2F186E}{}
{\noindent\textcolor{FuncColor}{$\triangleright$\enspace\texttt{RandomModule({\mdseries\slshape A, m, n})\index{RandomModule@\texttt{RandomModule}}
\label{RandomModule}
}\hfill{\scriptsize (operation)}}\\


 Arguments: \mbox{\texttt{\mdseries\slshape A, m, n}} \texttt{\symbol{45}}\texttt{\symbol{45}} a quiver algebra, two positive
integers. \\
\textbf{\indent Returns:\ }
a random module over the quiver algebra \mbox{\texttt{\mdseries\slshape A}} with a projective presentation $P_1 ---> P_0$ where they have \mbox{\texttt{\mdseries\slshape m}} and \mbox{\texttt{\mdseries\slshape n}} indecomposable projectives as direct summands.

}

 

\subsection{\textcolor{Chapter }{SocleSeries}}
\logpage{[ 6, 4, 41 ]}\nobreak
\hyperdef{L}{X84A724267E6F136D}{}
{\noindent\textcolor{FuncColor}{$\triangleright$\enspace\texttt{SocleSeries({\mdseries\slshape M})\index{SocleSeries@\texttt{SocleSeries}}
\label{SocleSeries}
}\hfill{\scriptsize (operation)}}\\


 Arguments: \mbox{\texttt{\mdseries\slshape M}} \texttt{\symbol{45}}\texttt{\symbol{45}} a path algebra module (\texttt{PathAlgebraMatModule}). \\
\textbf{\indent Returns:\ }
the socle series of the module \mbox{\texttt{\mdseries\slshape M}}.



 The function gives the socle series as a list of vectors \texttt{[n{\textunderscore}1,...,n{\textunderscore}s]}, where the algebra has $s$ isomorphism classes of simple modules and the numbers give the multiplicity of
each simple. The last vector listed corresponds to the socle layer, and so on
backwards. }

 

\subsection{\textcolor{Chapter }{SocleOfModule}}
\logpage{[ 6, 4, 42 ]}\nobreak
\hyperdef{L}{X79DF34618798E866}{}
{\noindent\textcolor{FuncColor}{$\triangleright$\enspace\texttt{SocleOfModule({\mdseries\slshape M})\index{SocleOfModule@\texttt{SocleOfModule}}
\label{SocleOfModule}
}\hfill{\scriptsize (operation)}}\\


 Arguments: \mbox{\texttt{\mdseries\slshape M}} \texttt{\symbol{45}}\texttt{\symbol{45}} a path algebra module (\texttt{PathAlgebraMatModule}). \\
\textbf{\indent Returns:\ }
the socle of the module \mbox{\texttt{\mdseries\slshape M}}.



 This operation only return the representation given by the socle of the module \mbox{\texttt{\mdseries\slshape M}}. The inclusion the socle of \mbox{\texttt{\mdseries\slshape M}} into \mbox{\texttt{\mdseries\slshape M}} can be computed using \texttt{SocleOfModuleInclusion} (\ref{SocleOfModuleInclusion}). }

 

\subsection{\textcolor{Chapter }{SubRepresentation}}
\logpage{[ 6, 4, 43 ]}\nobreak
\hyperdef{L}{X858AC23C83AC843E}{}
{\noindent\textcolor{FuncColor}{$\triangleright$\enspace\texttt{SubRepresentation({\mdseries\slshape M, gens})\index{SubRepresentation@\texttt{SubRepresentation}}
\label{SubRepresentation}
}\hfill{\scriptsize (operation)}}\\


 Arguments: \mbox{\texttt{\mdseries\slshape M}} \texttt{\symbol{45}}\texttt{\symbol{45}} a path algebra module (\texttt{PathAlgebraMatModule}), \mbox{\texttt{\mdseries\slshape gens}} \texttt{\symbol{45}}\texttt{\symbol{45}} elements in \mbox{\texttt{\mdseries\slshape M}}. \\
\textbf{\indent Returns:\ }
the submodule of the module \mbox{\texttt{\mdseries\slshape M}} generated by the elements \mbox{\texttt{\mdseries\slshape gens}}.



 The function checks if \mbox{\texttt{\mdseries\slshape gens}} are elements in \mbox{\texttt{\mdseries\slshape M}}, and returns an error message otherwise. The inclusion of the submodule
generated by the elements \mbox{\texttt{\mdseries\slshape gens}} into \mbox{\texttt{\mdseries\slshape M}} can be computed using \texttt{SubRepresentationInclusion} (\ref{SubRepresentationInclusion}). }

 

\subsection{\textcolor{Chapter }{SumOfSubmodules}}
\logpage{[ 6, 4, 44 ]}\nobreak
\hyperdef{L}{X8653599686499CD6}{}
{\noindent\textcolor{FuncColor}{$\triangleright$\enspace\texttt{SumOfSubmodules({\mdseries\slshape list})\index{SumOfSubmodules@\texttt{SumOfSubmodules}}
\label{SumOfSubmodules}
}\hfill{\scriptsize (operation)}}\\


 Arguments: \mbox{\texttt{\mdseries\slshape f, g}} or \mbox{\texttt{\mdseries\slshape list}} \texttt{\symbol{45}}\texttt{\symbol{45}} two inclusions of
PathAlgebraMatModules or a list of such. \\
\textbf{\indent Returns:\ }
the subrepresentation given by the sum of all the submodules given by the
inclusions \mbox{\texttt{\mdseries\slshape f, g}} or \mbox{\texttt{\mdseries\slshape list}}. 



The function checks if \mbox{\texttt{\mdseries\slshape list}} is non\texttt{\symbol{45}}empty and if $\mbox{\texttt{\mdseries\slshape f}}\colon M\to X$ and $\mbox{\texttt{\mdseries\slshape g}}\colon N\to X$ or all the homomorphism in \mbox{\texttt{\mdseries\slshape list}} have the same range and if they all are inclusions. If the function is given
two arguments \mbox{\texttt{\mdseries\slshape f}} and \mbox{\texttt{\mdseries\slshape g}}, then it returns $[h,f',g']$, where $h\colon M + N\to X$, $f'\colon M\to M + N$ and $g'\colon N\to M + N$. For a list of inclusions it returns a monomorphism from a module isomorphic
to the sum of the subrepresentations to $X$. }

 

\subsection{\textcolor{Chapter }{SupportModuleElement}}
\logpage{[ 6, 4, 45 ]}\nobreak
\hyperdef{L}{X856EA09A83A5A636}{}
{\noindent\textcolor{FuncColor}{$\triangleright$\enspace\texttt{SupportModuleElement({\mdseries\slshape m})\index{SupportModuleElement@\texttt{SupportModuleElement}}
\label{SupportModuleElement}
}\hfill{\scriptsize (operation)}}\\


 Arguments: \mbox{\texttt{\mdseries\slshape m}} \texttt{\symbol{45}}\texttt{\symbol{45}} an element of a path algebra module. \\
\textbf{\indent Returns:\ }
the primitive idempotents \mbox{\texttt{\mdseries\slshape v}} in the algebra over which the module containing the element \mbox{\texttt{\mdseries\slshape m}} is a module, such that \mbox{\texttt{\mdseries\slshape m\texttt{\symbol{94}}v}} is non\texttt{\symbol{45}}zero.



 The function checks if \mbox{\texttt{\mdseries\slshape m}} is an element in a module over a (quotient of a) path algebra, and returns
fail otherwise. }

 

\subsection{\textcolor{Chapter }{TopOfModule}}
\logpage{[ 6, 4, 46 ]}\nobreak
\hyperdef{L}{X87F571327E43AFB4}{}
{\noindent\textcolor{FuncColor}{$\triangleright$\enspace\texttt{TopOfModule({\mdseries\slshape M})\index{TopOfModule@\texttt{TopOfModule}}
\label{TopOfModule}
}\hfill{\scriptsize (operation)}}\\


 Arguments: \mbox{\texttt{\mdseries\slshape M}} or \mbox{\texttt{\mdseries\slshape f}} \texttt{\symbol{45}}\texttt{\symbol{45}} a path algebra module or a
homomorphism thereof. \\
\textbf{\indent Returns:\ }
the top of the module \mbox{\texttt{\mdseries\slshape M}} or the homomorphism induced on the top of the modules associated to the
homomorphism \mbox{\texttt{\mdseries\slshape f}}.



 This returns only the representation given by the top of the module \mbox{\texttt{\mdseries\slshape M}} or the homomorphism induced on the top of the modules associated to the
entered homomorphism. The operation \texttt{TopOfModuleProjection} (\ref{TopOfModuleProjection}) computes the projection of the module \mbox{\texttt{\mdseries\slshape M}} onto the top of the module \mbox{\texttt{\mdseries\slshape M}}. }

 }

 
\section{\textcolor{Chapter }{Special representations}}\logpage{[ 6, 5, 0 ]}
\hyperdef{L}{X7919F94382D9B38B}{}
{
 Here we collect the predefined representations/modules over a finite
dimensional quotient of a path algebra. 

\subsection{\textcolor{Chapter }{BasisOfProjectives}}
\logpage{[ 6, 5, 1 ]}\nobreak
\hyperdef{L}{X8048CD27796253CA}{}
{\noindent\textcolor{FuncColor}{$\triangleright$\enspace\texttt{BasisOfProjectives({\mdseries\slshape A})\index{BasisOfProjectives@\texttt{BasisOfProjectives}}
\label{BasisOfProjectives}
}\hfill{\scriptsize (attribute)}}\\


 Arguments: \mbox{\texttt{\mdseries\slshape A}} \texttt{\symbol{45}}\texttt{\symbol{45}} a finite dimensional (quotient of a)
path algebra. \\
\textbf{\indent Returns:\ }
a list of bases for all the indecomposable projective representations over \mbox{\texttt{\mdseries\slshape A}}. The basis for each indecomposable projective is given a list of elements in
nontips in \mbox{\texttt{\mdseries\slshape A}}. 



 The function checks if the algebra \mbox{\texttt{\mdseries\slshape A}} is a finite dimensional (quotient of a) path algebra, and returns an error
message otherwise. }

 

\subsection{\textcolor{Chapter }{ElementInIndecProjective}}
\logpage{[ 6, 5, 2 ]}\nobreak
\hyperdef{L}{X822401E583C75FCE}{}
{\noindent\textcolor{FuncColor}{$\triangleright$\enspace\texttt{ElementInIndecProjective({\mdseries\slshape A, m, s})\index{ElementInIndecProjective@\texttt{ElementInIndecProjective}}
\label{ElementInIndecProjective}
}\hfill{\scriptsize (operation)}}\\


 Arguments: \mbox{\texttt{\mdseries\slshape A}} \texttt{\symbol{45}}\texttt{\symbol{45}} a QuiverAlgebra, \mbox{\texttt{\mdseries\slshape m}} \texttt{\symbol{45}}\texttt{\symbol{45}} an element in an indecomposable
projective representation, \mbox{\texttt{\mdseries\slshape s}} \texttt{\symbol{45}}\texttt{\symbol{45}} an integer. \\
\textbf{\indent Returns:\ }
the element in the path algebra corresponding to \mbox{\texttt{\mdseries\slshape m}} in the right ideal from which the indecomposable projective representation is
constructed.

}

 

\subsection{\textcolor{Chapter }{ElementIn{\textunderscore}vA{\textunderscore}AsElementInIndecProj}}
\logpage{[ 6, 5, 3 ]}\nobreak
\hyperdef{L}{X7E298A807E5EB1A8}{}
{\noindent\textcolor{FuncColor}{$\triangleright$\enspace\texttt{ElementIn{\textunderscore}vA{\textunderscore}AsElementInIndecProj({\mdseries\slshape A, m})\index{ElementIn{\textunderscore}vA{\textunderscore}AsElementInIndecProj@\texttt{Element}\-\texttt{In{\textunderscore}v}\-\texttt{A{\textunderscore}}\-\texttt{As}\-\texttt{Element}\-\texttt{In}\-\texttt{Indec}\-\texttt{Proj}}
\label{ElementInuScorevAuScoreAsElementInIndecProj}
}\hfill{\scriptsize (operation)}}\\


 Arguments: \mbox{\texttt{\mdseries\slshape A}} \texttt{\symbol{45}}\texttt{\symbol{45}} a QuiverAlgebra, \mbox{\texttt{\mdseries\slshape m}} \texttt{\symbol{45}}\texttt{\symbol{45}} an element in \mbox{\texttt{\mdseries\slshape vA}} for a vertex \mbox{\texttt{\mdseries\slshape v}}. \\
\textbf{\indent Returns:\ }
the corresponding element in the indecomposable projective representation
assosicated with the vertex \mbox{\texttt{\mdseries\slshape v}}.

}

 

\subsection{\textcolor{Chapter }{IndecInjectiveModules}}
\logpage{[ 6, 5, 4 ]}\nobreak
\hyperdef{L}{X87741234871B1F5C}{}
{\noindent\textcolor{FuncColor}{$\triangleright$\enspace\texttt{IndecInjectiveModules({\mdseries\slshape A})\index{IndecInjectiveModules@\texttt{IndecInjectiveModules}}
\label{IndecInjectiveModules}
}\hfill{\scriptsize (attribute)}}\\


 Arguments: \mbox{\texttt{\mdseries\slshape A}} \texttt{\symbol{45}}\texttt{\symbol{45}} a finite dimensional (quotient of a)
path algebra. \\
\textbf{\indent Returns:\ }
a list of all the non\texttt{\symbol{45}}isomorphic indecomposable injective
representations over \mbox{\texttt{\mdseries\slshape A}}. 



 The function checks if the algebra \mbox{\texttt{\mdseries\slshape A}} is a finite dimensional (quotient of a) path algebra, and returns an error
message otherwise. }

 

\subsection{\textcolor{Chapter }{IndecProjectiveModules}}
\logpage{[ 6, 5, 5 ]}\nobreak
\hyperdef{L}{X85EDCFE27F66093F}{}
{\noindent\textcolor{FuncColor}{$\triangleright$\enspace\texttt{IndecProjectiveModules({\mdseries\slshape A})\index{IndecProjectiveModules@\texttt{IndecProjectiveModules}}
\label{IndecProjectiveModules}
}\hfill{\scriptsize (attribute)}}\\


 Arguments: \mbox{\texttt{\mdseries\slshape A}} \texttt{\symbol{45}}\texttt{\symbol{45}} a finite dimensional (quotient of a)
path algebra. \\
\textbf{\indent Returns:\ }
a list of all the non\texttt{\symbol{45}}isomorphic indecomposable projective
representations over \mbox{\texttt{\mdseries\slshape A}}. 



 The function checks if the algebra \mbox{\texttt{\mdseries\slshape A}} is a finite dimensional (quotient of a) path algebra, and returns an error
message otherwise. }

 

\subsection{\textcolor{Chapter }{SimpleModules}}
\logpage{[ 6, 5, 6 ]}\nobreak
\hyperdef{L}{X7C61261F7C5E53B8}{}
{\noindent\textcolor{FuncColor}{$\triangleright$\enspace\texttt{SimpleModules({\mdseries\slshape A})\index{SimpleModules@\texttt{SimpleModules}}
\label{SimpleModules}
}\hfill{\scriptsize (attribute)}}\\


 Arguments: \mbox{\texttt{\mdseries\slshape A}} \texttt{\symbol{45}}\texttt{\symbol{45}} a finite dimensional (admissible
quotient of a) path algebra. \\
\textbf{\indent Returns:\ }
a list of all the simple representations over \mbox{\texttt{\mdseries\slshape A}} . 



 The function checks if the algebra \mbox{\texttt{\mdseries\slshape A}} is a finite dimensional (admissible quotient of a) path algebra, and returns
an error message otherwise. }

 

\subsection{\textcolor{Chapter }{ZeroModule}}
\logpage{[ 6, 5, 7 ]}\nobreak
\hyperdef{L}{X7CCE2D12807AA35A}{}
{\noindent\textcolor{FuncColor}{$\triangleright$\enspace\texttt{ZeroModule({\mdseries\slshape A})\index{ZeroModule@\texttt{ZeroModule}}
\label{ZeroModule}
}\hfill{\scriptsize (attribute)}}\\


 Arguments: \mbox{\texttt{\mdseries\slshape A}} \texttt{\symbol{45}}\texttt{\symbol{45}} a finite dimensional (quotient of a)
path algebra. \\
\textbf{\indent Returns:\ }
the zero representation over \mbox{\texttt{\mdseries\slshape A}}. 



 The function checks if the algebra \mbox{\texttt{\mdseries\slshape A}} is a finite dimensional (quotient of a) path algebra, and returns an error
message otherwise. }

 }

 
\section{\textcolor{Chapter }{Functors on representations}}\logpage{[ 6, 6, 0 ]}
\hyperdef{L}{X7D99BF5A87DDC099}{}
{
 

\subsection{\textcolor{Chapter }{DualOfModule}}
\logpage{[ 6, 6, 1 ]}\nobreak
\hyperdef{L}{X82D7B50A7ACA47BF}{}
{\noindent\textcolor{FuncColor}{$\triangleright$\enspace\texttt{DualOfModule({\mdseries\slshape M})\index{DualOfModule@\texttt{DualOfModule}}
\label{DualOfModule}
}\hfill{\scriptsize (attribute)}}\\


 Arguments: \mbox{\texttt{\mdseries\slshape M}} \texttt{\symbol{45}}\texttt{\symbol{45}} a \texttt{PathAlgebraMatModule}. \\
\textbf{\indent Returns:\ }
 the dual of \mbox{\texttt{\mdseries\slshape M}} over the opposite algebra $A^\op$, if \mbox{\texttt{\mdseries\slshape M}} is a module over $A$. 

}

 

\subsection{\textcolor{Chapter }{DualOfModuleHomomorphism}}
\logpage{[ 6, 6, 2 ]}\nobreak
\hyperdef{L}{X847DDC417BFB8515}{}
{\noindent\textcolor{FuncColor}{$\triangleright$\enspace\texttt{DualOfModuleHomomorphism({\mdseries\slshape f})\index{DualOfModuleHomomorphism@\texttt{DualOfModuleHomomorphism}}
\label{DualOfModuleHomomorphism}
}\hfill{\scriptsize (attribute)}}\\


 Arguments: \mbox{\texttt{\mdseries\slshape f}} \texttt{\symbol{45}}\texttt{\symbol{45}} a map between two representations \mbox{\texttt{\mdseries\slshape M}} and \mbox{\texttt{\mdseries\slshape N}} over a path algebra $A$. \\
\textbf{\indent Returns:\ }
 the dual of this map over the opposite path algebra \mbox{\texttt{\mdseries\slshape A\texttt{\symbol{94}}\texttt{\symbol{92}}op}}. 

}

 

\subsection{\textcolor{Chapter }{DTr}}
\logpage{[ 6, 6, 3 ]}\nobreak
\hyperdef{L}{X82D31F887C14E921}{}
{\noindent\textcolor{FuncColor}{$\triangleright$\enspace\texttt{DTr({\mdseries\slshape M[, n], or, f[, n]})\index{DTr@\texttt{DTr}}
\label{DTr}
}\hfill{\scriptsize (operation)}}\\
\noindent\textcolor{FuncColor}{$\triangleright$\enspace\texttt{DualOfTranspose({\mdseries\slshape M[, n]})\index{DualOfTranspose@\texttt{DualOfTranspose}}
\label{DualOfTranspose}
}\hfill{\scriptsize (operation)}}\\


 Arguments: (first version) \mbox{\texttt{\mdseries\slshape M}} \texttt{\symbol{45}}\texttt{\symbol{45}} a path algebra module, (optional) \mbox{\texttt{\mdseries\slshape n}} \texttt{\symbol{45}}\texttt{\symbol{45}} an integer. \\
 (second version) \mbox{\texttt{\mdseries\slshape f}} \texttt{\symbol{45}}\texttt{\symbol{45}} a path algebra module homomorphism,
(optional \mbox{\texttt{\mdseries\slshape n}} \texttt{\symbol{45}}\texttt{\symbol{45}} an integer.\\
 (third version) \mbox{\texttt{\mdseries\slshape M}} \texttt{\symbol{45}}\texttt{\symbol{45}} a path algebra modules, (optional) \mbox{\texttt{\mdseries\slshape n}} \texttt{\symbol{45}}\texttt{\symbol{45}} an integer. \\
 \textbf{\indent Returns:\ }
(first and third version) the dual of the transpose of \mbox{\texttt{\mdseries\slshape M}} when called with only one argument, while it returns the dual of the transpose
applied to \mbox{\texttt{\mdseries\slshape M}} \mbox{\texttt{\mdseries\slshape n}} times otherwise. If \mbox{\texttt{\mdseries\slshape n}} is negative, then powers of \texttt{TrD} are computed. \texttt{DualOfTranspose} is a synonym for \texttt{DTr}. 

\textbf{\indent Returns:\ }
(second version) the \mbox{\texttt{\mdseries\slshape n}}\texttt{\symbol{45}}th power of the dual of the transpose of the homomorphism \mbox{\texttt{\mdseries\slshape f}}. 

}

 

\subsection{\textcolor{Chapter }{NakayamaFunctorOfModule}}
\logpage{[ 6, 6, 4 ]}\nobreak
\hyperdef{L}{X82ACC83D7EF5B32C}{}
{\noindent\textcolor{FuncColor}{$\triangleright$\enspace\texttt{NakayamaFunctorOfModule({\mdseries\slshape M})\index{NakayamaFunctorOfModule@\texttt{NakayamaFunctorOfModule}}
\label{NakayamaFunctorOfModule}
}\hfill{\scriptsize (attribute)}}\\


 Arguments: \mbox{\texttt{\mdseries\slshape M}} \texttt{\symbol{45}}\texttt{\symbol{45}} a \texttt{PathAlgebraMatModule}. \\
\textbf{\indent Returns:\ }
 the module $\Hom_K(\Hom_A(M,A), K)$ over $A$, when \mbox{\texttt{\mdseries\slshape M}} is a module over a $K$\texttt{\symbol{45}}algebra $A$. 

}

 

\subsection{\textcolor{Chapter }{NakayamaFunctorOfModuleHomomorphism}}
\logpage{[ 6, 6, 5 ]}\nobreak
\hyperdef{L}{X865C241C86D6168F}{}
{\noindent\textcolor{FuncColor}{$\triangleright$\enspace\texttt{NakayamaFunctorOfModuleHomomorphism({\mdseries\slshape f})\index{NakayamaFunctorOfModuleHomomorphism@\texttt{NakayamaFunctorOfModuleHomomorphism}}
\label{NakayamaFunctorOfModuleHomomorphism}
}\hfill{\scriptsize (attribute)}}\\


 Arguments: \mbox{\texttt{\mdseries\slshape f}} \texttt{\symbol{45}}\texttt{\symbol{45}} a map between two modules \mbox{\texttt{\mdseries\slshape M}} and \mbox{\texttt{\mdseries\slshape N}} over a path algebra $A$. \\
 \textbf{\indent Returns:\ }
 the homomorphism induced by $f\colon M\to N$ from the module $\Hom_K(\Hom_A(M,A),K)$ to $\Hom_K(\Hom_A(N,A), K)$, when \mbox{\texttt{\mdseries\slshape f}} is a module homomorphism over a $K$\texttt{\symbol{45}}algebra \mbox{\texttt{\mdseries\slshape A}}. 

}

 

\subsection{\textcolor{Chapter }{OppositeNakayamaFunctorOfModule}}
\logpage{[ 6, 6, 6 ]}\nobreak
\hyperdef{L}{X78E724307D9FE41D}{}
{\noindent\textcolor{FuncColor}{$\triangleright$\enspace\texttt{OppositeNakayamaFunctorOfModule({\mdseries\slshape M})\index{OppositeNakayamaFunctorOfModule@\texttt{OppositeNakayamaFunctorOfModule}}
\label{OppositeNakayamaFunctorOfModule}
}\hfill{\scriptsize (attribute)}}\\


 Arguments: \mbox{\texttt{\mdseries\slshape M}} \texttt{\symbol{45}}\texttt{\symbol{45}} a \texttt{PathAlgebraMatModule}. \\
\textbf{\indent Returns:\ }
 the module $\Hom_A(\Hom_K(M,K), A)$ over $A$, when \mbox{\texttt{\mdseries\slshape M}} is a module over a $K$\texttt{\symbol{45}}algebra $A$. 

}

 

\subsection{\textcolor{Chapter }{OppositeNakayamaFunctorOfModuleHomomorphism}}
\logpage{[ 6, 6, 7 ]}\nobreak
\hyperdef{L}{X82A057C5824917FA}{}
{\noindent\textcolor{FuncColor}{$\triangleright$\enspace\texttt{OppositeNakayamaFunctorOfModuleHomomorphism({\mdseries\slshape f})\index{OppositeNakayamaFunctorOfModuleHomomorphism@\texttt{Opposite}\-\texttt{Nakayama}\-\texttt{Functor}\-\texttt{Of}\-\texttt{Module}\-\texttt{Homomorphism}}
\label{OppositeNakayamaFunctorOfModuleHomomorphism}
}\hfill{\scriptsize (attribute)}}\\


 Arguments: \mbox{\texttt{\mdseries\slshape f}} \texttt{\symbol{45}}\texttt{\symbol{45}} a map between two modules \mbox{\texttt{\mdseries\slshape M}} and \mbox{\texttt{\mdseries\slshape N}} over a path algebra $A$. \\
 \textbf{\indent Returns:\ }
 the homomorphism induced by $f\colon M\to N$ from the module $\Hom_A(\Hom_K(M,K),A)$ to $\Hom_A(\Hom_K(N,K), A)$, when \mbox{\texttt{\mdseries\slshape f}} is a module homomorphism over a $K$\texttt{\symbol{45}}algebra \mbox{\texttt{\mdseries\slshape A}}. 

}

 

\subsection{\textcolor{Chapter }{RestrictionViaAlgebraHomomorphism}}
\logpage{[ 6, 6, 8 ]}\nobreak
\hyperdef{L}{X818DD1A67A5C03AB}{}
{\noindent\textcolor{FuncColor}{$\triangleright$\enspace\texttt{RestrictionViaAlgebraHomomorphism({\mdseries\slshape f, M})\index{RestrictionViaAlgebraHomomorphism@\texttt{RestrictionViaAlgebraHomomorphism}}
\label{RestrictionViaAlgebraHomomorphism}
}\hfill{\scriptsize (operation)}}\\


 Arguments: \mbox{\texttt{\mdseries\slshape f}} \texttt{\symbol{45}}\texttt{\symbol{45}} an IsAlgebraHomomorphism, \mbox{\texttt{\mdseries\slshape M}} \texttt{\symbol{45}}\texttt{\symbol{45}} an IsPathAlgebraMatModule. \\
\textbf{\indent Returns:\ }
Given an algebra homomorphism $f \colon A \hookrightarrow B$ and a module \mbox{\texttt{\mdseries\slshape M}} over $B$, this function returns \mbox{\texttt{\mdseries\slshape M}} as a module over $A$. 

}

 

\subsection{\textcolor{Chapter }{RestrictionViaAlgebraHomomorphismMap}}
\logpage{[ 6, 6, 9 ]}\nobreak
\hyperdef{L}{X86267BD982DB2221}{}
{\noindent\textcolor{FuncColor}{$\triangleright$\enspace\texttt{RestrictionViaAlgebraHomomorphismMap({\mdseries\slshape f, h})\index{RestrictionViaAlgebraHomomorphismMap@\texttt{Restriction}\-\texttt{Via}\-\texttt{Algebra}\-\texttt{HomomorphismMap}}
\label{RestrictionViaAlgebraHomomorphismMap}
}\hfill{\scriptsize (operation)}}\\


 Arguments: \mbox{\texttt{\mdseries\slshape f}} \texttt{\symbol{45}}\texttt{\symbol{45}} an IsAlgebraHomomorphism, \mbox{\texttt{\mdseries\slshape h}} \texttt{\symbol{45}}\texttt{\symbol{45}} an
IsPathAlgebraMatModuleHomomorphism. \\
\textbf{\indent Returns:\ }
Given an algebra homomorphism $f \colon A \hookrightarrow B$ and a homomorphism of modules \mbox{\texttt{\mdseries\slshape h}} from $M$ to $N$ over $B$, this function returns the induced homomorphism induced by \mbox{\texttt{\mdseries\slshape h}} as a homomorphism over $A$. 

}

 

\subsection{\textcolor{Chapter }{StarOfModule}}
\logpage{[ 6, 6, 10 ]}\nobreak
\hyperdef{L}{X7F07712F786AAEFD}{}
{\noindent\textcolor{FuncColor}{$\triangleright$\enspace\texttt{StarOfModule({\mdseries\slshape M})\index{StarOfModule@\texttt{StarOfModule}}
\label{StarOfModule}
}\hfill{\scriptsize (attribute)}}\\


 Arguments: \mbox{\texttt{\mdseries\slshape M}} \texttt{\symbol{45}}\texttt{\symbol{45}} a \texttt{PathAlgebraMatModule}. \\
\textbf{\indent Returns:\ }
 the module $\Hom_A(M,A)$ over the opposite of $A$, when \mbox{\texttt{\mdseries\slshape M}} is a module over an algebra $A$. 

}

 

\subsection{\textcolor{Chapter }{StarOfModuleHomomorphism}}
\logpage{[ 6, 6, 11 ]}\nobreak
\hyperdef{L}{X80AF678D795B6C57}{}
{\noindent\textcolor{FuncColor}{$\triangleright$\enspace\texttt{StarOfModuleHomomorphism({\mdseries\slshape f})\index{StarOfModuleHomomorphism@\texttt{StarOfModuleHomomorphism}}
\label{StarOfModuleHomomorphism}
}\hfill{\scriptsize (attribute)}}\\


 Arguments: \mbox{\texttt{\mdseries\slshape f}} \texttt{\symbol{45}}\texttt{\symbol{45}} a map between two modules \mbox{\texttt{\mdseries\slshape M}} and \mbox{\texttt{\mdseries\slshape N}} over a path algebra $A$. \\
 \textbf{\indent Returns:\ }
 the homomorphism induced by $f\colon M\to N$ from the module $\Hom_A(N,A)$ to $\Hom_A(M,A)$, when \mbox{\texttt{\mdseries\slshape f}} is a module homomorphism over an algebra \mbox{\texttt{\mdseries\slshape A}}. 

}

 

\subsection{\textcolor{Chapter }{TensorProductOfModules}}
\logpage{[ 6, 6, 12 ]}\nobreak
\hyperdef{L}{X814EE88D8474A99D}{}
{\noindent\textcolor{FuncColor}{$\triangleright$\enspace\texttt{TensorProductOfModules({\mdseries\slshape M, N})\index{TensorProductOfModules@\texttt{TensorProductOfModules}}
\label{TensorProductOfModules}
}\hfill{\scriptsize (operation)}}\\


 Arguments: \mbox{\texttt{\mdseries\slshape M}}, \mbox{\texttt{\mdseries\slshape N}} \texttt{\symbol{45}}\texttt{\symbol{45}} two path algebra modules\\
\textbf{\indent Returns:\ }
the tensor product $M\otimes_A N$ as a vector space and a function $M\times N \to M\otimes_A N$, given two representations \mbox{\texttt{\mdseries\slshape M}} and \mbox{\texttt{\mdseries\slshape N}}, where \mbox{\texttt{\mdseries\slshape M}} is a right module over $A$ and \mbox{\texttt{\mdseries\slshape N}} is a right module over the opposite of $A$. 

}

 

\subsection{\textcolor{Chapter }{TrD}}
\logpage{[ 6, 6, 13 ]}\nobreak
\hyperdef{L}{X7939949279208FA3}{}
{\noindent\textcolor{FuncColor}{$\triangleright$\enspace\texttt{TrD({\mdseries\slshape M[, n], or, f[, n]})\index{TrD@\texttt{TrD}}
\label{TrD}
}\hfill{\scriptsize (operation)}}\\
\noindent\textcolor{FuncColor}{$\triangleright$\enspace\texttt{TransposeOfDual({\mdseries\slshape M[, n]})\index{TransposeOfDual@\texttt{TransposeOfDual}}
\label{TransposeOfDual}
}\hfill{\scriptsize (operation)}}\\


 Arguments: (first version and third version) \mbox{\texttt{\mdseries\slshape M}} \texttt{\symbol{45}}\texttt{\symbol{45}} a path algebra module, (optional) \mbox{\texttt{\mdseries\slshape n}} \texttt{\symbol{45}}\texttt{\symbol{45}} an integer. \\
 (second version) \mbox{\texttt{\mdseries\slshape f}} \texttt{\symbol{45}}\texttt{\symbol{45}} a path algebra module homomorphism, \mbox{\texttt{\mdseries\slshape n}} \texttt{\symbol{45}}\texttt{\symbol{45}} an integer.\\
 \textbf{\indent Returns:\ }
(first and third version) the transpose of the dual of \mbox{\texttt{\mdseries\slshape M}} when called with only one argument, while it returns the transpose of the dual
applied to \mbox{\texttt{\mdseries\slshape M}} \mbox{\texttt{\mdseries\slshape n}} times otherwise. If \mbox{\texttt{\mdseries\slshape n}} is negative, then powers of \texttt{TrD} are computed. \texttt{TransposeOfDual} is a synonym for \texttt{TrD}.

\textbf{\indent Returns:\ }
(second version) the \mbox{\texttt{\mdseries\slshape n}}\texttt{\symbol{45}}th power of the transpose of the dual of the homomorphism \mbox{\texttt{\mdseries\slshape f}}. 

}

 

\subsection{\textcolor{Chapter }{TransposeOfModule}}
\logpage{[ 6, 6, 14 ]}\nobreak
\hyperdef{L}{X79C0B620842128AF}{}
{\noindent\textcolor{FuncColor}{$\triangleright$\enspace\texttt{TransposeOfModule({\mdseries\slshape M})\index{TransposeOfModule@\texttt{TransposeOfModule}}
\label{TransposeOfModule}
}\hfill{\scriptsize (attribute)}}\\


 Arguments: \mbox{\texttt{\mdseries\slshape M}} \texttt{\symbol{45}}\texttt{\symbol{45}} a path algebra module. \\
\textbf{\indent Returns:\ }
the transpose of the module \mbox{\texttt{\mdseries\slshape M}}.

}

 

\subsection{\textcolor{Chapter }{TransposeOfModuleHomomorphism}}
\logpage{[ 6, 6, 15 ]}\nobreak
\hyperdef{L}{X80C4C3BF80B39D66}{}
{\noindent\textcolor{FuncColor}{$\triangleright$\enspace\texttt{TransposeOfModuleHomomorphism({\mdseries\slshape f})\index{TransposeOfModuleHomomorphism@\texttt{TransposeOfModuleHomomorphism}}
\label{TransposeOfModuleHomomorphism}
}\hfill{\scriptsize (attribute)}}\\


 Arguments: \mbox{\texttt{\mdseries\slshape f}} \texttt{\symbol{45}}\texttt{\symbol{45}} a path algebra module homomorphism. \\
\textbf{\indent Returns:\ }
the transpose of the module homomorphism \mbox{\texttt{\mdseries\slshape f}}.

}

 }

 
\section{\textcolor{Chapter }{Vertex projective modules and submodules thereof}}\label{vertexproj}
\logpage{[ 6, 7, 0 ]}
\hyperdef{L}{X84F07A1579CBC26A}{}
{
 In general, if $R$ is a ring and $e$ is an idempotent of $R$, then $eR$ is a projective module of $R$. Then we can form a direct sum of these projective modules together to form
larger projective module. One can construct more general modules by providing
a \mbox{\texttt{\mdseries\slshape vertex projective presentation}}. In this case, $M$ is the cokernel as given by the following exact sequence: $\oplus_{j=1}^{r} w(j)R \rightarrow \oplus_{i=1}^{g} v(i)R \rightarrow{M}
\rightarrow 0$ for some map between $\oplus_{j=1}^{r} w(j)R$ and $\oplus_{i=1}^{g} v(i)R$. The maps $w$ and $v$ map the integers to some idempotent in $R$. 

\subsection{\textcolor{Chapter }{RightProjectiveModule}}
\logpage{[ 6, 7, 1 ]}\nobreak
\hyperdef{L}{X79175B097A0718FE}{}
{\noindent\textcolor{FuncColor}{$\triangleright$\enspace\texttt{RightProjectiveModule({\mdseries\slshape A, verts})\index{RightProjectiveModule@\texttt{RightProjectiveModule}}
\label{RightProjectiveModule}
}\hfill{\scriptsize (function)}}\\


 Arguments: \mbox{\texttt{\mdseries\slshape A}} \texttt{\symbol{45}}\texttt{\symbol{45}} a (quotient of a) path algebra, \mbox{\texttt{\mdseries\slshape verts}} \texttt{\symbol{45}}\texttt{\symbol{45}} a list of vertices. \\
\textbf{\indent Returns:\ }
 the right projective module over \mbox{\texttt{\mdseries\slshape A}} which is the direct sum of projective modules of the form \mbox{\texttt{\mdseries\slshape vA}} where the vertices are taken from \mbox{\texttt{\mdseries\slshape verts}}.



 The module created is in the category \texttt{IsPathAlgebraModule}. In this implementation the algebra can be a quotient of a path algebra. So
if the list was $[v,w]$ then the module created will be the direct sum $vA \oplus wA$, in that order. Elements of the modules are vectors of algebra elements, and
in each component, each path begins with the vertex in that position in the
list of vertices. Right projective modules are implemented as algebra modules
(see the main GAP\texttt{\symbol{45}}manual concerning "representations of
algebras") and all operations for algebra modules are applicable to right
projective modules. In particular, one can construct submodules using
GAP\texttt{\symbol{45}}command "SubAlgebraModule". }

 Here we create the right projective module $P = vA \oplus vA \oplus wA$. 
\begin{Verbatim}[commandchars=!@|,fontsize=\small,frame=single,label=Example]
  !gapprompt@gap>| !gapinput@F := GF(11);|
  GF(11)
  !gapprompt@gap>| !gapinput@Q := Quiver(["v","w", "x"],[["v","w","a"],["v","w","b"],|
  !gapprompt@>| !gapinput@["w","x","c"]]);|
  <quiver with 3 vertices and 3 arrows>
  !gapprompt@gap>| !gapinput@A := PathAlgebra(F,Q);|
  <GF(11)[<quiver with 3 vertices and 3 arrows>]>
  !gapprompt@gap>| !gapinput@P := RightProjectiveModule(A,[A.v,A.v,A.w]);|
  <right-module over <GF(11)[<quiver with 3 vertices and 3 arrows>]>>
  !gapprompt@gap>| !gapinput@Dimension(P);|
  12 
\end{Verbatim}
 

\subsection{\textcolor{Chapter }{CompletelyReduceGroebnerBasisForModule}}
\logpage{[ 6, 7, 2 ]}\nobreak
\hyperdef{L}{X855EFF36842AA3AE}{}
{\noindent\textcolor{FuncColor}{$\triangleright$\enspace\texttt{CompletelyReduceGroebnerBasisForModule({\mdseries\slshape GB})\index{CompletelyReduceGroebnerBasisForModule@\texttt{Completely}\-\texttt{Reduce}\-\texttt{Groebner}\-\texttt{Basis}\-\texttt{For}\-\texttt{Module}}
\label{CompletelyReduceGroebnerBasisForModule}
}\hfill{\scriptsize (function)}}\\


 Arguments: \mbox{\texttt{\mdseries\slshape GB}} \texttt{\symbol{45}}\texttt{\symbol{45}} an right Groebner basis for a
(submodule of a) vertex projective module over a path algebra. \\
\textbf{\indent Returns:\ }
a completely reduced right Groebner basis from the entered Groebner basis \mbox{\texttt{\mdseries\slshape GB}}.



 This function takes as input an right Groebner basis for a vertex projective
module or a submodule thereof, an constructs completely reduced right Groebner
from it. }

 

\subsection{\textcolor{Chapter }{IsLeftDivisible}}
\logpage{[ 6, 7, 3 ]}\nobreak
\hyperdef{L}{X7D931FBF7BF64C7C}{}
{\noindent\textcolor{FuncColor}{$\triangleright$\enspace\texttt{IsLeftDivisible({\mdseries\slshape x, y})\index{IsLeftDivisible@\texttt{IsLeftDivisible}}
\label{IsLeftDivisible}
}\hfill{\scriptsize (property)}}\\


 Arguments: \mbox{\texttt{\mdseries\slshape x, y}} \texttt{\symbol{45}}\texttt{\symbol{45}} two path algebra vectors. \\
\textbf{\indent Returns:\ }
true if the tip of \mbox{\texttt{\mdseries\slshape y}} left divides the tip of \mbox{\texttt{\mdseries\slshape x}}. False otherwise.



 Given two PathAlgebraVectors \mbox{\texttt{\mdseries\slshape x}} and \mbox{\texttt{\mdseries\slshape y}}, then \mbox{\texttt{\mdseries\slshape y}} is said to left divide \mbox{\texttt{\mdseries\slshape x}}, if the tip of \mbox{\texttt{\mdseries\slshape x}} and the tip of \mbox{\texttt{\mdseries\slshape y}} occur in the same coordinate, and the tipmonomial of the tip of \mbox{\texttt{\mdseries\slshape y}} leftdivides the tipmonomial of the tip of \mbox{\texttt{\mdseries\slshape x}}. }

 

\subsection{\textcolor{Chapter }{IsPathAlgebraModule}}
\logpage{[ 6, 7, 4 ]}\nobreak
\hyperdef{L}{X82A8398478788A5A}{}
{\noindent\textcolor{FuncColor}{$\triangleright$\enspace\texttt{IsPathAlgebraModule({\mdseries\slshape P})\index{IsPathAlgebraModule@\texttt{IsPathAlgebraModule}}
\label{IsPathAlgebraModule}
}\hfill{\scriptsize (property)}}\\


 Arguments: \mbox{\texttt{\mdseries\slshape P}} \texttt{\symbol{45}}\texttt{\symbol{45}} any object. \\
\textbf{\indent Returns:\ }
true if the argument \mbox{\texttt{\mdseries\slshape P}} is in the category \texttt{IsPathAlgebraModule}.

}

 

\subsection{\textcolor{Chapter }{IsPathAlgebraVector}}
\logpage{[ 6, 7, 5 ]}\nobreak
\hyperdef{L}{X83DC9F63800A8812}{}
{\noindent\textcolor{FuncColor}{$\triangleright$\enspace\texttt{IsPathAlgebraVector({\mdseries\slshape v})\index{IsPathAlgebraVector@\texttt{IsPathAlgebraVector}}
\label{IsPathAlgebraVector}
}\hfill{\scriptsize (property)}}\\


 Arguments: \mbox{\texttt{\mdseries\slshape v}} \texttt{\symbol{45}}\texttt{\symbol{45}} a path algebra vector. \\
\textbf{\indent Returns:\ }
true if \mbox{\texttt{\mdseries\slshape v}} has been constructed as a PathAlgebraVector. Otherwise it returns false.

}

 

\subsection{\textcolor{Chapter }{LeadingCoefficient (of PathAlgebraVector)}}
\logpage{[ 6, 7, 6 ]}\nobreak
\hyperdef{L}{X7935C4407BCB6F38}{}
{\noindent\textcolor{FuncColor}{$\triangleright$\enspace\texttt{LeadingCoefficient (of PathAlgebraVector)({\mdseries\slshape x})\index{LeadingCoefficient (of PathAlgebraVector)@\texttt{Leading}\-\texttt{Coefficient (of }\-\texttt{Path}\-\texttt{Algebra}\-\texttt{Vector)}}
\label{LeadingCoefficient (of PathAlgebraVector)}
}\hfill{\scriptsize (operation)}}\\


 Arguments: \mbox{\texttt{\mdseries\slshape x}} \texttt{\symbol{45}}\texttt{\symbol{45}} an element in a PathAlgebraModule. \\
\textbf{\indent Returns:\ }
the coefficient of the leading term/tip of a \texttt{PathAlgebraVector}.



 The tip of the element \mbox{\texttt{\mdseries\slshape x}} can by found by applying the command \texttt{LeadingTerm (of PathAlgebraVector)} (\ref{LeadingTerm (of PathAlgebraVector)}). }

 

\subsection{\textcolor{Chapter }{LeadingComponent}}
\logpage{[ 6, 7, 7 ]}\nobreak
\hyperdef{L}{X87DE0FE482A17ECC}{}
{\noindent\textcolor{FuncColor}{$\triangleright$\enspace\texttt{LeadingComponent({\mdseries\slshape v})\index{LeadingComponent@\texttt{LeadingComponent}}
\label{LeadingComponent}
}\hfill{\scriptsize (operation)}}\\


 Arguments: \mbox{\texttt{\mdseries\slshape v}} \texttt{\symbol{45}}\texttt{\symbol{45}} a path algebra vector. \\
\textbf{\indent Returns:\ }
\texttt{v[pos]}, where \texttt{pos} is the coordinate for the tip of the vector, whenever \mbox{\texttt{\mdseries\slshape v}} is non\texttt{\symbol{45}}zero. That is, it returns the coordinate of the
vector \mbox{\texttt{\mdseries\slshape v}} where the tip occors. It returns zero otherwise.

}

 

\subsection{\textcolor{Chapter }{LeadingPosition}}
\logpage{[ 6, 7, 8 ]}\nobreak
\hyperdef{L}{X839C070B7BAAE5DC}{}
{\noindent\textcolor{FuncColor}{$\triangleright$\enspace\texttt{LeadingPosition({\mdseries\slshape v})\index{LeadingPosition@\texttt{LeadingPosition}}
\label{LeadingPosition}
}\hfill{\scriptsize (operation)}}\\


 Arguments: \mbox{\texttt{\mdseries\slshape v}} \texttt{\symbol{45}}\texttt{\symbol{45}} a path algebra vector. \\
\textbf{\indent Returns:\ }
the coordinate in which the tip of the vector occurs.

}

 

\subsection{\textcolor{Chapter }{LeadingTerm (of PathAlgebraVector)}}
\logpage{[ 6, 7, 9 ]}\nobreak
\hyperdef{L}{X7E705C5C825D9187}{}
{\noindent\textcolor{FuncColor}{$\triangleright$\enspace\texttt{LeadingTerm (of PathAlgebraVector)({\mdseries\slshape x})\index{LeadingTerm (of PathAlgebraVector)@\texttt{LeadingTerm (of PathAlgebraVector)}}
\label{LeadingTerm (of PathAlgebraVector)}
}\hfill{\scriptsize (operation)}}\\


 Arguments: \mbox{\texttt{\mdseries\slshape x}} \texttt{\symbol{45}}\texttt{\symbol{45}} an element in a PathAlgebraModule. \\
\textbf{\indent Returns:\ }
the leading term/tip of a \texttt{PathAlgebraVector}.



 The tip of the element \mbox{\texttt{\mdseries\slshape x}} is computed using the following order: the tip is computed for each
coordinate, if the largest of these occur as a tip of several coordinates,
then the coordinate with the smallest index from 1 to the length of vector is
chosen. The position of the tip was computed when the \texttt{PathAlgebraVector} was created. }

 

\subsection{\textcolor{Chapter }{LeftDivision}}
\logpage{[ 6, 7, 10 ]}\nobreak
\hyperdef{L}{X8593BCDB8402C46C}{}
{\noindent\textcolor{FuncColor}{$\triangleright$\enspace\texttt{LeftDivision({\mdseries\slshape x, y})\index{LeftDivision@\texttt{LeftDivision}}
\label{LeftDivision}
}\hfill{\scriptsize (operation)}}\\


 Arguments: \mbox{\texttt{\mdseries\slshape x, y}} \texttt{\symbol{45}}\texttt{\symbol{45}} two path algebra vectors. \\
\textbf{\indent Returns:\ }
a scalar multiple of a path, say $\lambda$ such that the tips of $y*\lambda$ and $x$ are the same, if the tip of \mbox{\texttt{\mdseries\slshape y}} left divides the tip of \mbox{\texttt{\mdseries\slshape x}}. False otherwise.

}

 In the following example, we create two elements in $P$, perform some elementwise operations, and then construct a submodule using
the two elements as generators.

 
\begin{Verbatim}[commandchars=!@|,fontsize=\small,frame=single,label=Example]
  !gapprompt@gap>| !gapinput@p1 := Vectorize(P,[A.b*A.c,A.a*A.c,A.c]);|
  [ (Z(11)^0)*b*c, (Z(11)^0)*a*c, (Z(11)^0)*c ]
  !gapprompt@gap>| !gapinput@p2 := Vectorize(P,[A.a,A.b,A.w]);|
  [ (Z(11)^0)*a, (Z(11)^0)*b, (Z(11)^0)*w ]
  !gapprompt@gap>| !gapinput@2*p1 + p2;|
  [ (Z(11)^0)*a+(Z(11))*b*c, (Z(11)^0)*b+(Z(11))*a*c, 
    (Z(11)^0)*w+(Z(11))*c ]
  !gapprompt@gap>| !gapinput@S := SubAlgebraModule(P,[p1,p2]);|
  <right-module over <GF(11)[<quiver with 3 vertices and 3 arrows>]>>
  !gapprompt@gap>| !gapinput@Dimension(S);|
  3 
\end{Verbatim}
 

\subsection{\textcolor{Chapter }{\texttt{\symbol{94}} (a PathAlgebraModule element and a PathAlgebra element)}}
\logpage{[ 6, 7, 11 ]}\nobreak
\hyperdef{L}{X8749643879D32A01}{}
{\noindent\textcolor{FuncColor}{$\triangleright$\enspace\texttt{\texttt{\symbol{94}}({\mdseries\slshape m, a})\index{\texttt{\symbol{94}}@\texttt{\texttt{\symbol{94}}}!a PathAlgebraModule element and a PathAlgebra element}
\label{^:a PathAlgebraModule element and a PathAlgebra element}
}\hfill{\scriptsize (operation)}}\\


 Arguments: \mbox{\texttt{\mdseries\slshape m}} \texttt{\symbol{45}}\texttt{\symbol{45}} an element of a path algebra module, \mbox{\texttt{\mdseries\slshape a}} \texttt{\symbol{45}}\texttt{\symbol{45}} an element of a path algebra. \\
\textbf{\indent Returns:\ }
 the element \mbox{\texttt{\mdseries\slshape m}} multiplied with \mbox{\texttt{\mdseries\slshape a}}. 



 This action is defined by multiplying each component in \mbox{\texttt{\mdseries\slshape m}} by \mbox{\texttt{\mdseries\slshape a}} on the right. }

 
\begin{Verbatim}[commandchars=!@|,fontsize=\small,frame=single,label=Example]
  !gapprompt@gap>| !gapinput@p2^(A.c - A.w);|
  [ (Z(11)^5)*a+(Z(11)^0)*a*c, (Z(11)^5)*b+(Z(11)^0)*b*c, 
    (Z(11)^5)*w+(Z(11)^0)*c ] 
\end{Verbatim}
 

\subsection{\textcolor{Chapter }{{\textless} (for two elements in a PathAlgebraModule)}}
\logpage{[ 6, 7, 12 ]}\nobreak
\hyperdef{L}{X87B186CE868FCB30}{}
{\noindent\textcolor{FuncColor}{$\triangleright$\enspace\texttt{{\textless}({\mdseries\slshape m1, m2})\index{{\textless}@\texttt{{\textless}}!for two elements in a PathAlgebraModule}
\label{<:for two elements in a PathAlgebraModule}
}\hfill{\scriptsize (operation)}}\\


 Arguments: \mbox{\texttt{\mdseries\slshape m1, m2}} \texttt{\symbol{45}}\texttt{\symbol{45}} two elements of a PathAlgebraModule. \\
\textbf{\indent Returns:\ }
 `true' if \mbox{\texttt{\mdseries\slshape m1}} is less than \mbox{\texttt{\mdseries\slshape m2}} and false otherwise. 



 Elements are compared componentwise from left to right using the ordering of
the underlying algebra. The element \mbox{\texttt{\mdseries\slshape m1}} is less than \mbox{\texttt{\mdseries\slshape m2}} if the first time components are not equal, the component of \mbox{\texttt{\mdseries\slshape m1}} is less than the corresponding component of \mbox{\texttt{\mdseries\slshape m2}}. }

 
\begin{Verbatim}[commandchars=!@|,fontsize=\small,frame=single,label=Example]
  !gapprompt@gap>| !gapinput@p1 < p2;|
  false 
\end{Verbatim}
 

\subsection{\textcolor{Chapter }{/}}
\logpage{[ 6, 7, 13 ]}\nobreak
\hyperdef{L}{X7F51DF007F51DF00}{}
{\noindent\textcolor{FuncColor}{$\triangleright$\enspace\texttt{/({\mdseries\slshape M, N})\index{/@\texttt{/}}
\label{/}
}\hfill{\scriptsize (operation)}}\\


 Arguments: \mbox{\texttt{\mdseries\slshape M, N}} \texttt{\symbol{45}}\texttt{\symbol{45}} two finite dimensional \texttt{PathAlgebraModule}s. \\
\textbf{\indent Returns:\ }
 the factor module $M/N$. 



 This module is again a right algebra module, and all applicable methods and
operations are available for the resulting factor module. Furthermore, the
resulting module is a vector space, so operations for computing bases and
dimensions are also available. }

 
\begin{Verbatim}[commandchars=!@|,fontsize=\small,frame=single,label=Example]
  !gapprompt@gap>| !gapinput@PS := P/S;|
  <9-dimensional right-module over <GF(11)[<quiver with 3 vertices and 
  3 arrows>]>>
  !gapprompt@gap>| !gapinput@Basis(PS);|
  Basis( <9-dimensional right-module over <GF(11)[<quiver with 
  3 vertices and 3 arrows>]>>, 
  [ [ [ <zero> of ..., (Z(11)^0)*v, <zero> of ... ] ], 
    [ [ (Z(11)^0)*v, <zero> of ..., <zero> of ... ] ], 
    [ [ <zero> of ..., <zero> of ..., (Z(11)^0)*w ] ], 
    [ [ <zero> of ..., (Z(11)^0)*a, <zero> of ... ] ], 
    [ [ (Z(11)^0)*a, <zero> of ..., <zero> of ... ] ], 
    [ [ (Z(11)^0)*b, <zero> of ..., <zero> of ... ] ], 
    [ [ <zero> of ..., <zero> of ..., (Z(11)^0)*c ] ], 
    [ [ <zero> of ..., (Z(11)^0)*a*c, <zero> of ... ] ], 
    [ [ (Z(11)^0)*a*c, <zero> of ..., <zero> of ... ] ] ] ) 
\end{Verbatim}
 

\subsection{\textcolor{Chapter }{PathAlgebraVector}}
\logpage{[ 6, 7, 14 ]}\nobreak
\hyperdef{L}{X7B3C31F17E6FB3CD}{}
{\noindent\textcolor{FuncColor}{$\triangleright$\enspace\texttt{PathAlgebraVector({\mdseries\slshape fam, components})\index{PathAlgebraVector@\texttt{PathAlgebraVector}}
\label{PathAlgebraVector}
}\hfill{\scriptsize (operation)}}\\


 Arguments: \mbox{\texttt{\mdseries\slshape fam}} \texttt{\symbol{45}}\texttt{\symbol{45}} a PathAlgebraVectorFamily, \mbox{\texttt{\mdseries\slshape components}} \texttt{\symbol{45}}\texttt{\symbol{45}} a homogeneous list of elements. \\
\textbf{\indent Returns:\ }
a PathAlgebraVector in the PathAlgebraVectorFamily \mbox{\texttt{\mdseries\slshape fam}} with the components of the vector being equal to \mbox{\texttt{\mdseries\slshape components}}. 



 This function is typically used when constructing elements of a module
constructed by the command \texttt{RightProjectiveModule}. If \texttt{P} is constructed as say, \texttt{P := RightProjectiveModule(KQ, [KQ.v1, KQ.v1, KQ.v2])}, then \texttt{ExtRepOfObj(p)}, where \texttt{p} is an element if \texttt{P} is a \texttt{PathAlgebraVector}. The tip is computed using the following ordering: the tip is computed for
each coordinate, if the largest of these occur as a tip of several
coordinates, then the coordinate with the smallest index from 1 to the length
of vector is chosen. }

 

\subsection{\textcolor{Chapter }{ProjectivePathAlgebraPresentation}}
\logpage{[ 6, 7, 15 ]}\nobreak
\hyperdef{L}{X850B83A0801EE970}{}
{\noindent\textcolor{FuncColor}{$\triangleright$\enspace\texttt{ProjectivePathAlgebraPresentation({\mdseries\slshape M})\index{ProjectivePathAlgebraPresentation@\texttt{ProjectivePathAlgebraPresentation}}
\label{ProjectivePathAlgebraPresentation}
}\hfill{\scriptsize (operation)}}\\


 Arguments: \mbox{\texttt{\mdseries\slshape M}} \texttt{\symbol{45}}\texttt{\symbol{45}} a finite dimensional module over a
(quotient of a) path algebra. \\
\textbf{\indent Returns:\ }
a projective presentation of the entered module \mbox{\texttt{\mdseries\slshape M}} over a (qoutient of a) path algebra $A$. The projective presentation, or resolution is over the path algebra form
which $A$ was constructed. 



 This function takes as input a PathAlgebraMatModule and constructs a
projective presentation of this module over the path algebra over which it is
defined, i.e. a projective resolution of length 1. It returns a list of five
elements: (1) a projective module $P$ over the path algebra, which modulo the relations induced the projective cover
of \mbox{\texttt{\mdseries\slshape M}}, (2) a submodule $U$ of $P$ such that $P/U$ is isomorphic to \mbox{\texttt{\mdseries\slshape M}}, (3) module generators of $P$, (4) module generators for $U$ which forms a completely reduced right Groebner basis for $U$, and (5) a matrix with entries in the path algebra which gives the map from $U$ to $P$, if $U$ were considered a direct sum of vertex projective modules over the path
algebra. }

 

\subsection{\textcolor{Chapter }{RightGroebnerBasisOfModule}}
\logpage{[ 6, 7, 16 ]}\nobreak
\hyperdef{L}{X7F5109637D496354}{}
{\noindent\textcolor{FuncColor}{$\triangleright$\enspace\texttt{RightGroebnerBasisOfModule({\mdseries\slshape M})\index{RightGroebnerBasisOfModule@\texttt{RightGroebnerBasisOfModule}}
\label{RightGroebnerBasisOfModule}
}\hfill{\scriptsize (attribute)}}\\


 Arguments: \mbox{\texttt{\mdseries\slshape M}} \texttt{\symbol{45}}\texttt{\symbol{45}} a PathAlgebraModule. \\
\textbf{\indent Returns:\ }
a right Groebner basis for the module \mbox{\texttt{\mdseries\slshape M}}.



 It checks if the acting algebra on the module \mbox{\texttt{\mdseries\slshape M}} is a path algebra, and it returns an error message otherwise. The elements in
the right Groebner basis that is constructed, can be retrieved by the command \texttt{BasisVectors}. The underlying module is likewise returned by the command \texttt{UnderlyingModule}. The output of the function is satisfying the filter/category \texttt{IsRightPathAlgebraModuleGroebnerBasis}. }

 

\subsection{\textcolor{Chapter }{TargetVertex}}
\logpage{[ 6, 7, 17 ]}\nobreak
\hyperdef{L}{X7F5D460187F89CB7}{}
{\noindent\textcolor{FuncColor}{$\triangleright$\enspace\texttt{TargetVertex({\mdseries\slshape v})\index{TargetVertex@\texttt{TargetVertex}}
\label{TargetVertex}
}\hfill{\scriptsize (operation)}}\\


 Arguments: \mbox{\texttt{\mdseries\slshape v}} \texttt{\symbol{45}}\texttt{\symbol{45}} a PathAlgebraVector. \\
\textbf{\indent Returns:\ }
a vertex $w$ such that $v*w = v$, if such a vertex exists, and fail otherwise. 



 Given a PathAlgebraVector \mbox{\texttt{\mdseries\slshape v}}, if \mbox{\texttt{\mdseries\slshape v}} is right uniform, this function finds the vertex $w$ such that $v*w = v$ whenever \mbox{\texttt{\mdseries\slshape v}} is non\texttt{\symbol{45}}zero, and returns the zero path otherwise. If \mbox{\texttt{\mdseries\slshape v}} is not right uniform it returns fail. }

 

\subsection{\textcolor{Chapter }{UniformGeneratorsOfModule}}
\logpage{[ 6, 7, 18 ]}\nobreak
\hyperdef{L}{X7CDBA7818700F9D2}{}
{\noindent\textcolor{FuncColor}{$\triangleright$\enspace\texttt{UniformGeneratorsOfModule({\mdseries\slshape M})\index{UniformGeneratorsOfModule@\texttt{UniformGeneratorsOfModule}}
\label{UniformGeneratorsOfModule}
}\hfill{\scriptsize (attribute)}}\\


 Arguments: \mbox{\texttt{\mdseries\slshape M}} \texttt{\symbol{45}}\texttt{\symbol{45}} a \texttt{PathAlgebraModule}. \\
\textbf{\indent Returns:\ }
a set of right uniform generators of the mdoule \mbox{\texttt{\mdseries\slshape M}}. If \mbox{\texttt{\mdseries\slshape M}} is the zero module, then it returns an empty list. 

}

 

\subsection{\textcolor{Chapter }{Vectorize}}
\logpage{[ 6, 7, 19 ]}\nobreak
\hyperdef{L}{X78E05C8F7ADE2BCD}{}
{\noindent\textcolor{FuncColor}{$\triangleright$\enspace\texttt{Vectorize({\mdseries\slshape M, components})\index{Vectorize@\texttt{Vectorize}}
\label{Vectorize}
}\hfill{\scriptsize (function)}}\\


 Arguments: \mbox{\texttt{\mdseries\slshape M}} \texttt{\symbol{45}}\texttt{\symbol{45}} a module over a path algebra, \mbox{\texttt{\mdseries\slshape components}} \texttt{\symbol{45}}\texttt{\symbol{45}} a list of elements of \mbox{\texttt{\mdseries\slshape M}}. \\
\textbf{\indent Returns:\ }
a vector in \mbox{\texttt{\mdseries\slshape M}} from a list of path algebra elements \mbox{\texttt{\mdseries\slshape components}}, which defines the components in the resulting vector.



 The returned vector is normalized, so the vector's components may not match
the input components. }

 }

 }

 
\chapter{\textcolor{Chapter }{Homomorphisms of Right Modules over Path Algebras}}\label{Homomorphisms}
\logpage{[ 7, 0, 0 ]}
\hyperdef{L}{X7C049EFC82A7CAA7}{}
{
 This chapter describes the categories, representations, attributes, and
operations on homomorphisms between representations of quivers.

 Given two homorphisms $f\colon L\to M$ and $g\colon M\to N$, then the composition is written $f*g$. The elements in the modules or the representations of a quiver are row
vectors. Therefore the homomorphisms between two modules are acting on these
row vectors, that is, if $m_i$ is in $M[i]$ and $g_i\colon M[i]\to N[i]$ represents the linear map, then the value of $g$ applied to $m_i$ is the matrix product $m_i*g_i$.

 The example used throughout this chapter is the following. 
\begin{Verbatim}[commandchars=!@|,fontsize=\small,frame=single,label=Example]
  !gapprompt@gap>| !gapinput@Q := Quiver(3,[[1,2,"a"],[1,2,"b"],[2,2,"c"],[2,3,"d"],[3,1,"e"]]);;|
  !gapprompt@gap>| !gapinput@KQ := PathAlgebra(Rationals, Q);;|
  !gapprompt@gap>| !gapinput@AssignGeneratorVariables(KQ);;|
  !gapprompt@gap>| !gapinput@rels := [d*e,c^2,a*c*d-b*d,e*a];;|
  !gapprompt@gap>| !gapinput@A := KQ/rels;;|
  !gapprompt@gap>| !gapinput@mat :=[["a",[[1,2],[0,3],[1,5]]],["b",[[2,0],[3,0],[5,0]]],|
  !gapprompt@>| !gapinput@["c",[[0,0],[1,0]]],["d",[[1,2],[0,1]]],["e",[[0,0,0],[0,0,0]]]];;|
  !gapprompt@gap>| !gapinput@N := RightModuleOverPathAlgebra(A,mat);; |
\end{Verbatim}
 
\section{\textcolor{Chapter }{Categories and representation of homomorphisms}}\logpage{[ 7, 1, 0 ]}
\hyperdef{L}{X7B18E84678FA5EE0}{}
{
 

\subsection{\textcolor{Chapter }{IsPathAlgebraModuleHomomorphism}}
\logpage{[ 7, 1, 1 ]}\nobreak
\hyperdef{L}{X79AE9787877E0A28}{}
{\noindent\textcolor{FuncColor}{$\triangleright$\enspace\texttt{IsPathAlgebraModuleHomomorphism({\mdseries\slshape f})\index{IsPathAlgebraModuleHomomorphism@\texttt{IsPathAlgebraModuleHomomorphism}}
\label{IsPathAlgebraModuleHomomorphism}
}\hfill{\scriptsize (filter)}}\\


 Arguments: \mbox{\texttt{\mdseries\slshape f}} \texttt{\symbol{45}} any object in GAP.\\
 \textbf{\indent Returns:\ }
 true or false depending on if \mbox{\texttt{\mdseries\slshape f}} belongs to the categories \texttt{IsPathAlgebraModuleHomomorphism}. 



 This defines the category \texttt{IsPathAlgebraModuleHomomorphism}. }

 

\subsection{\textcolor{Chapter }{RightModuleHomOverAlgebra}}
\logpage{[ 7, 1, 2 ]}\nobreak
\hyperdef{L}{X8318ED607FE21F55}{}
{\noindent\textcolor{FuncColor}{$\triangleright$\enspace\texttt{RightModuleHomOverAlgebra({\mdseries\slshape M, N, mats})\index{RightModuleHomOverAlgebra@\texttt{RightModuleHomOverAlgebra}}
\label{RightModuleHomOverAlgebra}
}\hfill{\scriptsize (operation)}}\\


 Arguments: \mbox{\texttt{\mdseries\slshape M}}, \mbox{\texttt{\mdseries\slshape N}} \texttt{\symbol{45}} two modules over the same (quotient of a) path algebra, \mbox{\texttt{\mdseries\slshape mats}} \texttt{\symbol{45}} a list of matrices, one for each vertex in the quiver of
the path algebra.\\
 \textbf{\indent Returns:\ }
a homomorphism in the category \texttt{IsPathAlgebraModuleHomomorphism} from the module \mbox{\texttt{\mdseries\slshape M}} to the module \mbox{\texttt{\mdseries\slshape N}} given by the matrices \mbox{\texttt{\mdseries\slshape mats}}.



 The arguments \mbox{\texttt{\mdseries\slshape M}} and \mbox{\texttt{\mdseries\slshape N}} are two modules over the same algebra (this is checked), and \mbox{\texttt{\mdseries\slshape mats}} is a list of matrices \texttt{mats[i]}, where \texttt{mats[i]} represents the linear map from \texttt{M[i]} to \texttt{N[i]} with \texttt{i} running through all the vertices in the same order as when the underlying
quiver was created. If both \texttt{DimensionVector(M)[i]} and \texttt{DimensionVector(N)[i]} are non\texttt{\symbol{45}}zero, then \texttt{mats[i]} is a \texttt{DimensionVector(M)[i]} by \texttt{DimensionVector(N)[i]} matrix. If \texttt{DimensionVector(M)[i]} is zero and \texttt{DimensionVector(N)[i]} is non\texttt{\symbol{45}}zero, then \texttt{mats[i]} must be the zero \texttt{1} by \texttt{DimensionVector(N)[i]} matrix. Similarly for the other way around. If both \texttt{DimensionVector(M)[i]} and \texttt{DimensionVector(N)[i]} are zero, then \texttt{mats[i]} must be the \texttt{1} by \texttt{1} zero matrix. The function checks if \mbox{\texttt{\mdseries\slshape mats}} is a homomorphism from the module \mbox{\texttt{\mdseries\slshape M}} to the module \mbox{\texttt{\mdseries\slshape N}} by checking that the matrices given in \mbox{\texttt{\mdseries\slshape mats}} have the correct size and satisfy the appropriate commutativity conditions
with the matrices in the modules given by \mbox{\texttt{\mdseries\slshape M}} and \mbox{\texttt{\mdseries\slshape N}}. The source (or domain) and the range (or codomain) of the homomorphism
constructed can by obtained again by \texttt{Range} (\ref{Range}) and by \texttt{Source} (\ref{Source}), respectively. }

 
\begin{Verbatim}[commandchars=!@|,fontsize=\small,frame=single,label=Example]
  !gapprompt@gap>| !gapinput@L := RightModuleOverPathAlgebra(A,[["a",[0,1]],["b",[0,1]],|
  !gapprompt@>| !gapinput@["c",[[0]]],["d",[[1]]],["e",[1,0]]]);|
  <[ 0, 1, 1 ]>
  !gapprompt@gap>| !gapinput@DimensionVector(L);|
  [ 0, 1, 1 ]
  !gapprompt@gap>| !gapinput@f := RightModuleHomOverAlgebra(L,N,[[[0,0,0]], [[1,0]], |
  !gapprompt@>| !gapinput@[[1,2]]]);|
  <<[ 0, 1, 1 ]> ---> <[ 3, 2, 2 ]>>
  
  !gapprompt@gap>| !gapinput@IsPathAlgebraMatModuleHomomorphism(f);|
  true 
\end{Verbatim}
 

\subsection{\textcolor{Chapter }{UnderlyingLinearMap}}
\logpage{[ 7, 1, 3 ]}\nobreak
\hyperdef{L}{X8752DBF3814CC4D5}{}
{\noindent\textcolor{FuncColor}{$\triangleright$\enspace\texttt{UnderlyingLinearMap({\mdseries\slshape f})\index{UnderlyingLinearMap@\texttt{UnderlyingLinearMap}}
\label{UnderlyingLinearMap}
}\hfill{\scriptsize (operation)}}\\


 Arguments: \mbox{\texttt{\mdseries\slshape f}} \texttt{\symbol{45}} a homomorphism between two representations.\\
 \textbf{\indent Returns:\ }
the homomorphism \mbox{\texttt{\mdseries\slshape f}} as a matrix, where the homomorphism on a zero vectorspace is represented by a
zero matrix. The linear map/matrix operates on row vectors.

}

 }

 
\section{\textcolor{Chapter }{Generalities of homomorphisms}}\logpage{[ 7, 2, 0 ]}
\hyperdef{L}{X830529F9800BF688}{}
{
 

\subsection{\textcolor{Chapter }{\texttt{\symbol{92}}= (maps)}}
\logpage{[ 7, 2, 1 ]}\nobreak
\hyperdef{L}{X82F8641E84AD4922}{}
{\noindent\textcolor{FuncColor}{$\triangleright$\enspace\texttt{\texttt{\symbol{92}}= (maps)({\mdseries\slshape f, g})\index{\texttt{\symbol{92}}= (maps)@\texttt{\texttt{\symbol{92}}= (maps)}}
\label{bSlash= (maps)}
}\hfill{\scriptsize (operation)}}\\


 Arguments: \mbox{\texttt{\mdseries\slshape f}}, \mbox{\texttt{\mdseries\slshape g}} \texttt{\symbol{45}} two homomorphisms between two modules.\\
 \textbf{\indent Returns:\ }
true, if \texttt{Source(f) = Source(g)}, \texttt{Range(f) = Range(g)}, and the matrices defining the maps \mbox{\texttt{\mdseries\slshape f}} and \mbox{\texttt{\mdseries\slshape g}} coincide.

}

 

\subsection{\textcolor{Chapter }{\texttt{\symbol{92}}+ (maps)}}
\logpage{[ 7, 2, 2 ]}\nobreak
\hyperdef{L}{X7FA09A067BE00277}{}
{\noindent\textcolor{FuncColor}{$\triangleright$\enspace\texttt{\texttt{\symbol{92}}+ (maps)({\mdseries\slshape f, g})\index{\texttt{\symbol{92}}+ (maps)@\texttt{\texttt{\symbol{92}}+ (maps)}}
\label{bSlash+ (maps)}
}\hfill{\scriptsize (operation)}}\\


 Arguments: \mbox{\texttt{\mdseries\slshape f}}, \mbox{\texttt{\mdseries\slshape g}} \texttt{\symbol{45}} two homomorphisms between two modules.\\
 \textbf{\indent Returns:\ }
the sum \mbox{\texttt{\mdseries\slshape f+g}} of the maps \mbox{\texttt{\mdseries\slshape f}} and \mbox{\texttt{\mdseries\slshape g}}.



The function checks if the maps have the same source and the same range, and
returns an error message otherwise. }

 

\subsection{\textcolor{Chapter }{\texttt{\symbol{92}}* (maps)}}
\logpage{[ 7, 2, 3 ]}\nobreak
\hyperdef{L}{X8315573C7C90717E}{}
{\noindent\textcolor{FuncColor}{$\triangleright$\enspace\texttt{\texttt{\symbol{92}}* (maps)({\mdseries\slshape f, g})\index{\texttt{\symbol{92}}* (maps)@\texttt{\texttt{\symbol{92}}* (maps)}}
\label{bSlash* (maps)}
}\hfill{\scriptsize (operation)}}\\


 Arguments: \mbox{\texttt{\mdseries\slshape f}}, \mbox{\texttt{\mdseries\slshape g}} \texttt{\symbol{45}} two homomorphisms between two modules, or one scalar and
one homomorphism between modules.\\
 \textbf{\indent Returns:\ }
the composition \mbox{\texttt{\mdseries\slshape fg}} of the maps \mbox{\texttt{\mdseries\slshape f}} and \mbox{\texttt{\mdseries\slshape g}}, if the input are maps between representations of the same quivers. If \mbox{\texttt{\mdseries\slshape f}} or \mbox{\texttt{\mdseries\slshape g}} is a scalar, it returns the natural action of scalars on the maps between
representations.



The function checks if the maps are composable, in the first case and in the
second case it checks if the scalar is in the correct field, and returns an
error message otherwise. }

 

\subsection{\textcolor{Chapter }{CoKernelOfWhat}}
\logpage{[ 7, 2, 4 ]}\nobreak
\hyperdef{L}{X7F6E2378786AC02A}{}
{\noindent\textcolor{FuncColor}{$\triangleright$\enspace\texttt{CoKernelOfWhat({\mdseries\slshape f})\index{CoKernelOfWhat@\texttt{CoKernelOfWhat}}
\label{CoKernelOfWhat}
}\hfill{\scriptsize (attribute)}}\\


 Arguments: \mbox{\texttt{\mdseries\slshape f}} \texttt{\symbol{45}} a homomorphism between two modules.\\
 \textbf{\indent Returns:\ }
 a homomorphism \mbox{\texttt{\mdseries\slshape g}}, if \mbox{\texttt{\mdseries\slshape f}} has been computed as the cokernel of the homomorphism \mbox{\texttt{\mdseries\slshape g}}. 

}

 

\subsection{\textcolor{Chapter }{IdentityMapping}}
\logpage{[ 7, 2, 5 ]}\nobreak
\hyperdef{L}{X7EBAE0368470A603}{}
{\noindent\textcolor{FuncColor}{$\triangleright$\enspace\texttt{IdentityMapping({\mdseries\slshape M})\index{IdentityMapping@\texttt{IdentityMapping}}
\label{IdentityMapping}
}\hfill{\scriptsize (operation)}}\\


 Arguments: \mbox{\texttt{\mdseries\slshape M}} \texttt{\symbol{45}} a module.\\
 \textbf{\indent Returns:\ }
the identity map between \mbox{\texttt{\mdseries\slshape M}} and \mbox{\texttt{\mdseries\slshape M}}.

}

 

\subsection{\textcolor{Chapter }{ImageElm}}
\logpage{[ 7, 2, 6 ]}\nobreak
\hyperdef{L}{X7CFAB0157BFB1806}{}
{\noindent\textcolor{FuncColor}{$\triangleright$\enspace\texttt{ImageElm({\mdseries\slshape f, elem})\index{ImageElm@\texttt{ImageElm}}
\label{ImageElm}
}\hfill{\scriptsize (operation)}}\\


 Arguments: \mbox{\texttt{\mdseries\slshape f}} \texttt{\symbol{45}} a homomorphism between two modules, \mbox{\texttt{\mdseries\slshape elem}} \texttt{\symbol{45}} an element in the source of \mbox{\texttt{\mdseries\slshape f}}.\\
 \textbf{\indent Returns:\ }
the image of the element \mbox{\texttt{\mdseries\slshape elem}} in the source (or domain) of the homomorphism \mbox{\texttt{\mdseries\slshape f}}.



 The function checks if \mbox{\texttt{\mdseries\slshape elem}} is an element in the source of \mbox{\texttt{\mdseries\slshape f}}, and it returns an error message otherwise. }

 

\subsection{\textcolor{Chapter }{ImagesSet}}
\logpage{[ 7, 2, 7 ]}\nobreak
\hyperdef{L}{X8781348F7F5796A0}{}
{\noindent\textcolor{FuncColor}{$\triangleright$\enspace\texttt{ImagesSet({\mdseries\slshape f, elts})\index{ImagesSet@\texttt{ImagesSet}}
\label{ImagesSet}
}\hfill{\scriptsize (operation)}}\\


 Arguments: \mbox{\texttt{\mdseries\slshape f}} \texttt{\symbol{45}} a homomorphism between two modules, \mbox{\texttt{\mdseries\slshape elts}} \texttt{\symbol{45}} an element in the source of \mbox{\texttt{\mdseries\slshape f}}, or the source of \mbox{\texttt{\mdseries\slshape f}}.\\
 \textbf{\indent Returns:\ }
the non\texttt{\symbol{45}}zero images of a set of elements \mbox{\texttt{\mdseries\slshape elts}} in the source of the homomorphism \mbox{\texttt{\mdseries\slshape f}}, or if \mbox{\texttt{\mdseries\slshape elts}} is the source of \mbox{\texttt{\mdseries\slshape f}}, it returns a basis of the image.



 The function checks if the set of elements \mbox{\texttt{\mdseries\slshape elts}} consists of elements in the source of \mbox{\texttt{\mdseries\slshape f}}, and it returns an error message otherwise. }

 
\begin{Verbatim}[commandchars=@|A,fontsize=\small,frame=single,label=Example]
  @gapprompt|gap>A @gapinput|B := BasisVectors(Basis(N)); A
  [ [ [ 1, 0, 0 ], [ 0, 0 ], [ 0, 0 ] ], 
    [ [ 0, 1, 0 ], [ 0, 0 ], [ 0, 0 ] ], 
    [ [ 0, 0, 1 ], [ 0, 0 ], [ 0, 0 ] ], 
    [ [ 0, 0, 0 ], [ 1, 0 ], [ 0, 0 ] ], 
    [ [ 0, 0, 0 ], [ 0, 1 ], [ 0, 0 ] ], 
    [ [ 0, 0, 0 ], [ 0, 0 ], [ 1, 0 ] ], 
    [ [ 0, 0, 0 ], [ 0, 0 ], [ 0, 1 ] ] ]
  @gapprompt|gap>A @gapinput|PreImagesRepresentativeNC(f,B[4]);     A
  [ [ 0 ], [ 1 ], [ 0 ] ]
  @gapprompt|gap>A @gapinput|PreImagesRepresentativeNC(f,B[5]);A
  fail
  @gapprompt|gap>A @gapinput|BL := BasisVectors(Basis(L));A
  [ [ [ 0 ], [ 1 ], [ 0 ] ], [ [ 0 ], [ 0 ], [ 1 ] ] ]
  @gapprompt|gap>A @gapinput|ImageElm(f,BL[1]);A
  [ [ 0, 0, 0 ], [ 1, 0 ], [ 0, 0 ] ]
  @gapprompt|gap>A @gapinput|ImagesSet(f,L);A
  [ [ [ 0, 0, 0 ], [ 1, 0 ], [ 0, 0 ] ], 
    [ [ 0, 0, 0 ], [ 0, 0 ], [ 1, 2 ] ] ]
  @gapprompt|gap>A @gapinput|ImagesSet(f,BL);A
  [ [ [ 0, 0, 0 ], [ 1, 0 ], [ 0, 0 ] ], 
    [ [ 0, 0, 0 ], [ 0, 0 ], [ 1, 2 ] ] ]
  @gapprompt|gap>A @gapinput|z := Zero(f);;A
  @gapprompt|gap>A @gapinput|f = z;A
  false
  @gapprompt|gap>A @gapinput|Range(f) = Range(z);A
  true
  @gapprompt|gap>A @gapinput|y := ZeroMapping(L,N);;A
  @gapprompt|gap>A @gapinput|y = z;            A
  true
  @gapprompt|gap>A @gapinput|id := IdentityMapping(N);;A
  @gapprompt|gap>A @gapinput|f*id;;A
  @gapprompt|gap>A @gapinput|#This causes an error!A
  @gapprompt|gap>A @gapinput|id*f;A
  Error, codomain of the first argument is not equal to the domain of th\
  e second argument,  called from
  <function>( <arguments> ) called from read-eval-loop
  Entering break read-eval-print loop ...
  you can 'quit;' to quit to outer loop, or
  you can 'return;' to continue
  @gapbrkprompt|brk>A @gapinput|quit;;A
  @gapprompt|gap>A @gapinput|2*f + z;A
  <<[ 0, 1, 1 ]> ---> <[ 3, 2, 2 ]>> 
\end{Verbatim}
 

\subsection{\textcolor{Chapter }{ImageOfWhat}}
\logpage{[ 7, 2, 8 ]}\nobreak
\hyperdef{L}{X78EE24857C79789E}{}
{\noindent\textcolor{FuncColor}{$\triangleright$\enspace\texttt{ImageOfWhat({\mdseries\slshape f})\index{ImageOfWhat@\texttt{ImageOfWhat}}
\label{ImageOfWhat}
}\hfill{\scriptsize (attribute)}}\\


 Arguments: \mbox{\texttt{\mdseries\slshape f}} \texttt{\symbol{45}} a homomorphism between two modules.\\
 \textbf{\indent Returns:\ }
 a homomorphism \mbox{\texttt{\mdseries\slshape g}}, if \mbox{\texttt{\mdseries\slshape f}} has been computed as the image projection or the image inclusion of the
homomorphism \mbox{\texttt{\mdseries\slshape g}}. 

}

 

\subsection{\textcolor{Chapter }{IsInjective}}
\logpage{[ 7, 2, 9 ]}\nobreak
\hyperdef{L}{X7F065FD7822C0A12}{}
{\noindent\textcolor{FuncColor}{$\triangleright$\enspace\texttt{IsInjective({\mdseries\slshape f})\index{IsInjective@\texttt{IsInjective}}
\label{IsInjective}
}\hfill{\scriptsize (property)}}\\


 Arguments: \mbox{\texttt{\mdseries\slshape f}} \texttt{\symbol{45}} a homomorphism between two modules.\\
 \textbf{\indent Returns:\ }
 \texttt{true} if the homomorphism \mbox{\texttt{\mdseries\slshape f}} is one\texttt{\symbol{45}}to\texttt{\symbol{45}}one. 

}

 

\subsection{\textcolor{Chapter }{IsIsomorphism}}
\logpage{[ 7, 2, 10 ]}\nobreak
\hyperdef{L}{X7E07BBF57B92BA56}{}
{\noindent\textcolor{FuncColor}{$\triangleright$\enspace\texttt{IsIsomorphism({\mdseries\slshape f})\index{IsIsomorphism@\texttt{IsIsomorphism}}
\label{IsIsomorphism}
}\hfill{\scriptsize (operation)}}\\


 Arguments: \mbox{\texttt{\mdseries\slshape f}} \texttt{\symbol{45}} a homomorphism between two modules.\\
 \textbf{\indent Returns:\ }
 \texttt{true} if the homomorphism \mbox{\texttt{\mdseries\slshape f}} is an isomorphism. 

}

 

\subsection{\textcolor{Chapter }{IsLeftMinimal}}
\logpage{[ 7, 2, 11 ]}\nobreak
\hyperdef{L}{X7E5D33B8853B9490}{}
{\noindent\textcolor{FuncColor}{$\triangleright$\enspace\texttt{IsLeftMinimal({\mdseries\slshape f})\index{IsLeftMinimal@\texttt{IsLeftMinimal}}
\label{IsLeftMinimal}
}\hfill{\scriptsize (property)}}\\


 Arguments: \mbox{\texttt{\mdseries\slshape f}} \texttt{\symbol{45}} a homomorphism between two modules.\\
 \textbf{\indent Returns:\ }
 \texttt{true} if the homomorphism \mbox{\texttt{\mdseries\slshape f}} is left minimal. 

}

 

\subsection{\textcolor{Chapter }{IsRightMinimal}}
\logpage{[ 7, 2, 12 ]}\nobreak
\hyperdef{L}{X876706C77FB707E5}{}
{\noindent\textcolor{FuncColor}{$\triangleright$\enspace\texttt{IsRightMinimal({\mdseries\slshape f})\index{IsRightMinimal@\texttt{IsRightMinimal}}
\label{IsRightMinimal}
}\hfill{\scriptsize (property)}}\\


 Arguments: \mbox{\texttt{\mdseries\slshape f}} \texttt{\symbol{45}} a homomorphism between two modules.\\
 \textbf{\indent Returns:\ }
 \texttt{true} if the homomorphism \mbox{\texttt{\mdseries\slshape f}} is right minimal. 

}

 
\begin{Verbatim}[commandchars=!@|,fontsize=\small,frame=single,label=Example]
  !gapprompt@gap>| !gapinput@L := RightModuleOverPathAlgebra(A,[["a",[0,1]],["b",[0,1]],|
  !gapprompt@>| !gapinput@["c",[[0]]],["d",[[1]]],["e",[1,0]]]);;|
  !gapprompt@gap>| !gapinput@f := RightModuleHomOverAlgebra(L,N,[[[0,0,0]], [[1,0]], |
  !gapprompt@>| !gapinput@[[1,2]]]);|
  <<[ 0, 1, 1 ]> ---> <[ 3, 2, 2 ]>>
  
  !gapprompt@gap>| !gapinput@g := CoKernelProjection(f);|
  <<[ 3, 2, 2 ]> ---> <[ 3, 1, 1 ]>>
  
  !gapprompt@gap>| !gapinput@CoKernelOfWhat(g) = f;|
  true
  !gapprompt@gap>| !gapinput@h := ImageProjection(f);|
  <<[ 0, 1, 1 ]> ---> <[ 0, 1, 1 ]>>
  
  !gapprompt@gap>| !gapinput@ImageOfWhat(h) = f;|
  true
  !gapprompt@gap>| !gapinput@IsInjective(f); IsSurjective(f); IsIsomorphism(f); |
  true
  false
  false
  !gapprompt@gap>| !gapinput@IsIsomorphism(h);|
  true 
\end{Verbatim}
 

\subsection{\textcolor{Chapter }{IsSplitEpimorphism}}
\logpage{[ 7, 2, 13 ]}\nobreak
\hyperdef{L}{X80A66EFA862E56BC}{}
{\noindent\textcolor{FuncColor}{$\triangleright$\enspace\texttt{IsSplitEpimorphism({\mdseries\slshape f})\index{IsSplitEpimorphism@\texttt{IsSplitEpimorphism}}
\label{IsSplitEpimorphism}
}\hfill{\scriptsize (property)}}\\


 Arguments: \mbox{\texttt{\mdseries\slshape f}} \texttt{\symbol{45}} a homomorphism between two modules.\\
 \textbf{\indent Returns:\ }
 \texttt{true} if the homomorphism \mbox{\texttt{\mdseries\slshape f}} is a splittable epimorphism, otherwise \texttt{false}. 

}

 

\subsection{\textcolor{Chapter }{IsSplitMonomorphism}}
\logpage{[ 7, 2, 14 ]}\nobreak
\hyperdef{L}{X7DFACF1F7D7F7EE9}{}
{\noindent\textcolor{FuncColor}{$\triangleright$\enspace\texttt{IsSplitMonomorphism({\mdseries\slshape f})\index{IsSplitMonomorphism@\texttt{IsSplitMonomorphism}}
\label{IsSplitMonomorphism}
}\hfill{\scriptsize (property)}}\\


 Arguments: \mbox{\texttt{\mdseries\slshape f}} \texttt{\symbol{45}} a homomorphism between two modules.\\
 \textbf{\indent Returns:\ }
 \texttt{true} if the homomorphism \mbox{\texttt{\mdseries\slshape f}} is a splittable monomorphism, otherwise \texttt{false}. 

}

 
\begin{Verbatim}[commandchars=!@|,fontsize=\small,frame=single,label=Example]
  !gapprompt@gap>| !gapinput@S := SimpleModules(A)[1];;|
  !gapprompt@gap>| !gapinput@H := HomOverAlgebra(N,S);; |
  !gapprompt@gap>| !gapinput@IsSplitMonomorphism(H[1]);  |
  false
  !gapprompt@gap>| !gapinput@IsSplitEpimorphism(H[1]);|
  true
\end{Verbatim}
 

\subsection{\textcolor{Chapter }{IsSurjective}}
\logpage{[ 7, 2, 15 ]}\nobreak
\hyperdef{L}{X784ECE847E005B8F}{}
{\noindent\textcolor{FuncColor}{$\triangleright$\enspace\texttt{IsSurjective({\mdseries\slshape f})\index{IsSurjective@\texttt{IsSurjective}}
\label{IsSurjective}
}\hfill{\scriptsize (property)}}\\


 Arguments: \mbox{\texttt{\mdseries\slshape f}} \texttt{\symbol{45}} a homomorphism between two modules.\\
 \textbf{\indent Returns:\ }
 \texttt{true} if the homomorphism \mbox{\texttt{\mdseries\slshape f}} is onto. 

}

 

\subsection{\textcolor{Chapter }{IsZero(homomorphism)}}
\logpage{[ 7, 2, 16 ]}\nobreak
\hyperdef{L}{X7C910CBD7C225C08}{}
{\noindent\textcolor{FuncColor}{$\triangleright$\enspace\texttt{IsZero(homomorphism)({\mdseries\slshape f})\index{IsZero(homomorphism)@\texttt{IsZero(homomorphism)}}
\label{IsZero(homomorphism)}
}\hfill{\scriptsize (property)}}\\


 Arguments: \mbox{\texttt{\mdseries\slshape f}} \texttt{\symbol{45}} a homomorphism between two modules.\\
 \textbf{\indent Returns:\ }
 \texttt{true} if the homomorphism \mbox{\texttt{\mdseries\slshape f}} is a zero homomorphism. 

}

 

\subsection{\textcolor{Chapter }{KernelOfWhat}}
\logpage{[ 7, 2, 17 ]}\nobreak
\hyperdef{L}{X7EF520F67BA7F082}{}
{\noindent\textcolor{FuncColor}{$\triangleright$\enspace\texttt{KernelOfWhat({\mdseries\slshape f})\index{KernelOfWhat@\texttt{KernelOfWhat}}
\label{KernelOfWhat}
}\hfill{\scriptsize (attribute)}}\\


 Arguments: \mbox{\texttt{\mdseries\slshape f}} \texttt{\symbol{45}} a homomorphism between two modules.\\
 \textbf{\indent Returns:\ }
 a homomorphism \mbox{\texttt{\mdseries\slshape g}}, if \mbox{\texttt{\mdseries\slshape f}} has been computed as the kernel of the homomorphism \mbox{\texttt{\mdseries\slshape g}}. 

}

 
\begin{Verbatim}[commandchars=!@|,fontsize=\small,frame=single,label=Example]
  !gapprompt@gap>| !gapinput@L := RightModuleOverPathAlgebra(A,[["a",[0,1]],["b",[0,1]],|
  !gapprompt@>| !gapinput@["c",[[0]]],["d",[[1]]],["e",[1,0]]]);|
  <[ 0, 1, 1 ]>
  !gapprompt@gap>| !gapinput@f := RightModuleHomOverAlgebra(L,N,[[[0,0,0]], [[1,0]], |
  !gapprompt@>| !gapinput@[[1,2]]]);;|
  !gapprompt@gap>| !gapinput@IsZero(0*f);|
  true
  !gapprompt@gap>| !gapinput@g := KernelInclusion(f);|
  <<[ 0, 0, 0 ]> ---> <[ 0, 1, 1 ]>>
  
  !gapprompt@gap>| !gapinput@KnownAttributesOfObject(g);|
  [ "Range", "Source", "PathAlgebraOfMatModuleMap", "KernelOfWhat" ]
  !gapprompt@gap>| !gapinput@KernelOfWhat(g) = f;|
  true 
\end{Verbatim}
 

\subsection{\textcolor{Chapter }{LeftInverseOfHomomorphism}}
\logpage{[ 7, 2, 18 ]}\nobreak
\hyperdef{L}{X7A40540E79DBD804}{}
{\noindent\textcolor{FuncColor}{$\triangleright$\enspace\texttt{LeftInverseOfHomomorphism({\mdseries\slshape f})\index{LeftInverseOfHomomorphism@\texttt{LeftInverseOfHomomorphism}}
\label{LeftInverseOfHomomorphism}
}\hfill{\scriptsize (attribute)}}\\


 Arguments: \mbox{\texttt{\mdseries\slshape f}} \texttt{\symbol{45}} a homomorphism between two modules.\\
 \textbf{\indent Returns:\ }
\texttt{false} if the homomorphism \mbox{\texttt{\mdseries\slshape f}} is not a splittable epimorphism, otherwise it returns a splitting of the split
epimorphism \mbox{\texttt{\mdseries\slshape f}}. 

}

 

\subsection{\textcolor{Chapter }{MatricesOfPathAlgebraMatModuleHomomorphism}}
\logpage{[ 7, 2, 19 ]}\nobreak
\hyperdef{L}{X7B63FBAF84533D75}{}
{\noindent\textcolor{FuncColor}{$\triangleright$\enspace\texttt{MatricesOfPathAlgebraMatModuleHomomorphism({\mdseries\slshape f})\index{MatricesOfPathAlgebraMatModuleHomomorphism@\texttt{Matrices}\-\texttt{Of}\-\texttt{Path}\-\texttt{Algebra}\-\texttt{Mat}\-\texttt{Module}\-\texttt{Homomorphism}}
\label{MatricesOfPathAlgebraMatModuleHomomorphism}
}\hfill{\scriptsize (operation)}}\\


 Arguments: \mbox{\texttt{\mdseries\slshape f}} \texttt{\symbol{45}} a homomorphism between two modules.\\
 \textbf{\indent Returns:\ }
the matrices defining the homomorphism \mbox{\texttt{\mdseries\slshape f}}.

}

 
\begin{Verbatim}[commandchars=!@|,fontsize=\small,frame=single,label=Example]
  !gapprompt@gap>| !gapinput@MatricesOfPathAlgebraMatModuleHomomorphism(f);|
  [ [ [ 0, 0, 0 ] ], [ [ 1, 0 ] ], [ [ 1, 2 ] ] ]
  !gapprompt@gap>| !gapinput@Range(f);|
  <[ 3, 2, 2 ]>
  !gapprompt@gap>| !gapinput@Source(f);|
  <[ 0, 1, 1 ]>
  !gapprompt@gap>| !gapinput@Source(f) = L;|
  true 
\end{Verbatim}
 

\subsection{\textcolor{Chapter }{PathAlgebraOfMatModuleMap}}
\logpage{[ 7, 2, 20 ]}\nobreak
\hyperdef{L}{X7D3B488586DA3938}{}
{\noindent\textcolor{FuncColor}{$\triangleright$\enspace\texttt{PathAlgebraOfMatModuleMap({\mdseries\slshape f})\index{PathAlgebraOfMatModuleMap@\texttt{PathAlgebraOfMatModuleMap}}
\label{PathAlgebraOfMatModuleMap}
}\hfill{\scriptsize (attribute)}}\\


 Arguments: \mbox{\texttt{\mdseries\slshape f}} \texttt{\symbol{45}}\texttt{\symbol{45}} a homomorphism between two path
algebra modules (\texttt{PathAlgebraMatModule}). \\
\textbf{\indent Returns:\ }
the algebra over which the range and the source of the homomorphism \mbox{\texttt{\mdseries\slshape f}} is defined.

}

 

\subsection{\textcolor{Chapter }{PreImagesRepresentativeNC}}
\logpage{[ 7, 2, 21 ]}\nobreak
\hyperdef{L}{X7981A59C80E21FD8}{}
{\noindent\textcolor{FuncColor}{$\triangleright$\enspace\texttt{PreImagesRepresentativeNC({\mdseries\slshape f, elem})\index{PreImagesRepresentativeNC@\texttt{PreImagesRepresentativeNC}}
\label{PreImagesRepresentativeNC}
}\hfill{\scriptsize (operation)}}\\


 Arguments: \mbox{\texttt{\mdseries\slshape f}} \texttt{\symbol{45}} a homomorphism between two modules, \mbox{\texttt{\mdseries\slshape elem}} \texttt{\symbol{45}} an element in the range of \mbox{\texttt{\mdseries\slshape f}}.\\
 \textbf{\indent Returns:\ }
a preimage of the element \mbox{\texttt{\mdseries\slshape elem}} in the range (or codomain) the homomorphism \mbox{\texttt{\mdseries\slshape f}} if a preimage exists, otherwise it returns \texttt{fail}.



 The function checks if \mbox{\texttt{\mdseries\slshape elem}} is an element in the range of \mbox{\texttt{\mdseries\slshape f}} and returns an error message if not. }

 

\subsection{\textcolor{Chapter }{Range}}
\logpage{[ 7, 2, 22 ]}\nobreak
\hyperdef{L}{X829F76BB80BD55DB}{}
{\noindent\textcolor{FuncColor}{$\triangleright$\enspace\texttt{Range({\mdseries\slshape f})\index{Range@\texttt{Range}}
\label{Range}
}\hfill{\scriptsize (attribute)}}\\


 Arguments: \mbox{\texttt{\mdseries\slshape f}} \texttt{\symbol{45}} a homomorphism between two modules.\\
 \textbf{\indent Returns:\ }
the range (or codomain) the homomorphism \mbox{\texttt{\mdseries\slshape f}}.

}

 

\subsection{\textcolor{Chapter }{RightInverseOfHomomorphism}}
\logpage{[ 7, 2, 23 ]}\nobreak
\hyperdef{L}{X8105F85B8260C4F9}{}
{\noindent\textcolor{FuncColor}{$\triangleright$\enspace\texttt{RightInverseOfHomomorphism({\mdseries\slshape f})\index{RightInverseOfHomomorphism@\texttt{RightInverseOfHomomorphism}}
\label{RightInverseOfHomomorphism}
}\hfill{\scriptsize (attribute)}}\\


 Arguments: \mbox{\texttt{\mdseries\slshape f}} \texttt{\symbol{45}} a homomorphism between two modules.\\
 \textbf{\indent Returns:\ }
\texttt{false} if the homomorphism \mbox{\texttt{\mdseries\slshape f}} is not a splittable monomorphism, otherwise it returns a splitting of the
split monomorphism \mbox{\texttt{\mdseries\slshape f}}. 

}

 

\subsection{\textcolor{Chapter }{Source}}
\logpage{[ 7, 2, 24 ]}\nobreak
\hyperdef{L}{X7DE8173F80E07AB1}{}
{\noindent\textcolor{FuncColor}{$\triangleright$\enspace\texttt{Source({\mdseries\slshape f})\index{Source@\texttt{Source}}
\label{Source}
}\hfill{\scriptsize (attribute)}}\\


 Arguments: \mbox{\texttt{\mdseries\slshape f}} \texttt{\symbol{45}} a homomorphism between two modules.\\
 \textbf{\indent Returns:\ }
the source (or domain) the homomorphism \mbox{\texttt{\mdseries\slshape f}}.

}

 

\subsection{\textcolor{Chapter }{Zero}}
\logpage{[ 7, 2, 25 ]}\nobreak
\hyperdef{L}{X804B376481243046}{}
{\noindent\textcolor{FuncColor}{$\triangleright$\enspace\texttt{Zero({\mdseries\slshape f})\index{Zero@\texttt{Zero}}
\label{Zero}
}\hfill{\scriptsize (operation)}}\\


 Arguments: \mbox{\texttt{\mdseries\slshape f}} \texttt{\symbol{45}} a homomorphism between two modules.\\
 \textbf{\indent Returns:\ }
the zero map between \texttt{Source(f)} and \texttt{Range(f)}.

}

 

\subsection{\textcolor{Chapter }{ZeroMapping}}
\logpage{[ 7, 2, 26 ]}\nobreak
\hyperdef{L}{X795FF8DC785F110A}{}
{\noindent\textcolor{FuncColor}{$\triangleright$\enspace\texttt{ZeroMapping({\mdseries\slshape M, N})\index{ZeroMapping@\texttt{ZeroMapping}}
\label{ZeroMapping}
}\hfill{\scriptsize (operation)}}\\


 Arguments: \mbox{\texttt{\mdseries\slshape M}}, \mbox{\texttt{\mdseries\slshape N}} \texttt{\symbol{45}} two modules.\\
 \textbf{\indent Returns:\ }
the zero map between \mbox{\texttt{\mdseries\slshape M}} and \mbox{\texttt{\mdseries\slshape N}}.

}

 

\subsection{\textcolor{Chapter }{HomomorphismFromImages}}
\logpage{[ 7, 2, 27 ]}\nobreak
\hyperdef{L}{X81CA4D9E7D50A9A8}{}
{\noindent\textcolor{FuncColor}{$\triangleright$\enspace\texttt{HomomorphismFromImages({\mdseries\slshape M, N, genImages})\index{HomomorphismFromImages@\texttt{HomomorphismFromImages}}
\label{HomomorphismFromImages}
}\hfill{\scriptsize (operation)}}\\


 Arguments: \mbox{\texttt{\mdseries\slshape M}}, \mbox{\texttt{\mdseries\slshape N}} \texttt{\symbol{45}}\texttt{\symbol{45}} two modules, \mbox{\texttt{\mdseries\slshape genImages}} \texttt{\symbol{45}}\texttt{\symbol{45}} a list.\\
 \textbf{\indent Returns:\ }
A map $f$ between \mbox{\texttt{\mdseries\slshape M}} and \mbox{\texttt{\mdseries\slshape N}}, given by \mbox{\texttt{\mdseries\slshape genImages}}. 



 Let \texttt{B} be the basis \texttt{BasisVectors( Basis( M ) )} of \mbox{\texttt{\mdseries\slshape M}}. Then the number of elements of \texttt{genImages} should be equal to the number of elements of \texttt{B}, and \texttt{genImages[i]} is an element of \texttt{N} and the image of \texttt{B[i]} under \texttt{f}. The method fails if \texttt{f} is not a homomorphism, or if \texttt{B[i]} and \texttt{genImages[i]} are supported in different vertices. }

 }

 
\section{\textcolor{Chapter }{Homomorphisms and modules constructed from homomorphisms and modules}}\logpage{[ 7, 3, 0 ]}
\hyperdef{L}{X7E8D1A3F7C03CFF1}{}
{
 

\subsection{\textcolor{Chapter }{ADRAlgebraOfAlgebra}}
\logpage{[ 7, 3, 1 ]}\nobreak
\hyperdef{L}{X79A45BC87F1AA033}{}
{\noindent\textcolor{FuncColor}{$\triangleright$\enspace\texttt{ADRAlgebraOfAlgebra({\mdseries\slshape A})\index{ADRAlgebraOfAlgebra@\texttt{ADRAlgebraOfAlgebra}}
\label{ADRAlgebraOfAlgebra}
}\hfill{\scriptsize (operation)}}\\


 Arguments: \mbox{\texttt{\mdseries\slshape A}} \texttt{\symbol{45}} an quiver algebra over a field.\\
 \textbf{\indent Returns:\ }
Auslander\texttt{\symbol{45}}Dlab\texttt{\symbol{45}}Ringel algebra of the
algebra \mbox{\texttt{\mdseries\slshape A}}. 



 Given a finite dimensional quiver algebra \mbox{\texttt{\mdseries\slshape A}} this function constructs the endomorphism algebra of the \mbox{\texttt{\mdseries\slshape A}}\texttt{\symbol{45}}modules $\oplus_{i=1}^n A/\textrm{rad}(A)^i,$ where $n$ is the Loewy length of \mbox{\texttt{\mdseries\slshape A}}, which is known to have finite global dimension less or equal to $n + 1$. }

 

\subsection{\textcolor{Chapter }{ADRAlgebraOfBasicModule}}
\logpage{[ 7, 3, 2 ]}\nobreak
\hyperdef{L}{X825762D5795A419B}{}
{\noindent\textcolor{FuncColor}{$\triangleright$\enspace\texttt{ADRAlgebraOfBasicModule({\mdseries\slshape L})\index{ADRAlgebraOfBasicModule@\texttt{ADRAlgebraOfBasicModule}}
\label{ADRAlgebraOfBasicModule}
}\hfill{\scriptsize (operation)}}\\


 Arguments: \mbox{\texttt{\mdseries\slshape L}} \texttt{\symbol{45}} a list of PathAlgebraMatModules over a quiver algebra
over a finite field.\\
 \textbf{\indent Returns:\ }
Auslander\texttt{\symbol{45}}Dlab\texttt{\symbol{45}}Ringel algebra of the
module given as the direct of the list of non\texttt{\symbol{45}}isomorphic
modules \mbox{\texttt{\mdseries\slshape L}}. 



 Given a list of non\texttt{\symbol{45}}isomorphic modules \mbox{\texttt{\mdseries\slshape L}} this function constructs the endomorphism algebra of the basic module $M$ which generate the same additive subcategory as the module $\oplus_{i=1}^n U/\textrm{rad}^i(U),$ where $U$ is the direct sum of the modules on the list \mbox{\texttt{\mdseries\slshape L}} and $n$ is the Loewy length of \mbox{\texttt{\mdseries\slshape U}}. }

 

\subsection{\textcolor{Chapter }{AllIndecModulesOfLengthAtMost}}
\logpage{[ 7, 3, 3 ]}\nobreak
\hyperdef{L}{X7864A84A80553958}{}
{\noindent\textcolor{FuncColor}{$\triangleright$\enspace\texttt{AllIndecModulesOfLengthAtMost({\mdseries\slshape A, n})\index{AllIndecModulesOfLengthAtMost@\texttt{AllIndecModulesOfLengthAtMost}}
\label{AllIndecModulesOfLengthAtMost}
}\hfill{\scriptsize (operation)}}\\


 Arguments: \mbox{\texttt{\mdseries\slshape A, n}} \texttt{\symbol{45}} an algebra over a finite field, an integer.\\
 \textbf{\indent Returns:\ }
all the different indecomposable modules over the algebra \mbox{\texttt{\mdseries\slshape A}} of length at most \mbox{\texttt{\mdseries\slshape n}}. 



 This function is only implemented for algebras over a finite field. }

 

\subsection{\textcolor{Chapter }{AllModulesOfLengthAtMost}}
\logpage{[ 7, 3, 4 ]}\nobreak
\hyperdef{L}{X83B5C7D484D98A34}{}
{\noindent\textcolor{FuncColor}{$\triangleright$\enspace\texttt{AllModulesOfLengthAtMost({\mdseries\slshape A, n})\index{AllModulesOfLengthAtMost@\texttt{AllModulesOfLengthAtMost}}
\label{AllModulesOfLengthAtMost}
}\hfill{\scriptsize (operation)}}\\


 Arguments: \mbox{\texttt{\mdseries\slshape A, n}} \texttt{\symbol{45}} an algebra over a finite field, an integer.\\
 \textbf{\indent Returns:\ }
all the different modules over the algebra \mbox{\texttt{\mdseries\slshape A}} of length at most \mbox{\texttt{\mdseries\slshape n}}. 



 This function is only implemented for algebras over a finite field. }

 

\subsection{\textcolor{Chapter }{AllSimpleSubmodulesOfModule}}
\logpage{[ 7, 3, 5 ]}\nobreak
\hyperdef{L}{X84F2D28D7F8694DC}{}
{\noindent\textcolor{FuncColor}{$\triangleright$\enspace\texttt{AllSimpleSubmodulesOfModule({\mdseries\slshape M})\index{AllSimpleSubmodulesOfModule@\texttt{AllSimpleSubmodulesOfModule}}
\label{AllSimpleSubmodulesOfModule}
}\hfill{\scriptsize (operation)}}\\


 Arguments: \mbox{\texttt{\mdseries\slshape M}} \texttt{\symbol{45}} a module.\\
 \textbf{\indent Returns:\ }
all the different simple submodules of a module given as inclusions into the
module \mbox{\texttt{\mdseries\slshape M}}. 



 This function is only implemented for algebras over a finite field. }

 

\subsection{\textcolor{Chapter }{AllSubmodulesOfModule}}
\logpage{[ 7, 3, 6 ]}\nobreak
\hyperdef{L}{X7C3D5CF978CF5058}{}
{\noindent\textcolor{FuncColor}{$\triangleright$\enspace\texttt{AllSubmodulesOfModule({\mdseries\slshape M})\index{AllSubmodulesOfModule@\texttt{AllSubmodulesOfModule}}
\label{AllSubmodulesOfModule}
}\hfill{\scriptsize (operation)}}\\


 Arguments: \mbox{\texttt{\mdseries\slshape M}} \texttt{\symbol{45}} a module.\\
 \textbf{\indent Returns:\ }
all the different submodules of a module given as inclusions into the module \mbox{\texttt{\mdseries\slshape M}}. It returns the list of submodules as a list of lists according to the length
of the submodules, namely, first a list of the zero module, second a list of
all simple submodules, third a list of all submodules of length 2, and so on. 



 This function is only implemented for algebras over a finite field. }

 

\subsection{\textcolor{Chapter }{CoKernel}}
\logpage{[ 7, 3, 7 ]}\nobreak
\hyperdef{L}{X875F177A82BF9B8B}{}
{\noindent\textcolor{FuncColor}{$\triangleright$\enspace\texttt{CoKernel({\mdseries\slshape f})\index{CoKernel@\texttt{CoKernel}}
\label{CoKernel}
}\hfill{\scriptsize (attribute)}}\\


 Arguments: \mbox{\texttt{\mdseries\slshape f}} \texttt{\symbol{45}} a homomorphism between two modules.\\
 \textbf{\indent Returns:\ }
the cokernel of a homomorphism \mbox{\texttt{\mdseries\slshape f}} between two modules.



 This function returns the cokernel of the homomorphism \mbox{\texttt{\mdseries\slshape f}} as a module. }

 

\subsection{\textcolor{Chapter }{CoKernelProjection}}
\logpage{[ 7, 3, 8 ]}\nobreak
\hyperdef{L}{X8267B6477A8F808F}{}
{\noindent\textcolor{FuncColor}{$\triangleright$\enspace\texttt{CoKernelProjection({\mdseries\slshape f})\index{CoKernelProjection@\texttt{CoKernelProjection}}
\label{CoKernelProjection}
}\hfill{\scriptsize (attribute)}}\\


 Arguments: \mbox{\texttt{\mdseries\slshape f}} \texttt{\symbol{45}} a homomorphism between two modules.\\
 \textbf{\indent Returns:\ }
the cokernel of a homomorphism \mbox{\texttt{\mdseries\slshape f}} between two modules.



 This function returns the cokernel of the homomorphism \mbox{\texttt{\mdseries\slshape f}} as the projection homomorphism from the range of the homomorphism \mbox{\texttt{\mdseries\slshape f}} to the cokernel of the homomorphism \mbox{\texttt{\mdseries\slshape f}}. }

 

\subsection{\textcolor{Chapter }{EndModuloProjOverAlgebra}}
\logpage{[ 7, 3, 9 ]}\nobreak
\hyperdef{L}{X7D302484872EBCA5}{}
{\noindent\textcolor{FuncColor}{$\triangleright$\enspace\texttt{EndModuloProjOverAlgebra({\mdseries\slshape M})\index{EndModuloProjOverAlgebra@\texttt{EndModuloProjOverAlgebra}}
\label{EndModuloProjOverAlgebra}
}\hfill{\scriptsize (operation)}}\\


 Arguments: \mbox{\texttt{\mdseries\slshape M}} \texttt{\symbol{45}} a module.\\
 \textbf{\indent Returns:\ }
the natural homomorphism from the endomorphism ring of \mbox{\texttt{\mdseries\slshape M}} to the endomorphism ring of \mbox{\texttt{\mdseries\slshape M}} modulo the ideal generated by those endomorphisms of \mbox{\texttt{\mdseries\slshape M}} which factor through a projective module.



 The operation returns an error message if the zero module is entered as an
argument. }

 

\subsection{\textcolor{Chapter }{EndOfModuleAsQuiverAlgebra}}
\logpage{[ 7, 3, 10 ]}\nobreak
\hyperdef{L}{X8179107081C47D81}{}
{\noindent\textcolor{FuncColor}{$\triangleright$\enspace\texttt{EndOfModuleAsQuiverAlgebra({\mdseries\slshape M})\index{EndOfModuleAsQuiverAlgebra@\texttt{EndOfModuleAsQuiverAlgebra}}
\label{EndOfModuleAsQuiverAlgebra}
}\hfill{\scriptsize (operation)}}\\


 Arguments: \mbox{\texttt{\mdseries\slshape M}} \texttt{\symbol{45}} a PathAlgebraMatModule.\\
 \textbf{\indent Returns:\ }
 a list of three elements, (i) the endomorphism ring of \mbox{\texttt{\mdseries\slshape M}}, (ii) the adjacency matrix of the quiver of the endomorphism ring and (iii)
the endomorphism ring as a quiver algebra. 



 Suppose \mbox{\texttt{\mdseries\slshape M}} is a module over a quiver algebra over a field $K$. The function checks if the endomorphism ring of \mbox{\texttt{\mdseries\slshape M}} is K\texttt{\symbol{45}}elementary (not necessary for it to be a quiver
algebra, but this is a TODO improvement), and returns error message otherwise. }

 

\subsection{\textcolor{Chapter }{EndOfBasicModuleAsQuiverAlgebra}}
\logpage{[ 7, 3, 11 ]}\nobreak
\hyperdef{L}{X7E0B0F807933DEEE}{}
{\noindent\textcolor{FuncColor}{$\triangleright$\enspace\texttt{EndOfBasicModuleAsQuiverAlgebra({\mdseries\slshape L})\index{EndOfBasicModuleAsQuiverAlgebra@\texttt{EndOfBasicModuleAsQuiverAlgebra}}
\label{EndOfBasicModuleAsQuiverAlgebra}
}\hfill{\scriptsize (operation)}}\\


 Arguments: \mbox{\texttt{\mdseries\slshape L}} \texttt{\symbol{45}} a list of PathAlgebraMatModule's.\\
 \textbf{\indent Returns:\ }
 the endomorphism ring as a quiver algebra. 



 The argument of this function is a list of PathAlgebraMatModules over a quiver
algebra over a field K and optionally a positive integer. Hence the input
should either be (i) a list of non\texttt{\symbol{45}}isomorphic modules or
(ii) a list of two elements where the first entry is a list of
non\texttt{\symbol{45}}isomorphic modules and the second entry is a positive
integer. It assumes that the direct sum of the modules in the list is basic.
It checks if the endomorphism ring of the direct sum of the modules on the
list is K\texttt{\symbol{45}}elementary. If so, then it finds the path algebra
the endomorphism ring of the direct sum $M$ of the modules on the list \mbox{\texttt{\mdseries\slshape L}} is a quotient of, and quotient out the relations that are found. If the
argument is only a list of modules, then the function returns the quiver
algebra isomorphic to the endomorphism ring of $M$. If the argument is a list of modules and a positive integer, then not all
relations might not been found and the function returns a warning if that is
the case. This function was first constructed mainly by Daniel Owens and with
help from Rene Marczinzik. Later modified by Oeyvind Solberg. }

 

\subsection{\textcolor{Chapter }{EndOverAlgebra}}
\logpage{[ 7, 3, 12 ]}\nobreak
\hyperdef{L}{X79FDBE1B795308A9}{}
{\noindent\textcolor{FuncColor}{$\triangleright$\enspace\texttt{EndOverAlgebra({\mdseries\slshape M})\index{EndOverAlgebra@\texttt{EndOverAlgebra}}
\label{EndOverAlgebra}
}\hfill{\scriptsize (attribute)}}\\


 Arguments: \mbox{\texttt{\mdseries\slshape M}} \texttt{\symbol{45}} a module.\\
 \textbf{\indent Returns:\ }
the endomorphism ring of \mbox{\texttt{\mdseries\slshape M}} as a subalgebra of the direct sum of the full matrix rings of \texttt{DimensionVector(M)[i] x DimensionVector(M)[i]}, where \mbox{\texttt{\mdseries\slshape i}} runs over all vertices where \texttt{DimensionVector(M)[i]} is non\texttt{\symbol{45}}zero.



 The endomorphism is an algebra with one, and one can apply for example \texttt{RadicalOfAlgebra} to find the radical of the endomorphism ring. }

 

\subsection{\textcolor{Chapter }{FromEndMToHomMM}}
\logpage{[ 7, 3, 13 ]}\nobreak
\hyperdef{L}{X7E9554FC7A4616E1}{}
{\noindent\textcolor{FuncColor}{$\triangleright$\enspace\texttt{FromEndMToHomMM({\mdseries\slshape f})\index{FromEndMToHomMM@\texttt{FromEndMToHomMM}}
\label{FromEndMToHomMM}
}\hfill{\scriptsize (operation)}}\\


 Arguments: \mbox{\texttt{\mdseries\slshape f}} \texttt{\symbol{45}}\texttt{\symbol{45}} an element in \texttt{EndOverAlgebra(M)}.\\
 \textbf{\indent Returns:\ }
the homomorphism from \mbox{\texttt{\mdseries\slshape M}} to \mbox{\texttt{\mdseries\slshape M}} corresponding to the element \mbox{\texttt{\mdseries\slshape f}} in the endomorphism ring \texttt{EndOverAlgebra(M)} of \mbox{\texttt{\mdseries\slshape M}}.

}

 

\subsection{\textcolor{Chapter }{FromHomMMToEndM}}
\logpage{[ 7, 3, 14 ]}\nobreak
\hyperdef{L}{X7F4EECD880D88DC8}{}
{\noindent\textcolor{FuncColor}{$\triangleright$\enspace\texttt{FromHomMMToEndM({\mdseries\slshape f})\index{FromHomMMToEndM@\texttt{FromHomMMToEndM}}
\label{FromHomMMToEndM}
}\hfill{\scriptsize (operation)}}\\


 Arguments: \mbox{\texttt{\mdseries\slshape f}} \texttt{\symbol{45}}\texttt{\symbol{45}} an element in \texttt{HomOverAlgebra(M,M)}.\\
 \textbf{\indent Returns:\ }
the element \mbox{\texttt{\mdseries\slshape f}} in the endomorphism ring \texttt{EndOverAlgebra(M)} of \mbox{\texttt{\mdseries\slshape M}} corresponding to the the homomorphism from \mbox{\texttt{\mdseries\slshape M}} to \mbox{\texttt{\mdseries\slshape M}} given by \mbox{\texttt{\mdseries\slshape f}}. 

}

 

\subsection{\textcolor{Chapter }{HomFactoringThroughProjOverAlgebra}}
\logpage{[ 7, 3, 15 ]}\nobreak
\hyperdef{L}{X78B00ECD7C33C43C}{}
{\noindent\textcolor{FuncColor}{$\triangleright$\enspace\texttt{HomFactoringThroughProjOverAlgebra({\mdseries\slshape M, N})\index{HomFactoringThroughProjOverAlgebra@\texttt{HomFactoringThroughProjOverAlgebra}}
\label{HomFactoringThroughProjOverAlgebra}
}\hfill{\scriptsize (operation)}}\\


 Arguments: \mbox{\texttt{\mdseries\slshape M}}, \mbox{\texttt{\mdseries\slshape N}} \texttt{\symbol{45}} two modules.\\
 \textbf{\indent Returns:\ }
a basis for the vector space of homomorphisms from \mbox{\texttt{\mdseries\slshape M}} to \mbox{\texttt{\mdseries\slshape N}} which factors through a projective module.



 The function checks if \mbox{\texttt{\mdseries\slshape M}} and \mbox{\texttt{\mdseries\slshape N}} are modules over the same algebra, and returns an error message otherwise. }

 

\subsection{\textcolor{Chapter }{HomFromProjective}}
\logpage{[ 7, 3, 16 ]}\nobreak
\hyperdef{L}{X844682C07989D181}{}
{\noindent\textcolor{FuncColor}{$\triangleright$\enspace\texttt{HomFromProjective({\mdseries\slshape m, M})\index{HomFromProjective@\texttt{HomFromProjective}}
\label{HomFromProjective}
}\hfill{\scriptsize (operation)}}\\


 Arguments: \mbox{\texttt{\mdseries\slshape m}}, \mbox{\texttt{\mdseries\slshape M}} \texttt{\symbol{45}} an element and a module.\\
 \textbf{\indent Returns:\ }
the homomorphism from the indecomposable projective module defined by the
support of the element \mbox{\texttt{\mdseries\slshape m}} to the module \mbox{\texttt{\mdseries\slshape M}}.



 The function checks if \mbox{\texttt{\mdseries\slshape m}} is an element in \mbox{\texttt{\mdseries\slshape M}} and if the element \mbox{\texttt{\mdseries\slshape m}} is supported in only one vertex. Otherwise it returns fail. }

 

\subsection{\textcolor{Chapter }{HomOverAlgebra}}
\logpage{[ 7, 3, 17 ]}\nobreak
\hyperdef{L}{X8681E72F7FD4BFCE}{}
{\noindent\textcolor{FuncColor}{$\triangleright$\enspace\texttt{HomOverAlgebra({\mdseries\slshape M, N})\index{HomOverAlgebra@\texttt{HomOverAlgebra}}
\label{HomOverAlgebra}
}\hfill{\scriptsize (operation)}}\\
\noindent\textcolor{FuncColor}{$\triangleright$\enspace\texttt{HomOverAlgebraWithBasisFunction({\mdseries\slshape M, N})\index{HomOverAlgebraWithBasisFunction@\texttt{HomOverAlgebraWithBasisFunction}}
\label{HomOverAlgebraWithBasisFunction}
}\hfill{\scriptsize (operation)}}\\


 Arguments: \mbox{\texttt{\mdseries\slshape M}}, \mbox{\texttt{\mdseries\slshape N}} \texttt{\symbol{45}} two modules.\\
 \textbf{\indent Returns:\ }
a basis for the vector space of homomorphisms from \mbox{\texttt{\mdseries\slshape M}} to \mbox{\texttt{\mdseries\slshape N}} in the first version. In the second version it also returns a list of length
two, where the first entry is the basis found by \texttt{HomOverAlgebra} and the second entry is a function from the space of homomorphisms from \mbox{\texttt{\mdseries\slshape M}} to \mbox{\texttt{\mdseries\slshape N}} to the vector space with the basis given by the first entry.



 The function checks if \mbox{\texttt{\mdseries\slshape M}} and \mbox{\texttt{\mdseries\slshape N}} are modules over the same algebra, and returns an error message and fail
otherwise. }

 

\subsection{\textcolor{Chapter }{Image}}
\logpage{[ 7, 3, 18 ]}\nobreak
\hyperdef{L}{X87F4D35A826599C6}{}
{\noindent\textcolor{FuncColor}{$\triangleright$\enspace\texttt{Image({\mdseries\slshape f})\index{Image@\texttt{Image}}
\label{Image}
}\hfill{\scriptsize (attribute)}}\\


 Arguments: \mbox{\texttt{\mdseries\slshape f}} \texttt{\symbol{45}} a homomorphism between two modules.\\
 \textbf{\indent Returns:\ }
the image of a homomorphism \mbox{\texttt{\mdseries\slshape f}} as a module.

}

 

\subsection{\textcolor{Chapter }{ImageInclusion}}
\logpage{[ 7, 3, 19 ]}\nobreak
\hyperdef{L}{X7B076247877382A9}{}
{\noindent\textcolor{FuncColor}{$\triangleright$\enspace\texttt{ImageInclusion({\mdseries\slshape f})\index{ImageInclusion@\texttt{ImageInclusion}}
\label{ImageInclusion}
}\hfill{\scriptsize (attribute)}}\\


 Arguments: \mbox{\texttt{\mdseries\slshape f}} \texttt{\symbol{45}} a homomorphism between two modules.\\
 \textbf{\indent Returns:\ }
the inclusion of the image of a homomorphism \mbox{\texttt{\mdseries\slshape f}} into the range of \mbox{\texttt{\mdseries\slshape f}}.

}

 

\subsection{\textcolor{Chapter }{ImageProjection}}
\logpage{[ 7, 3, 20 ]}\nobreak
\hyperdef{L}{X7F336D357FE717EC}{}
{\noindent\textcolor{FuncColor}{$\triangleright$\enspace\texttt{ImageProjection({\mdseries\slshape f})\index{ImageProjection@\texttt{ImageProjection}}
\label{ImageProjection}
}\hfill{\scriptsize (attribute)}}\\


 Arguments: \mbox{\texttt{\mdseries\slshape f}} \texttt{\symbol{45}} a homomorphism between two modules.\\
 \textbf{\indent Returns:\ }
the projection from the source of \mbox{\texttt{\mdseries\slshape f}} to the image of the homomorphism \mbox{\texttt{\mdseries\slshape f}}.

}

 

\subsection{\textcolor{Chapter }{ImageProjectionInclusion}}
\logpage{[ 7, 3, 21 ]}\nobreak
\hyperdef{L}{X83B9B54A841BE792}{}
{\noindent\textcolor{FuncColor}{$\triangleright$\enspace\texttt{ImageProjectionInclusion({\mdseries\slshape f})\index{ImageProjectionInclusion@\texttt{ImageProjectionInclusion}}
\label{ImageProjectionInclusion}
}\hfill{\scriptsize (attribute)}}\\


 Arguments: \mbox{\texttt{\mdseries\slshape f}} \texttt{\symbol{45}} a homomorphism between two modules.\\
 \textbf{\indent Returns:\ }
both the projection from the source of \mbox{\texttt{\mdseries\slshape f}} to the image of the homomorphism \mbox{\texttt{\mdseries\slshape f}} and the inclusion of the image of a homomorphism \mbox{\texttt{\mdseries\slshape f}} into the range of \mbox{\texttt{\mdseries\slshape f}} as a list of two elements (first the projection and then the inclusion).

}

 

\subsection{\textcolor{Chapter }{IsomorphismOfModules}}
\logpage{[ 7, 3, 22 ]}\nobreak
\hyperdef{L}{X8212964C83E70122}{}
{\noindent\textcolor{FuncColor}{$\triangleright$\enspace\texttt{IsomorphismOfModules({\mdseries\slshape M, N})\index{IsomorphismOfModules@\texttt{IsomorphismOfModules}}
\label{IsomorphismOfModules}
}\hfill{\scriptsize (operation)}}\\


 Arguments: \mbox{\texttt{\mdseries\slshape M, N}} \texttt{\symbol{45}} two PathAlgebraMatModules.\\
 \textbf{\indent Returns:\ }
false if \mbox{\texttt{\mdseries\slshape M}} and \mbox{\texttt{\mdseries\slshape N}} are two non\texttt{\symbol{45}}isomorphic modules, otherwise it returns an
isomorphism from \mbox{\texttt{\mdseries\slshape M}} to \mbox{\texttt{\mdseries\slshape N}}.



 The function checks if \mbox{\texttt{\mdseries\slshape M}} and \mbox{\texttt{\mdseries\slshape N}} are modules over the same algebra, and returns an error message otherwise. }

 

\subsection{\textcolor{Chapter }{Kernel}}
\logpage{[ 7, 3, 23 ]}\nobreak
\hyperdef{L}{X7DCD99628504B810}{}
{\noindent\textcolor{FuncColor}{$\triangleright$\enspace\texttt{Kernel({\mdseries\slshape f})\index{Kernel@\texttt{Kernel}}
\label{Kernel}
}\hfill{\scriptsize (attribute)}}\\
\noindent\textcolor{FuncColor}{$\triangleright$\enspace\texttt{KernelInclusion({\mdseries\slshape f})\index{KernelInclusion@\texttt{KernelInclusion}}
\label{KernelInclusion}
}\hfill{\scriptsize (attribute)}}\\


 Arguments: \mbox{\texttt{\mdseries\slshape f}} \texttt{\symbol{45}} a homomorphism between two modules.\\
 \textbf{\indent Returns:\ }
the kernel of a homomorphism \mbox{\texttt{\mdseries\slshape f}} between two modules.



 The first variant \texttt{Kernel} returns the kernel of the homomorphism \mbox{\texttt{\mdseries\slshape f}} as a module, while the latter one returns the inclusion homomorphism of the
kernel into the source of the homomorphism \mbox{\texttt{\mdseries\slshape f}}. }

 
\begin{Verbatim}[commandchars=!@|,fontsize=\small,frame=single,label=Example]
  !gapprompt@gap>| !gapinput@hom := HomOverAlgebra(N,N);|
  [ <<[ 3, 2, 2 ]> ---> <[ 3, 2, 2 ]>>
      , <<[ 3, 2, 2 ]> ---> <[ 3, 2, 2 ]>>
      , <<[ 3, 2, 2 ]> ---> <[ 3, 2, 2 ]>>
      , <<[ 3, 2, 2 ]> ---> <[ 3, 2, 2 ]>>
      , <<[ 3, 2, 2 ]> ---> <[ 3, 2, 2 ]>>
       ]
  !gapprompt@gap>| !gapinput@g := hom[1];|
  <<[ 3, 2, 2 ]> ---> <[ 3, 2, 2 ]>>
  
  !gapprompt@gap>| !gapinput@M := CoKernel(g);|
  <[ 2, 2, 2 ]>
  !gapprompt@gap>| !gapinput@f := CoKernelProjection(g);|
  <<[ 3, 2, 2 ]> ---> <[ 2, 2, 2 ]>>
  
  !gapprompt@gap>| !gapinput@Range(f) = M;|
  true
  !gapprompt@gap>| !gapinput@endo := EndOverAlgebra(N);|
  <algebra-with-one of dimension 5 over Rationals>
  !gapprompt@gap>| !gapinput@RadicalOfAlgebra(endo);|
  <algebra of dimension 3 over Rationals>
  !gapprompt@gap>| !gapinput@B := BasisVectors(Basis(N));|
  [ [ [ 1, 0, 0 ], [ 0, 0 ], [ 0, 0 ] ], 
    [ [ 0, 1, 0 ], [ 0, 0 ], [ 0, 0 ] ], 
    [ [ 0, 0, 1 ], [ 0, 0 ], [ 0, 0 ] ], 
    [ [ 0, 0, 0 ], [ 1, 0 ], [ 0, 0 ] ], 
    [ [ 0, 0, 0 ], [ 0, 1 ], [ 0, 0 ] ], 
    [ [ 0, 0, 0 ], [ 0, 0 ], [ 1, 0 ] ], 
    [ [ 0, 0, 0 ], [ 0, 0 ], [ 0, 1 ] ] ]
  !gapprompt@gap>| !gapinput@p := HomFromProjective(B[1],N);|
  <<[ 1, 4, 3 ]> ---> <[ 3, 2, 2 ]>>
  
  !gapprompt@gap>| !gapinput@U := Image(p);|
  <[ 1, 2, 2 ]>
  !gapprompt@gap>| !gapinput@projinc := ImageProjectionInclusion(p);|
  [ <<[ 1, 4, 3 ]> ---> <[ 1, 2, 2 ]>>
      , <<[ 1, 2, 2 ]> ---> <[ 3, 2, 2 ]>>
       ]
  !gapprompt@gap>| !gapinput@U = Range(projinc[1]);                                      |
  true
  !gapprompt@gap>| !gapinput@Kernel(p);|
  <[ 0, 2, 1 ]> 
\end{Verbatim}
 

\subsection{\textcolor{Chapter }{LeftMinimalVersion}}
\logpage{[ 7, 3, 24 ]}\nobreak
\hyperdef{L}{X7F2D09C07F6DDBF8}{}
{\noindent\textcolor{FuncColor}{$\triangleright$\enspace\texttt{LeftMinimalVersion({\mdseries\slshape f})\index{LeftMinimalVersion@\texttt{LeftMinimalVersion}}
\label{LeftMinimalVersion}
}\hfill{\scriptsize (attribute)}}\\


 Arguments: \mbox{\texttt{\mdseries\slshape f}} \texttt{\symbol{45}} a homomorphism between two modules.\\
 \textbf{\indent Returns:\ }
 the left minimal version \mbox{\texttt{\mdseries\slshape f'}} of the homomorphism \mbox{\texttt{\mdseries\slshape f}} together with the a list \texttt{B} of modules such that the direct sum of the modules, \texttt{Range(f')} and the modules in the list \texttt{B} is isomorphic to \texttt{Range(f)}. 

}

 

\subsection{\textcolor{Chapter }{MatrixOfHomomorphismBetweenProjectives}}
\logpage{[ 7, 3, 25 ]}\nobreak
\hyperdef{L}{X7CC14E9C804AF143}{}
{\noindent\textcolor{FuncColor}{$\triangleright$\enspace\texttt{MatrixOfHomomorphismBetweenProjectives({\mdseries\slshape f})\index{MatrixOfHomomorphismBetweenProjectives@\texttt{Matrix}\-\texttt{Of}\-\texttt{Homomorphism}\-\texttt{Between}\-\texttt{Projectives}}
\label{MatrixOfHomomorphismBetweenProjectives}
}\hfill{\scriptsize (operation)}}\\


 Arguments: \mbox{\texttt{\mdseries\slshape f}} \texttt{\symbol{45}}\texttt{\symbol{45}} a homomorphism between two projective
modules.\\
 \textbf{\indent Returns:\ }
 for a homomorphism \mbox{\texttt{\mdseries\slshape f}} of projective $A$\texttt{\symbol{45}}modules from $P = \oplus v_iA$ to $P' = \oplus w_iA$, where $v_i$ and $w_i$ are vertices, the homomorphism as a matrix in $\oplus v_iAw_i$. 

}

 

\subsection{\textcolor{Chapter }{FromMatrixToHomomorphismOfProjectives}}
\logpage{[ 7, 3, 26 ]}\nobreak
\hyperdef{L}{X7F6CFBA086101B54}{}
{\noindent\textcolor{FuncColor}{$\triangleright$\enspace\texttt{FromMatrixToHomomorphismOfProjectives({\mdseries\slshape A, mat, vert1, vert0})\index{FromMatrixToHomomorphismOfProjectives@\texttt{From}\-\texttt{Matrix}\-\texttt{To}\-\texttt{Homomorphism}\-\texttt{Of}\-\texttt{Projectives}}
\label{FromMatrixToHomomorphismOfProjectives}
}\hfill{\scriptsize (operation)}}\\


 Arguments: \mbox{\texttt{\mdseries\slshape A}} \texttt{\symbol{45}}\texttt{\symbol{45}} a QuiverAlgebra, \mbox{\texttt{\mdseries\slshape mat}} \texttt{\symbol{45}}\texttt{\symbol{45}} a matrix over \mbox{\texttt{\mdseries\slshape A}}, \mbox{\texttt{\mdseries\slshape vert1, vert0}} \texttt{\symbol{45}}\texttt{\symbol{45}} two lists of vertex indices\\
 \textbf{\indent Returns:\ }
 a homomorphism of projective \mbox{\texttt{\mdseries\slshape A}}\texttt{\symbol{45}}modules from $P_1 = \oplus w_iA$ to $P_0 = \oplus v_iA$, where $w_i$ and $v_i$ are vertices determined by the two last arguments. 

}

 

\subsection{\textcolor{Chapter }{RightMinimalVersion}}
\logpage{[ 7, 3, 27 ]}\nobreak
\hyperdef{L}{X7D479D1C7C4DF55B}{}
{\noindent\textcolor{FuncColor}{$\triangleright$\enspace\texttt{RightMinimalVersion({\mdseries\slshape f})\index{RightMinimalVersion@\texttt{RightMinimalVersion}}
\label{RightMinimalVersion}
}\hfill{\scriptsize (attribute)}}\\


 Arguments: \mbox{\texttt{\mdseries\slshape f}} \texttt{\symbol{45}} a homomorphism between two modules.\\
 \textbf{\indent Returns:\ }
 the right minimal version \mbox{\texttt{\mdseries\slshape f'}} of the homomorphism \mbox{\texttt{\mdseries\slshape f}} together with the a list \texttt{B} of modules such that the direct sum of the modules, \texttt{Source(f')} and the modules on the list \texttt{B} is isomorphic to \texttt{Source(f)}. 

}

 
\begin{Verbatim}[commandchars=!@|,fontsize=\small,frame=single,label=Example]
  !gapprompt@gap>| !gapinput@H:= HomOverAlgebra(N,N);;|
  !gapprompt@gap>| !gapinput@RightMinimalVersion(H[1]);   |
  [ <<[ 1, 0, 0 ]> ---> <[ 3, 2, 2 ]>>
      , [ <[ 2, 2, 2 ]> ] ]
  !gapprompt@gap>| !gapinput@LeftMinimalVersion(H[1]);             |
  [ <<[ 3, 2, 2 ]> ---> <[ 1, 0, 0 ]>>
      , [ <[ 2, 2, 2 ]> ] ]
  !gapprompt@gap>| !gapinput@S := SimpleModules(A)[1];;|
  !gapprompt@gap>| !gapinput@MinimalRightApproximation(N,S);|
  <<[ 1, 0, 0 ]> ---> <[ 1, 0, 0 ]>>
  
  !gapprompt@gap>| !gapinput@S := SimpleModules(A)[3];;|
  !gapprompt@gap>| !gapinput@MinimalLeftApproximation(S,N);   |
  <<[ 0, 0, 1 ]> ---> <[ 2, 2, 2 ]>> 
\end{Verbatim}
 

\subsection{\textcolor{Chapter }{RadicalOfModuleInclusion}}
\logpage{[ 7, 3, 28 ]}\nobreak
\hyperdef{L}{X7BF04DAE78E98D9C}{}
{\noindent\textcolor{FuncColor}{$\triangleright$\enspace\texttt{RadicalOfModuleInclusion({\mdseries\slshape M})\index{RadicalOfModuleInclusion@\texttt{RadicalOfModuleInclusion}}
\label{RadicalOfModuleInclusion}
}\hfill{\scriptsize (attribute)}}\\


 Arguments: \mbox{\texttt{\mdseries\slshape M}} \texttt{\symbol{45}} a module.\\
 \textbf{\indent Returns:\ }
the inclusion of the radical of the module \mbox{\texttt{\mdseries\slshape M}} into \mbox{\texttt{\mdseries\slshape M}}.



 The radical of \mbox{\texttt{\mdseries\slshape M}} can be accessed using \texttt{Source}, or it can be computed directly via the command \texttt{RadicalOfModule} (\ref{RadicalOfModule}). If the algebra over which \mbox{\texttt{\mdseries\slshape M}} is a module is not a finite dimensional path algebra or an admissible quotient
of a path algebra, then it will search for other methods. }

 

\subsection{\textcolor{Chapter }{RejectOfModule}}
\logpage{[ 7, 3, 29 ]}\nobreak
\hyperdef{L}{X7A103F4F8169D8C5}{}
{\noindent\textcolor{FuncColor}{$\triangleright$\enspace\texttt{RejectOfModule({\mdseries\slshape M, N})\index{RejectOfModule@\texttt{RejectOfModule}}
\label{RejectOfModule}
}\hfill{\scriptsize (operation)}}\\


 Arguments: \mbox{\texttt{\mdseries\slshape N}}, \mbox{\texttt{\mdseries\slshape M}} \texttt{\symbol{45}}\texttt{\symbol{45}} two path algebra modules (\texttt{PathAlgebraMatModule}).\\
 \textbf{\indent Returns:\ }
the reject of the module \mbox{\texttt{\mdseries\slshape M}} in the module \mbox{\texttt{\mdseries\slshape N}} as an inclusion homomorhpism from the reject of \mbox{\texttt{\mdseries\slshape M}} into \mbox{\texttt{\mdseries\slshape N}}.

}

 

\subsection{\textcolor{Chapter }{SocleOfModuleInclusion}}
\logpage{[ 7, 3, 30 ]}\nobreak
\hyperdef{L}{X82EB23337C5F4DBB}{}
{\noindent\textcolor{FuncColor}{$\triangleright$\enspace\texttt{SocleOfModuleInclusion({\mdseries\slshape M})\index{SocleOfModuleInclusion@\texttt{SocleOfModuleInclusion}}
\label{SocleOfModuleInclusion}
}\hfill{\scriptsize (operation)}}\\


 Arguments: \mbox{\texttt{\mdseries\slshape M}} \texttt{\symbol{45}} a module.\\
 \textbf{\indent Returns:\ }
the inclusion of the socle of the module \mbox{\texttt{\mdseries\slshape M}} into \mbox{\texttt{\mdseries\slshape M}}.



 The socle of \mbox{\texttt{\mdseries\slshape M}} can be accessed using \texttt{Source}, or it can be computed directly via the command \texttt{SocleOfModule} (\ref{SocleOfModule}). }

 

\subsection{\textcolor{Chapter }{SubRepresentationInclusion}}
\logpage{[ 7, 3, 31 ]}\nobreak
\hyperdef{L}{X7E9BF05587D4A86A}{}
{\noindent\textcolor{FuncColor}{$\triangleright$\enspace\texttt{SubRepresentationInclusion({\mdseries\slshape M, gens})\index{SubRepresentationInclusion@\texttt{SubRepresentationInclusion}}
\label{SubRepresentationInclusion}
}\hfill{\scriptsize (operation)}}\\


 Arguments: \mbox{\texttt{\mdseries\slshape M}} \texttt{\symbol{45}} a module, \mbox{\texttt{\mdseries\slshape gens}} \texttt{\symbol{45}} a list of elements in \mbox{\texttt{\mdseries\slshape M}}.\\
 \textbf{\indent Returns:\ }
the inclusion of the submodule generated by the generators \mbox{\texttt{\mdseries\slshape gens}} into the module \mbox{\texttt{\mdseries\slshape M}}.



 The function checks if \mbox{\texttt{\mdseries\slshape gens}} consists of elements in \mbox{\texttt{\mdseries\slshape M}}, and returns an error message otherwise. The module given by the submodule
generated by the generators \mbox{\texttt{\mdseries\slshape gens}} can be accessed using \texttt{Source}. }

 

\subsection{\textcolor{Chapter }{TopOfModuleProjection}}
\logpage{[ 7, 3, 32 ]}\nobreak
\hyperdef{L}{X804BD7CD804E63C4}{}
{\noindent\textcolor{FuncColor}{$\triangleright$\enspace\texttt{TopOfModuleProjection({\mdseries\slshape M})\index{TopOfModuleProjection@\texttt{TopOfModuleProjection}}
\label{TopOfModuleProjection}
}\hfill{\scriptsize (operation)}}\\


 Arguments: \mbox{\texttt{\mdseries\slshape M}} \texttt{\symbol{45}} a module.\\
 \textbf{\indent Returns:\ }
the projection from the module \mbox{\texttt{\mdseries\slshape M}} to the top of the module \mbox{\texttt{\mdseries\slshape M}}.



 The module given by the top of the module \mbox{\texttt{\mdseries\slshape M}} can be accessed using \texttt{Range} of the homomorphism. }

 
\begin{Verbatim}[commandchars=!@|,fontsize=\small,frame=single,label=Example]
  !gapprompt@gap>| !gapinput@f := RadicalOfModuleInclusion(N);|
  <<[ 0, 2, 2 ]> ---> <[ 3, 2, 2 ]>>
  
  !gapprompt@gap>| !gapinput@radN := Source(f);|
  <[ 0, 2, 2 ]>
  !gapprompt@gap>| !gapinput@g := SocleOfModuleInclusion(N);|
  <<[ 1, 0, 2 ]> ---> <[ 3, 2, 2 ]>>
  
  !gapprompt@gap>| !gapinput@U := SubRepresentationInclusion(N,[B[5]+B[6],B[7]]);|
  <<[ 0, 2, 2 ]> ---> <[ 3, 2, 2 ]>>
  
  !gapprompt@gap>| !gapinput@h := TopOfModuleProjection(N);|
  <<[ 3, 2, 2 ]> ---> <[ 3, 0, 0 ]>> 
\end{Verbatim}
 

\subsection{\textcolor{Chapter }{TraceOfModule}}
\logpage{[ 7, 3, 33 ]}\nobreak
\hyperdef{L}{X7DB504E184361BBE}{}
{\noindent\textcolor{FuncColor}{$\triangleright$\enspace\texttt{TraceOfModule({\mdseries\slshape M, N})\index{TraceOfModule@\texttt{TraceOfModule}}
\label{TraceOfModule}
}\hfill{\scriptsize (operation)}}\\


 Arguments: \mbox{\texttt{\mdseries\slshape M}}, \mbox{\texttt{\mdseries\slshape C}} \texttt{\symbol{45}}\texttt{\symbol{45}} two path algebra modules (\texttt{PathAlgebraMatModule}).\\
 \textbf{\indent Returns:\ }
the trace of the module \mbox{\texttt{\mdseries\slshape M}} in the module \mbox{\texttt{\mdseries\slshape N}} as an inclusion homomorhpism from the trace of \mbox{\texttt{\mdseries\slshape M}} to \mbox{\texttt{\mdseries\slshape N}}.

}

 }

 }

 
\chapter{\textcolor{Chapter }{Homological algebra}}\label{HomologicalAlgebra}
\logpage{[ 8, 0, 0 ]}
\hyperdef{L}{X7D310EDE847865C1}{}
{
 This chapter describes the homological algebra that is implemented in QPA.

 The example used throughout this chapter is the following. 
\begin{Verbatim}[commandchars=!@|,fontsize=\small,frame=single,label=Example]
  !gapprompt@gap>| !gapinput@Q := Quiver(3,[[1,2,"a"],[1,2,"b"],[2,2,"c"],[2,3,"d"],|
  !gapprompt@>| !gapinput@[3,1,"e"]]);;|
  !gapprompt@gap>| !gapinput@KQ := PathAlgebra(Rationals, Q);;|
  !gapprompt@gap>| !gapinput@AssignGeneratorVariables(KQ);;|
  !gapprompt@gap>| !gapinput@rels := [d*e,c^2,a*c*d-b*d,e*a];;|
  !gapprompt@gap>| !gapinput@A := KQ/rels;;|
  !gapprompt@gap>| !gapinput@mat := [["a",[[1,2],[0,3],[1,5]]],["b",[[2,0],[3,0],[5,0]]],|
  !gapprompt@>| !gapinput@["c",[[0,0],[1,0]]],["d",[[1,2],[0,1]]],["e",[[0,0,0],[0,0,0]]]];;|
  !gapprompt@gap>| !gapinput@N := RightModuleOverPathAlgebra(A,mat);;|
\end{Verbatim}
 
\section{\textcolor{Chapter }{Homological algebra}}\logpage{[ 8, 1, 0 ]}
\hyperdef{L}{X7D310EDE847865C1}{}
{
 

\subsection{\textcolor{Chapter }{1stSyzygy}}
\logpage{[ 8, 1, 1 ]}\nobreak
\hyperdef{L}{X7EC01D588141BB96}{}
{\noindent\textcolor{FuncColor}{$\triangleright$\enspace\texttt{1stSyzygy({\mdseries\slshape M})\index{1stSyzygy@\texttt{1stSyzygy}}
\label{1stSyzygy}
}\hfill{\scriptsize (attribute)}}\\


 Arguments: \mbox{\texttt{\mdseries\slshape M}} \texttt{\symbol{45}}\texttt{\symbol{45}} a path algebra module (\texttt{PathAlgebraMatModule}). \\
\textbf{\indent Returns:\ }
the first syzygy of the representation \mbox{\texttt{\mdseries\slshape M}} as a representation. 

}

 

\subsection{\textcolor{Chapter }{AllComplementsOfAlmostCompleteTiltingModule}}
\logpage{[ 8, 1, 2 ]}\nobreak
\hyperdef{L}{X7902D1477B9FF116}{}
{\noindent\textcolor{FuncColor}{$\triangleright$\enspace\texttt{AllComplementsOfAlmostCompleteTiltingModule({\mdseries\slshape M, X})\index{AllComplementsOfAlmostCompleteTiltingModule@\texttt{All}\-\texttt{Complements}\-\texttt{Of}\-\texttt{Almost}\-\texttt{Complete}\-\texttt{Tilting}\-\texttt{Module}}
\label{AllComplementsOfAlmostCompleteTiltingModule}
}\hfill{\scriptsize (operation)}}\\
\noindent\textcolor{FuncColor}{$\triangleright$\enspace\texttt{AllComplementsOfAlmostCompleteCotiltingModule({\mdseries\slshape M, X})\index{AllComplementsOfAlmostCompleteCotiltingModule@\texttt{All}\-\texttt{Complements}\-\texttt{Of}\-\texttt{Almost}\-\texttt{Complete}\-\texttt{Cotilting}\-\texttt{Module}}
\label{AllComplementsOfAlmostCompleteCotiltingModule}
}\hfill{\scriptsize (operation)}}\\


 Arguments: \mbox{\texttt{\mdseries\slshape M}}, \mbox{\texttt{\mdseries\slshape X}} \texttt{\symbol{45}} two PathAlgebraMatModule's.\\
 \textbf{\indent Returns:\ }
all the complements of the almost complete (co\texttt{\symbol{45}})tilting
module \mbox{\texttt{\mdseries\slshape M}} as two exact sequences, the first is all complements which are gotten as an \mbox{\texttt{\mdseries\slshape add M}}\texttt{\symbol{45}}resolution of \mbox{\texttt{\mdseries\slshape X}} and the second is all complements which are gotten as an \mbox{\texttt{\mdseries\slshape add M}}\texttt{\symbol{45}}coresolution of \mbox{\texttt{\mdseries\slshape X}}. If there are no complements to the left of \mbox{\texttt{\mdseries\slshape X}}, then an empty list is returned. Similarly for to the right of \mbox{\texttt{\mdseries\slshape X}}. In particular, if \mbox{\texttt{\mdseries\slshape X}} has no other complements the list \texttt{[[],[]]} is returned. 

}

 

\subsection{\textcolor{Chapter }{CotiltingModule}}
\logpage{[ 8, 1, 3 ]}\nobreak
\hyperdef{L}{X81479679781E04FA}{}
{\noindent\textcolor{FuncColor}{$\triangleright$\enspace\texttt{CotiltingModule({\mdseries\slshape M, n})\index{CotiltingModule@\texttt{CotiltingModule}}
\label{CotiltingModule}
}\hfill{\scriptsize (operation)}}\\


 Arguments: \mbox{\texttt{\mdseries\slshape M}}, \mbox{\texttt{\mdseries\slshape n}} \texttt{\symbol{45}} a PathAlgebraMatModule and a positive integer.\\
 \textbf{\indent Returns:\ }
false if \mbox{\texttt{\mdseries\slshape M}} is not a cotilting module of injective dimension at most \mbox{\texttt{\mdseries\slshape a}}. Otherwise, it returns the injective dimension of \mbox{\texttt{\mdseries\slshape M}} and the resolution of all indecomposable injective modules in \mbox{\texttt{\mdseries\slshape add M}}. 

}

 

\subsection{\textcolor{Chapter }{DominantDimensionOfAlgebra}}
\logpage{[ 8, 1, 4 ]}\nobreak
\hyperdef{L}{X7D6AD2AB87865333}{}
{\noindent\textcolor{FuncColor}{$\triangleright$\enspace\texttt{DominantDimensionOfAlgebra({\mdseries\slshape A, n})\index{DominantDimensionOfAlgebra@\texttt{DominantDimensionOfAlgebra}}
\label{DominantDimensionOfAlgebra}
}\hfill{\scriptsize (operation)}}\\


 Arguments: \mbox{\texttt{\mdseries\slshape A}}, \mbox{\texttt{\mdseries\slshape n}} \texttt{\symbol{45}} a quiver algebra, a positive integer.\\
 \textbf{\indent Returns:\ }
 the dominant dimension of the algebra \mbox{\texttt{\mdseries\slshape A}} if the dominant dimension is less or equal to \mbox{\texttt{\mdseries\slshape n}}. If the function can decide that the dominant dimension is infinite, it
returns \texttt{infinity}. Otherwise, if the dominant dimension is larger than \mbox{\texttt{\mdseries\slshape n}}, then it returns \texttt{false}. 

}

 

\subsection{\textcolor{Chapter }{DominantDimensionOfModule}}
\logpage{[ 8, 1, 5 ]}\nobreak
\hyperdef{L}{X8582841B79A11466}{}
{\noindent\textcolor{FuncColor}{$\triangleright$\enspace\texttt{DominantDimensionOfModule({\mdseries\slshape M, n})\index{DominantDimensionOfModule@\texttt{DominantDimensionOfModule}}
\label{DominantDimensionOfModule}
}\hfill{\scriptsize (operation)}}\\


 Arguments: \mbox{\texttt{\mdseries\slshape M}}, \mbox{\texttt{\mdseries\slshape n}} \texttt{\symbol{45}} a PathAlgebraMatModule, a positive integer.\\
 \textbf{\indent Returns:\ }
 the dominant dimension of the module \mbox{\texttt{\mdseries\slshape M}} if the dominant dimension is less or equal to \mbox{\texttt{\mdseries\slshape n}}. If the function can decide that the dominant dimension is infinite, it
returns \texttt{infinity}. Otherwise, if the dominant dimension is larger than \mbox{\texttt{\mdseries\slshape n}}, then it returns \texttt{false}. 

}

 

\subsection{\textcolor{Chapter }{ExtAlgebraGenerators}}
\logpage{[ 8, 1, 6 ]}\nobreak
\hyperdef{L}{X7B9209CD7A46F544}{}
{\noindent\textcolor{FuncColor}{$\triangleright$\enspace\texttt{ExtAlgebraGenerators({\mdseries\slshape M, n})\index{ExtAlgebraGenerators@\texttt{ExtAlgebraGenerators}}
\label{ExtAlgebraGenerators}
}\hfill{\scriptsize (operation)}}\\


 Arguments: \mbox{\texttt{\mdseries\slshape M}} \texttt{\symbol{45}} a module, \mbox{\texttt{\mdseries\slshape n}} \texttt{\symbol{45}} a positive integer.\\
 \textbf{\indent Returns:\ }
a list of three elements, where the first element is the dimensions of
Ext\texttt{\symbol{94}}[0..n](M,M), the second element is the number of
minimal generators in the degrees [0..n], and the third element is the
generators in these degrees.



 This function computes the generators of the Ext\texttt{\symbol{45}}algebra $Ext^*(M,M)$ up to degree \mbox{\texttt{\mdseries\slshape n}}. }

 

\subsection{\textcolor{Chapter }{ExtOverAlgebra}}
\logpage{[ 8, 1, 7 ]}\nobreak
\hyperdef{L}{X787217547958E9C3}{}
{\noindent\textcolor{FuncColor}{$\triangleright$\enspace\texttt{ExtOverAlgebra({\mdseries\slshape M, N})\index{ExtOverAlgebra@\texttt{ExtOverAlgebra}}
\label{ExtOverAlgebra}
}\hfill{\scriptsize (operation)}}\\


 Arguments: \mbox{\texttt{\mdseries\slshape M}}, \mbox{\texttt{\mdseries\slshape N}} \texttt{\symbol{45}} two modules.\\
 \textbf{\indent Returns:\ }
a list of three elements \texttt{ExtOverAlgebra}, where the first element is the map from the first syzygy, $\Omega(M)$ to the projective cover, $P(M)$ of the module \mbox{\texttt{\mdseries\slshape M}}, the second element is a basis of $\Ext^1(M,N)$ in terms of elements in $\Hom(\Omega(M),N)$ and the third element is a function that takes as an argument a homomorphism
in \texttt{Hom(Omega(M),N)} and returns the coefficients of this element when written in terms of the
basis of $\Ext^1(M,N)$.



 The function checks if the arguments \mbox{\texttt{\mdseries\slshape M}} and \mbox{\texttt{\mdseries\slshape N}} are modules of the same algebra, and returns an error message otherwise. It $\Ext^1(M,N)$ is zero, an empty list is returned. }

 

\subsection{\textcolor{Chapter }{FaithfulDimension}}
\logpage{[ 8, 1, 8 ]}\nobreak
\hyperdef{L}{X7A7BF8178390DB7A}{}
{\noindent\textcolor{FuncColor}{$\triangleright$\enspace\texttt{FaithfulDimension({\mdseries\slshape M})\index{FaithfulDimension@\texttt{FaithfulDimension}}
\label{FaithfulDimension}
}\hfill{\scriptsize (attribute)}}\\


 Arguments: \mbox{\texttt{\mdseries\slshape M}} \texttt{\symbol{45}} a PathAlgebraMatModule.\\
 \textbf{\indent Returns:\ }
 the faithful dimension of the module \mbox{\texttt{\mdseries\slshape M}}. 

}

 

\subsection{\textcolor{Chapter }{GlobalDimensionOfAlgebra}}
\logpage{[ 8, 1, 9 ]}\nobreak
\hyperdef{L}{X7B2001337C48B3C9}{}
{\noindent\textcolor{FuncColor}{$\triangleright$\enspace\texttt{GlobalDimensionOfAlgebra({\mdseries\slshape A, n})\index{GlobalDimensionOfAlgebra@\texttt{GlobalDimensionOfAlgebra}}
\label{GlobalDimensionOfAlgebra}
}\hfill{\scriptsize (operation)}}\\


 Arguments: \mbox{\texttt{\mdseries\slshape A}}, \mbox{\texttt{\mdseries\slshape n}} \texttt{\symbol{45}} a quiver algebra, a positive integer.\\
 \textbf{\indent Returns:\ }
 the global dimension of \mbox{\texttt{\mdseries\slshape A}} if the global dimension is less or equal to \mbox{\texttt{\mdseries\slshape n}}. If the function can decide that the global dimension is infinite, it returns \texttt{infinity}. Otherwise, if the global dimension is larger than \mbox{\texttt{\mdseries\slshape n}}, then it returns \texttt{false}. 

}

 

\subsection{\textcolor{Chapter }{GorensteinDimension}}
\logpage{[ 8, 1, 10 ]}\nobreak
\hyperdef{L}{X8380C47D80FAAED8}{}
{\noindent\textcolor{FuncColor}{$\triangleright$\enspace\texttt{GorensteinDimension({\mdseries\slshape A})\index{GorensteinDimension@\texttt{GorensteinDimension}}
\label{GorensteinDimension}
}\hfill{\scriptsize (attribute)}}\\


 Arguments: \mbox{\texttt{\mdseries\slshape A}} \texttt{\symbol{45}} a quiver algebra.\\
 \textbf{\indent Returns:\ }
 the Gorenstein dimension of \mbox{\texttt{\mdseries\slshape A}}, if the Gorenstein dimension has been computed. Otherwise it returns an error
message. 

}

 

\subsection{\textcolor{Chapter }{GorensteinDimensionOfAlgebra}}
\logpage{[ 8, 1, 11 ]}\nobreak
\hyperdef{L}{X84BC33C07B4BBBC0}{}
{\noindent\textcolor{FuncColor}{$\triangleright$\enspace\texttt{GorensteinDimensionOfAlgebra({\mdseries\slshape A, n})\index{GorensteinDimensionOfAlgebra@\texttt{GorensteinDimensionOfAlgebra}}
\label{GorensteinDimensionOfAlgebra}
}\hfill{\scriptsize (operation)}}\\


 Arguments: \mbox{\texttt{\mdseries\slshape A}}, \mbox{\texttt{\mdseries\slshape n}} \texttt{\symbol{45}} a quiver algebra, a positive integer.\\
 \textbf{\indent Returns:\ }
 the Gorenstein dimension of \mbox{\texttt{\mdseries\slshape A}} if the Gorenstein dimension is less or equal to \mbox{\texttt{\mdseries\slshape n}}. Otherwise, if the Gorenstein dimension is larger than \mbox{\texttt{\mdseries\slshape n}}, then it returns \texttt{false}. 

}

 

\subsection{\textcolor{Chapter }{HaveFiniteCoresolutionInAddM}}
\logpage{[ 8, 1, 12 ]}\nobreak
\hyperdef{L}{X7B1A0308813BA30F}{}
{\noindent\textcolor{FuncColor}{$\triangleright$\enspace\texttt{HaveFiniteCoresolutionInAddM({\mdseries\slshape N, M, n})\index{HaveFiniteCoresolutionInAddM@\texttt{HaveFiniteCoresolutionInAddM}}
\label{HaveFiniteCoresolutionInAddM}
}\hfill{\scriptsize (operation)}}\\


 Arguments: \mbox{\texttt{\mdseries\slshape N}}, \mbox{\texttt{\mdseries\slshape M}}, \mbox{\texttt{\mdseries\slshape n}} \texttt{\symbol{45}} two PathAlgebraMatModule's and an integer.\\
 \textbf{\indent Returns:\ }
false if \mbox{\texttt{\mdseries\slshape N}} does not have a coresolution of length at most \mbox{\texttt{\mdseries\slshape n}} in \mbox{\texttt{\mdseries\slshape add M}}, otherwise it returns the coresolution of \mbox{\texttt{\mdseries\slshape N}} of length at most \mbox{\texttt{\mdseries\slshape n}}. 

}

 

\subsection{\textcolor{Chapter }{HaveFiniteResolutionInAddM}}
\logpage{[ 8, 1, 13 ]}\nobreak
\hyperdef{L}{X83271D587CD605E1}{}
{\noindent\textcolor{FuncColor}{$\triangleright$\enspace\texttt{HaveFiniteResolutionInAddM({\mdseries\slshape N, M, n})\index{HaveFiniteResolutionInAddM@\texttt{HaveFiniteResolutionInAddM}}
\label{HaveFiniteResolutionInAddM}
}\hfill{\scriptsize (operation)}}\\


 Arguments: \mbox{\texttt{\mdseries\slshape N}}, \mbox{\texttt{\mdseries\slshape M}}, \mbox{\texttt{\mdseries\slshape n}} \texttt{\symbol{45}} two PathAlgebraMatModule's and an integer.\\
 \textbf{\indent Returns:\ }
false if \mbox{\texttt{\mdseries\slshape N}} does not have a resolution of length at most \mbox{\texttt{\mdseries\slshape n}} in \mbox{\texttt{\mdseries\slshape add M}}, otherwise it returns the resolution of \mbox{\texttt{\mdseries\slshape N}} of length at most \mbox{\texttt{\mdseries\slshape n}}. 

}

 

\subsection{\textcolor{Chapter }{InjDimension}}
\logpage{[ 8, 1, 14 ]}\nobreak
\hyperdef{L}{X80C6D02F7B543CC9}{}
{\noindent\textcolor{FuncColor}{$\triangleright$\enspace\texttt{InjDimension({\mdseries\slshape M})\index{InjDimension@\texttt{InjDimension}}
\label{InjDimension}
}\hfill{\scriptsize (attribute)}}\\


 Arguments: \mbox{\texttt{\mdseries\slshape M}} \texttt{\symbol{45}} a PathAlgebraMatModule.\\
 

 If the injective dimension of the module \mbox{\texttt{\mdseries\slshape M}} has been computed, then the projective dimension is returned. }

 

\subsection{\textcolor{Chapter }{InjDimensionOfModule}}
\logpage{[ 8, 1, 15 ]}\nobreak
\hyperdef{L}{X7C16FB658383C36F}{}
{\noindent\textcolor{FuncColor}{$\triangleright$\enspace\texttt{InjDimensionOfModule({\mdseries\slshape M, n})\index{InjDimensionOfModule@\texttt{InjDimensionOfModule}}
\label{InjDimensionOfModule}
}\hfill{\scriptsize (operation)}}\\


 Arguments: \mbox{\texttt{\mdseries\slshape M, n}} \texttt{\symbol{45}} a PathAlgebraMatModule, a positive integer.\\
 \textbf{\indent Returns:\ }
Returns the injective dimension of the module \mbox{\texttt{\mdseries\slshape M}} if it is less or equal to \mbox{\texttt{\mdseries\slshape n}}. Otherwise it returns false.

}

 

\subsection{\textcolor{Chapter }{InjectiveEnvelope}}
\logpage{[ 8, 1, 16 ]}\nobreak
\hyperdef{L}{X8434C5FC87773723}{}
{\noindent\textcolor{FuncColor}{$\triangleright$\enspace\texttt{InjectiveEnvelope({\mdseries\slshape M})\index{InjectiveEnvelope@\texttt{InjectiveEnvelope}}
\label{InjectiveEnvelope}
}\hfill{\scriptsize (attribute)}}\\


 Arguments: \mbox{\texttt{\mdseries\slshape M}} \texttt{\symbol{45}} a module.\\
 \textbf{\indent Returns:\ }
the injective envelope of \mbox{\texttt{\mdseries\slshape M}}, that is, returns the map $M\to I(M)$.



 If the module \mbox{\texttt{\mdseries\slshape M}} is zero, then the zero map from \mbox{\texttt{\mdseries\slshape M}} is returned. }

 

\subsection{\textcolor{Chapter }{IsCotiltingModule}}
\logpage{[ 8, 1, 17 ]}\nobreak
\hyperdef{L}{X81910CD67E69BBF4}{}
{\noindent\textcolor{FuncColor}{$\triangleright$\enspace\texttt{IsCotiltingModule({\mdseries\slshape M})\index{IsCotiltingModule@\texttt{IsCotiltingModule}}
\label{IsCotiltingModule}
}\hfill{\scriptsize (attribute)}}\\


 Arguments: \mbox{\texttt{\mdseries\slshape M}} \texttt{\symbol{45}} a PathAlgebraMatModule.\\
 \textbf{\indent Returns:\ }
 true if the module \mbox{\texttt{\mdseries\slshape M}} has been checked to be a cotilting mdoule, otherwise it returns an error
message. 

}

 

\subsection{\textcolor{Chapter }{IsNthSyzygy}}
\logpage{[ 8, 1, 18 ]}\nobreak
\hyperdef{L}{X7AA1E69C7A1009DA}{}
{\noindent\textcolor{FuncColor}{$\triangleright$\enspace\texttt{IsNthSyzygy({\mdseries\slshape M, n})\index{IsNthSyzygy@\texttt{IsNthSyzygy}}
\label{IsNthSyzygy}
}\hfill{\scriptsize (operation)}}\\


 Arguments: \mbox{\texttt{\mdseries\slshape M}} \texttt{\symbol{45}}\texttt{\symbol{45}} a path algebra module (\texttt{PathAlgebraMatModule}), \mbox{\texttt{\mdseries\slshape n}} \texttt{\symbol{45}}\texttt{\symbol{45}} a positive integer. \\
\textbf{\indent Returns:\ }
\texttt{true}, if the representation \mbox{\texttt{\mdseries\slshape M}} is isomorphic to a direct summand of an \mbox{\texttt{\mdseries\slshape n}}\texttt{\symbol{45}}th syzygy of some module or is a projective module, and
false otherwise. 

}

 

\subsection{\textcolor{Chapter }{IsOmegaPeriodic}}
\logpage{[ 8, 1, 19 ]}\nobreak
\hyperdef{L}{X7BE6A5B586463067}{}
{\noindent\textcolor{FuncColor}{$\triangleright$\enspace\texttt{IsOmegaPeriodic({\mdseries\slshape M, n})\index{IsOmegaPeriodic@\texttt{IsOmegaPeriodic}}
\label{IsOmegaPeriodic}
}\hfill{\scriptsize (operation)}}\\


 Arguments: \mbox{\texttt{\mdseries\slshape M}} \texttt{\symbol{45}}\texttt{\symbol{45}} a path algebra module (\texttt{PathAlgebraMatModule}), \mbox{\texttt{\mdseries\slshape n}} \texttt{\symbol{45}}\texttt{\symbol{45}} a positive integer. \\
\textbf{\indent Returns:\ }
\texttt{i}, where \texttt{i} is the smallest positive integer less or equal \texttt{n} such that the representation \mbox{\texttt{\mdseries\slshape M}} is isomorphic to the \texttt{i}\texttt{\symbol{45}}th syzygy of \mbox{\texttt{\mdseries\slshape M}}, and false otherwise. 

}

 

\subsection{\textcolor{Chapter }{IsTtiltingModule}}
\logpage{[ 8, 1, 20 ]}\nobreak
\hyperdef{L}{X7CA63DA785F30E99}{}
{\noindent\textcolor{FuncColor}{$\triangleright$\enspace\texttt{IsTtiltingModule({\mdseries\slshape M})\index{IsTtiltingModule@\texttt{IsTtiltingModule}}
\label{IsTtiltingModule}
}\hfill{\scriptsize (attribute)}}\\


 Arguments: \mbox{\texttt{\mdseries\slshape M}} \texttt{\symbol{45}} a PathAlgebraMatModule.\\
 \textbf{\indent Returns:\ }
 true if the module \mbox{\texttt{\mdseries\slshape M}} has been checked to be a tilting mdoule, otherwise it returns an error
message. 

}

 

\subsection{\textcolor{Chapter }{IyamaGenerator}}
\logpage{[ 8, 1, 21 ]}\nobreak
\hyperdef{L}{X7B433C86806B60A4}{}
{\noindent\textcolor{FuncColor}{$\triangleright$\enspace\texttt{IyamaGenerator({\mdseries\slshape M})\index{IyamaGenerator@\texttt{IyamaGenerator}}
\label{IyamaGenerator}
}\hfill{\scriptsize (operation)}}\\


 Arguments: \mbox{\texttt{\mdseries\slshape M}} \texttt{\symbol{45}}\texttt{\symbol{45}} a path algebra module (\texttt{PathAlgebraMatModule}). \\
\textbf{\indent Returns:\ }
a module $N$ such that \mbox{\texttt{\mdseries\slshape M}} is a direct summand of $N$ and such that the global dimension of the endomorphism ring of $N$ is finite using the algorithm provided by Osamu Iyama (add reference here). 

}

 

\subsection{\textcolor{Chapter }{LeftApproximationByAddTHat}}
\logpage{[ 8, 1, 22 ]}\nobreak
\hyperdef{L}{X80800F4F80398356}{}
{\noindent\textcolor{FuncColor}{$\triangleright$\enspace\texttt{LeftApproximationByAddTHat({\mdseries\slshape T, M})\index{LeftApproximationByAddTHat@\texttt{LeftApproximationByAddTHat}}
\label{LeftApproximationByAddTHat}
}\hfill{\scriptsize (operation)}}\\


 Arguments: \mbox{\texttt{\mdseries\slshape T}}, \mbox{\texttt{\mdseries\slshape M}} \texttt{\symbol{45}}\texttt{\symbol{45}} two path algebra modules (\texttt{PathAlgebraMatModule}). \\
\textbf{\indent Returns:\ }
the minimal left $\widehat{\add T}$\texttt{\symbol{45}}approximation of \mbox{\texttt{\mdseries\slshape M}}. 



 The function checks if the first argument is a cotilting module, that is,
checks if the attribute of \texttt{IsCotiltingModule} is set. This attribute can be set by giving the command \texttt{CotiltingModule( T, n )} for some positive integer \texttt{n} which is at least the injective dimension of the module \mbox{\texttt{\mdseries\slshape T}}. }

 

\subsection{\textcolor{Chapter }{LeftFacMApproximation}}
\logpage{[ 8, 1, 23 ]}\nobreak
\hyperdef{L}{X806A7DF7864D9431}{}
{\noindent\textcolor{FuncColor}{$\triangleright$\enspace\texttt{LeftFacMApproximation({\mdseries\slshape C, M})\index{LeftFacMApproximation@\texttt{LeftFacMApproximation}}
\label{LeftFacMApproximation}
}\hfill{\scriptsize (operation)}}\\
\noindent\textcolor{FuncColor}{$\triangleright$\enspace\texttt{MinimalLeftFacMApproximation({\mdseries\slshape C, M})\index{MinimalLeftFacMApproximation@\texttt{MinimalLeftFacMApproximation}}
\label{MinimalLeftFacMApproximation}
}\hfill{\scriptsize (operation)}}\\


 Arguments: \mbox{\texttt{\mdseries\slshape C}}, \mbox{\texttt{\mdseries\slshape M}} \texttt{\symbol{45}}\texttt{\symbol{45}} two path algebra modules (\texttt{PathAlgebraMatModule}). \\
\textbf{\indent Returns:\ }
a left $\operatorname{Fac} M$\texttt{\symbol{45}}approximation of the module $C$, where the first version returns a not necessarily minimal left $\operatorname{Fac} M$\texttt{\symbol{45}}approximation and the second returns a minimal
approximation. 

}

 

\subsection{\textcolor{Chapter }{LeftMutationOfTiltingModuleComplement}}
\logpage{[ 8, 1, 24 ]}\nobreak
\hyperdef{L}{X82720AAD7E80714B}{}
{\noindent\textcolor{FuncColor}{$\triangleright$\enspace\texttt{LeftMutationOfTiltingModuleComplement({\mdseries\slshape M, N})\index{LeftMutationOfTiltingModuleComplement@\texttt{Left}\-\texttt{Mutation}\-\texttt{Of}\-\texttt{Tilting}\-\texttt{Module}\-\texttt{Complement}}
\label{LeftMutationOfTiltingModuleComplement}
}\hfill{\scriptsize (operation)}}\\
\noindent\textcolor{FuncColor}{$\triangleright$\enspace\texttt{LeftMutationOfCotiltingModuleComplement({\mdseries\slshape M, N})\index{LeftMutationOfCotiltingModuleComplement@\texttt{Left}\-\texttt{Mutation}\-\texttt{Of}\-\texttt{Cotilting}\-\texttt{Module}\-\texttt{Complement}}
\label{LeftMutationOfCotiltingModuleComplement}
}\hfill{\scriptsize (operation)}}\\


 Arguments: \mbox{\texttt{\mdseries\slshape M}}, \mbox{\texttt{\mdseries\slshape N}} \texttt{\symbol{45}}\texttt{\symbol{45}} two path algebra modules (\texttt{PathAlgebraMatModule}). \\
\textbf{\indent Returns:\ }
a left mutation of the complement \mbox{\texttt{\mdseries\slshape N}} of the almost complete tilting/cotilting module \mbox{\texttt{\mdseries\slshape M}}, if such a complement exists. Otherwise it returns false. 

}

 

\subsection{\textcolor{Chapter }{LeftSubMApproximation}}
\logpage{[ 8, 1, 25 ]}\nobreak
\hyperdef{L}{X86EB50377F2E910D}{}
{\noindent\textcolor{FuncColor}{$\triangleright$\enspace\texttt{LeftSubMApproximation({\mdseries\slshape C, M})\index{LeftSubMApproximation@\texttt{LeftSubMApproximation}}
\label{LeftSubMApproximation}
}\hfill{\scriptsize (operation)}}\\
\noindent\textcolor{FuncColor}{$\triangleright$\enspace\texttt{MinimalLeftSubMApproximation({\mdseries\slshape C, M})\index{MinimalLeftSubMApproximation@\texttt{MinimalLeftSubMApproximation}}
\label{MinimalLeftSubMApproximation}
}\hfill{\scriptsize (operation)}}\\


 Arguments: \mbox{\texttt{\mdseries\slshape C}}, \mbox{\texttt{\mdseries\slshape M}} \texttt{\symbol{45}}\texttt{\symbol{45}} two path algebra modules (\texttt{PathAlgebraMatModule}). \\
\textbf{\indent Returns:\ }
a minimal left $\operatorname{Sub} M$\texttt{\symbol{45}}approximation of the module $C$. 

}

 

\subsection{\textcolor{Chapter }{LiftingInclusionMorphisms}}
\logpage{[ 8, 1, 26 ]}\nobreak
\hyperdef{L}{X848B51017E269802}{}
{\noindent\textcolor{FuncColor}{$\triangleright$\enspace\texttt{LiftingInclusionMorphisms({\mdseries\slshape f, g})\index{LiftingInclusionMorphisms@\texttt{LiftingInclusionMorphisms}}
\label{LiftingInclusionMorphisms}
}\hfill{\scriptsize (operation)}}\\


 Arguments: \mbox{\texttt{\mdseries\slshape f}}, \mbox{\texttt{\mdseries\slshape g}} \texttt{\symbol{45}} two homomorphisms with common range.\\
 \textbf{\indent Returns:\ }
 a factorization of \mbox{\texttt{\mdseries\slshape g}} in terms of \mbox{\texttt{\mdseries\slshape f}}, whenever possible and \texttt{fail} otherwise. 



 Given two inclusions $f\colon B\to C$ and $g\colon A\to C$, this function constructs a morphism from $A$ to $B$, whenever the image of \mbox{\texttt{\mdseries\slshape g}} is contained in the image of \mbox{\texttt{\mdseries\slshape f}}. Otherwise the function returns fail. The function checks if \mbox{\texttt{\mdseries\slshape f}} and \mbox{\texttt{\mdseries\slshape g}} are one\texttt{\symbol{45}}to\texttt{\symbol{45}}one, if they have the same
range and if the image of \mbox{\texttt{\mdseries\slshape g}} is contained in the image of \mbox{\texttt{\mdseries\slshape f}}. }

 

\subsection{\textcolor{Chapter }{LiftingMorphismFromProjective}}
\logpage{[ 8, 1, 27 ]}\nobreak
\hyperdef{L}{X7A10EF0587ADA535}{}
{\noindent\textcolor{FuncColor}{$\triangleright$\enspace\texttt{LiftingMorphismFromProjective({\mdseries\slshape f, g})\index{LiftingMorphismFromProjective@\texttt{LiftingMorphismFromProjective}}
\label{LiftingMorphismFromProjective}
}\hfill{\scriptsize (operation)}}\\


 Arguments: \mbox{\texttt{\mdseries\slshape f}}, \mbox{\texttt{\mdseries\slshape g}} \texttt{\symbol{45}} two homomorphisms with common range.\\
 \textbf{\indent Returns:\ }
 a factorization of \mbox{\texttt{\mdseries\slshape g}} in terms of \mbox{\texttt{\mdseries\slshape f}}, whenever possible and \texttt{fail} otherwise. 



 Given two morphisms $f\colon B\to C$ and $g\colon P\to C$, where $P$ is a direct sum of indecomposable projective modules constructed via \texttt{DirectSumOfQPAModules} and \mbox{\texttt{\mdseries\slshape f}} an epimorphism, this function finds a lifting of \mbox{\texttt{\mdseries\slshape g}} to $B$. The function checks if $P$ is a direct sum of indecomposable projective modules, if \mbox{\texttt{\mdseries\slshape f}} is onto and if \mbox{\texttt{\mdseries\slshape f}} and \mbox{\texttt{\mdseries\slshape g}} have the same range. }

 
\begin{Verbatim}[commandchars=!@|,fontsize=\small,frame=single,label=Example]
  !gapprompt@gap>| !gapinput@B := BasisVectors(Basis(N));|
  [ [ [ 1, 0, 0 ], [ 0, 0 ], [ 0, 0 ] ], 
    [ [ 0, 1, 0 ], [ 0, 0 ], [ 0, 0 ] ], 
    [ [ 0, 0, 1 ], [ 0, 0 ], [ 0, 0 ] ], 
    [ [ 0, 0, 0 ], [ 1, 0 ], [ 0, 0 ] ], 
    [ [ 0, 0, 0 ], [ 0, 1 ], [ 0, 0 ] ], 
    [ [ 0, 0, 0 ], [ 0, 0 ], [ 1, 0 ] ], 
    [ [ 0, 0, 0 ], [ 0, 0 ], [ 0, 1 ] ] ]
  !gapprompt@gap>| !gapinput@g := SubRepresentationInclusion(N,[B[1],B[4]]);|
  <<[ 1, 2, 2 ]> ---> <[ 3, 2, 2 ]>>
  
  !gapprompt@gap>| !gapinput@f := SubRepresentationInclusion(N,[B[1],B[2]]);|
  <<[ 2, 2, 2 ]> ---> <[ 3, 2, 2 ]>>
  
  !gapprompt@gap>| !gapinput@LiftingInclusionMorphisms(f,g);|
  <<[ 1, 2, 2 ]> ---> <[ 2, 2, 2 ]>>
  
  !gapprompt@gap>| !gapinput@S := SimpleModules(A); |
  [ <[ 1, 0, 0 ]>, <[ 0, 1, 0 ]>, <[ 0, 0, 1 ]> ]
  !gapprompt@gap>| !gapinput@homNS := HomOverAlgebra(N,S[1]);|
  [ <<[ 3, 2, 2 ]> ---> <[ 1, 0, 0 ]>>
      , <<[ 3, 2, 2 ]> ---> <[ 1, 0, 0 ]>>
      , <<[ 3, 2, 2 ]> ---> <[ 1, 0, 0 ]>>
       ]
  !gapprompt@gap>| !gapinput@f := homNS[1];|
  <<[ 3, 2, 2 ]> ---> <[ 1, 0, 0 ]>>
  
  !gapprompt@gap>| !gapinput@p := ProjectiveCover(S[1]);|
  <<[ 1, 4, 3 ]> ---> <[ 1, 0, 0 ]>>
  
  !gapprompt@gap>| !gapinput@LiftingMorphismFromProjective(f,p);|
  <<[ 1, 4, 3 ]> ---> <[ 3, 2, 2 ]>>
  [ [(1)*v1], [(1)*v2], [(1)*v3], [(1)*a], [(1)*b], [(1)*c], [(1)*d], [(1)*e] 
   ] )> > 
\end{Verbatim}
 

\subsection{\textcolor{Chapter }{LeftApproximationByAddM}}
\logpage{[ 8, 1, 28 ]}\nobreak
\hyperdef{L}{X853E593B83DDD4B8}{}
{\noindent\textcolor{FuncColor}{$\triangleright$\enspace\texttt{LeftApproximationByAddM({\mdseries\slshape C, M})\index{LeftApproximationByAddM@\texttt{LeftApproximationByAddM}}
\label{LeftApproximationByAddM}
}\hfill{\scriptsize (operation)}}\\
\noindent\textcolor{FuncColor}{$\triangleright$\enspace\texttt{MinimalLeftAddMApproximation({\mdseries\slshape C, M})\index{MinimalLeftAddMApproximation@\texttt{MinimalLeftAddMApproximation}}
\label{MinimalLeftAddMApproximation}
}\hfill{\scriptsize (attribute)}}\\
\noindent\textcolor{FuncColor}{$\triangleright$\enspace\texttt{MinimalLeftApproximation({\mdseries\slshape C, M})\index{MinimalLeftApproximation@\texttt{MinimalLeftApproximation}}
\label{MinimalLeftApproximation}
}\hfill{\scriptsize (attribute)}}\\


 Arguments: \mbox{\texttt{\mdseries\slshape C}}, \mbox{\texttt{\mdseries\slshape M}} \texttt{\symbol{45}} two modules.\\
 \textbf{\indent Returns:\ }
 the minimal left $\add M$\texttt{\symbol{45}}approximation in the two last versions of the module \mbox{\texttt{\mdseries\slshape C}}. In the first it returns some left approximation. Note: The order of the
arguments is opposite of the order for minimal right approximations. 

}

 

\subsection{\textcolor{Chapter }{RightApproximationByAddM}}
\logpage{[ 8, 1, 29 ]}\nobreak
\hyperdef{L}{X7F689558823AC784}{}
{\noindent\textcolor{FuncColor}{$\triangleright$\enspace\texttt{RightApproximationByAddM({\mdseries\slshape M, C/modulelist, C})\index{RightApproximationByAddM@\texttt{RightApproximationByAddM}}
\label{RightApproximationByAddM}
}\hfill{\scriptsize (operation)}}\\
\noindent\textcolor{FuncColor}{$\triangleright$\enspace\texttt{MinimalRightApproximation({\mdseries\slshape M, C})\index{MinimalRightApproximation@\texttt{MinimalRightApproximation}}
\label{MinimalRightApproximation}
}\hfill{\scriptsize (attribute)}}\\
\noindent\textcolor{FuncColor}{$\triangleright$\enspace\texttt{MinimalRightAddMApproximation({\mdseries\slshape M, C})\index{MinimalRightAddMApproximation@\texttt{MinimalRightAddMApproximation}}
\label{MinimalRightAddMApproximation}
}\hfill{\scriptsize (attribute)}}\\


 Arguments: \mbox{\texttt{\mdseries\slshape M}}, \mbox{\texttt{\mdseries\slshape C}} \texttt{\symbol{45}} two modules.\\
 \textbf{\indent Returns:\ }
 the minimal right $\add M$\texttt{\symbol{45}}approximation in the two last versions of the module \mbox{\texttt{\mdseries\slshape C}}. In the two first it returns some right approximation, where in the first
version the input is two modules, while in the second version the input is a
list of modules and a module. Note: The order of the arguments is opposite of
the order for minimal left approximations. 

}

 

\subsection{\textcolor{Chapter }{RadicalRightApproximationByAddM}}
\logpage{[ 8, 1, 30 ]}\nobreak
\hyperdef{L}{X7947619F7C8539E0}{}
{\noindent\textcolor{FuncColor}{$\triangleright$\enspace\texttt{RadicalRightApproximationByAddM({\mdseries\slshape modulelist, t})\index{RadicalRightApproximationByAddM@\texttt{RadicalRightApproximationByAddM}}
\label{RadicalRightApproximationByAddM}
}\hfill{\scriptsize (operation)}}\\


 Arguments: \mbox{\texttt{\mdseries\slshape modulelist}}, \mbox{\texttt{\mdseries\slshape t}} \texttt{\symbol{45}} a list of modules and an index of this list.\\
 \textbf{\indent Returns:\ }
 a radical right approximation of \texttt{moduleslist[ t ]} by the additive closure of the modules in the list of modules \mbox{\texttt{\mdseries\slshape modulelist}}, that is, returns a homomorphism $f\colon M_{M_t} \to M_t$, where $M_t$ is the t\texttt{\symbol{45}}th module in the \mbox{\texttt{\mdseries\slshape modulelist}}. 

}

 

\subsection{\textcolor{Chapter }{MorphismOnKernel}}
\logpage{[ 8, 1, 31 ]}\nobreak
\hyperdef{L}{X87AB8D2E8660869D}{}
{\noindent\textcolor{FuncColor}{$\triangleright$\enspace\texttt{MorphismOnKernel({\mdseries\slshape f, g, alpha, beta})\index{MorphismOnKernel@\texttt{MorphismOnKernel}}
\label{MorphismOnKernel}
}\hfill{\scriptsize (operation)}}\\
\noindent\textcolor{FuncColor}{$\triangleright$\enspace\texttt{MorphismOnImage({\mdseries\slshape f, g, alpha, beta})\index{MorphismOnImage@\texttt{MorphismOnImage}}
\label{MorphismOnImage}
}\hfill{\scriptsize (operation)}}\\
\noindent\textcolor{FuncColor}{$\triangleright$\enspace\texttt{MorphismOnCoKernel({\mdseries\slshape f, g, alpha, beta})\index{MorphismOnCoKernel@\texttt{MorphismOnCoKernel}}
\label{MorphismOnCoKernel}
}\hfill{\scriptsize (operation)}}\\


 Arguments: \mbox{\texttt{\mdseries\slshape f}}, \mbox{\texttt{\mdseries\slshape g}}, \mbox{\texttt{\mdseries\slshape alpha}}, \mbox{\texttt{\mdseries\slshape beta}} \texttt{\symbol{45}} four homomorphisms of modules.\\
 \textbf{\indent Returns:\ }
the morphism induced on the kernels, the images or the cokernels of the
morphisms \mbox{\texttt{\mdseries\slshape f}} and \mbox{\texttt{\mdseries\slshape g}}, respectively, whenever $f\colon A\to B$, $\beta\colon B\to B'$, $\alpha\colon A\to A'$ and $g\colon A'\to B'$ forms a commutative diagram.



 It is checked if \mbox{\texttt{\mdseries\slshape f}}, \mbox{\texttt{\mdseries\slshape g}}, \mbox{\texttt{\mdseries\slshape alpha}}, \mbox{\texttt{\mdseries\slshape beta}} forms a commutative diagram, that is, if $f \beta - \alpha g = 0$. }

 
\begin{Verbatim}[commandchars=!@|,fontsize=\small,frame=single,label=Example]
  !gapprompt@gap>| !gapinput@hom := HomOverAlgebra(N,N);|
  [ <<[ 3, 2, 2 ]> ---> <[ 3, 2, 2 ]>>
      , <<[ 3, 2, 2 ]> ---> <[ 3, 2, 2 ]>>
      , <<[ 3, 2, 2 ]> ---> <[ 3, 2, 2 ]>>
      , <<[ 3, 2, 2 ]> ---> <[ 3, 2, 2 ]>>
      , <<[ 3, 2, 2 ]> ---> <[ 3, 2, 2 ]>>
       ]
  !gapprompt@gap>| !gapinput@g := MorphismOnKernel(hom[1],hom[2],hom[1],hom[2]);|
  <<[ 2, 2, 2 ]> ---> <[ 2, 2, 2 ]>>
  
  !gapprompt@gap>| !gapinput@IsomorphicModules(Source(g),Range(g));|
  true
  !gapprompt@gap>| !gapinput@p := ProjectiveCover(N);|
  <<[ 3, 12, 9 ]> ---> <[ 3, 2, 2 ]>>
  
  !gapprompt@gap>| !gapinput@N1 := Kernel(p);|
  <[ 0, 10, 7 ]>
  !gapprompt@gap>| !gapinput@pullback := PullBack(p,hom[1]);|
  [ <<[ 3, 12, 9 ]> ---> <[ 3, 2, 2 ]>>
      , <<[ 3, 12, 9 ]> ---> <[ 3, 12, 9 ]>>
       ]
  !gapprompt@gap>| !gapinput@Kernel(pullback[1]);|
  <[ 0, 10, 7 ]>
  !gapprompt@gap>| !gapinput@IsomorphicModules(N1,Kernel(pullback[1]));|
  true
  !gapprompt@gap>| !gapinput@t := LiftingMorphismFromProjective(p,p*hom[1]);|
  <<[ 3, 12, 9 ]> ---> <[ 3, 12, 9 ]>>
  
  !gapprompt@gap>| !gapinput@s := MorphismOnKernel(p,p,t,hom[1]);    |
  <<[ 0, 10, 7 ]> ---> <[ 0, 10, 7 ]>>
  
  !gapprompt@gap>| !gapinput@Source(s) = N1;|
  true
  !gapprompt@gap>| !gapinput@q := KernelInclusion(p);|
  <<[ 0, 10, 7 ]> ---> <[ 3, 12, 9 ]>>
  
  !gapprompt@gap>| !gapinput@pushout := PushOut(q,s);|
  [ <<[ 0, 10, 7 ]> ---> <[ 3, 12, 9 ]>>
      , <<[ 3, 12, 9 ]> ---> <[ 3, 12, 9 ]>>
       ]
  !gapprompt@gap>| !gapinput@U := CoKernel(pushout[1]);|
  <[ 3, 2, 2 ]>
  !gapprompt@gap>| !gapinput@IsomorphicModules(U,N);|
  true 
\end{Verbatim}
 

\subsection{\textcolor{Chapter }{NthSyzygy}}
\logpage{[ 8, 1, 32 ]}\nobreak
\hyperdef{L}{X7EB51A487B5B239D}{}
{\noindent\textcolor{FuncColor}{$\triangleright$\enspace\texttt{NthSyzygy({\mdseries\slshape M, n})\index{NthSyzygy@\texttt{NthSyzygy}}
\label{NthSyzygy}
}\hfill{\scriptsize (operation)}}\\


 Arguments: \mbox{\texttt{\mdseries\slshape M}} \texttt{\symbol{45}}\texttt{\symbol{45}} a path algebra module (\texttt{PathAlgebraMatModule}), \mbox{\texttt{\mdseries\slshape n}} \texttt{\symbol{45}}\texttt{\symbol{45}} a non\texttt{\symbol{45}}negative
integer. \\
\textbf{\indent Returns:\ }
the \mbox{\texttt{\mdseries\slshape n}}\texttt{\symbol{45}}th syzygy of \mbox{\texttt{\mdseries\slshape M}}. 



 This functions computes a projective resolution of \mbox{\texttt{\mdseries\slshape M}} and finds the \mbox{\texttt{\mdseries\slshape n}}\texttt{\symbol{45}}th syzygy of the module \mbox{\texttt{\mdseries\slshape M}}. }

 

\subsection{\textcolor{Chapter }{NumberOfComplementsOfAlmostCompleteTiltingModule}}
\logpage{[ 8, 1, 33 ]}\nobreak
\hyperdef{L}{X788CD7B282582D93}{}
{\noindent\textcolor{FuncColor}{$\triangleright$\enspace\texttt{NumberOfComplementsOfAlmostCompleteTiltingModule({\mdseries\slshape M})\index{NumberOfComplementsOfAlmostCompleteTiltingModule@\texttt{Number}\-\texttt{Of}\-\texttt{Complements}\-\texttt{Of}\-\texttt{Almost}\-\texttt{Complete}\-\texttt{Tilting}\-\texttt{Module}}
\label{NumberOfComplementsOfAlmostCompleteTiltingModule}
}\hfill{\scriptsize (operation)}}\\
\noindent\textcolor{FuncColor}{$\triangleright$\enspace\texttt{NumberOfComplementsOfAlmostCompleteCotiltingModule({\mdseries\slshape M})\index{NumberOfComplementsOfAlmostCompleteCotiltingModule@\texttt{Number}\-\texttt{Of}\-\texttt{Complements}\-\texttt{Of}\-\texttt{Almost}\-\texttt{Complete}\-\texttt{Cotilting}\-\texttt{Module}}
\label{NumberOfComplementsOfAlmostCompleteCotiltingModule}
}\hfill{\scriptsize (operation)}}\\


 Arguments: \mbox{\texttt{\mdseries\slshape M}} \texttt{\symbol{45}}\texttt{\symbol{45}} a PathAlgebraMatModule.\\
\textbf{\indent Returns:\ }
the number complements of an almost complete tilting/cotilting module \mbox{\texttt{\mdseries\slshape M}}, assuming that \mbox{\texttt{\mdseries\slshape M}} is an almost complete tilting module. 

}

 

\subsection{\textcolor{Chapter }{ProjDimension}}
\logpage{[ 8, 1, 34 ]}\nobreak
\hyperdef{L}{X7D3A779D7EFF32D1}{}
{\noindent\textcolor{FuncColor}{$\triangleright$\enspace\texttt{ProjDimension({\mdseries\slshape M})\index{ProjDimension@\texttt{ProjDimension}}
\label{ProjDimension}
}\hfill{\scriptsize (attribute)}}\\


 Arguments: \mbox{\texttt{\mdseries\slshape M}} \texttt{\symbol{45}} a PathAlgebraMatModule.\\
 \textbf{\indent Returns:\ }
the projective dimension of the module \mbox{\texttt{\mdseries\slshape M}}, if it has been computed. 

}

 

\subsection{\textcolor{Chapter }{ProjDimensionOfModule}}
\logpage{[ 8, 1, 35 ]}\nobreak
\hyperdef{L}{X86DA085884BAF968}{}
{\noindent\textcolor{FuncColor}{$\triangleright$\enspace\texttt{ProjDimensionOfModule({\mdseries\slshape M, n})\index{ProjDimensionOfModule@\texttt{ProjDimensionOfModule}}
\label{ProjDimensionOfModule}
}\hfill{\scriptsize (operation)}}\\


 Arguments: \mbox{\texttt{\mdseries\slshape M, n}} \texttt{\symbol{45}} a PathAlgebraMatModule, a positive integer.\\
 \textbf{\indent Returns:\ }
Returns the projective dimension of the module \mbox{\texttt{\mdseries\slshape M}} if it is less or equal to \mbox{\texttt{\mdseries\slshape n}}. Otherwise it returns false.

}

 

\subsection{\textcolor{Chapter }{ProjectiveCover}}
\logpage{[ 8, 1, 36 ]}\nobreak
\hyperdef{L}{X83257E2F7F8E0068}{}
{\noindent\textcolor{FuncColor}{$\triangleright$\enspace\texttt{ProjectiveCover({\mdseries\slshape M})\index{ProjectiveCover@\texttt{ProjectiveCover}}
\label{ProjectiveCover}
}\hfill{\scriptsize (attribute)}}\\


 Arguments: \mbox{\texttt{\mdseries\slshape M}} \texttt{\symbol{45}} a module.\\
 \textbf{\indent Returns:\ }
the projective cover of \mbox{\texttt{\mdseries\slshape M}}, that is, returns the map $P(M)\to M$.



 If the module \mbox{\texttt{\mdseries\slshape M}} is zero, then the zero map to \mbox{\texttt{\mdseries\slshape M}} is returned. }

 

\subsection{\textcolor{Chapter }{ProjectiveResolutionOfPathAlgebraModule}}
\logpage{[ 8, 1, 37 ]}\nobreak
\hyperdef{L}{X8346B3997D9AB706}{}
{\noindent\textcolor{FuncColor}{$\triangleright$\enspace\texttt{ProjectiveResolutionOfPathAlgebraModule({\mdseries\slshape M, n})\index{ProjectiveResolutionOfPathAlgebraModule@\texttt{Projective}\-\texttt{Resolution}\-\texttt{Of}\-\texttt{Path}\-\texttt{Algebra}\-\texttt{Module}}
\label{ProjectiveResolutionOfPathAlgebraModule}
}\hfill{\scriptsize (operation)}}\\


 Arguments: \mbox{\texttt{\mdseries\slshape M}} \texttt{\symbol{45}} a path algebra module (\texttt{PathAlgebraMatModule}), \mbox{\texttt{\mdseries\slshape n}} \texttt{\symbol{45}} a positive integer.\\
 \textbf{\indent Returns:\ }
in terms of attributes \texttt{RProjectives}, \texttt{ProjectivesFList} and \texttt{Maps} a projective resolution of \mbox{\texttt{\mdseries\slshape M}} out to stage \mbox{\texttt{\mdseries\slshape n}}, where \texttt{RProjectives} are the projectives in the resolution lifted up to projectives over the path
algebra, \texttt{ProjectivesFList} are the generators of the projective modules given in \texttt{RProjectives} in terms of elements in the first projective in the resolution and \texttt{Maps} contains the information about the maps in the resolution.



 The algorithm for computing this projective resolution is based on the paper \cite{GreenSolbergZacharia}. In addition, the algebra over which the modules are defined is available via
the attribute \texttt{ParentAlgebra}. }

 

\subsection{\textcolor{Chapter }{ProjectiveResolutionOfSimpleModuleOverEndo}}
\logpage{[ 8, 1, 38 ]}\nobreak
\hyperdef{L}{X80678790799AEC16}{}
{\noindent\textcolor{FuncColor}{$\triangleright$\enspace\texttt{ProjectiveResolutionOfSimpleModuleOverEndo({\mdseries\slshape modulelist, t, length})\index{ProjectiveResolutionOfSimpleModuleOverEndo@\texttt{Projective}\-\texttt{Resolution}\-\texttt{Of}\-\texttt{Simple}\-\texttt{Module}\-\texttt{Over}\-\texttt{Endo}}
\label{ProjectiveResolutionOfSimpleModuleOverEndo}
}\hfill{\scriptsize (operation)}}\\


 Arguments: \mbox{\texttt{\mdseries\slshape modulelist}} \texttt{\symbol{45}} a list of module, \mbox{\texttt{\mdseries\slshape t}} \texttt{\symbol{45}} an index of the list of modules, \mbox{\texttt{\mdseries\slshape length}} \texttt{\symbol{45}} length of the resolution.\\
 \textbf{\indent Returns:\ }
information about the projective dimension and
non\texttt{\symbol{45}}projective summands of the syzygies of the simple
module corresponding to the \mbox{\texttt{\mdseries\slshape t}}\texttt{\symbol{45}}th indecomposable projective module over the endomorphism
ring of the direct sum of all the modules in \mbox{\texttt{\mdseries\slshape modulelist}} (all assumed to be indecomposable). The non\texttt{\symbol{45}}projective
summands in the syzygies from the second syzygy up to the \mbox{\texttt{\mdseries\slshape length}}\texttt{\symbol{45}}syzygy are always returned. If the projective dimension is
less or equal to \mbox{\texttt{\mdseries\slshape length}}, the projective dimension is returned. Otherwise, it returns that the
projective dimension is bigger that \mbox{\texttt{\mdseries\slshape length}}. The output has the format \texttt{[ info on projective dimension, syzygies ]}. 

}

 

\subsection{\textcolor{Chapter }{PullBack}}
\logpage{[ 8, 1, 39 ]}\nobreak
\hyperdef{L}{X7C705F2A79F8E43C}{}
{\noindent\textcolor{FuncColor}{$\triangleright$\enspace\texttt{PullBack({\mdseries\slshape f, g})\index{PullBack@\texttt{PullBack}}
\label{PullBack}
}\hfill{\scriptsize (operation)}}\\


 Arguments: \mbox{\texttt{\mdseries\slshape f}}, \mbox{\texttt{\mdseries\slshape g}} \texttt{\symbol{45}} two homomorphisms with a common range.\\
 \textbf{\indent Returns:\ }
the pullback of the maps \mbox{\texttt{\mdseries\slshape f}} and \mbox{\texttt{\mdseries\slshape g}}.



 It is checked if \mbox{\texttt{\mdseries\slshape f}} and \mbox{\texttt{\mdseries\slshape g}} have the same range. Given the input $f\colon A\to B$ (horizontal map) and $g\colon C\to B$ (vertical map), the pullback $E$ is returned as the two homomorphisms $[f',g']$, where $f'\colon E\to C$ (horizontal map) and $g'\colon E\to A$ (vertical map). }

 

\subsection{\textcolor{Chapter }{PushOut}}
\logpage{[ 8, 1, 40 ]}\nobreak
\hyperdef{L}{X81A2D49D85923894}{}
{\noindent\textcolor{FuncColor}{$\triangleright$\enspace\texttt{PushOut({\mdseries\slshape f, g})\index{PushOut@\texttt{PushOut}}
\label{PushOut}
}\hfill{\scriptsize (operation)}}\\


 Arguments: \mbox{\texttt{\mdseries\slshape f}}, \mbox{\texttt{\mdseries\slshape g}} \texttt{\symbol{45}} two homomorphisms between modules with a common source.\\
 \textbf{\indent Returns:\ }
the pushout of the maps \mbox{\texttt{\mdseries\slshape f}} and \mbox{\texttt{\mdseries\slshape g}}.



 It is checked if \mbox{\texttt{\mdseries\slshape f}} and \mbox{\texttt{\mdseries\slshape g}} have the same source. Given the input $f\colon A\to B$ (horizontal map) and $g\colon A\to C$ (vertical map), the pushout $E$ is returned as the two homomorphisms $[f',g']$, where $f'\colon C\to E$ (horizontal map) and $g'\colon B\to E$ (vertical map). }

 
\begin{Verbatim}[commandchars=!@|,fontsize=\small,frame=single,label=Example]
  !gapprompt@gap>| !gapinput@S := SimpleModules(A);|
  [ <[ 1, 0, 0 ]>, <[ 0, 1, 0 ]>, <[ 0, 0, 1 ]> ]
  !gapprompt@gap>| !gapinput@Ext := ExtOverAlgebra(S[2],S[2]);|
  [ <<[ 0, 1, 2 ]> ---> <[ 0, 2, 2 ]>>
      , [ <<[ 0, 1, 2 ]> ---> <[ 0, 1, 0 ]>>
           ], function( map ) ... end ]
  !gapprompt@gap>| !gapinput@Length(Ext[2]);|
  1
  !gapprompt@gap>| !gapinput@# i.e. Ext^1(S[2],S[2]) is 1-dimensional|
  !gapprompt@gap>| !gapinput@pushout := PushOut(Ext[2][1],Ext[1]);   |
  [ <<[ 0, 2, 2 ]> ---> <[ 0, 2, 0 ]>>
      , <<[ 0, 1, 0 ]> ---> <[ 0, 2, 0 ]>>
       ]
  !gapprompt@gap>| !gapinput@f := CoKernelProjection(pushout[1]);|
  <<[ 0, 2, 0 ]> ---> <[ 0, 0, 0 ]>>
  
  !gapprompt@gap>| !gapinput@U := Range(pushout[1]); |
  <[ 0, 2, 0 ]> 
\end{Verbatim}
 

\subsection{\textcolor{Chapter }{RightApproximationByPerpT}}
\logpage{[ 8, 1, 41 ]}\nobreak
\hyperdef{L}{X7C8B72157BDE0957}{}
{\noindent\textcolor{FuncColor}{$\triangleright$\enspace\texttt{RightApproximationByPerpT({\mdseries\slshape T, M})\index{RightApproximationByPerpT@\texttt{RightApproximationByPerpT}}
\label{RightApproximationByPerpT}
}\hfill{\scriptsize (operation)}}\\


 Arguments: \mbox{\texttt{\mdseries\slshape T}}, \mbox{\texttt{\mdseries\slshape M}} \texttt{\symbol{45}}\texttt{\symbol{45}} two path algebra modules (\texttt{PathAlgebraMatModule}). \\
\textbf{\indent Returns:\ }
the minimal right $^\perp T$\texttt{\symbol{45}}approximation of \mbox{\texttt{\mdseries\slshape M}}. 



 The function checks if the first argument is a cotilting module, that is,
checks if the attribute of \texttt{IsCotiltingModule} is set. This attribute can be set by giving the command \texttt{CotiltingModule( T, n )} for some positive integer \texttt{n} which is at least the injective dimension of the module \mbox{\texttt{\mdseries\slshape T}}. }

 

\subsection{\textcolor{Chapter }{RightFacMApproximation}}
\logpage{[ 8, 1, 42 ]}\nobreak
\hyperdef{L}{X873C632A83DF48E2}{}
{\noindent\textcolor{FuncColor}{$\triangleright$\enspace\texttt{RightFacMApproximation({\mdseries\slshape M, C})\index{RightFacMApproximation@\texttt{RightFacMApproximation}}
\label{RightFacMApproximation}
}\hfill{\scriptsize (operation)}}\\
\noindent\textcolor{FuncColor}{$\triangleright$\enspace\texttt{MinimalRightFacMApproximation({\mdseries\slshape M, C})\index{MinimalRightFacMApproximation@\texttt{MinimalRightFacMApproximation}}
\label{MinimalRightFacMApproximation}
}\hfill{\scriptsize (operation)}}\\


 Arguments: \mbox{\texttt{\mdseries\slshape M}}, \mbox{\texttt{\mdseries\slshape C}} \texttt{\symbol{45}}\texttt{\symbol{45}} two path algebra modules (\texttt{PathAlgebraMatModule}). \\
\textbf{\indent Returns:\ }
a minimal right $\operatorname{Fac} M$\texttt{\symbol{45}}approximation of the module $C$. 

}

 

\subsection{\textcolor{Chapter }{RightMutationOfTiltingModuleComplement}}
\logpage{[ 8, 1, 43 ]}\nobreak
\hyperdef{L}{X8018953F7951F8F2}{}
{\noindent\textcolor{FuncColor}{$\triangleright$\enspace\texttt{RightMutationOfTiltingModuleComplement({\mdseries\slshape M, N})\index{RightMutationOfTiltingModuleComplement@\texttt{Right}\-\texttt{Mutation}\-\texttt{Of}\-\texttt{Tilting}\-\texttt{Module}\-\texttt{Complement}}
\label{RightMutationOfTiltingModuleComplement}
}\hfill{\scriptsize (operation)}}\\
\noindent\textcolor{FuncColor}{$\triangleright$\enspace\texttt{RightMutationOfCotiltingModuleComplement({\mdseries\slshape M, N})\index{RightMutationOfCotiltingModuleComplement@\texttt{Right}\-\texttt{Mutation}\-\texttt{Of}\-\texttt{Cotilting}\-\texttt{Module}\-\texttt{Complement}}
\label{RightMutationOfCotiltingModuleComplement}
}\hfill{\scriptsize (operation)}}\\


 Arguments: \mbox{\texttt{\mdseries\slshape M}}, \mbox{\texttt{\mdseries\slshape N}} \texttt{\symbol{45}}\texttt{\symbol{45}} two path algebra modules (\texttt{PathAlgebraMatModule}). \\
\textbf{\indent Returns:\ }
a right mutation of the complement \mbox{\texttt{\mdseries\slshape N}} of the almost complete tilting/cotilting module \mbox{\texttt{\mdseries\slshape M}}, if such a complement exists. Otherwise it returns false. 

}

 

\subsection{\textcolor{Chapter }{RightSubMApproximation}}
\logpage{[ 8, 1, 44 ]}\nobreak
\hyperdef{L}{X81BD4EEA831D4A54}{}
{\noindent\textcolor{FuncColor}{$\triangleright$\enspace\texttt{RightSubMApproximation({\mdseries\slshape M, C})\index{RightSubMApproximation@\texttt{RightSubMApproximation}}
\label{RightSubMApproximation}
}\hfill{\scriptsize (operation)}}\\
\noindent\textcolor{FuncColor}{$\triangleright$\enspace\texttt{MinimalRightSubMApproximation({\mdseries\slshape M, C})\index{MinimalRightSubMApproximation@\texttt{MinimalRightSubMApproximation}}
\label{MinimalRightSubMApproximation}
}\hfill{\scriptsize (operation)}}\\


 Arguments: \mbox{\texttt{\mdseries\slshape M}}, \mbox{\texttt{\mdseries\slshape C}} \texttt{\symbol{45}}\texttt{\symbol{45}} two path algebra modules (\texttt{PathAlgebraMatModule}). \\
\textbf{\indent Returns:\ }
a right $\operatorname{Sub} M$\texttt{\symbol{45}}approximation of the module $C$, where the first version returns a not necessarily minimal right $\operatorname{Sub} M$\texttt{\symbol{45}}approximation and the second returns a minimal
approximation. 

}

 

\subsection{\textcolor{Chapter }{N{\textunderscore}RigidModule}}
\logpage{[ 8, 1, 45 ]}\nobreak
\hyperdef{L}{X78CA656D7D4F2446}{}
{\noindent\textcolor{FuncColor}{$\triangleright$\enspace\texttt{N{\textunderscore}RigidModule({\mdseries\slshape M, n})\index{N{\textunderscore}RigidModule@\texttt{N{\textunderscore}RigidModule}}
\label{NuScoreRigidModule}
}\hfill{\scriptsize (operation)}}\\


 Arguments: \mbox{\texttt{\mdseries\slshape M}}, \mbox{\texttt{\mdseries\slshape n}} \texttt{\symbol{45}} a PathAlgebraMatModule, an integer.\\
 \textbf{\indent Returns:\ }
true if \mbox{\texttt{\mdseries\slshape M}} is a \mbox{\texttt{\mdseries\slshape n}}\texttt{\symbol{45}}rigid module. Otherwise it returns false.

}

 

\subsection{\textcolor{Chapter }{TiltingModule}}
\logpage{[ 8, 1, 46 ]}\nobreak
\hyperdef{L}{X824D52847879AB91}{}
{\noindent\textcolor{FuncColor}{$\triangleright$\enspace\texttt{TiltingModule({\mdseries\slshape M, n})\index{TiltingModule@\texttt{TiltingModule}}
\label{TiltingModule}
}\hfill{\scriptsize (operation)}}\\


 Arguments: \mbox{\texttt{\mdseries\slshape M}}, \mbox{\texttt{\mdseries\slshape n}} \texttt{\symbol{45}} a PathAlgebraMatModule and a positive integer.\\
 \textbf{\indent Returns:\ }
false if \mbox{\texttt{\mdseries\slshape M}} is not a tilting module of projective dimension at most \mbox{\texttt{\mdseries\slshape n}}. Otherwise, it returns the projective dimension of \mbox{\texttt{\mdseries\slshape M}} and the coresolution of all indecomposable projective modules in $\add M$. 

}

 }

 }

 
\chapter{\textcolor{Chapter }{Auslander\texttt{\symbol{45}}Reiten theory}}\label{AR-theory}
\logpage{[ 9, 0, 0 ]}
\hyperdef{L}{X855427278501E7FB}{}
{
 This chapter describes the functions implemented for almost split sequences
and Auslander\texttt{\symbol{45}}Reiten theory in QPA.

 
\section{\textcolor{Chapter }{Almost split sequences and AR\texttt{\symbol{45}}quivers}}\logpage{[ 9, 1, 0 ]}
\hyperdef{L}{X79B0EA987E050C6D}{}
{
 

\subsection{\textcolor{Chapter }{AlmostSplitSequence}}
\logpage{[ 9, 1, 1 ]}\nobreak
\hyperdef{L}{X87BADA287BD9972C}{}
{\noindent\textcolor{FuncColor}{$\triangleright$\enspace\texttt{AlmostSplitSequence({\mdseries\slshape M, or, M, e})\index{AlmostSplitSequence@\texttt{AlmostSplitSequence}}
\label{AlmostSplitSequence}
}\hfill{\scriptsize (attribute)}}\\


 Arguments: (first version, only the first argument) \mbox{\texttt{\mdseries\slshape M}} \texttt{\symbol{45}} an indecomposable non\texttt{\symbol{45}}projective
module, (second version, this additional argument) \mbox{\texttt{\mdseries\slshape e}} \texttt{\symbol{45}} either l = left or r = right\\
 \textbf{\indent Returns:\ }
the almost split sequence ending in the module \mbox{\texttt{\mdseries\slshape M}} if it is indecomposable and not projective, for the first variant. The second
variant finds the almost split sequence starting or ending in the module \mbox{\texttt{\mdseries\slshape M}} depending on whether the second argument \mbox{\texttt{\mdseries\slshape e}} is l or r (l = almost split sequence starting with \mbox{\texttt{\mdseries\slshape M}}, or r = almost split sequence ending in \mbox{\texttt{\mdseries\slshape M}}), if the module is indecomposable and not injective or not projective,
respectively. It returns fail if the module is injective (l) or projective
(r). 



 The almost split sequence is returned as a pair of maps, the monomorphism and
the epimorphism. The function assumes that the module \mbox{\texttt{\mdseries\slshape M}} is indecomposable, and the source of the monomorphism (l) or the range of the
epimorphism (r) is a module that is isomorphic to \mbox{\texttt{\mdseries\slshape M}}, not necessarily identical. }

 

\subsection{\textcolor{Chapter }{AlmostSplitSequenceInPerpT}}
\logpage{[ 9, 1, 2 ]}\nobreak
\hyperdef{L}{X7E8DBA647B3EC71B}{}
{\noindent\textcolor{FuncColor}{$\triangleright$\enspace\texttt{AlmostSplitSequenceInPerpT({\mdseries\slshape T, M})\index{AlmostSplitSequenceInPerpT@\texttt{AlmostSplitSequenceInPerpT}}
\label{AlmostSplitSequenceInPerpT}
}\hfill{\scriptsize (operation)}}\\


 Arguments: \mbox{\texttt{\mdseries\slshape T}} \texttt{\symbol{45}} a cotilting module, \mbox{\texttt{\mdseries\slshape M}} \texttt{\symbol{45}} an indecomposable non\texttt{\symbol{45}}projective
module\\
 \textbf{\indent Returns:\ }
the almost split sequence in $^\perp T$ ending in the module \mbox{\texttt{\mdseries\slshape M}}, if the module is indecomposable and not projective (that is, not projective
object in $^\perp T$). It returns fail if the module \mbox{\texttt{\mdseries\slshape M}} is in $\add T$ projective. The almost split sequence is returned as a pair of maps, the
monomorphism and the epimorphism, and the range of the epimorphism is a module
that is isomorphic to the input, not necessarily identical. 



 The function assumes that the module \mbox{\texttt{\mdseries\slshape M}} is indecomposable and in $^\perp T$, and the range of the epimorphism is a module that is isomorphic to the
input, not necessarily identical. }

 

\subsection{\textcolor{Chapter }{IrreducibleMorphismsEndingIn}}
\logpage{[ 9, 1, 3 ]}\nobreak
\hyperdef{L}{X8774261F862EC761}{}
{\noindent\textcolor{FuncColor}{$\triangleright$\enspace\texttt{IrreducibleMorphismsEndingIn({\mdseries\slshape M})\index{IrreducibleMorphismsEndingIn@\texttt{IrreducibleMorphismsEndingIn}}
\label{IrreducibleMorphismsEndingIn}
}\hfill{\scriptsize (attribute)}}\\
\noindent\textcolor{FuncColor}{$\triangleright$\enspace\texttt{IrreducibleMorphismsStartingIn({\mdseries\slshape M})\index{IrreducibleMorphismsStartingIn@\texttt{IrreducibleMorphismsStartingIn}}
\label{IrreducibleMorphismsStartingIn}
}\hfill{\scriptsize (attribute)}}\\


 Arguments: \mbox{\texttt{\mdseries\slshape M}} \texttt{\symbol{45}} an indecomposable module\\
 \textbf{\indent Returns:\ }
the collection of irreducible morphisms ending and starting in the module \mbox{\texttt{\mdseries\slshape M}}, respectively. The argument is assumed to be an indecomposable module. 



 The irreducible morphisms are returned as a list of maps. Even in the case of
only one irreducible morphism, it is returned as a list. The function assumes
that the module \mbox{\texttt{\mdseries\slshape M}} is indecomposable over a quiver algebra with a finite field as the ground
ring. }

 

\subsection{\textcolor{Chapter }{IsTauPeriodic}}
\logpage{[ 9, 1, 4 ]}\nobreak
\hyperdef{L}{X873085DD8534F600}{}
{\noindent\textcolor{FuncColor}{$\triangleright$\enspace\texttt{IsTauPeriodic({\mdseries\slshape M, n})\index{IsTauPeriodic@\texttt{IsTauPeriodic}}
\label{IsTauPeriodic}
}\hfill{\scriptsize (operation)}}\\


 Arguments: \mbox{\texttt{\mdseries\slshape M}} \texttt{\symbol{45}}\texttt{\symbol{45}} a path algebra module (\texttt{PathAlgebraMatModule}), \mbox{\texttt{\mdseries\slshape n}} \texttt{\symbol{45}}\texttt{\symbol{45}} be a positive integer. \\
\textbf{\indent Returns:\ }
\texttt{i}, where \texttt{i} is the smallest positive integer less or equal \texttt{n} such that the representation \mbox{\texttt{\mdseries\slshape M}} is isomorphic to the $\tau^i(M)$, and false otherwise. 

}

 

\subsection{\textcolor{Chapter }{PredecessorOfModule}}
\logpage{[ 9, 1, 5 ]}\nobreak
\hyperdef{L}{X81A363848166449D}{}
{\noindent\textcolor{FuncColor}{$\triangleright$\enspace\texttt{PredecessorOfModule({\mdseries\slshape M, n})\index{PredecessorOfModule@\texttt{PredecessorOfModule}}
\label{PredecessorOfModule}
}\hfill{\scriptsize (operation)}}\\


 Arguments: \mbox{\texttt{\mdseries\slshape M}} \texttt{\symbol{45}} an indecomposable non\texttt{\symbol{45}}projective
module and \mbox{\texttt{\mdseries\slshape n}} \texttt{\symbol{45}} a positive integer.\\
 \textbf{\indent Returns:\ }
the predecessors of the module \mbox{\texttt{\mdseries\slshape M}} in the AR\texttt{\symbol{45}}quiver of the algebra \mbox{\texttt{\mdseries\slshape M}} is given over of distance less or equal to \mbox{\texttt{\mdseries\slshape n}}. 



 It returns two lists, the first is the indecomposable modules in the different
layers and the second is the valuations for the arrows in the
AR\texttt{\symbol{45}}quiver. The different entries in the first list are the
modules at distance zero, one, two, three, and so on, until layer \mbox{\texttt{\mdseries\slshape n}}. The \texttt{m}\texttt{\symbol{45}}th entry in the second list is the valuations of the
irreducible morphism from indecomposable module number \texttt{i} in layer \texttt{m+1} to indecomposable module number \texttt{j} in layer \texttt{m} for the values of \texttt{i} and \texttt{j} there is an irreducible morphism. Whenever \texttt{false} occur in the output, it means that this valuation has not been computed. The
function assumes that the module \mbox{\texttt{\mdseries\slshape M}} is indecomposable and that the quotient of the path algebra is given over a
finite field. }

 
\begin{Verbatim}[commandchars=!@|,fontsize=\small,frame=single,label=Example]
  !gapprompt@gap>| !gapinput@A := KroneckerAlgebra(GF(4),2);       |
  <GF(2^2)[<quiver with 2 vertices and 2 arrows>]>
  !gapprompt@gap>| !gapinput@S := SimpleModules(A)[1];             |
  <[ 1, 0 ]>
  !gapprompt@gap>| !gapinput@ass := AlmostSplitSequence(S);   |
  [ <<[ 3, 2 ]> ---> <[ 4, 2 ]>>
      , <<[ 4, 2 ]> ---> <[ 1, 0 ]>>
       ]
  !gapprompt@gap>| !gapinput@DecomposeModule(Range(ass[1]));|
  [ <[ 2, 1 ]>, <[ 2, 1 ]> ]
  !gapprompt@gap>| !gapinput@PredecessorsOfModule(S,5);   |
  [ [ [ <[ 1, 0 ]> ], [ <[ 2, 1 ]> ], [ <[ 3, 2 ]> ], [ <[ 4, 3 ]> ], 
        [ <[ 5, 4 ]> ], [ <[ 6, 5 ]> ] ], 
    [ [ [ 1, 1, [ 2, false ] ] ], [ [ 1, 1, [ 2, 2 ] ] ], 
        [ [ 1, 1, [ 2, 2 ] ] ], [ [ 1, 1, [ 2, 2 ] ] ], 
        [ [ 1, 1, [ false, 2 ] ] ] ] ]
  !gapprompt@gap>| !gapinput@A:=NakayamaAlgebra([5,4,3,2,1],GF(4));|
  <GF(2^2)[<quiver with 5 vertices and 4 arrows>]>
  !gapprompt@gap>| !gapinput@S := SimpleModules(A)[1];             |
  <[ 1, 0, 0, 0, 0 ]>
  !gapprompt@gap>| !gapinput@PredecessorsOfModule(S,5);|
  [ [ [ <[ 1, 0, 0, 0, 0 ]> ], [ <[ 1, 1, 0, 0, 0 ]> ], 
        [ <[ 0, 1, 0, 0, 0 ]>, <[ 1, 1, 1, 0, 0 ]> ], 
        [ <[ 0, 1, 1, 0, 0 ]>, <[ 1, 1, 1, 1, 0 ]> ], 
        [ <[ 0, 0, 1, 0, 0 ]>, <[ 0, 1, 1, 1, 0 ]>, <[ 1, 1, 1, 1, 1 ]> 
           ], [ <[ 0, 0, 1, 1, 0 ]>, <[ 0, 1, 1, 1, 1 ]> ] ], 
    [ [ [ 1, 1, [ 1, false ] ] ], 
        [ [ 1, 1, [ 1, 1 ] ], [ 2, 1, [ 1, false ] ] ], 
        [ [ 1, 1, [ 1, 1 ] ], [ 1, 2, [ 1, 1 ] ], 
            [ 2, 2, [ 1, false ] ] ], 
        [ [ 1, 1, [ 1, 1 ] ], [ 2, 1, [ 1, 1 ] ], [ 2, 2, [ 1, 1 ] ], 
            [ 3, 2, [ 1, false ] ] ], 
        [ [ 1, 1, [ false, 1 ] ], [ 1, 2, [ false, 1 ] ], 
            [ 2, 2, [ false, 1 ] ], [ 2, 3, [ false, 1 ] ] ] ] ] 
\end{Verbatim}
 }

 }

 
\chapter{\textcolor{Chapter }{Chain complexes}}\logpage{[ 10, 0, 0 ]}
\hyperdef{L}{X7A06103979B92808}{}
{
  
\section{\textcolor{Chapter }{Introduction}}\logpage{[ 10, 1, 0 ]}
\hyperdef{L}{X7DFB63A97E67C0A1}{}
{
  If $\mathcal{A}$ is an abelian category, then a chain complex of objects of $\mathcal A$ is a sequence 
\[ \cdots \longrightarrow C_{i+1} \stackrel{d_{i+1}}{\longrightarrow} C_i
\stackrel{d_i}{\longrightarrow} C_{i-1} \stackrel{d_{i-1}}{\longrightarrow}
\cdots \]
 where $C_i$ is an object of $\mathcal A$ for all $i$, and $d_i$ is a morphism of $\mathcal A$ for all $i$ such that the composition of two consecutive maps of the complex is zero. The
maps are called the differentials of the complex. A complex is called \emph{bounded above} (resp. below) if there is a bound $b$ such that $C_i = 0$ for all $i>b$ (resp. $i<b$). A complex is \emph{bounded} if it is both bounded below and bounded above. 

 The challenge when representing chain complexes in software is to handle their
infinite nature. If a complex is not bounded, or not known to be bounded, how
can we represent it in an immutable way? Our solution is to use a category
called \texttt{InfList} (for ``infinite list'') to store the differentials of the complex. The
properties of the \texttt{IsInfList} category is described in \ref{InfLists}. An \texttt{IsQPAComplex} object consists of one \texttt{IsInfList} for the differentials, and it also has an \texttt{IsCat} object as an attribute. The \texttt{IsCat} category is a representation of an abelian category, see \ref{Cats}. 

 To work with bounded complexes one does not need to know much about the \texttt{IsInfList} category. A bounded complex can be created by simply giving a list of the
differentials and the degree of the first differential as input (see \texttt{FiniteComplex} (\ref{FiniteComplex})), and to create a stalk complex the stalk object and its degree suffice as
input (see \texttt{StalkComplex} (\ref{StalkComplex})). In both cases an \texttt{IsCat} object is also needed. 
\begin{Verbatim}[commandchars=!@|,fontsize=\small,frame=single,label=Example]
  !gapprompt@gap>| !gapinput@C := FiniteComplex(cat, 1, [g,f]);|
  0 -> 2:(1,0) -> 1:(2,2) -> 0:(1,1) -> 0 
  !gapprompt@gap>| !gapinput@Ms := StalkComplex(cat, M, 3);|
  0 -> 3:(2,2) -> 0 
\end{Verbatim}
 }

 
\section{\textcolor{Chapter }{Infinite lists}}\label{InfLists}
\logpage{[ 10, 2, 0 ]}
\hyperdef{L}{X7AC5660E80079755}{}
{
  In this section we give documentation for the \texttt{IsInfList} category. We start by giving a representation of $\pm \infty$. Then we quickly describe the \texttt{IsInfList} category, before we turn to the underlying structure of the infinite lists
\texttt{\symbol{45}}\texttt{\symbol{45}} the half infinite lists (\texttt{IsHalfInfList}). Most of the functionality of the infinite lists come from this category.
Finally, we give the constructors for infinite lists, and some methods for
manipulating such objects. 

\subsection{\textcolor{Chapter }{IsInfiniteNumber}}
\logpage{[ 10, 2, 1 ]}\nobreak
\hyperdef{L}{X8466C7117CEC99D2}{}
{\noindent\textcolor{FuncColor}{$\triangleright$\enspace\texttt{IsInfiniteNumber\index{IsInfiniteNumber@\texttt{IsInfiniteNumber}}
\label{IsInfiniteNumber}
}\hfill{\scriptsize (Category)}}\\


 A category for infinite numbers. }

 

\subsection{\textcolor{Chapter }{PositiveInfinity}}
\logpage{[ 10, 2, 2 ]}\nobreak
\hyperdef{L}{X7F12EE4479A03527}{}
{\noindent\textcolor{FuncColor}{$\triangleright$\enspace\texttt{PositiveInfinity\index{PositiveInfinity@\texttt{PositiveInfinity}}
\label{PositiveInfinity}
}\hfill{\scriptsize (Var)}}\\


 A global variable representing the number $\infty$. It is greater than any integer, but it can not be compared to numbers which
are not integers. It belongs to the \texttt{IsInfiniteNumber} category. }

 

\subsection{\textcolor{Chapter }{NegativeInfinity}}
\logpage{[ 10, 2, 3 ]}\nobreak
\hyperdef{L}{X8759931784478A4D}{}
{\noindent\textcolor{FuncColor}{$\triangleright$\enspace\texttt{NegativeInfinity\index{NegativeInfinity@\texttt{NegativeInfinity}}
\label{NegativeInfinity}
}\hfill{\scriptsize (Var)}}\\


 A global variable representing the number $-\infty$. It is smaller than any integer, but it can not be compared to numbers which
are not integers. It belongs to the \texttt{IsInfiniteNumber} category. }

 

\subsection{\textcolor{Chapter }{IsInfList}}
\logpage{[ 10, 2, 4 ]}\nobreak
\hyperdef{L}{X84731DA279E797A5}{}
{\noindent\textcolor{FuncColor}{$\triangleright$\enspace\texttt{IsInfList\index{IsInfList@\texttt{IsInfList}}
\label{IsInfList}
}\hfill{\scriptsize (Category)}}\\


 An infinite list is an immutable representation of a list with possibly
infinite range of indices. It consists of three parts: The ``middle part'' is
finite and covers some range $[a,b]$ of indices, the ``positive part'' covers the range $[b+1,\infty)$ of indices, and the ``negative part'' covers the range $(-\infty,a-1]$ of indices. Note that none of the three parts are mandatory: The middle part
may be an empty list, and the positive part may be set to \texttt{fail} to achieve index range ending at $b < \infty$. Similarly, if the index range has lower bound $a < \infty$, put the negative part to be \texttt{fail}. 

 Each of the two infinite parts are described in one of the following ways: (1)
A finite list which is repeated indefinitely; (2) A function which takes an
index in the list as argument and returns the corresponding list item; (3) A
function which takes an item from the list as argument and returns the next
item. 

 The two infinite parts are represented as ``half infinite lists'', see \ref{IsHalfInfListsubsect}. An infinite list can be constructed in the following ways: 
\begin{itemize}
\item  From two half infinite lists and a middle part, \texttt{MakeInfListFromHalfInfLists} (\ref{MakeInfListFromHalfInfLists}).
\item  Directly, by giving the same input as when constructing the above, \texttt{MakeInfList} (\ref{MakeInfList}).
\item  If all values of the infinite list are the image of the index under a function $f$, one can use \texttt{FunctionInfList} (\ref{FunctionInfList}).
\item  If all values of the infinite list are the same, one can use \texttt{ConstantInfList} (\ref{ConstantInfList}).
\item  If the infinite list has a finite range, one can use \texttt{FiniteInfList} (\ref{FiniteInfList}).
\end{itemize}
 In addition, new infinite lists can be constructed from others by shift,
splice, concatenation, extracting parts or applying a function to the
elements. }

 

\subsection{\textcolor{Chapter }{IsHalfInfList}}
\logpage{[ 10, 2, 5 ]}\nobreak
\label{IsHalfInfListsubsect}
\hyperdef{L}{X87524BBB7C2F9FCB}{}
{\noindent\textcolor{FuncColor}{$\triangleright$\enspace\texttt{IsHalfInfList\index{IsHalfInfList@\texttt{IsHalfInfList}}
\label{IsHalfInfList}
}\hfill{\scriptsize (Category)}}\\


 A half infinite list is a representation of a list with indices in the range $[a,\infty)$ or $(-\infty,b]$. An infinite list is typically made from two half infinite lists, and half
infinite lists can be extracted from an infinite list. Hence, the half
infinite list stores much of the information about an infinite list. One main
difference between an infinite list and a half infinite list is that the half
infinite list does not have any finite part, as the ``middle'' part of an
infinite list. }

 

\subsection{\textcolor{Chapter }{\texttt{\symbol{92}}\texttt{\symbol{94}}}}
\logpage{[ 10, 2, 6 ]}\nobreak
\hyperdef{L}{X7D21FB1A7D21FB1A}{}
{\noindent\textcolor{FuncColor}{$\triangleright$\enspace\texttt{\texttt{\symbol{92}}\texttt{\symbol{94}}({\mdseries\slshape list, pos})\index{\texttt{\symbol{92}}\texttt{\symbol{94}}@\texttt{\texttt{\symbol{92}}\texttt{\symbol{94}}}}
\label{bSlash^}
}\hfill{\scriptsize (operation)}}\\


 Arguments: \mbox{\texttt{\mdseries\slshape list}} \texttt{\symbol{45}}\texttt{\symbol{45}} either an infinite list or a half
infinite list, \mbox{\texttt{\mdseries\slshape pos}} \texttt{\symbol{45}}\texttt{\symbol{45}} a valid index for \mbox{\texttt{\mdseries\slshape list}}. \textbf{\indent Returns:\ }
 The value at position \mbox{\texttt{\mdseries\slshape pos}} of \mbox{\texttt{\mdseries\slshape list}}.

}

 

\subsection{\textcolor{Chapter }{MakeHalfInfList}}
\logpage{[ 10, 2, 7 ]}\nobreak
\label{MakeHalfInfListsubsect}
\hyperdef{L}{X7E01944C83F87EED}{}
{\noindent\textcolor{FuncColor}{$\triangleright$\enspace\texttt{MakeHalfInfList({\mdseries\slshape start, direction, typeWithArgs, callback, repeatifyCallback})\index{MakeHalfInfList@\texttt{MakeHalfInfList}}
\label{MakeHalfInfList}
}\hfill{\scriptsize (function)}}\\


Arguments: \mbox{\texttt{\mdseries\slshape start}} \texttt{\symbol{45}}\texttt{\symbol{45}} an integer, \mbox{\texttt{\mdseries\slshape direction}} \texttt{\symbol{45}}\texttt{\symbol{45}} either $1$ or $-1$, \mbox{\texttt{\mdseries\slshape typeWithArgs}} \texttt{\symbol{45}}\texttt{\symbol{45}} a list which may have different
formats, \mbox{\texttt{\mdseries\slshape callback}} \texttt{\symbol{45}}\texttt{\symbol{45}} a function, \mbox{\texttt{\mdseries\slshape repeatifyCallback}} \texttt{\symbol{45}}\texttt{\symbol{45}} a function. \\
 \textbf{\indent Returns:\ }
A newly created half infinite list with index range from \mbox{\texttt{\mdseries\slshape start}} to $\infty$, or from $-\infty$ to \mbox{\texttt{\mdseries\slshape start}}.



 If the range should be $[\mathtt{start},\infty)$ then the value of \mbox{\texttt{\mdseries\slshape direction}} is $1$. if the range should be $(-\infty,\mathtt{start}]$, then the value of \mbox{\texttt{\mdseries\slshape direction}} is $-1$. 

 The argument \mbox{\texttt{\mdseries\slshape typeWithArgs}} can take one of the following forms: 
\begin{itemize}
\item \texttt{[ "repeat", repeatList ]}
\item \texttt{[ "next", nextFunction, initialValue ]}
\item \texttt{[ "next/repeat", nextFunction, initialValue ]}
\item \texttt{[ "pos", posFunction ]}
\item \texttt{[ "pos", posFunction, storeValues ]}
\end{itemize}
 \texttt{repeatList} is a list of values that should be repeated in the half infinite list. \texttt{nextFunction} returns the value at position $i$, given the value at the previous position as argument. Here \texttt{initialValue} is the value at position \texttt{start}. Similarly, \texttt{posFunction} returns the value at any position $i$, and it may or may not store the values between the previous computed indices
and the newly computed index. The default value of \texttt{storeValues} is \texttt{true} for \texttt{"next"} and \texttt{"pos"}, and \texttt{false} for \texttt{"repeat"}. The type \texttt{"next/repeat"} works exactly like the type \texttt{"next"}, except that when values in the list are computed, the list will try to
discover if the values are repeating. If this happens, the function \mbox{\texttt{\mdseries\slshape repeatifyCallback}} is called with two arguments: the non\texttt{\symbol{45}}repeating part at the
beginning as a normal list (this might be empty) and a new HalfInfList of type \texttt{"repeat"} for the repeating part. 

 The argument \texttt{callback} is a function that is called whenever a new value of the list is computed. It
takes three arguments: The current position, the direction and the type (that
is, \texttt{typeWithArgs[1]}). If no callback function is needed, use \texttt{false}. 

 All the information given to create the list is stored, and can be retrieved
later by the operations listed in \ref{StartPositionsubsect}\texttt{\symbol{45}}\texttt{\symbol{45}}\ref{HighestKnownValue}. }

 
\begin{Verbatim}[commandchars=!@|,fontsize=\small,frame=single,label=Example]
  !gapprompt@gap>| !gapinput@# make a HalfInfList from 0 to inf which repeats the list [ 2, 4, 6 ]|
  !gapprompt@gap>| !gapinput@list1 := MakeHalfInfList( 0, 1, [ "repeat", [ 2, 4, 6 ] ], false );|
  <object>
  !gapprompt@gap>| !gapinput@list1^0;|
  2
  !gapprompt@gap>| !gapinput@list1^5;|
  6
  !gapprompt@gap>| !gapinput@# make a HalfInfList from 0 to inf with x^2 in position x|
  !gapprompt@gap>| !gapinput@f := function(x) return x^2; end;;|
  !gapprompt@gap>| !gapinput@list2 := MakeHalfInfList( 0, 1, [ "pos", f, false ], false );|
  <object>
  !gapprompt@gap>| !gapinput@list2^0;|
  0
  !gapprompt@gap>| !gapinput@list2^10;|
  100 
  !gapprompt@gap>| !gapinput@# make a HalfInfList from 0 to -inf where each new value adds 3|
  !gapprompt@gap>| !gapinput@# to the previous and the value in position 0 is 10|
  !gapprompt@gap>| !gapinput@g := function(x) return x+3; end;;|
  !gapprompt@gap>| !gapinput@list3 := MakeHalfInfList( 0, -1, [ "next", g, 7 ], false );|
  <object>
  !gapprompt@gap>| !gapinput@list3^0;|
  10
  !gapprompt@gap>| !gapinput@list3^-10;|
  40 
\end{Verbatim}
 

\subsection{\textcolor{Chapter }{StartPosition}}
\logpage{[ 10, 2, 8 ]}\nobreak
\label{StartPositionsubsect}
\hyperdef{L}{X8458C2698505C12B}{}
{\noindent\textcolor{FuncColor}{$\triangleright$\enspace\texttt{StartPosition({\mdseries\slshape list})\index{StartPosition@\texttt{StartPosition}}
\label{StartPosition}
}\hfill{\scriptsize (operation)}}\\


\mbox{\texttt{\mdseries\slshape list}} \texttt{\symbol{45}}\texttt{\symbol{45}} a half infinite list.\\
 \textbf{\indent Returns:\ }
The start position of \mbox{\texttt{\mdseries\slshape list}}.

}

 

\subsection{\textcolor{Chapter }{Direction}}
\logpage{[ 10, 2, 9 ]}\nobreak
\hyperdef{L}{X812437C47A208D4A}{}
{\noindent\textcolor{FuncColor}{$\triangleright$\enspace\texttt{Direction({\mdseries\slshape list})\index{Direction@\texttt{Direction}}
\label{Direction}
}\hfill{\scriptsize (operation)}}\\


\mbox{\texttt{\mdseries\slshape list}} \texttt{\symbol{45}}\texttt{\symbol{45}} a half infinite list.\\
 \textbf{\indent Returns:\ }
The direction of \mbox{\texttt{\mdseries\slshape list}} (either $1$ or $-1$).

}

 

\subsection{\textcolor{Chapter }{InfListType}}
\logpage{[ 10, 2, 10 ]}\nobreak
\hyperdef{L}{X7C3728BF7ED7149A}{}
{\noindent\textcolor{FuncColor}{$\triangleright$\enspace\texttt{InfListType({\mdseries\slshape list})\index{InfListType@\texttt{InfListType}}
\label{InfListType}
}\hfill{\scriptsize (operation)}}\\


\mbox{\texttt{\mdseries\slshape list}} \texttt{\symbol{45}}\texttt{\symbol{45}} a half infinite list.\\
 \textbf{\indent Returns:\ }
The type of \mbox{\texttt{\mdseries\slshape list}} (either \texttt{"pos"}, \texttt{"repeat"} or \texttt{"next"}).

}

 

\subsection{\textcolor{Chapter }{RepeatingList}}
\logpage{[ 10, 2, 11 ]}\nobreak
\hyperdef{L}{X82AA0979868A6967}{}
{\noindent\textcolor{FuncColor}{$\triangleright$\enspace\texttt{RepeatingList({\mdseries\slshape list})\index{RepeatingList@\texttt{RepeatingList}}
\label{RepeatingList}
}\hfill{\scriptsize (operation)}}\\


\mbox{\texttt{\mdseries\slshape list}} \texttt{\symbol{45}}\texttt{\symbol{45}} a half infinite list.\\
 \textbf{\indent Returns:\ }
The repeating list of \mbox{\texttt{\mdseries\slshape list}} if \mbox{\texttt{\mdseries\slshape list}} is of type \texttt{"repeat"}, and \texttt{fail} otherwise.

}

 

\subsection{\textcolor{Chapter }{ElementFunction}}
\logpage{[ 10, 2, 12 ]}\nobreak
\hyperdef{L}{X7AF21D62848B270B}{}
{\noindent\textcolor{FuncColor}{$\triangleright$\enspace\texttt{ElementFunction({\mdseries\slshape list})\index{ElementFunction@\texttt{ElementFunction}}
\label{ElementFunction}
}\hfill{\scriptsize (operation)}}\\


\mbox{\texttt{\mdseries\slshape list}} \texttt{\symbol{45}}\texttt{\symbol{45}} a half infinite list.\\
 \textbf{\indent Returns:\ }
The element function of \mbox{\texttt{\mdseries\slshape list}} if \mbox{\texttt{\mdseries\slshape list}} is of type \texttt{"next"} or \texttt{"pos"}, and \texttt{fail} otherwise.

}

 

\subsection{\textcolor{Chapter }{IsStoringValues}}
\logpage{[ 10, 2, 13 ]}\nobreak
\hyperdef{L}{X8728A2107F27D68F}{}
{\noindent\textcolor{FuncColor}{$\triangleright$\enspace\texttt{IsStoringValues({\mdseries\slshape list})\index{IsStoringValues@\texttt{IsStoringValues}}
\label{IsStoringValues}
}\hfill{\scriptsize (operation)}}\\


\mbox{\texttt{\mdseries\slshape list}} \texttt{\symbol{45}}\texttt{\symbol{45}} a half infinite list.\\
 \textbf{\indent Returns:\ }
\texttt{true} if all elements of the list are stored, \texttt{false} otherwise.

}

 

\subsection{\textcolor{Chapter }{NewValueCallback}}
\logpage{[ 10, 2, 14 ]}\nobreak
\hyperdef{L}{X86888FBD8559B29B}{}
{\noindent\textcolor{FuncColor}{$\triangleright$\enspace\texttt{NewValueCallback({\mdseries\slshape list})\index{NewValueCallback@\texttt{NewValueCallback}}
\label{NewValueCallback}
}\hfill{\scriptsize (operation)}}\\


\mbox{\texttt{\mdseries\slshape list}} \texttt{\symbol{45}}\texttt{\symbol{45}} a half infinite list.\\
 \textbf{\indent Returns:\ }
The callback function of the list.

}

 

\subsection{\textcolor{Chapter }{IsRepeating}}
\logpage{[ 10, 2, 15 ]}\nobreak
\hyperdef{L}{X856E1541834688C6}{}
{\noindent\textcolor{FuncColor}{$\triangleright$\enspace\texttt{IsRepeating({\mdseries\slshape list})\index{IsRepeating@\texttt{IsRepeating}}
\label{IsRepeating}
}\hfill{\scriptsize (operation)}}\\


\mbox{\texttt{\mdseries\slshape list}} \texttt{\symbol{45}}\texttt{\symbol{45}} a half infinite list.\\
 \textbf{\indent Returns:\ }
\texttt{true} if the type of the list is \texttt{"repeat"}.

}

 

\subsection{\textcolor{Chapter }{InitialValue}}
\logpage{[ 10, 2, 16 ]}\nobreak
\hyperdef{L}{X8208162D7F3CC139}{}
{\noindent\textcolor{FuncColor}{$\triangleright$\enspace\texttt{InitialValue({\mdseries\slshape list})\index{InitialValue@\texttt{InitialValue}}
\label{InitialValue}
}\hfill{\scriptsize (operation)}}\\


\mbox{\texttt{\mdseries\slshape list}} \texttt{\symbol{45}}\texttt{\symbol{45}} a half infinite list.\\
 \textbf{\indent Returns:\ }
If the list is of type \texttt{"next"} then the initial value is returned, otherwise it fails.

}

 

\subsection{\textcolor{Chapter }{LowestKnownPosition}}
\logpage{[ 10, 2, 17 ]}\nobreak
\hyperdef{L}{X7F12274E794510CA}{}
{\noindent\textcolor{FuncColor}{$\triangleright$\enspace\texttt{LowestKnownPosition({\mdseries\slshape list})\index{LowestKnownPosition@\texttt{LowestKnownPosition}}
\label{LowestKnownPosition}
}\hfill{\scriptsize (operation)}}\\


\mbox{\texttt{\mdseries\slshape list}} \texttt{\symbol{45}}\texttt{\symbol{45}} a half infinite list or an infinite
list.\\
 \textbf{\indent Returns:\ }
The lowest index $i$ such that the value at position $i$ is known without computation (that is, it is either stored, or the list has
type \texttt{"repeat"}).

}

 

\subsection{\textcolor{Chapter }{HighestKnownValue}}
\logpage{[ 10, 2, 18 ]}\nobreak
\label{HighestKnownPositionsubsect}
\hyperdef{L}{X7B5E57F0859BEE17}{}
{\noindent\textcolor{FuncColor}{$\triangleright$\enspace\texttt{HighestKnownValue({\mdseries\slshape list})\index{HighestKnownValue@\texttt{HighestKnownValue}}
\label{HighestKnownValue}
}\hfill{\scriptsize (operation)}}\\


\mbox{\texttt{\mdseries\slshape list}} \texttt{\symbol{45}}\texttt{\symbol{45}} a half infinite list.\\
 \textbf{\indent Returns:\ }
The highest index $i$ such that the value at position $i$ is known without computation (that is, it is either stored, or the list has
type \texttt{"repeat"}).

}

 
\begin{Verbatim}[commandchars=!@|,fontsize=\small,frame=single,label=Example]
  !gapprompt@gap>| !gapinput@# we reuse the IsHalfInfLists from the previous example|
  !gapprompt@gap>| !gapinput@HighestKnownPosition(list1);|
  +inf
  !gapprompt@gap>| !gapinput@HighestKnownPosition(list2);|
  "none"
  !gapprompt@gap>| !gapinput@HighestKnownPosition(list3);|
  0 
\end{Verbatim}
 For a description of the command \texttt{Shift} of a half infinite list, see \ref{Shiftsubsect}. 

\subsection{\textcolor{Chapter }{Cut}}
\logpage{[ 10, 2, 19 ]}\nobreak
\hyperdef{L}{X7D8067A7842BC4A4}{}
{\noindent\textcolor{FuncColor}{$\triangleright$\enspace\texttt{Cut({\mdseries\slshape list, pos})\index{Cut@\texttt{Cut}}
\label{Cut}
}\hfill{\scriptsize (operation)}}\\


 Arguments: \mbox{\texttt{\mdseries\slshape list}} \texttt{\symbol{45}}\texttt{\symbol{45}} a half infinite list, \mbox{\texttt{\mdseries\slshape pos}} \texttt{\symbol{45}}\texttt{\symbol{45}} an integer within the range of \mbox{\texttt{\mdseries\slshape list}}.\\
 \textbf{\indent Returns:\ }
A new half infinite list which is \mbox{\texttt{\mdseries\slshape list}} with some part cut off. 



If the direction of \mbox{\texttt{\mdseries\slshape list}} is positive, then the new list has range from \texttt{cut} to $\infty$. If the direction of \mbox{\texttt{\mdseries\slshape list}} is negative, then the new list has range from $-\infty$ to \texttt{cut}. The values at position $i$ of the new half infinite list is the same as the value at position $i$ of \mbox{\texttt{\mdseries\slshape list}}. }

 

\subsection{\textcolor{Chapter }{HalfInfList}}
\logpage{[ 10, 2, 20 ]}\nobreak
\hyperdef{L}{X79A67E377B71C540}{}
{\noindent\textcolor{FuncColor}{$\triangleright$\enspace\texttt{HalfInfList({\mdseries\slshape list, func})\index{HalfInfList@\texttt{HalfInfList}}
\label{HalfInfList}
}\hfill{\scriptsize (operation)}}\\


 Arguments: \mbox{\texttt{\mdseries\slshape list}} \texttt{\symbol{45}}\texttt{\symbol{45}} a half infinite list, \mbox{\texttt{\mdseries\slshape func}} \texttt{\symbol{45}}\texttt{\symbol{45}} a function which takes an element of
the list as argument. \\
 \textbf{\indent Returns:\ }
A half infinite list with the same range as \mbox{\texttt{\mdseries\slshape list}}, where the value at position $i$ is the image of the value at position $i$ of \mbox{\texttt{\mdseries\slshape list}} under \mbox{\texttt{\mdseries\slshape func}}. 

}

 

\subsection{\textcolor{Chapter }{MakeInfListFromHalfInfLists}}
\logpage{[ 10, 2, 21 ]}\nobreak
\hyperdef{L}{X7CE66EBA844DBB57}{}
{\noindent\textcolor{FuncColor}{$\triangleright$\enspace\texttt{MakeInfListFromHalfInfLists({\mdseries\slshape basePosition, middle, positive, negative})\index{MakeInfListFromHalfInfLists@\texttt{MakeInfListFromHalfInfLists}}
\label{MakeInfListFromHalfInfLists}
}\hfill{\scriptsize (function)}}\\


 Arguments: \mbox{\texttt{\mdseries\slshape basePosition}} \texttt{\symbol{45}}\texttt{\symbol{45}} an integer, \mbox{\texttt{\mdseries\slshape middle}} \texttt{\symbol{45}}\texttt{\symbol{45}} a list, \mbox{\texttt{\mdseries\slshape positive}} \texttt{\symbol{45}}\texttt{\symbol{45}} a half infinite list, \mbox{\texttt{\mdseries\slshape negative}} \texttt{\symbol{45}}\texttt{\symbol{45}} a half infinite list. \\
 \textbf{\indent Returns:\ }
An infinite list with \mbox{\texttt{\mdseries\slshape middle}} as is middle part, \mbox{\texttt{\mdseries\slshape positive}} as its positive part and \mbox{\texttt{\mdseries\slshape negative}} as its negative part.



 The starting position of \mbox{\texttt{\mdseries\slshape positive}} must be \texttt{basePosition + Length( middle )}, and the starting position of \mbox{\texttt{\mdseries\slshape negative}} must be \texttt{basePosition \texttt{\symbol{45}} 1}. The returned list has \texttt{middle[1]} in position \mbox{\texttt{\mdseries\slshape basePosition}}, \texttt{middle[2]} in position \mbox{\texttt{\mdseries\slshape basePosition + 1}} and so on. Note that one probably wants the \mbox{\texttt{\mdseries\slshape positive}} half infinite list to have direction $1$, and the \mbox{\texttt{\mdseries\slshape negative}} half infinite list to have direction $-1$. }

 
\begin{Verbatim}[commandchars=!@|,fontsize=\small,frame=single,label=Example]
  !gapprompt@gap>| !gapinput@# we want to construct an infinite list with 0 in position|
  !gapprompt@gap>| !gapinput@# 0 to 5, and x^2 in position x where x goes from 6 to inf, |
  !gapprompt@gap>| !gapinput@# and alternatingly 1 and -1 in position -1 to -inf.|
  !gapprompt@gap>| !gapinput@#|
  !gapprompt@gap>| !gapinput@basePosition := 0;;|
  !gapprompt@gap>| !gapinput@middle := [0,0,0,0,0,0];;|
  !gapprompt@gap>| !gapinput@f := function(x) return x^2; end;;|
  !gapprompt@gap>| !gapinput@positive := MakeHalfInfList( 6, 1, [ "pos", f, false ], false );|
  <object>
  !gapprompt@gap>| !gapinput@altList := [ 1, -1 ];;|
  !gapprompt@gap>| !gapinput@negative := MakeHalfInfList( -1, -1, [ "repeat", altList ], false );|
  <object>
  !gapprompt@gap>| !gapinput@inflist := MakeInfListFromHalfInfLists( basePosition, middle,|
  !gapprompt@>| !gapinput@                                           positive, negative );|
  <object>
  !gapprompt@gap>| !gapinput@inflist^0; inflist^5; inflist^6; inflist^-1; inflist^-4;|
  0
  0
  36
  1
  -1 
\end{Verbatim}
 

\subsection{\textcolor{Chapter }{MakeInfList}}
\logpage{[ 10, 2, 22 ]}\nobreak
\hyperdef{L}{X801DE2517D08C19E}{}
{\noindent\textcolor{FuncColor}{$\triangleright$\enspace\texttt{MakeInfList({\mdseries\slshape basePosition, middle, positive, negative, callback})\index{MakeInfList@\texttt{MakeInfList}}
\label{MakeInfList}
}\hfill{\scriptsize (function)}}\\


 Arguments: \mbox{\texttt{\mdseries\slshape basePosition}} \texttt{\symbol{45}}\texttt{\symbol{45}} an integer, \mbox{\texttt{\mdseries\slshape middle}} \texttt{\symbol{45}}\texttt{\symbol{45}} a list, \mbox{\texttt{\mdseries\slshape positive}} \texttt{\symbol{45}}\texttt{\symbol{45}} a list describing the positive part, \mbox{\texttt{\mdseries\slshape negative}} \texttt{\symbol{45}}\texttt{\symbol{45}} a list describing the negative part. \\
 \textbf{\indent Returns:\ }
 An infinite list with \mbox{\texttt{\mdseries\slshape middle}} as is middle part, \mbox{\texttt{\mdseries\slshape positive}} as its positive part and \mbox{\texttt{\mdseries\slshape negative}} as its negative part.



 The major difference between this construction and the previous is that here
the half infinite lists that will make the positive and negative parts are not
entered directly as arguments. Instead, one enters ``description lists'',
which are of the same format as the argument \mbox{\texttt{\mdseries\slshape typeWithArgs}} of \texttt{MakeHalfInfList} (\ref{MakeHalfInfList}). If the positive and/or negative part is specified with type \texttt{"next/repeat"}, then it will initially be of type \texttt{"next"}, but will be replaced by a HalfInfList of type \texttt{"repeat"} if it is discovered that the values are repeating. }

 
\begin{Verbatim}[commandchars=!@|,fontsize=\small,frame=single,label=Example]
  !gapprompt@gap>| !gapinput@# we construct the same infinite list as in the previous example|
  !gapprompt@gap>| !gapinput@basePosition := 0;;|
  !gapprompt@gap>| !gapinput@middle := [0,0,0,0,0,0];;|
  !gapprompt@gap>| !gapinput@f := function(x) return x^2; end;;|
  !gapprompt@gap>| !gapinput@altList := [ 1, -1 ];;|
  !gapprompt@gap>| !gapinput@inflist2 := MakeInfList( 0, middle, [ "pos", f, false ], [ "repeat", |
  !gapprompt@>| !gapinput@                            altList ], false );|
  <object>
  !gapprompt@gap>| !gapinput@inflist2^0; inflist2^5; inflist2^6; inflist2^-1; inflist2^-4;|
  0
  0
  36
  1
  -1
  !gapprompt@gap>| !gapinput@n := function( x ) return ( x + 1 ) mod 5; end;;|
  !gapprompt@gap>| !gapinput@list := MakeInfList( 0, [ 0 ], [ "next/repeat", n, 0 ],|
  !gapprompt@>| !gapinput@                        [ "repeat", [ 0 ] ], false );;|
  !gapprompt@gap>| !gapinput@list^2;|
  2
  !gapprompt@gap>| !gapinput@IsRepeating( PositivePart( list ) );|
  false
  !gapprompt@gap>| !gapinput@list^11;|
  1
  !gapprompt@gap>| !gapinput@IsRepeating( PositivePart( list ) );|
  true 
\end{Verbatim}
 

\subsection{\textcolor{Chapter }{FunctionInfList}}
\logpage{[ 10, 2, 23 ]}\nobreak
\hyperdef{L}{X7EEB276380E1584F}{}
{\noindent\textcolor{FuncColor}{$\triangleright$\enspace\texttt{FunctionInfList({\mdseries\slshape func})\index{FunctionInfList@\texttt{FunctionInfList}}
\label{FunctionInfList}
}\hfill{\scriptsize (function)}}\\


 Arguments: \mbox{\texttt{\mdseries\slshape func}} \texttt{\symbol{45}}\texttt{\symbol{45}} a function that takes an integer as
argument.\\
 \textbf{\indent Returns:\ }
An infinite list where the value at position $i$ is the function \mbox{\texttt{\mdseries\slshape func}} applied to $i$. 



 }

 

\subsection{\textcolor{Chapter }{ConstantInfList}}
\logpage{[ 10, 2, 24 ]}\nobreak
\hyperdef{L}{X78C7FBAE83914790}{}
{\noindent\textcolor{FuncColor}{$\triangleright$\enspace\texttt{ConstantInfList({\mdseries\slshape value})\index{ConstantInfList@\texttt{ConstantInfList}}
\label{ConstantInfList}
}\hfill{\scriptsize (function)}}\\


 Arguments: \mbox{\texttt{\mdseries\slshape value}} \texttt{\symbol{45}}\texttt{\symbol{45}} an object.\\
 \textbf{\indent Returns:\ }
An infinite list which has the object \mbox{\texttt{\mdseries\slshape value}} in every position. 



 }

 

\subsection{\textcolor{Chapter }{FiniteInfList}}
\logpage{[ 10, 2, 25 ]}\nobreak
\hyperdef{L}{X7D68E0627F556E14}{}
{\noindent\textcolor{FuncColor}{$\triangleright$\enspace\texttt{FiniteInfList({\mdseries\slshape basePosition, list})\index{FiniteInfList@\texttt{FiniteInfList}}
\label{FiniteInfList}
}\hfill{\scriptsize (function)}}\\


 Arguments: \mbox{\texttt{\mdseries\slshape basePosition}} \texttt{\symbol{45}}\texttt{\symbol{45}} an integer, \mbox{\texttt{\mdseries\slshape list}} \texttt{\symbol{45}}\texttt{\symbol{45}} a list of length $n$.\\
 \textbf{\indent Returns:\ }
An infinite list with $\mathtt{list[1]},\ldots,\mathtt{list[n]}$ in positions $\mathtt{basePosition},\ldots,$ $\mathtt{basePosition + n}$. 



The range of this list is $[\mathtt{basePosition}, \mathtt{basePosition + n}]$. }

 

\subsection{\textcolor{Chapter }{MiddleStart}}
\logpage{[ 10, 2, 26 ]}\nobreak
\hyperdef{L}{X8134072B7C7309CD}{}
{\noindent\textcolor{FuncColor}{$\triangleright$\enspace\texttt{MiddleStart({\mdseries\slshape list})\index{MiddleStart@\texttt{MiddleStart}}
\label{MiddleStart}
}\hfill{\scriptsize (operation)}}\\


 Arguments: \mbox{\texttt{\mdseries\slshape list}} \texttt{\symbol{45}}\texttt{\symbol{45}} an infinite list. \\
 \textbf{\indent Returns:\ }
The starting position of the "middle" part of \mbox{\texttt{\mdseries\slshape list}}. 



 }

 

\subsection{\textcolor{Chapter }{MiddleEnd}}
\logpage{[ 10, 2, 27 ]}\nobreak
\hyperdef{L}{X7EB933657AF3D723}{}
{\noindent\textcolor{FuncColor}{$\triangleright$\enspace\texttt{MiddleEnd({\mdseries\slshape list})\index{MiddleEnd@\texttt{MiddleEnd}}
\label{MiddleEnd}
}\hfill{\scriptsize (operation)}}\\


 Arguments: \mbox{\texttt{\mdseries\slshape list}} \texttt{\symbol{45}}\texttt{\symbol{45}} an infinite list. \\
 \textbf{\indent Returns:\ }
The ending position of the middle part of \mbox{\texttt{\mdseries\slshape list}}. 



 }

 

\subsection{\textcolor{Chapter }{MiddlePart}}
\logpage{[ 10, 2, 28 ]}\nobreak
\hyperdef{L}{X7AC3ECBB7BC93594}{}
{\noindent\textcolor{FuncColor}{$\triangleright$\enspace\texttt{MiddlePart({\mdseries\slshape list})\index{MiddlePart@\texttt{MiddlePart}}
\label{MiddlePart}
}\hfill{\scriptsize (operation)}}\\


 Arguments: \mbox{\texttt{\mdseries\slshape list}} \texttt{\symbol{45}}\texttt{\symbol{45}} an infinite list. \\
 \textbf{\indent Returns:\ }
The middle part (as a list) of \mbox{\texttt{\mdseries\slshape list}}. 



 }

 

\subsection{\textcolor{Chapter }{PositivePart}}
\logpage{[ 10, 2, 29 ]}\nobreak
\hyperdef{L}{X7896D6DD788046E7}{}
{\noindent\textcolor{FuncColor}{$\triangleright$\enspace\texttt{PositivePart({\mdseries\slshape list})\index{PositivePart@\texttt{PositivePart}}
\label{PositivePart}
}\hfill{\scriptsize (operation)}}\\


 Arguments: \mbox{\texttt{\mdseries\slshape list}} \texttt{\symbol{45}}\texttt{\symbol{45}} an infinite list. \\
 \textbf{\indent Returns:\ }
The positive part (as a half infinite list) of \mbox{\texttt{\mdseries\slshape list}}. 



 }

 

\subsection{\textcolor{Chapter }{NegativePart}}
\logpage{[ 10, 2, 30 ]}\nobreak
\hyperdef{L}{X81B768838567F98D}{}
{\noindent\textcolor{FuncColor}{$\triangleright$\enspace\texttt{NegativePart({\mdseries\slshape list})\index{NegativePart@\texttt{NegativePart}}
\label{NegativePart}
}\hfill{\scriptsize (operation)}}\\


 Arguments: \mbox{\texttt{\mdseries\slshape list}} \texttt{\symbol{45}}\texttt{\symbol{45}} an infinite list. \\
 \textbf{\indent Returns:\ }
The negative part (as a halft infinite list) of \mbox{\texttt{\mdseries\slshape list}}. 



 }

 

\subsection{\textcolor{Chapter }{HighestKnownPosition}}
\logpage{[ 10, 2, 31 ]}\nobreak
\hyperdef{L}{X8511A0B88003B989}{}
{\noindent\textcolor{FuncColor}{$\triangleright$\enspace\texttt{HighestKnownPosition({\mdseries\slshape list})\index{HighestKnownPosition@\texttt{HighestKnownPosition}}
\label{HighestKnownPosition}
}\hfill{\scriptsize (operation)}}\\


 Arguments: \mbox{\texttt{\mdseries\slshape list}} \texttt{\symbol{45}}\texttt{\symbol{45}} an infinite list. \\
 \textbf{\indent Returns:\ }
The highest index $i$ such that the value at position $i$ is known without computation. 



 }

 

\subsection{\textcolor{Chapter }{UpperBound(list)}}
\logpage{[ 10, 2, 32 ]}\nobreak
\hyperdef{L}{X81CD94E284DE7E67}{}
{\noindent\textcolor{FuncColor}{$\triangleright$\enspace\texttt{UpperBound(list)({\mdseries\slshape list})\index{UpperBound(list)@\texttt{UpperBound(list)}}
\label{UpperBound(list)}
}\hfill{\scriptsize (operation)}}\\


 Arguments: \mbox{\texttt{\mdseries\slshape list}} \texttt{\symbol{45}}\texttt{\symbol{45}} an infinite list. \\
 \textbf{\indent Returns:\ }
The highest index in the range of the list. 



 }

 

\subsection{\textcolor{Chapter }{LowerBound(list)}}
\logpage{[ 10, 2, 33 ]}\nobreak
\hyperdef{L}{X8639C48C7E19956A}{}
{\noindent\textcolor{FuncColor}{$\triangleright$\enspace\texttt{LowerBound(list)({\mdseries\slshape list})\index{LowerBound(list)@\texttt{LowerBound(list)}}
\label{LowerBound(list)}
}\hfill{\scriptsize (operation)}}\\


 Arguments: \mbox{\texttt{\mdseries\slshape list}} \texttt{\symbol{45}}\texttt{\symbol{45}} an infinite list. \\
 \textbf{\indent Returns:\ }
The lowest index in the range of the list. 



 }

 

\subsection{\textcolor{Chapter }{FinitePartAsList}}
\logpage{[ 10, 2, 34 ]}\nobreak
\hyperdef{L}{X7A311EBE79F31AEA}{}
{\noindent\textcolor{FuncColor}{$\triangleright$\enspace\texttt{FinitePartAsList({\mdseries\slshape list, startPos, endPos})\index{FinitePartAsList@\texttt{FinitePartAsList}}
\label{FinitePartAsList}
}\hfill{\scriptsize (operation)}}\\


 Arguments: \mbox{\texttt{\mdseries\slshape list}} \texttt{\symbol{45}}\texttt{\symbol{45}} an infinite list, \mbox{\texttt{\mdseries\slshape startPos}} \texttt{\symbol{45}}\texttt{\symbol{45}} an integer, \mbox{\texttt{\mdseries\slshape endPos}} \texttt{\symbol{45}}\texttt{\symbol{45}} an integer.\\
 \textbf{\indent Returns:\ }
A list containing the values of \mbox{\texttt{\mdseries\slshape list}} in positions $\mathtt{endPos},\ldots,\mathtt{startPos}$. 



Note that both integers in the input must be within the index range of \mbox{\texttt{\mdseries\slshape list}}. }

 

\subsection{\textcolor{Chapter }{PositivePartFrom}}
\logpage{[ 10, 2, 35 ]}\nobreak
\hyperdef{L}{X8028ED6883CFA39C}{}
{\noindent\textcolor{FuncColor}{$\triangleright$\enspace\texttt{PositivePartFrom({\mdseries\slshape list, pos})\index{PositivePartFrom@\texttt{PositivePartFrom}}
\label{PositivePartFrom}
}\hfill{\scriptsize (operation)}}\\


 Arguments: \mbox{\texttt{\mdseries\slshape list}} \texttt{\symbol{45}}\texttt{\symbol{45}} an infinite list, \mbox{\texttt{\mdseries\slshape pos}} \texttt{\symbol{45}}\texttt{\symbol{45}} an integer.\\
 \textbf{\indent Returns:\ }
An infinite list (\emph{not} a half infinite list) with index range from \texttt{pos} to \texttt{UpperBound(list)}. 



The value at position $i$ of the new infinite list is the same as the value at position $i$ of \mbox{\texttt{\mdseries\slshape list}}. }

 

\subsection{\textcolor{Chapter }{NegativePartFrom}}
\logpage{[ 10, 2, 36 ]}\nobreak
\hyperdef{L}{X7863903B7E281CF6}{}
{\noindent\textcolor{FuncColor}{$\triangleright$\enspace\texttt{NegativePartFrom({\mdseries\slshape list, pos})\index{NegativePartFrom@\texttt{NegativePartFrom}}
\label{NegativePartFrom}
}\hfill{\scriptsize (operation)}}\\


 Arguments: \mbox{\texttt{\mdseries\slshape list}} \texttt{\symbol{45}}\texttt{\symbol{45}} an infinite list, \mbox{\texttt{\mdseries\slshape pos}} \texttt{\symbol{45}}\texttt{\symbol{45}} an integer.\\
 \textbf{\indent Returns:\ }
An infinite list (\emph{not} a half infinite list) with index range from \texttt{LowerBound(list)} to \texttt{pos}. 



The value at position $i$ of the new infinite list is the same as the value at position $i$ of \mbox{\texttt{\mdseries\slshape list}}. }

 

\subsection{\textcolor{Chapter }{Splice}}
\logpage{[ 10, 2, 37 ]}\nobreak
\hyperdef{L}{X80153272857B2EDE}{}
{\noindent\textcolor{FuncColor}{$\triangleright$\enspace\texttt{Splice({\mdseries\slshape positiveList, negativeList, joinPosition})\index{Splice@\texttt{Splice}}
\label{Splice}
}\hfill{\scriptsize (operation)}}\\


 Arguments: \mbox{\texttt{\mdseries\slshape positiveList}} \texttt{\symbol{45}}\texttt{\symbol{45}} an infinite list, \mbox{\texttt{\mdseries\slshape negativeList}} \texttt{\symbol{45}}\texttt{\symbol{45}} an infinite list, \mbox{\texttt{\mdseries\slshape joinPosition}} \texttt{\symbol{45}}\texttt{\symbol{45}} an integer.\\
 \textbf{\indent Returns:\ }
A new infinite list which is identical to \mbox{\texttt{\mdseries\slshape positiveList}} for indices greater than \mbox{\texttt{\mdseries\slshape joinPosition}} and identical to \mbox{\texttt{\mdseries\slshape negativeList}} for indices smaller than or equal to \mbox{\texttt{\mdseries\slshape joinPosition}}. 



 }

 

\subsection{\textcolor{Chapter }{InfConcatenation}}
\logpage{[ 10, 2, 38 ]}\nobreak
\hyperdef{L}{X84FC63C183F0A977}{}
{\noindent\textcolor{FuncColor}{$\triangleright$\enspace\texttt{InfConcatenation({\mdseries\slshape arg})\index{InfConcatenation@\texttt{InfConcatenation}}
\label{InfConcatenation}
}\hfill{\scriptsize (function)}}\\


 Arguments: \mbox{\texttt{\mdseries\slshape arg}} \texttt{\symbol{45}}\texttt{\symbol{45}} a number of infinite lists. \\
 \textbf{\indent Returns:\ }
A new infinite list.



 If the length of \mbox{\texttt{\mdseries\slshape arg}} is greater than or equal to $2$, then the new infinite list consists of the following parts: It has the
positive part of \texttt{arg[1]}, and the middle part is the concatenation of the middle parts of all lists in \mbox{\texttt{\mdseries\slshape arg}}, such that \texttt{MiddleEnd} of the new list is the same as \texttt{MiddleEnd( arg[1] )}. The negative part of the new list is the negative part of \texttt{arg[Length(arg)]}, although shiftet so that it starts in the correct position. }

 
\begin{Verbatim}[commandchars=!@|,fontsize=\small,frame=single,label=Example]
  !gapprompt@gap>| !gapinput@# we do an InfConcatenation of three lists.|
  !gapprompt@gap>| !gapinput@f := function(x) return x; end;;|
  !gapprompt@gap>| !gapinput@g := function(x) return x+1; end;;|
  !gapprompt@gap>| !gapinput@h := function(x) return x^2; end;;|
  !gapprompt@gap>| !gapinput@InfList1 := MakeInfList( 0, [ 10 ], [ "pos", f, false ], |
  !gapprompt@>| !gapinput@                            [ "repeat", [ 10, 15 ] ], false );|
  <object>
  !gapprompt@gap>| !gapinput@InfList2 := MakeInfList( 0, [ 20 ], [ "pos", g, false ], |
  !gapprompt@>| !gapinput@                            [ "repeat", [ 20, 25 ] ], false );|
  <object>
  !gapprompt@gap>| !gapinput@InfList3 := MakeInfList( 0, [ 30 ], [ "pos", h, false ], |
  !gapprompt@>| !gapinput@                            [ "repeat", [ 30, 35 ] ], false );|
  <object>
  !gapprompt@gap>| !gapinput@concList := InfConcatenation( InfList1, InfList2, InfList3 );|
  <object> 
  !gapprompt@gap>| !gapinput@MiddlePart(concList);|
  [ 30, 20, 10 ] 
\end{Verbatim}
 The newly created \texttt{concList} looks as follows around the middle part: 
\[ \begin{array}{lrrrrrrrrrrr} \text{position} & \cdots & 3 & 2 & 1 & 0 & -1 & -2
& -3 & -4 & -5 & \cdots \\ \text{value} & \cdots & 3 & 2 & 1 & 10 & 20 & 30 &
30 & 35 & 30 & \cdots \\ \end{array} \]
 

\subsection{\textcolor{Chapter }{InfList}}
\logpage{[ 10, 2, 39 ]}\nobreak
\hyperdef{L}{X7F4CF360858E80A6}{}
{\noindent\textcolor{FuncColor}{$\triangleright$\enspace\texttt{InfList({\mdseries\slshape list, func})\index{InfList@\texttt{InfList}}
\label{InfList}
}\hfill{\scriptsize (operation)}}\\


 Arguments: \mbox{\texttt{\mdseries\slshape list}} \texttt{\symbol{45}}\texttt{\symbol{45}} an infinite list, \mbox{\texttt{\mdseries\slshape func}} \texttt{\symbol{45}}\texttt{\symbol{45}} a function which takes an element of
the list as argument. \\
 \textbf{\indent Returns:\ }
An infinite list with the same range as \mbox{\texttt{\mdseries\slshape list}}, where the value at position $i$ is the image of the value at position $i$ of \mbox{\texttt{\mdseries\slshape list}} under \mbox{\texttt{\mdseries\slshape func}}. 

}

 

\subsection{\textcolor{Chapter }{IntegersList}}
\logpage{[ 10, 2, 40 ]}\nobreak
\hyperdef{L}{X791D879483C59410}{}
{\noindent\textcolor{FuncColor}{$\triangleright$\enspace\texttt{IntegersList\index{IntegersList@\texttt{IntegersList}}
\label{IntegersList}
}\hfill{\scriptsize (global variable)}}\\


 An infinite list with range $(-\infty,\infty)$ where the value at position $i$ is the number $i$ (that is, a representation of the integers). }

 }

 
\section{\textcolor{Chapter }{Representation of categories}}\label{Cats}
\logpage{[ 10, 3, 0 ]}
\hyperdef{L}{X7CAF603281B94AC8}{}
{
  A chain complex consists of objects and morphisms from some category. In QPA,
this category will usually be the category of right modules over some quotient
of a path algebra. 

\subsection{\textcolor{Chapter }{IsCat}}
\logpage{[ 10, 3, 1 ]}\nobreak
\hyperdef{L}{X7F0E546178977963}{}
{\noindent\textcolor{FuncColor}{$\triangleright$\enspace\texttt{IsCat\index{IsCat@\texttt{IsCat}}
\label{IsCat}
}\hfill{\scriptsize (Category)}}\\


 The category for categories. A category is a record, storing a number of
properties that is specified within each category. Two categories can be
compared using \texttt{=}. Currently, the only implemented category is the one of right modules over a
(quotient of a) path algebra. }

 

\subsection{\textcolor{Chapter }{CatOfRightAlgebraModules}}
\logpage{[ 10, 3, 2 ]}\nobreak
\hyperdef{L}{X7BB9B02E7E76FF9C}{}
{\noindent\textcolor{FuncColor}{$\triangleright$\enspace\texttt{CatOfRightAlgebraModules({\mdseries\slshape A})\index{CatOfRightAlgebraModules@\texttt{CatOfRightAlgebraModules}}
\label{CatOfRightAlgebraModules}
}\hfill{\scriptsize (operation)}}\\


Arguments: \mbox{\texttt{\mdseries\slshape A}} \texttt{\symbol{45}}\texttt{\symbol{45}} a (quotient of a) path algebra.\\
 \textbf{\indent Returns:\ }
The category mod $A$. 



 mod $A$ has several properties, which can be accessed using the \texttt{.} mark. Some of the properties store functions. All properties are demonstrated
in the following example. 
\begin{itemize}
\item \texttt{zeroObj} \texttt{\symbol{45}}\texttt{\symbol{45}} returns the zero module of mod $A$.
\item \texttt{isZeroObj} \texttt{\symbol{45}}\texttt{\symbol{45}} returns true if the given module is
zero.
\item \texttt{zeroMap} \texttt{\symbol{45}}\texttt{\symbol{45}} returns the ZeroMapping function.
\item \texttt{isZeroMapping} \texttt{\symbol{45}}\texttt{\symbol{45}} returns the IsZero test.
\item \texttt{composeMaps} \texttt{\symbol{45}}\texttt{\symbol{45}} returns the composition of the two
given maps.
\item \texttt{ker} \texttt{\symbol{45}}\texttt{\symbol{45}} returns the Kernel function.
\item \texttt{im} \texttt{\symbol{45}}\texttt{\symbol{45}} returns the Image function.
\item \texttt{isExact} \texttt{\symbol{45}}\texttt{\symbol{45}} returns true if two consecutive maps
are exact. 
\end{itemize}
 }

 
\begin{Verbatim}[commandchars=!@|,fontsize=\small,frame=single,label=Example]
  !gapprompt@gap>| !gapinput@alg;|
  <algebra-with-one over Rationals, with 7 generators>
  !gapprompt@gap>| !gapinput@# L, M, and N are alg-modules|
  !gapprompt@gap>| !gapinput@# f: L --> M and g: M --> N are non-zero morphisms|
  !gapprompt@gap>| !gapinput@cat := CatOfRightAlgebraModules(alg);|
  <cat: right modules over algebra>
  !gapprompt@gap>| !gapinput@cat.zeroObj;|
  <right-module over <algebra-with-one over Rationals, with 7 generators>>
  !gapprompt@gap>| !gapinput@cat.isZeroObj(M);|
  false
  !gapprompt@gap>| !gapinput@cat.zeroMap(M,N);|
  <mapping: <3-dimensional right-module over AlgebraWithOne( Rationals, 
  [ [(1)*v1], [(1)*v2], [(1)*v3], [(1)*v4], [(1)*a], [(1)*b], [(1)*c] ])> -> 
    <1-dimensional right-module over AlgebraWithOne( Rationals,
    [ [(1)*v1], [(1)*v2], [(1)*v3], [(1)*v4], [(1)*a], [(1)*b], [(1)*c] ] )> >
  !gapprompt@gap>| !gapinput@cat.composeMaps(g,f);|
  <mapping: <1-dimensional right-module over AlgebraWithOne( Rationals, 
    [ [(1)*v1], [(1)*v2], [(1)*v3], [(1)*v4], [(1)*a], [(1)*b], [(1)*c]]
    -> <1-dimensional right-module over AlgebraWithOne( Rationals,
    [ [(1)*v1], [(1)*v2], [(1)*v3], [(1)*v4], [(1)*a], [(1)*b], [(1)*c] ] )> >
  !gapprompt@gap>| !gapinput@cat.ker(g);|
  <2-dimensional right-module over <algebra-with-one over Rationals,
    with 7 generators>>
  !gapprompt@gap>| !gapinput@cat.isExact(g,f);|
  false 
\end{Verbatim}
 }

 
\section{\textcolor{Chapter }{Making a complex}}\logpage{[ 10, 4, 0 ]}
\hyperdef{L}{X7EC3D95F7C791F7E}{}
{
  The most general constructor for complexes is the function \texttt{Complex} (\ref{Complex}). In addition to this, there are constructors for common special cases: 
\begin{itemize}
\item \texttt{ZeroComplex} (\ref{ZeroComplex})
\item \texttt{StalkComplex} (\ref{StalkComplex})
\item \texttt{FiniteComplex} (\ref{FiniteComplex})
\item \texttt{ShortExactSequence} (\ref{ShortExactSequence})
\end{itemize}
  

\subsection{\textcolor{Chapter }{IsQPAComplex}}
\logpage{[ 10, 4, 1 ]}\nobreak
\hyperdef{L}{X838C9EBD87B8CB9D}{}
{\noindent\textcolor{FuncColor}{$\triangleright$\enspace\texttt{IsQPAComplex\index{IsQPAComplex@\texttt{IsQPAComplex}}
\label{IsQPAComplex}
}\hfill{\scriptsize (Category)}}\\


 The category for chain complexes. }

 

\subsection{\textcolor{Chapter }{IsZeroComplex}}
\logpage{[ 10, 4, 2 ]}\nobreak
\hyperdef{L}{X7E0CE8AC87533991}{}
{\noindent\textcolor{FuncColor}{$\triangleright$\enspace\texttt{IsZeroComplex\index{IsZeroComplex@\texttt{IsZeroComplex}}
\label{IsZeroComplex}
}\hfill{\scriptsize (Category)}}\\


 Category for zero complexes, subcategory of \texttt{IsQPAComplex} (\ref{IsQPAComplex}). }

 

\subsection{\textcolor{Chapter }{Complex}}
\logpage{[ 10, 4, 3 ]}\nobreak
\hyperdef{L}{X7E51958C86D73D17}{}
{\noindent\textcolor{FuncColor}{$\triangleright$\enspace\texttt{Complex({\mdseries\slshape cat, baseDegree, middle, positive, negative})\index{Complex@\texttt{Complex}}
\label{Complex}
}\hfill{\scriptsize (function)}}\\
\textbf{\indent Returns:\ }
A newly created chain complex



 The first argument, \mbox{\texttt{\mdseries\slshape cat}} is an \texttt{IsCat} (\ref{IsCat}) object describing the category to create a chain complex over.

 The rest of the arguments describe the differentials of the complex. These are
divided into three parts: one finite (``middle'') and two infinite (``positive'' and ``negative''). The positive part contains all differentials in degrees higher than those
in the middle part, and the negative part contains all differentials in
degrees lower than those in the middle part. (The middle part may be placed
anywhere, so the positive part can \texttt{\symbol{45}}\texttt{\symbol{45}}
despite its name \texttt{\symbol{45}}\texttt{\symbol{45}} contain some
differentials of negative degree. Conversely, the negative part can contain
some differentials of positive degree.)

 The argument \mbox{\texttt{\mdseries\slshape middle}} is a list containing the differentials for the middle part. The argument \mbox{\texttt{\mdseries\slshape baseDegree}} gives the degree of the first differential in this list. The second
differential is placed in degree $\mbox{\texttt{\mdseries\slshape baseDegree}}+1$, and so on. Thus, the middle part consists of the degrees 
\[ \mbox{\texttt{\mdseries\slshape baseDegree}},\quad \mbox{\texttt{\mdseries\slshape baseDegree}} + 1,\quad \ldots\quad \mbox{\texttt{\mdseries\slshape baseDegree}} + \text{Length}(\mbox{\texttt{\mdseries\slshape middle}}). \]
 Each of the arguments \mbox{\texttt{\mdseries\slshape positive}} and \mbox{\texttt{\mdseries\slshape negative}} can be one of the following: 
\begin{itemize}
\item The string \texttt{"zero"}, meaning that the part contains only zero objects and zero morphisms.
\item A list of the form \texttt{[ "repeat", L ]}, where \texttt{L} is a list of morphisms. The part will contain the differentials in \texttt{L} repeated infinitely many times. The convention for the order of elements in \texttt{L} is that \texttt{L[1]} is the differential which is closest to the middle part, and \texttt{L[Length(L)]} is farthest away from the middle part.
\item A list of the form \texttt{[ "pos", f ]} or \texttt{[ "pos", f, store ]}, where \texttt{f} is a function of two arguments, and \texttt{store} (if included) is a boolean. The function \texttt{f} is used to compute the differentials in this part. The function \texttt{f} is not called immediately by the \texttt{Complex} constructor, but will be called later as the differentials in this part are
needed. The function call \texttt{f(C,i)} (where \texttt{C} is the complex and \texttt{i} an integer) should produce the differential in degree \texttt{i}. The function may use \texttt{C} to look up other differentials in the complex, as long as this does not cause
an infinite loop. If \texttt{store} is \texttt{true} (or not specified), each computed differential is stored, and they are
computed in order from the one closest to the middle part, regardless of which
order they are requested in.
\item A list of the form \texttt{[ "next", f, init ]}, where \texttt{f} is a function of one argument, and \texttt{init} is a morphism. The function \texttt{f} is used to compute the differentials in this part. For the first differential
in the part (that is, the one closest to the middle part), \texttt{f} is called with \texttt{init} as argument. For the next differential, \texttt{f} is called with the first differential as argument, and so on. Thus, the
differentials are 
\[ f(\text{init}),\quad f^2(\text{init}),\quad f^3(\text{init}),\quad \ldots \]
 Each differential is stored when it has been computed. 
\item A list of the form \texttt{[ "next/repeat", f, init ]}. This works like the type \texttt{"next"}, but may be automatically converted to type \texttt{"repeat"} later, if it is discovered that the differentials are repeating. 
\end{itemize}
 }

 
\begin{Verbatim}[commandchars=!@|,fontsize=\small,frame=single,label=Example]
  !gapprompt@gap>| !gapinput@A := PathAlgebra( Rationals, Quiver( 2, [ [ 1, 2, "a" ] ] ) );;|
  !gapprompt@gap>| !gapinput@M := RightModuleOverPathAlgebra( A, [ 2, 2 ], [ [ "a", [ [ 1, 0 ], [ 0, 1 ] ] ] ] );;|
  !gapprompt@gap>| !gapinput@d := RightModuleHomOverAlgebra( M, M, [ [ [ 0, 0 ], [ 1, 0 ] ], [ [ 0, 0 ], [ 1, 0 ] ] ] );;|
  !gapprompt@gap>| !gapinput@IsZero( d * d );|
  true
  !gapprompt@gap>| !gapinput@C := Complex( CatOfRightAlgebraModules( A ), 0, [ d ],|
  !gapprompt@>| !gapinput@                 [ "next/repeat", function( x ) return d; end, d ], "zero" );|
  --- -> 0:(2,2) -> -1:(2,2) -> 0
  !gapprompt@gap>| !gapinput@ObjectOfComplex( C, 3 );|
  <[ 2, 2 ]>
  !gapprompt@gap>| !gapinput@C;|
  --- -> [ 1:(2,2) -> ] 0:(2,2) -> -1:(2,2) -> 0 
\end{Verbatim}
 

\subsection{\textcolor{Chapter }{ZeroComplex}}
\logpage{[ 10, 4, 4 ]}\nobreak
\hyperdef{L}{X80F8DD20789BBBC0}{}
{\noindent\textcolor{FuncColor}{$\triangleright$\enspace\texttt{ZeroComplex({\mdseries\slshape cat})\index{ZeroComplex@\texttt{ZeroComplex}}
\label{ZeroComplex}
}\hfill{\scriptsize (function)}}\\
\textbf{\indent Returns:\ }
A newly created zero complex



 This function creates a zero complex (a complex consisting of only zero
objects and zero morphisms) over the category described by the \texttt{IsCat} (\ref{IsCat}) object \mbox{\texttt{\mdseries\slshape cat}}. }

 

\subsection{\textcolor{Chapter }{FiniteComplex}}
\logpage{[ 10, 4, 5 ]}\nobreak
\hyperdef{L}{X7C75868E82A831C8}{}
{\noindent\textcolor{FuncColor}{$\triangleright$\enspace\texttt{FiniteComplex({\mdseries\slshape cat, baseDegree, differentials})\index{FiniteComplex@\texttt{FiniteComplex}}
\label{FiniteComplex}
}\hfill{\scriptsize (function)}}\\
\textbf{\indent Returns:\ }
A newly created complex



 This function creates a complex where all but finitely many objects are the
zero object.

 The argument \mbox{\texttt{\mdseries\slshape cat}} is an \texttt{IsCat} (\ref{IsCat}) object describing the category to create a chain complex over.

 The argument \mbox{\texttt{\mdseries\slshape differentials}} is a list of morphisms. The argument \mbox{\texttt{\mdseries\slshape baseDegree}} gives the degree for the first differential in this list. The subsequent
differentials are placed in degrees $\mbox{\texttt{\mdseries\slshape baseDegree}}+1$, and so on.

 This means that the \mbox{\texttt{\mdseries\slshape differentials}} argument specifies the differentials in degrees 
\[ \mbox{\texttt{\mdseries\slshape baseDegree}},\quad \mbox{\texttt{\mdseries\slshape baseDegree}} + 1,\quad \ldots \quad \mbox{\texttt{\mdseries\slshape baseDegree}} + \text{Length}(\mbox{\texttt{\mdseries\slshape differentials}}); \]
 and thus implicitly the objects in degrees 
\[ \mbox{\texttt{\mdseries\slshape baseDegree}} - 1,\quad \mbox{\texttt{\mdseries\slshape baseDegree}},\quad \ldots \quad \mbox{\texttt{\mdseries\slshape baseDegree}} + \text{Length}(\mbox{\texttt{\mdseries\slshape differentials}}). \]
 All other objects in the complex are zero. }

 
\begin{Verbatim}[commandchars=!@|,fontsize=\small,frame=single,label=Example]
  !gapprompt@gap>| !gapinput@# L, M and N are modules over the same algebra A|
  !gapprompt@gap>| !gapinput@# cat is the category mod A|
  !gapprompt@gap>| !gapinput@# f: L --> M and g: M --> N maps|
  !gapprompt@gap>| !gapinput@C := FiniteComplex(cat, 1, [g,f]);|
  0 -> 2:(1,0) -> 1:(2,2) -> 0:(1,1) -> 0 
\end{Verbatim}
 

\subsection{\textcolor{Chapter }{StalkComplex}}
\logpage{[ 10, 4, 6 ]}\nobreak
\hyperdef{L}{X80B47D78785F6859}{}
{\noindent\textcolor{FuncColor}{$\triangleright$\enspace\texttt{StalkComplex({\mdseries\slshape cat, obj, degree})\index{StalkComplex@\texttt{StalkComplex}}
\label{StalkComplex}
}\hfill{\scriptsize (function)}}\\


 Arguments: \mbox{\texttt{\mdseries\slshape cat}} \texttt{\symbol{45}}\texttt{\symbol{45}} a category, \mbox{\texttt{\mdseries\slshape obj}} \texttt{\symbol{45}}\texttt{\symbol{45}} an object in \mbox{\texttt{\mdseries\slshape cat}}, \mbox{\texttt{\mdseries\slshape degree}} \texttt{\symbol{45}}\texttt{\symbol{45}} the degree \mbox{\texttt{\mdseries\slshape obj}} should be placed in.\\
 \textbf{\indent Returns:\ }
a newly created complex. 



 The new complex is a stalk complex with \mbox{\texttt{\mdseries\slshape obj}} in position \mbox{\texttt{\mdseries\slshape degree}}, and zero elsewhere. }

 
\begin{Verbatim}[commandchars=!@|,fontsize=\small,frame=single,label=Example]
  !gapprompt@gap>| !gapinput@Ms := StalkComplex(cat, M, 3);|
  0 -> 3:(2,2) -> 0 
\end{Verbatim}
 

\subsection{\textcolor{Chapter }{ShortExactSequence}}
\logpage{[ 10, 4, 7 ]}\nobreak
\hyperdef{L}{X7A685484784087FE}{}
{\noindent\textcolor{FuncColor}{$\triangleright$\enspace\texttt{ShortExactSequence({\mdseries\slshape cat, f, g})\index{ShortExactSequence@\texttt{ShortExactSequence}}
\label{ShortExactSequence}
}\hfill{\scriptsize (function)}}\\


 Arguments: \mbox{\texttt{\mdseries\slshape cat}} \texttt{\symbol{45}}\texttt{\symbol{45}} a category, \mbox{\texttt{\mdseries\slshape f}} and \mbox{\texttt{\mdseries\slshape g}} \texttt{\symbol{45}}\texttt{\symbol{45}} maps in \mbox{\texttt{\mdseries\slshape cat}}, where \mbox{\texttt{\mdseries\slshape f}}: $A \rightarrow B$ and \mbox{\texttt{\mdseries\slshape g}}: $B \rightarrow C$.\\
 \textbf{\indent Returns:\ }
a newly created complex. 



If the sequence $0 \rightarrow A \rightarrow B \rightarrow C \rightarrow 0$ is exact, this complex (with $B$ in degree 0) is returned. }

 
\begin{Verbatim}[commandchars=!@|,fontsize=\small,frame=single,label=Example]
  !gapprompt@gap>| !gapinput@ses := ShortExactSequence(cat, f, g);|
  0 -> 1:(0,0,1,0) -> 0:(0,1,1,1) -> -1:(0,1,0,1) -> 0 
\end{Verbatim}
 }

 
\section{\textcolor{Chapter }{Information about a complex}}\logpage{[ 10, 5, 0 ]}
\hyperdef{L}{X7F7E0C197FA31B29}{}
{
  

\subsection{\textcolor{Chapter }{CatOfComplex}}
\logpage{[ 10, 5, 1 ]}\nobreak
\hyperdef{L}{X7BBF131485F3C792}{}
{\noindent\textcolor{FuncColor}{$\triangleright$\enspace\texttt{CatOfComplex({\mdseries\slshape C})\index{CatOfComplex@\texttt{CatOfComplex}}
\label{CatOfComplex}
}\hfill{\scriptsize (attribute)}}\\
\textbf{\indent Returns:\ }
 The category the objects of the complex \mbox{\texttt{\mdseries\slshape C}} live in. 



 }

 

\subsection{\textcolor{Chapter }{ObjectOfComplex}}
\logpage{[ 10, 5, 2 ]}\nobreak
\hyperdef{L}{X8239596E87830A51}{}
{\noindent\textcolor{FuncColor}{$\triangleright$\enspace\texttt{ObjectOfComplex({\mdseries\slshape C, i})\index{ObjectOfComplex@\texttt{ObjectOfComplex}}
\label{ObjectOfComplex}
}\hfill{\scriptsize (operation)}}\\


 Arguments: \mbox{\texttt{\mdseries\slshape C}} \texttt{\symbol{45}}\texttt{\symbol{45}} a complex, \mbox{\texttt{\mdseries\slshape i}} \texttt{\symbol{45}}\texttt{\symbol{45}} an integer.\\
 \textbf{\indent Returns:\ }
 The object at position \mbox{\texttt{\mdseries\slshape i}} in the complex. 



 }

 

\subsection{\textcolor{Chapter }{DifferentialOfComplex}}
\logpage{[ 10, 5, 3 ]}\nobreak
\hyperdef{L}{X7B597AC77A3D71F9}{}
{\noindent\textcolor{FuncColor}{$\triangleright$\enspace\texttt{DifferentialOfComplex({\mdseries\slshape C, i})\index{DifferentialOfComplex@\texttt{DifferentialOfComplex}}
\label{DifferentialOfComplex}
}\hfill{\scriptsize (operation)}}\\


 Arguments: \mbox{\texttt{\mdseries\slshape C}} \texttt{\symbol{45}}\texttt{\symbol{45}} a complex, \mbox{\texttt{\mdseries\slshape i}} \texttt{\symbol{45}}\texttt{\symbol{45}} an integer.\\
 \textbf{\indent Returns:\ }
 The map in \mbox{\texttt{\mdseries\slshape C}} between objects at positions $i$ and $i-1$. 



 }

 

\subsection{\textcolor{Chapter }{DifferentialsOfComplex}}
\logpage{[ 10, 5, 4 ]}\nobreak
\hyperdef{L}{X858D2D1A8674D2EC}{}
{\noindent\textcolor{FuncColor}{$\triangleright$\enspace\texttt{DifferentialsOfComplex({\mdseries\slshape C})\index{DifferentialsOfComplex@\texttt{DifferentialsOfComplex}}
\label{DifferentialsOfComplex}
}\hfill{\scriptsize (attribute)}}\\


 Arguments: \mbox{\texttt{\mdseries\slshape C}} \texttt{\symbol{45}}\texttt{\symbol{45}} a complex\\
 \textbf{\indent Returns:\ }
 The differentials of the complex, stored as an \texttt{IsInfList} object. 

}

 

\subsection{\textcolor{Chapter }{CyclesOfComplex}}
\logpage{[ 10, 5, 5 ]}\nobreak
\hyperdef{L}{X85FBFD9C8382C93B}{}
{\noindent\textcolor{FuncColor}{$\triangleright$\enspace\texttt{CyclesOfComplex({\mdseries\slshape C, i})\index{CyclesOfComplex@\texttt{CyclesOfComplex}}
\label{CyclesOfComplex}
}\hfill{\scriptsize (operation)}}\\


 Arguments: \mbox{\texttt{\mdseries\slshape C}} \texttt{\symbol{45}}\texttt{\symbol{45}} a complex, \mbox{\texttt{\mdseries\slshape i}} \texttt{\symbol{45}}\texttt{\symbol{45}} an integer. \\
 \textbf{\indent Returns:\ }
The $i$\texttt{\symbol{45}}cycle of the complex, that is the subobject $Ker(d_i)$ of \texttt{ObjectOfComplex(C,i)}. 



 }

 

\subsection{\textcolor{Chapter }{BoundariesOfComplex}}
\logpage{[ 10, 5, 6 ]}\nobreak
\hyperdef{L}{X8168D8CC7805FE30}{}
{\noindent\textcolor{FuncColor}{$\triangleright$\enspace\texttt{BoundariesOfComplex({\mdseries\slshape C, i})\index{BoundariesOfComplex@\texttt{BoundariesOfComplex}}
\label{BoundariesOfComplex}
}\hfill{\scriptsize (operation)}}\\


 Arguments: \mbox{\texttt{\mdseries\slshape C}} \texttt{\symbol{45}}\texttt{\symbol{45}} a complex, \mbox{\texttt{\mdseries\slshape i}} \texttt{\symbol{45}}\texttt{\symbol{45}} an integer. \\
 \textbf{\indent Returns:\ }
 The $i$\texttt{\symbol{45}}boundary of the complex, that is the subobject $Im(d_{i+1})$ of \texttt{ObjectOfComplex(C,i)}. 



 }

 

\subsection{\textcolor{Chapter }{HomologyOfComplex}}
\logpage{[ 10, 5, 7 ]}\nobreak
\hyperdef{L}{X7F8D61447AB57EF3}{}
{\noindent\textcolor{FuncColor}{$\triangleright$\enspace\texttt{HomologyOfComplex({\mdseries\slshape C, i})\index{HomologyOfComplex@\texttt{HomologyOfComplex}}
\label{HomologyOfComplex}
}\hfill{\scriptsize (operation)}}\\


 Arguments: \mbox{\texttt{\mdseries\slshape C}} \texttt{\symbol{45}}\texttt{\symbol{45}} a complex, \mbox{\texttt{\mdseries\slshape i}} \texttt{\symbol{45}}\texttt{\symbol{45}} an integer.\\
 \textbf{\indent Returns:\ }
 The $i$th homology of the complex, that is, $Ker(d_i)/Im(d_{i+1})$. 



Note: this operation is currently not available. When working in the category
of right $kQ/I$\texttt{\symbol{45}}modules, it is possible to "cheat" and use the following
procedure to compute the homology of a complex: }

 
\begin{Verbatim}[commandchars=!@|,fontsize=\small,frame=single,label=Example]
  !gapprompt@gap>| !gapinput@C;                     |
  0 -> 4:(0,1) -> 3:(1,0) -> 2:(2,2) -> 1:(1,1) -> 0:(2,2) -> 0
  !gapprompt@gap>| !gapinput@# Want to compute the homology in degree 2|
  !gapprompt@gap>| !gapinput@f := DifferentialOfComplex(C,3);|
  <mapping: <1-dimensional right-module over AlgebraWithOne( Rationals, 
    [ [(1)*v1], [(1)*v2], [(1)*a], [(1)*b] ] )> ->
    < 4-dimensional right-module over AlgebraWithOne( Rationals, 
    [ [(1)*v1], [(1)*v2], [(1)*a], [(1)*b] ] )> >
  !gapprompt@gap>| !gapinput@g := KernelInclusion(DifferentialOfComplex(C,2));|
    <mapping: <2-dimensional right-module over AlgebraWithOne( Rationals, 
    [ [(1)*v1], [(1)*v2], [(1)*a], [(1)*b] ] )> ->
    < 4-dimensional right-module over AlgebraWithOne( Rationals, 
    [ [(1)*v1], [(1)*v2], [(1)*a], [(1)*b] ] )> >
  !gapprompt@gap>| !gapinput@# We know that Im f is included in Ker g, so can find the|
  !gapprompt@gap>| !gapinput@# lifting morphism h from C_3 to Ker g.|
  !gapprompt@gap>| !gapinput@h := LiftingInclusionMorphisms(g,f);|
    <mapping: <1-dimensional right-module over AlgebraWithOne( Rationals, 
    [ [(1)*v1], [(1)*v2], [(1)*a], [(1)*b] ] )> ->
    < 2-dimensional right-module over AlgebraWithOne( Rationals, 
    [ [(1)*v1], [(1)*v2], [(1)*a], [(1)*b] ] )> >
  !gapprompt@gap>| !gapinput@# The cokernel of h is Ker g / Im f |
  !gapprompt@gap>| !gapinput@Homology := CoKernel(h);|
  <1-dimensional right-module over <algebra-with-one over Rationals, with 
    4 generators>> 
\end{Verbatim}
 

\subsection{\textcolor{Chapter }{IsFiniteComplex}}
\logpage{[ 10, 5, 8 ]}\nobreak
\hyperdef{L}{X7DCBCF7682DC87FC}{}
{\noindent\textcolor{FuncColor}{$\triangleright$\enspace\texttt{IsFiniteComplex({\mdseries\slshape C})\index{IsFiniteComplex@\texttt{IsFiniteComplex}}
\label{IsFiniteComplex}
}\hfill{\scriptsize (operation)}}\\


 Arguments: \mbox{\texttt{\mdseries\slshape C}} \texttt{\symbol{45}}\texttt{\symbol{45}} a complex.\\
 \textbf{\indent Returns:\ }
true if \mbox{\texttt{\mdseries\slshape C}} is a finite complex, false otherwise. 



 }

 

\subsection{\textcolor{Chapter }{UpperBound(complex)}}
\logpage{[ 10, 5, 9 ]}\nobreak
\hyperdef{L}{X7D41C1AE7B4FA417}{}
{\noindent\textcolor{FuncColor}{$\triangleright$\enspace\texttt{UpperBound(complex)({\mdseries\slshape C})\index{UpperBound(complex)@\texttt{UpperBound(complex)}}
\label{UpperBound(complex)}
}\hfill{\scriptsize (operation)}}\\


 Arguments: \mbox{\texttt{\mdseries\slshape C}} \texttt{\symbol{45}}\texttt{\symbol{45}} a complex.\\
 \textbf{\indent Returns:\ }
If it exists: The smallest integer $i$ such that the object at position $i$ is non\texttt{\symbol{45}}zero, but for all $j > i$ the object at position $j$ is zero. 



If \mbox{\texttt{\mdseries\slshape C}} is not a finite complex, the operation will return fail or infinity, depending
on how \mbox{\texttt{\mdseries\slshape C}} was defined. }

 

\subsection{\textcolor{Chapter }{LowerBound(complex)}}
\logpage{[ 10, 5, 10 ]}\nobreak
\hyperdef{L}{X82E960CC81884F1A}{}
{\noindent\textcolor{FuncColor}{$\triangleright$\enspace\texttt{LowerBound(complex)({\mdseries\slshape C})\index{LowerBound(complex)@\texttt{LowerBound(complex)}}
\label{LowerBound(complex)}
}\hfill{\scriptsize (operation)}}\\


 Arguments: \mbox{\texttt{\mdseries\slshape C}} \texttt{\symbol{45}}\texttt{\symbol{45}} a complex.\\
 \textbf{\indent Returns:\ }
If it exists: The greatest integer $i$ such that the object at position $i$ is non\texttt{\symbol{45}}zero, but for all $j < i$ the object at position $j$ is zero. 



If \mbox{\texttt{\mdseries\slshape C}} is not a finite complex, the operation will return fail or negative infinity,
depending on how \mbox{\texttt{\mdseries\slshape C}} was defined. }

 

\subsection{\textcolor{Chapter }{LengthOfComplex}}
\logpage{[ 10, 5, 11 ]}\nobreak
\hyperdef{L}{X7C1F81277A5AC8B4}{}
{\noindent\textcolor{FuncColor}{$\triangleright$\enspace\texttt{LengthOfComplex({\mdseries\slshape C})\index{LengthOfComplex@\texttt{LengthOfComplex}}
\label{LengthOfComplex}
}\hfill{\scriptsize (operation)}}\\


 Arguments: \mbox{\texttt{\mdseries\slshape C}} \texttt{\symbol{45}}\texttt{\symbol{45}} a complex.\\
 \textbf{\indent Returns:\ }
the length of the complex. 



The length is defined as follows: If \mbox{\texttt{\mdseries\slshape C}} is a zero complex, the length is zero. If \mbox{\texttt{\mdseries\slshape C}} is a finite complex, the length is the upper bound
\texttt{\symbol{45}}\texttt{\symbol{45}} the lower bound + 1. If \mbox{\texttt{\mdseries\slshape C}} is an infinite complex, the length is infinity. }

 

\subsection{\textcolor{Chapter }{HighestKnownDegree}}
\logpage{[ 10, 5, 12 ]}\nobreak
\hyperdef{L}{X82D615CD796C9C0A}{}
{\noindent\textcolor{FuncColor}{$\triangleright$\enspace\texttt{HighestKnownDegree({\mdseries\slshape C})\index{HighestKnownDegree@\texttt{HighestKnownDegree}}
\label{HighestKnownDegree}
}\hfill{\scriptsize (operation)}}\\


 Arguments: \mbox{\texttt{\mdseries\slshape C}} \texttt{\symbol{45}}\texttt{\symbol{45}} a complex.\\
 \textbf{\indent Returns:\ }
The greatest integer $i$ such that the object at position $i$ is known (or computed). 



For a finite complex, this will be infinity. }

 

\subsection{\textcolor{Chapter }{LowestKnownDegree}}
\logpage{[ 10, 5, 13 ]}\nobreak
\hyperdef{L}{X868105FF86FA1B6E}{}
{\noindent\textcolor{FuncColor}{$\triangleright$\enspace\texttt{LowestKnownDegree({\mdseries\slshape C})\index{LowestKnownDegree@\texttt{LowestKnownDegree}}
\label{LowestKnownDegree}
}\hfill{\scriptsize (operation)}}\\


 Arguments: \mbox{\texttt{\mdseries\slshape C}} \texttt{\symbol{45}}\texttt{\symbol{45}} a complex.\\
 \textbf{\indent Returns:\ }
The smallest integer $i$ such that the object at position $i$ is known (or computed). 



For a finite complex, this will be negative infinity. }

 
\begin{Verbatim}[commandchars=!@|,fontsize=\small,frame=single,label=Example]
  !gapprompt@gap>| !gapinput@C;|
  0 -> 4:(0,1) -> 3:(1,0) -> 2:(2,2) -> 1:(1,1) -> 0:(2,2) -> 0
  !gapprompt@gap>| !gapinput@IsFiniteComplex(C);|
  true
  !gapprompt@gap>| !gapinput@UpperBound(C);|
  4
  !gapprompt@gap>| !gapinput@LowerBound(C);|
  0
  !gapprompt@gap>| !gapinput@LengthOfComplex(C);|
  5
  !gapprompt@gap>| !gapinput@HighestKnownDegree(C);|
  +inf
  !gapprompt@gap>| !gapinput@LowestKnownDegree(C);|
  -inf 
\end{Verbatim}
 

\subsection{\textcolor{Chapter }{IsExactSequence}}
\logpage{[ 10, 5, 14 ]}\nobreak
\hyperdef{L}{X793465497B435197}{}
{\noindent\textcolor{FuncColor}{$\triangleright$\enspace\texttt{IsExactSequence({\mdseries\slshape C})\index{IsExactSequence@\texttt{IsExactSequence}}
\label{IsExactSequence}
}\hfill{\scriptsize (property)}}\\


 Arguments: \mbox{\texttt{\mdseries\slshape C}} \texttt{\symbol{45}}\texttt{\symbol{45}} a complex.\\
 \textbf{\indent Returns:\ }
true if \mbox{\texttt{\mdseries\slshape C}} is exact at every position. 



If the complex is not finite and not repeating, the function fails. }

 

\subsection{\textcolor{Chapter }{IsExactInDegree}}
\logpage{[ 10, 5, 15 ]}\nobreak
\hyperdef{L}{X7A6AE1C378E2EEE2}{}
{\noindent\textcolor{FuncColor}{$\triangleright$\enspace\texttt{IsExactInDegree({\mdseries\slshape C, i})\index{IsExactInDegree@\texttt{IsExactInDegree}}
\label{IsExactInDegree}
}\hfill{\scriptsize (operation)}}\\


 Arguments: \mbox{\texttt{\mdseries\slshape C}} \texttt{\symbol{45}}\texttt{\symbol{45}} a complex, \mbox{\texttt{\mdseries\slshape i}} \texttt{\symbol{45}}\texttt{\symbol{45}} an integer.\\
 \textbf{\indent Returns:\ }
true if \mbox{\texttt{\mdseries\slshape C}} is exact at position \mbox{\texttt{\mdseries\slshape i}}. 



 }

 

\subsection{\textcolor{Chapter }{IsShortExactSequence}}
\logpage{[ 10, 5, 16 ]}\nobreak
\hyperdef{L}{X87ADD4F685457000}{}
{\noindent\textcolor{FuncColor}{$\triangleright$\enspace\texttt{IsShortExactSequence({\mdseries\slshape C})\index{IsShortExactSequence@\texttt{IsShortExactSequence}}
\label{IsShortExactSequence}
}\hfill{\scriptsize (property)}}\\


 Arguments: \mbox{\texttt{\mdseries\slshape C}} \texttt{\symbol{45}}\texttt{\symbol{45}} a complex.\\
 \textbf{\indent Returns:\ }
true if \mbox{\texttt{\mdseries\slshape C}} is exact and of the form 
\[ \ldots \rightarrow 0 \rightarrow A \rightarrow B \rightarrow C \rightarrow 0
\rightarrow \ldots \]
 This could be positioned in any degree (as opposed to the construction of a
short exact sequence, where $B$ will be put in degree zero). 



 }

 
\begin{Verbatim}[commandchars=!@|,fontsize=\small,frame=single,label=Example]
  !gapprompt@gap>| !gapinput@C;                |
  0 -> 4:(0,1) -> 3:(1,0) -> 2:(2,2) -> 1:(1,1) -> 0:(2,2) -> 0
  !gapprompt@gap>| !gapinput@IsExactSequence(C);|
  false
  !gapprompt@gap>| !gapinput@IsExactInDegree(C,1);|
  true
  !gapprompt@gap>| !gapinput@IsExactInDegree(C,2);|
  false 
\end{Verbatim}
 

\subsection{\textcolor{Chapter }{ForEveryDegree}}
\logpage{[ 10, 5, 17 ]}\nobreak
\hyperdef{L}{X849C525A7F6EBF6D}{}
{\noindent\textcolor{FuncColor}{$\triangleright$\enspace\texttt{ForEveryDegree({\mdseries\slshape C, func})\index{ForEveryDegree@\texttt{ForEveryDegree}}
\label{ForEveryDegree}
}\hfill{\scriptsize (operation)}}\\


 Arguments: \mbox{\texttt{\mdseries\slshape C}} \texttt{\symbol{45}}\texttt{\symbol{45}} a complex, \mbox{\texttt{\mdseries\slshape func}} \texttt{\symbol{45}}\texttt{\symbol{45}} a function operating on two
consecutive maps.\\
 \textbf{\indent Returns:\ }
true if \mbox{\texttt{\mdseries\slshape func}} returns true for any two consecutive differentials, fail if this can not be
decided, false otherwise. 

}

 }

 
\section{\textcolor{Chapter }{Transforming and combining complexes}}\logpage{[ 10, 6, 0 ]}
\hyperdef{L}{X8764E5C88284301B}{}
{
  

\subsection{\textcolor{Chapter }{Shift}}
\logpage{[ 10, 6, 1 ]}\nobreak
\label{Shiftsubsect}
\hyperdef{L}{X81F7B0AB84D6319B}{}
{\noindent\textcolor{FuncColor}{$\triangleright$\enspace\texttt{Shift({\mdseries\slshape C, i, or, list, shift})\index{Shift@\texttt{Shift}}
\label{Shift}
}\hfill{\scriptsize (operation)}}\\


 Arguments: (first version) \mbox{\texttt{\mdseries\slshape C}} \texttt{\symbol{45}}\texttt{\symbol{45}} a complex, \mbox{\texttt{\mdseries\slshape i}} \texttt{\symbol{45}}\texttt{\symbol{45}} an integer.\\
 (second version) \mbox{\texttt{\mdseries\slshape list}} \texttt{\symbol{45}}\texttt{\symbol{45}} an infinite list, \mbox{\texttt{\mdseries\slshape shift}} \texttt{\symbol{45}}\texttt{\symbol{45}} an integer.\\
 \textbf{\indent Returns:\ }
(first version) A new complex, which is a shift of \mbox{\texttt{\mdseries\slshape C}}. 



If \mbox{\texttt{\mdseries\slshape i}} {\textgreater} 0, the complex is shifted to the left. If \mbox{\texttt{\mdseries\slshape i}} {\textless} 0, the complex is shifted to the right. Note that shifting might
change the differentials: In the shifted complex, $d_{new}$ is defined to be $(-1)^i d_{old}$. \textbf{\indent Returns:\ }
( second version) A new infinite list which is \mbox{\texttt{\mdseries\slshape list}} with all values shifted \mbox{\texttt{\mdseries\slshape shift}} positions to the right if \mbox{\texttt{\mdseries\slshape shift}} is positive, and to the left if \mbox{\texttt{\mdseries\slshape shift}} is negative. 

}

 
\begin{Verbatim}[commandchars=@|B,fontsize=\small,frame=single,label=Example]
  @gapprompt|gap>B @gapinput|C;B
  0 -> 4:(0,1) -> 3:(1,0) -> 2:(2,2) -> 1:(1,1) -> 0:(2,2) -> 0
  @gapprompt|gap>B @gapinput|Shift(C,1);B
  0 -> 3:(0,1) -> 2:(1,0) -> 1:(2,2) -> 0:(1,1) -> -1:(2,2) -> 0
  @gapprompt|gap>B @gapinput|D := Shift(C,-1);B
  0 -> 5:(0,1) -> 4:(1,0) -> 3:(2,2) -> 2:(1,1) -> 1:(2,2) -> 0
  @gapprompt|gap>B @gapinput|dc := DifferentialOfComplex(C,3)!.maps;B
  [ [ [ 1, 0 ] ], [ [ 0, 0 ] ] ]
  @gapprompt|gap>B @gapinput|dd := DifferentialOfComplex(D,4)!.maps;B
  [ [ [ -1, 0 ] ], [ [ 0, 0 ] ] ]
  @gapprompt|gap>B @gapinput|MatricesOfPathAlgebraMatModuleHomomorphism(dc);B
  [ [ [ 1, 0 ] ], [ [ 0, 0 ] ] ]
  @gapprompt|gap>B @gapinput|MatricesOfPathAlgebraMatModuleHomomorphism(dd);B
  [ [ [ -1, 0 ] ], [ [ 0, 0 ] ] ] 
\end{Verbatim}
 

\subsection{\textcolor{Chapter }{ShiftUnsigned}}
\logpage{[ 10, 6, 2 ]}\nobreak
\hyperdef{L}{X86B56AC87A433E5B}{}
{\noindent\textcolor{FuncColor}{$\triangleright$\enspace\texttt{ShiftUnsigned({\mdseries\slshape C, i})\index{ShiftUnsigned@\texttt{ShiftUnsigned}}
\label{ShiftUnsigned}
}\hfill{\scriptsize (operation)}}\\


 Arguments: \mbox{\texttt{\mdseries\slshape C}} \texttt{\symbol{45}}\texttt{\symbol{45}} a complex, \mbox{\texttt{\mdseries\slshape i}} \texttt{\symbol{45}}\texttt{\symbol{45}} an integer.\\
 \textbf{\indent Returns:\ }
A new complex, which is a shift of \mbox{\texttt{\mdseries\slshape C}}. 



Does the same as \texttt{Shift} change the sign of the differential. Although this is a
non\texttt{\symbol{45}}mathematical definition of shift, it is still useful
for technical purposes, when manipulating and creating complexes. }

 

\subsection{\textcolor{Chapter }{YonedaProduct}}
\logpage{[ 10, 6, 3 ]}\nobreak
\hyperdef{L}{X7897A65085171913}{}
{\noindent\textcolor{FuncColor}{$\triangleright$\enspace\texttt{YonedaProduct({\mdseries\slshape C, D})\index{YonedaProduct@\texttt{YonedaProduct}}
\label{YonedaProduct}
}\hfill{\scriptsize (operation)}}\\


 Arguments: \mbox{\texttt{\mdseries\slshape C}}, \mbox{\texttt{\mdseries\slshape D}} \texttt{\symbol{45}}\texttt{\symbol{45}} complexes.\\
 \textbf{\indent Returns:\ }
The Yoneda product of the two complexes, which is a complex. 



To compute the Yoneda product, \mbox{\texttt{\mdseries\slshape C}} and \mbox{\texttt{\mdseries\slshape D}} must be such that the object in degree \texttt{LowerBound(C)} equals the object in degree \texttt{UpperBound(D)}, that is 
\[ \ldots \rightarrow C_{i+1} \rightarrow C_{i} \rightarrow A \rightarrow 0
\rightarrow \ldots \]
 
\[ \ldots \rightarrow 0 \rightarrow A \rightarrow D_{j} \rightarrow D_{j-1}
\rightarrow \ldots \]
 The product is of this form: 
\[ \ldots \rightarrow C_{i+1} \rightarrow C_{i} \rightarrow D_{j} \rightarrow
D_{j-1} \rightarrow \ldots \]
 where the map $C_{i} \rightarrow D_{j}$ is the composition of the maps $C_{i} \rightarrow A$ and $A \rightarrow D_{j}$. Also, the object $D_{j}$ is in degree $j$. }

 
\begin{Verbatim}[commandchars=!@|,fontsize=\small,frame=single,label=Example]
  !gapprompt@gap>| !gapinput@C2;|
  0 -> 4:(0,1) -> 3:(1,0) -> 2:(2,2) -> 1:(1,1) -> 0:(0,0) -> 0
  !gapprompt@gap>| !gapinput@C3;|
  0 -> -1:(1,1) -> -2:(2,2) -> -3:(1,1) -> 0
  !gapprompt@gap>| !gapinput@YonedaProduct(C2,C3);|
  0 -> 1:(0,1) -> 0:(1,0) -> -1:(2,2) -> -2:(2,2) -> -3:(1,1) -> 0 
\end{Verbatim}
 

\subsection{\textcolor{Chapter }{BrutalTruncationBelow}}
\logpage{[ 10, 6, 4 ]}\nobreak
\hyperdef{L}{X87E1CA507EB6B86E}{}
{\noindent\textcolor{FuncColor}{$\triangleright$\enspace\texttt{BrutalTruncationBelow({\mdseries\slshape C, i})\index{BrutalTruncationBelow@\texttt{BrutalTruncationBelow}}
\label{BrutalTruncationBelow}
}\hfill{\scriptsize (operation)}}\\


 Arguments: \mbox{\texttt{\mdseries\slshape C}} \texttt{\symbol{45}}\texttt{\symbol{45}} a complex, \mbox{\texttt{\mdseries\slshape i}} \texttt{\symbol{45}}\texttt{\symbol{45}} an integer.\\
 \textbf{\indent Returns:\ }
A newly created complex. 



Replace all objects with degree $j$ {\textless} $i$ with zero. The differentials affected will also become zero. }

 

\subsection{\textcolor{Chapter }{BrutalTruncationAbove}}
\logpage{[ 10, 6, 5 ]}\nobreak
\hyperdef{L}{X7CE9481684DD304E}{}
{\noindent\textcolor{FuncColor}{$\triangleright$\enspace\texttt{BrutalTruncationAbove({\mdseries\slshape C, i})\index{BrutalTruncationAbove@\texttt{BrutalTruncationAbove}}
\label{BrutalTruncationAbove}
}\hfill{\scriptsize (operation)}}\\


 Arguments: \mbox{\texttt{\mdseries\slshape C}} \texttt{\symbol{45}}\texttt{\symbol{45}} a complex, \mbox{\texttt{\mdseries\slshape i}} \texttt{\symbol{45}}\texttt{\symbol{45}} an integer.\\
 \textbf{\indent Returns:\ }
A newly created complex. 



Replace all objects with degree $j$ {\textgreater} $i$ with zero. The differentials affected will also become zero. }

 

\subsection{\textcolor{Chapter }{BrutalTruncation}}
\logpage{[ 10, 6, 6 ]}\nobreak
\hyperdef{L}{X8547C51781BBB48F}{}
{\noindent\textcolor{FuncColor}{$\triangleright$\enspace\texttt{BrutalTruncation({\mdseries\slshape C, i, j})\index{BrutalTruncation@\texttt{BrutalTruncation}}
\label{BrutalTruncation}
}\hfill{\scriptsize (operation)}}\\


 Arguments: \mbox{\texttt{\mdseries\slshape C}} \texttt{\symbol{45}}\texttt{\symbol{45}} a complex, \mbox{\texttt{\mdseries\slshape i, j}} \texttt{\symbol{45}}\texttt{\symbol{45}} integers.\\
 \textbf{\indent Returns:\ }
A newly created complex. 



Brutally truncates in both ends. The integer arguments must be ordered such
that \mbox{\texttt{\mdseries\slshape i}} {\textgreater} \mbox{\texttt{\mdseries\slshape j}}. }

 

\subsection{\textcolor{Chapter }{SyzygyTruncation}}
\logpage{[ 10, 6, 7 ]}\nobreak
\hyperdef{L}{X7E3E34167B79248A}{}
{\noindent\textcolor{FuncColor}{$\triangleright$\enspace\texttt{SyzygyTruncation({\mdseries\slshape C, i})\index{SyzygyTruncation@\texttt{SyzygyTruncation}}
\label{SyzygyTruncation}
}\hfill{\scriptsize (operation)}}\\


 Arguments: \mbox{\texttt{\mdseries\slshape C}} \texttt{\symbol{45}}\texttt{\symbol{45}} a complex, \mbox{\texttt{\mdseries\slshape i}} \texttt{\symbol{45}}\texttt{\symbol{45}} an integer.\\
 \textbf{\indent Returns:\ }
A newly created complex. 



Replace the object in degree $i$ with the kernel of $d_i$, and $d_{i+1}$ with the natural inclusion. All objects in degree $j > i+1$ are replaced with zero. }

 

\subsection{\textcolor{Chapter }{CosyzygyTruncation}}
\logpage{[ 10, 6, 8 ]}\nobreak
\hyperdef{L}{X7C729A2587E5DA8D}{}
{\noindent\textcolor{FuncColor}{$\triangleright$\enspace\texttt{CosyzygyTruncation({\mdseries\slshape C, i})\index{CosyzygyTruncation@\texttt{CosyzygyTruncation}}
\label{CosyzygyTruncation}
}\hfill{\scriptsize (operation)}}\\


 Arguments: \mbox{\texttt{\mdseries\slshape C}} \texttt{\symbol{45}}\texttt{\symbol{45}} a complex, \mbox{\texttt{\mdseries\slshape i}} \texttt{\symbol{45}}\texttt{\symbol{45}} an integer.\\
 \textbf{\indent Returns:\ }
A newly created complex. 



Replace the object in degree $i-2$ with the cokernel of $d_i$, and $d_{i-1}$ with the natural projection. All objects in degree $j < i-2$ are replaced with zero. }

 

\subsection{\textcolor{Chapter }{SyzygyCosyzygyTruncation}}
\logpage{[ 10, 6, 9 ]}\nobreak
\hyperdef{L}{X797FC6037F6605E7}{}
{\noindent\textcolor{FuncColor}{$\triangleright$\enspace\texttt{SyzygyCosyzygyTruncation({\mdseries\slshape C, i, j})\index{SyzygyCosyzygyTruncation@\texttt{SyzygyCosyzygyTruncation}}
\label{SyzygyCosyzygyTruncation}
}\hfill{\scriptsize (operation)}}\\


 Arguments: \mbox{\texttt{\mdseries\slshape C}} \texttt{\symbol{45}}\texttt{\symbol{45}} a complex, \mbox{\texttt{\mdseries\slshape i}} \texttt{\symbol{45}}\texttt{\symbol{45}} an integer.\\
 \textbf{\indent Returns:\ }
A newly created complex. 



Performs both the above truncations. The integer arguments must be ordered
such that \mbox{\texttt{\mdseries\slshape i}} {\textgreater} \mbox{\texttt{\mdseries\slshape j}}. }

 }

 
\section{\textcolor{Chapter }{Chain maps}}\logpage{[ 10, 7, 0 ]}
\hyperdef{L}{X85F418EB859E7597}{}
{
  An \texttt{IsChainMap} (\ref{IsChainMap}) object represents a chain map between two complexes over the same category. 

\subsection{\textcolor{Chapter }{IsChainMap}}
\logpage{[ 10, 7, 1 ]}\nobreak
\hyperdef{L}{X7F76E62D837F0E66}{}
{\noindent\textcolor{FuncColor}{$\triangleright$\enspace\texttt{IsChainMap\index{IsChainMap@\texttt{IsChainMap}}
\label{IsChainMap}
}\hfill{\scriptsize (Category)}}\\


 The category for chain maps. }

 

\subsection{\textcolor{Chapter }{ChainMap}}
\logpage{[ 10, 7, 2 ]}\nobreak
\hyperdef{L}{X7BCD94877DF261C4}{}
{\noindent\textcolor{FuncColor}{$\triangleright$\enspace\texttt{ChainMap({\mdseries\slshape source, range, basePosition, middle, positive, negative})\index{ChainMap@\texttt{ChainMap}}
\label{ChainMap}
}\hfill{\scriptsize (function)}}\\


 Arguments: \mbox{\texttt{\mdseries\slshape source}}, \mbox{\texttt{\mdseries\slshape range}} \texttt{\symbol{45}}\texttt{\symbol{45}} complexes, \mbox{\texttt{\mdseries\slshape basePosition}} \texttt{\symbol{45}}\texttt{\symbol{45}} an integer, \mbox{\texttt{\mdseries\slshape middle}} \texttt{\symbol{45}}\texttt{\symbol{45}} a list of morphisms, \mbox{\texttt{\mdseries\slshape positive}} \texttt{\symbol{45}}\texttt{\symbol{45}} a list or the string \texttt{"zero"}, \mbox{\texttt{\mdseries\slshape negative}} \texttt{\symbol{45}}\texttt{\symbol{45}} a list or the string \texttt{"zero"}. \\
 \textbf{\indent Returns:\ }
A newly created chain map



 The arguments \mbox{\texttt{\mdseries\slshape source}} and \mbox{\texttt{\mdseries\slshape range}} are the complexes which the new chain map should map between.

 The rest of the arguments describe the individual morphisms which constitute
the chain map, in a similar way to the last four arguments to the \texttt{Complex} (\ref{Complex}) function.

 The morphisms of the chain map are divided into three parts: one finite (``middle'') and two infinite (``positive'' and ``negative''). The positive part contains all morphisms in degrees higher than those in
the middle part, and the negative part contains all morphisms in degrees lower
than those in the middle part. (The middle part may be placed anywhere, so the
positive part can \texttt{\symbol{45}}\texttt{\symbol{45}} despite its name
\texttt{\symbol{45}}\texttt{\symbol{45}} contain some morphisms of negative
degree. Conversely, the negative part can contain some morphisms of positive
degree.)

 The argument \mbox{\texttt{\mdseries\slshape middle}} is a list containing the morphisms for the middle part. The argument \mbox{\texttt{\mdseries\slshape baseDegree}} gives the degree of the first morphism in this list. The second morphism is
placed in degree $\mbox{\texttt{\mdseries\slshape baseDegree}}+1$, and so on. Thus, the middle part consists of the degrees 
\[ \mbox{\texttt{\mdseries\slshape baseDegree}},\quad \mbox{\texttt{\mdseries\slshape baseDegree}} + 1,\quad \ldots\quad \mbox{\texttt{\mdseries\slshape baseDegree}} + \text{Length}(\mbox{\texttt{\mdseries\slshape middle}}) - 1. \]
 Each of the arguments \mbox{\texttt{\mdseries\slshape positive}} and \mbox{\texttt{\mdseries\slshape negative}} can be one of the following: 
\begin{itemize}
\item The string \texttt{"zero"}, meaning that the part contains only zero morphisms.
\item A list of the form \texttt{[ "repeat", L ]}, where \texttt{L} is a list of morphisms. The part will contain the morphisms in \texttt{L} repeated infinitely many times. The convention for the order of elements in \texttt{L} is that \texttt{L[1]} is the morphism which is closest to the middle part, and \texttt{L[Length(L)]} is farthest away from the middle part. (Using this only makes sense if the
objects of both the source and range complex repeat in a compatible way.)
\item A list of the form \texttt{[ "pos", f ]} or \texttt{[ "pos", f, store ]}, where \texttt{f} is a function of two arguments, and \texttt{store} (if included) is a boolean. The function \texttt{f} is used to compute the morphisms in this part. The function \texttt{f} is not called immediately by the \texttt{ChainMap} constructor, but will be called later as the morphisms in this part are
needed. The function call \texttt{f(M,i)} (where \texttt{M} is the chain map and \texttt{i} an integer) should produce the morphism in degree \texttt{i}. The function may use \texttt{M} to look up other morphisms in the chain map (and to access the source and
range complexes), as long as this does not cause an infinite loop. If \texttt{store} is \texttt{true} (or not specified), each computed morphism is stored, and they are computed in
order from the one closest to the middle part, regardless of which order they
are requested in.
\item A list of the form \texttt{[ "next", f, init ]}, where \texttt{f} is a function of one argument, and \texttt{init} is a morphism. The function \texttt{f} is used to compute the morphisms in this part. For the first morphism in the
part (that is, the one closest to the middle part), \texttt{f} is called with \texttt{init} as argument. For the next morphism, \texttt{f} is called with the first morphism as argument, and so on. Thus, the morphisms
are 
\[ f(\text{init}),\quad f^2(\text{init}),\quad f^3(\text{init}),\quad \ldots \]
 Each morphism is stored when it has been computed. 
\end{itemize}
 }

 

\subsection{\textcolor{Chapter }{ZeroChainMap}}
\logpage{[ 10, 7, 3 ]}\nobreak
\hyperdef{L}{X7C9597F6810F195E}{}
{\noindent\textcolor{FuncColor}{$\triangleright$\enspace\texttt{ZeroChainMap({\mdseries\slshape source, range})\index{ZeroChainMap@\texttt{ZeroChainMap}}
\label{ZeroChainMap}
}\hfill{\scriptsize (function)}}\\
\textbf{\indent Returns:\ }
A newly created zero chain map



 This function creates a zero chain map (a chain map in which every morphism is
zero) from the complex \mbox{\texttt{\mdseries\slshape source}} to the complex \mbox{\texttt{\mdseries\slshape range}}.

 }

 

\subsection{\textcolor{Chapter }{FiniteChainMap}}
\logpage{[ 10, 7, 4 ]}\nobreak
\hyperdef{L}{X8699B48A783E1900}{}
{\noindent\textcolor{FuncColor}{$\triangleright$\enspace\texttt{FiniteChainMap({\mdseries\slshape source, range, baseDegree, morphisms})\index{FiniteChainMap@\texttt{FiniteChainMap}}
\label{FiniteChainMap}
}\hfill{\scriptsize (function)}}\\
\textbf{\indent Returns:\ }
A newly created chain map



 This function creates a complex where all but finitely many morphisms are
zero.

 The arguments \mbox{\texttt{\mdseries\slshape source}} and \mbox{\texttt{\mdseries\slshape range}} are the complexes which the new chain map should map between.

 The argument \mbox{\texttt{\mdseries\slshape morphisms}} is a list of morphisms. The argument \mbox{\texttt{\mdseries\slshape baseDegree}} gives the degree for the first morphism in this list. The subsequent morphisms
are placed in degrees $\mbox{\texttt{\mdseries\slshape baseDegree}}+1$, and so on.

 This means that the \mbox{\texttt{\mdseries\slshape morphisms}} argument specifies the morphisms in degrees 
\[ \mbox{\texttt{\mdseries\slshape baseDegree}},\quad \mbox{\texttt{\mdseries\slshape baseDegree}} + 1,\quad \ldots \quad \mbox{\texttt{\mdseries\slshape baseDegree}} + \text{Length}(\mbox{\texttt{\mdseries\slshape morphisms}}) - 1. \]
 All other morphisms in the chain map are zero.

 }

 

\subsection{\textcolor{Chapter }{ComplexAndChainMaps}}
\logpage{[ 10, 7, 5 ]}\nobreak
\hyperdef{L}{X805BA1698394DE4C}{}
{\noindent\textcolor{FuncColor}{$\triangleright$\enspace\texttt{ComplexAndChainMaps({\mdseries\slshape sourceComplexes, rangeComplexes, basePosition, middle, positive, negative})\index{ComplexAndChainMaps@\texttt{ComplexAndChainMaps}}
\label{ComplexAndChainMaps}
}\hfill{\scriptsize (function)}}\\


 Arguments: \mbox{\texttt{\mdseries\slshape sourceComplexes}} \texttt{\symbol{45}}\texttt{\symbol{45}} a list of complexes, \mbox{\texttt{\mdseries\slshape rangeComplexes}} \texttt{\symbol{45}}\texttt{\symbol{45}} a list of complexes, \mbox{\texttt{\mdseries\slshape basePosition}} \texttt{\symbol{45}}\texttt{\symbol{45}} an integer, \mbox{\texttt{\mdseries\slshape middle}} \texttt{\symbol{45}}\texttt{\symbol{45}} a list of morphisms, \mbox{\texttt{\mdseries\slshape positive}} \texttt{\symbol{45}}\texttt{\symbol{45}} a list or the string \texttt{"zero"}, \mbox{\texttt{\mdseries\slshape negative}} \texttt{\symbol{45}}\texttt{\symbol{45}} a list or the string \texttt{"zero"}. \\
 \textbf{\indent Returns:\ }
A list consisting of a newly created complex, and one or more newly created
chain maps.



 This is a combined constructor to make one complex and a set of chain maps at
the same time. All the chain maps will have the new complex as either source
or range.

 The argument \mbox{\texttt{\mdseries\slshape sourceComplexes}} is a list of the complexes to be sources of the chain maps which have the new
complex as range. The argument \mbox{\texttt{\mdseries\slshape rangeComplexes}} is a list of the complexes to be ranges of the chain maps which have the new
complex as source.

 Let $S$ and $R$ stand for the lengths of the lists \mbox{\texttt{\mdseries\slshape sourceComplexes}} and \mbox{\texttt{\mdseries\slshape rangeComplexes}}, respectively. Then the number of new chain maps which are created is $S+R$.

 The last four arguments describe the individual differentials of the new
complex, as well as the individual morphisms which constitute each of the new
chain maps. These arguments are treated in a similar way to the last four
arguments to the \texttt{Complex} (\ref{Complex}) and \texttt{ChainMap} (\ref{ChainMap}) constructors. In those constructors, the last four arguments describe, for
each degree, how to get the differential or morphism for that degree. Here, we
for each degree need both a differential for the complex, and one morphism for
each chain map. So for each degree $i$, we will have a list 
\[ L_i = [ d_i, m_i^1, \ldots, m_i^S, n_i^1, \ldots, n_i^R ], \]
 where $d_i$ is the differential for the new complex in degree $i$, $m_i^j$ is the morphism in degree $i$ of the chain map from \texttt{sourceComplexes[j]} to the new complex, and $n_i^j$ is the morphism in degree $i$ of the chain map from the new complex to \texttt{rangeComplexes[j]}.

 The degrees of the new complex and chain maps are divided into three parts:
one finite (``middle'') and two infinite (``positive'' and ``negative''). The positive part contains all degrees higher than those in the middle
part, and the negative part contains all degrees lower than those in the
middle part.

 The argument \mbox{\texttt{\mdseries\slshape middle}} is a list containing the lists $L_i$ for the middle part. The argument \mbox{\texttt{\mdseries\slshape baseDegree}} gives the degree of the first morphism in this list. The second morphism is
placed in degree $\mbox{\texttt{\mdseries\slshape baseDegree}}+1$, and so on. Thus, the middle part consists of the degrees 
\[ \mbox{\texttt{\mdseries\slshape baseDegree}},\quad \mbox{\texttt{\mdseries\slshape baseDegree}} + 1,\quad \ldots\quad \mbox{\texttt{\mdseries\slshape baseDegree}} + \text{Length}(\mbox{\texttt{\mdseries\slshape middle}}) - 1. \]
 Each of the arguments \mbox{\texttt{\mdseries\slshape positive}} and \mbox{\texttt{\mdseries\slshape negative}} can be one of the following: 
\begin{itemize}
\item The string \texttt{"zero"}, meaning that the part contains only zero morphisms.
\item A list of the form \texttt{[ "repeat", L ]}, where \texttt{L} is a list of morphisms. The part will contain the morphisms in \texttt{L} repeated infinitely many times. The convention for the order of elements in \texttt{L} is that \texttt{L[1]} is the morphism which is closest to the middle part, and \texttt{L[Length(L)]} is farthest away from the middle part. (Using this only makes sense if the
objects of both the source and range complex repeat in a compatible way.)
\item A list of the form \texttt{[ "pos", f ]} or \texttt{[ "pos", f, store ]}, where \texttt{f} is a function of two arguments, and \texttt{store} (if included) is a boolean. The function \texttt{f} is used to compute the morphisms in this part. The function \texttt{f} is not called immediately by the \texttt{ChainMap} constructor, but will be called later as the morphisms in this part are
needed. The function call \texttt{f(M,i)} (where \texttt{M} is the chain map and \texttt{i} an integer) should produce the morphism in degree \texttt{i}. The function may use \texttt{M} to look up other morphisms in the chain map (and to access the source and
range complexes), as long as this does not cause an infinite loop. If \texttt{store} is \texttt{true} (or not specified), each computed morphism is stored, and they are computed in
order from the one closest to the middle part, regardless of which order they
are requested in.
\item A list of the form \texttt{[ "next", f, init ]}, where \texttt{f} is a function of one argument, and \texttt{init} is a morphism. The function \texttt{f} is used to compute the morphisms in this part. For the first morphism in the
part (that is, the one closest to the middle part), \texttt{f} is called with \texttt{init} as argument. For the next morphism, \texttt{f} is called with the first morphism as argument, and so on. Thus, the morphisms
are 
\[ f(\text{init}),\quad f^2(\text{init}),\quad f^3(\text{init}),\quad \ldots \]
 Each morphism is stored when it has been computed. 
\end{itemize}
 The return value of the \texttt{ComplexAndChainMaps} constructor is a list 
\[ [ C, M_1, \ldots, M_S, N_1, \ldots, N_R ], \]
 where $C$ is the new complex, $M_1,\ldots,M_S$ are the new chain maps with $C$ as range, and $N_1,\ldots,N_R$ are the new chain maps with $C$ as source. }

 

\subsection{\textcolor{Chapter }{MorphismOfChainMap}}
\logpage{[ 10, 7, 6 ]}\nobreak
\hyperdef{L}{X7F20B8807ECE404E}{}
{\noindent\textcolor{FuncColor}{$\triangleright$\enspace\texttt{MorphismOfChainMap({\mdseries\slshape M, i})\index{MorphismOfChainMap@\texttt{MorphismOfChainMap}}
\label{MorphismOfChainMap}
}\hfill{\scriptsize (operation)}}\\


 Arguments: \mbox{\texttt{\mdseries\slshape M}} \texttt{\symbol{45}}\texttt{\symbol{45}} a chain map, \mbox{\texttt{\mdseries\slshape i}} \texttt{\symbol{45}}\texttt{\symbol{45}} an integer.\\
 \textbf{\indent Returns:\ }
 The morphism at position \mbox{\texttt{\mdseries\slshape i}} in the chain map. 



 }

 

\subsection{\textcolor{Chapter }{MorphismsOfChainMap}}
\logpage{[ 10, 7, 7 ]}\nobreak
\hyperdef{L}{X792B5B8684A75F72}{}
{\noindent\textcolor{FuncColor}{$\triangleright$\enspace\texttt{MorphismsOfChainMap({\mdseries\slshape M})\index{MorphismsOfChainMap@\texttt{MorphismsOfChainMap}}
\label{MorphismsOfChainMap}
}\hfill{\scriptsize (attribute)}}\\


 Arguments: \mbox{\texttt{\mdseries\slshape M}} \texttt{\symbol{45}}\texttt{\symbol{45}} a chain map.\\
 \textbf{\indent Returns:\ }
 The morphisms of the chain map, stored as an \texttt{IsInfList} (\ref{IsInfList}) object. 

}

 

\subsection{\textcolor{Chapter }{ComparisonLifting}}
\logpage{[ 10, 7, 8 ]}\nobreak
\hyperdef{L}{X78EC487B79B20DDE}{}
{\noindent\textcolor{FuncColor}{$\triangleright$\enspace\texttt{ComparisonLifting({\mdseries\slshape f, PC, EC})\index{ComparisonLifting@\texttt{ComparisonLifting}}
\label{ComparisonLifting}
}\hfill{\scriptsize (operation)}}\\


 Arguments: \mbox{\texttt{\mdseries\slshape f}} \texttt{\symbol{45}}\texttt{\symbol{45}} a map between modules $M$ and $N$, \mbox{\texttt{\mdseries\slshape PC}} \texttt{\symbol{45}}\texttt{\symbol{45}} a chain complex, \mbox{\texttt{\mdseries\slshape EC}} \texttt{\symbol{45}}\texttt{\symbol{45}} a chain complex.\\
 \textbf{\indent Returns:\ }
 The map \mbox{\texttt{\mdseries\slshape f}} lifted to a chain map from \mbox{\texttt{\mdseries\slshape PC}} to \mbox{\texttt{\mdseries\slshape EC}}. 



 The complex \mbox{\texttt{\mdseries\slshape PC}} must have $M$ in some fixed degree $i$, it should be bounded with only zero objects in degrees smaller than $i$, and it should have only projective objects in degrees greater than $i$ (or projective objects in degrees $[i+1,j]$ and zero in degrees greater than $j$). The complex \mbox{\texttt{\mdseries\slshape EC}} should also have zero in degrees smaller than $i$, it should have $N$ in degree $i$ and it should be exact for all degrees. The returned chain map has \mbox{\texttt{\mdseries\slshape f}} in degree $i$. }

 

\subsection{\textcolor{Chapter }{ComparisonLiftingToProjectiveResolution}}
\logpage{[ 10, 7, 9 ]}\nobreak
\hyperdef{L}{X7A95C5E78298E0C6}{}
{\noindent\textcolor{FuncColor}{$\triangleright$\enspace\texttt{ComparisonLiftingToProjectiveResolution({\mdseries\slshape f})\index{ComparisonLiftingToProjectiveResolution@\texttt{Comparison}\-\texttt{Lifting}\-\texttt{To}\-\texttt{Projective}\-\texttt{Resolution}}
\label{ComparisonLiftingToProjectiveResolution}
}\hfill{\scriptsize (operation)}}\\


 Arguments: \mbox{\texttt{\mdseries\slshape f}} \texttt{\symbol{45}}\texttt{\symbol{45}} a map between modules $M$ and $N$.\\
 \textbf{\indent Returns:\ }
 The map \mbox{\texttt{\mdseries\slshape f}} lifted to a chain map from the projective resolution of $M$ to the projective resolution of $N$. 



 The returned chain map has \mbox{\texttt{\mdseries\slshape f}} in degree $-1$ (the projective resolution of a module includes the module itself in degree $-1$). }

 

\subsection{\textcolor{Chapter }{MappingCone}}
\logpage{[ 10, 7, 10 ]}\nobreak
\hyperdef{L}{X7EF42CF48165D831}{}
{\noindent\textcolor{FuncColor}{$\triangleright$\enspace\texttt{MappingCone({\mdseries\slshape f})\index{MappingCone@\texttt{MappingCone}}
\label{MappingCone}
}\hfill{\scriptsize (operation)}}\\


 Arguments: \mbox{\texttt{\mdseries\slshape f}} \texttt{\symbol{45}}\texttt{\symbol{45}} a chain map between chain complexes $A$ and $B$. \\
 \textbf{\indent Returns:\ }
 A list with the mapping cone of \mbox{\texttt{\mdseries\slshape f}} and the inclusion of $B$ into the cone, and the projection of the cone onto $A[-1]$. 



 }

 
\begin{Verbatim}[commandchars=!@|,fontsize=\small,frame=single,label=Example]
  !gapprompt@gap>| !gapinput@# Constructs a quiver and a quotient of a path algebra|
  !gapprompt@gap>| !gapinput@Q := Quiver( 4, [ [1,2,"a"], [2,3,"b"], [3,4,"c"] ] );;|
  !gapprompt@gap>| !gapinput@PA := PathAlgebra( Rationals, Q );;|
  !gapprompt@gap>| !gapinput@rels := [ PA.a*PA.b ];;|
  !gapprompt@gap>| !gapinput@gb := GBNPGroebnerBasis( rels, PA );;|
  !gapprompt@gap>| !gapinput@I := Ideal( PA, gb );;|
  !gapprompt@gap>| !gapinput@grb := GroebnerBasis( I, gb );;|
  !gapprompt@gap>| !gapinput@alg := PA/I;;|
  !gapprompt@gap>| !gapinput@|
  !gapprompt@gap>| !gapinput@# Two modules M and N, and a map between them|
  !gapprompt@gap>| !gapinput@M := RightModuleOverPathAlgebra( alg, [0,1,1,0], [["b", [[1]] ]] );;|
  !gapprompt@gap>| !gapinput@N := RightModuleOverPathAlgebra( alg, [0,1,0,0], [] );;|
  !gapprompt@gap>| !gapinput@f := RightModuleHomOverAlgebra(M, N, [ [[0]],[[1]],[[0]],[[0]] ]);;|
  !gapprompt@gap>| !gapinput@|
  !gapprompt@gap>| !gapinput@# Lifts f to a map between the projective resolutions of M and N|
  !gapprompt@gap>| !gapinput@lf := ComparisonLiftingToProjectiveResolution(f);|
  <chain map>
  !gapprompt@gap>| !gapinput@|
  !gapprompt@gap>| !gapinput@# Computes the mapping cone of the chain map|
  !gapprompt@gap>| !gapinput@H := MappingCone(lf);|
  [ --- -> -1:(0,1,0,0) -> ---, <chain map>, <chain map> ]
  !gapprompt@gap>| !gapinput@cone := H[1];|
  --- -> -1:(0,1,0,0) -> ---
  !gapprompt@gap>| !gapinput@ObjectOfComplex(Source(lf),0);|
  <[ 0, 1, 1, 1 ]>
  !gapprompt@gap>| !gapinput@ObjectOfComplex(Range(lf),1);|
  <[ 0, 0, 1, 1 ]>
  !gapprompt@gap>| !gapinput@ObjectOfComplex(cone,1);|
  <[ 0, 1, 2, 2 ]> 
  !gapprompt@gap>| !gapinput@Source(H[2]) = Range(lf);|
  true
\end{Verbatim}
 }

 }

 
\chapter{\textcolor{Chapter }{Projective resolutions and the bounded derived category}}\logpage{[ 11, 0, 0 ]}
\hyperdef{L}{X7EDB390B82D9E644}{}
{
  What is implemented so far for working with the bounded derived category $\mathcal{D}^{b}( \modc A )$. We use the isomorphism $\mathcal{D}^{b}( \modc A ) \cong \mathcal{K}^{-,b}(\proj A)$, and will hence need a way to describe complexes where all objectives are
projective (or, dually, injective). 
\section{\textcolor{Chapter }{Projective and injective complexes}}\logpage{[ 11, 1, 0 ]}
\hyperdef{L}{X7C2FFD3E7D0D5D7F}{}
{
  

\subsection{\textcolor{Chapter }{InjectiveResolution}}
\logpage{[ 11, 1, 1 ]}\nobreak
\hyperdef{L}{X854F304E7A7F6FF2}{}
{\noindent\textcolor{FuncColor}{$\triangleright$\enspace\texttt{InjectiveResolution({\mdseries\slshape M})\index{InjectiveResolution@\texttt{InjectiveResolution}}
\label{InjectiveResolution}
}\hfill{\scriptsize (operation)}}\\


 Arguments: \mbox{\texttt{\mdseries\slshape M}} \texttt{\symbol{45}}\texttt{\symbol{45}} a module.\\
 \textbf{\indent Returns:\ }
The injective resolution of \mbox{\texttt{\mdseries\slshape M}} with \mbox{\texttt{\mdseries\slshape M}} in degree $-1$. 



 }

 

\subsection{\textcolor{Chapter }{IsProjectiveComplex}}
\logpage{[ 11, 1, 2 ]}\nobreak
\hyperdef{L}{X8581C1F37EA425DE}{}
{\noindent\textcolor{FuncColor}{$\triangleright$\enspace\texttt{IsProjectiveComplex({\mdseries\slshape C})\index{IsProjectiveComplex@\texttt{IsProjectiveComplex}}
\label{IsProjectiveComplex}
}\hfill{\scriptsize (property)}}\\


 Arguments: \mbox{\texttt{\mdseries\slshape C}} \texttt{\symbol{45}}\texttt{\symbol{45}} a complex.\\
 \textbf{\indent Returns:\ }
true if \mbox{\texttt{\mdseries\slshape C}} is either a finite complex of projectives or an infinite complex of
projectives constructed as a projective resolution (\texttt{ProjectiveResolutionOfComplex} (\ref{ProjectiveResolutionOfComplex})), false otherwise. 



A complex for which this property is true, will be printed in a different
manner than ordinary complexes. Instead of writing the dimension vector of the
objects in each degree, they are listed as \texttt{n1P1 + n2P2 + . . . + ntPt}, where \texttt{ni} is the number of copies of each indecomposable projective module $P_i$ corresponding to vertex $i$ of the quiver). Note that if a complex is both projective and injective, it is
printed as a projective complex. }

 

\subsection{\textcolor{Chapter }{IsInjectiveComplex}}
\logpage{[ 11, 1, 3 ]}\nobreak
\hyperdef{L}{X7BDFFF987E99A5EA}{}
{\noindent\textcolor{FuncColor}{$\triangleright$\enspace\texttt{IsInjectiveComplex({\mdseries\slshape C})\index{IsInjectiveComplex@\texttt{IsInjectiveComplex}}
\label{IsInjectiveComplex}
}\hfill{\scriptsize (property)}}\\


 Arguments: \mbox{\texttt{\mdseries\slshape C}} \texttt{\symbol{45}}\texttt{\symbol{45}} a complex.\\
 \textbf{\indent Returns:\ }
true if \mbox{\texttt{\mdseries\slshape C}} is either a finite complex of injectives or an infinite complex of injectives
constructed as $D\mathrm{Hom}_{A}(-,A)$ of a projective complex (\texttt{ProjectiveToInjectiveComplex} (\ref{ProjectiveToInjectiveComplex})), false otherwise. 



A complex for which this property is true, will be printed in a different
manner than ordinary complexes. Instead of writing the dimension vector of the
objects in each degree, they are listed as \texttt{n1I1 + n2I2 + . . . + ntIt}, where \texttt{ni} is the number of copies of each indecomposable projective module $I_i$ corresponding to vertex $i$ of the quiver). Note that if a complex is both projective and injective, it is
printed as a projective complex. }

 

\subsection{\textcolor{Chapter }{ProjectiveResolution}}
\logpage{[ 11, 1, 4 ]}\nobreak
\hyperdef{L}{X8273999C7B352F22}{}
{\noindent\textcolor{FuncColor}{$\triangleright$\enspace\texttt{ProjectiveResolution({\mdseries\slshape M})\index{ProjectiveResolution@\texttt{ProjectiveResolution}}
\label{ProjectiveResolution}
}\hfill{\scriptsize (operation)}}\\


 Arguments: \mbox{\texttt{\mdseries\slshape M}} \texttt{\symbol{45}}\texttt{\symbol{45}} a module.\\
 \textbf{\indent Returns:\ }
The projective resolution of \mbox{\texttt{\mdseries\slshape M}} with \mbox{\texttt{\mdseries\slshape M}} in degree $-1$. 



 }

 }

 
\section{\textcolor{Chapter }{The bounded derived category}}\logpage{[ 11, 2, 0 ]}
\hyperdef{L}{X83D4593C80C2C4F6}{}
{
  Let $\mathcal{D}^{b}( \modc A )$ denote the bounded derived category. If $C$ is an element of $\mathcal{D}^{b}( \modc A )$, that is, a bounded complex of $A$\texttt{\symbol{45}}modules, there exists a projective resolution $P$ of $C$ which is a complex of projective $A$\texttt{\symbol{45}}modules quasi\texttt{\symbol{45}}isomorphic to $C$. Moreover, there exists such a $P$ with the following properties: 
\begin{itemize}
\item $P$ is minimal (in the homotopy category).
\item $C$ is bounded, so $C_i = 0$ for $i < k$ for a lower bound $k$ and $C_i = 0$ for $i > j$ for an upper bound $j$. Then $P_i = 0$ for $i < k$, and $P$ is exact in degree $i$ for $i > j$.
\end{itemize}
 The function \texttt{ProjectiveResolutionOfComplex} computes such a projective resolution of any bounded complex. If $A$ has finite global dimension, then $\mathcal{D}^{b}( \modc A )$ has AR\texttt{\symbol{45}}triangles, and there exists an algorithm for
computing the AR\texttt{\symbol{45}}translation of a complex $C \in \mathcal{D}^{b}( \modc A )$: 
\begin{itemize}
\item Compute a projective resolution $P'$ of $C$.
\item Shift $P'$ one degree to the right.
\item Compute $I = D\mathrm{Hom}_{A}(P',A)$ to get a complex of injectives.
\item Compute a projective resolution $P$ of $I$.
\end{itemize}
 Then $P$ is the AR\texttt{\symbol{45}}translation of $C$, sometimes written $\tau(C)$. A projective complex $P$ is displayed/printed such that each indecomposable projective module occurring
in the complex is displayed as $Pi$ for the corresponding vertex $i$ in the quiver of the algebra for the acting path algebra. The index $i$ is given by the ording of the vertices in the quiver. A similar scheme is
followed when displaying/printing an injective complex. The following
documents the \textsf{QPA} functions for working with complexes in the derived category. 

\subsection{\textcolor{Chapter }{ProjectiveResolutionOfComplex}}
\logpage{[ 11, 2, 1 ]}\nobreak
\hyperdef{L}{X7BC3EE977FA24151}{}
{\noindent\textcolor{FuncColor}{$\triangleright$\enspace\texttt{ProjectiveResolutionOfComplex({\mdseries\slshape C})\index{ProjectiveResolutionOfComplex@\texttt{ProjectiveResolutionOfComplex}}
\label{ProjectiveResolutionOfComplex}
}\hfill{\scriptsize (operation)}}\\


 Arguments: \mbox{\texttt{\mdseries\slshape C}} \texttt{\symbol{45}}\texttt{\symbol{45}} a finite complex.\\
 \textbf{\indent Returns:\ }
A projective complex $P$ which is the projective resolution of $C$, as described in the introduction to this section. 



If the algebra has infinite global dimension, the projective resolution of $C$ could possibly be infinite. }

 

\subsection{\textcolor{Chapter }{ProjectiveToInjectiveComplex}}
\logpage{[ 11, 2, 2 ]}\nobreak
\hyperdef{L}{X84F050D07B19ABC4}{}
{\noindent\textcolor{FuncColor}{$\triangleright$\enspace\texttt{ProjectiveToInjectiveComplex({\mdseries\slshape P})\index{ProjectiveToInjectiveComplex@\texttt{ProjectiveToInjectiveComplex}}
\label{ProjectiveToInjectiveComplex}
}\hfill{\scriptsize (operation)}}\\
\noindent\textcolor{FuncColor}{$\triangleright$\enspace\texttt{ProjectiveToInjectiveFiniteComplex({\mdseries\slshape P})\index{ProjectiveToInjectiveFiniteComplex@\texttt{ProjectiveToInjectiveFiniteComplex}}
\label{ProjectiveToInjectiveFiniteComplex}
}\hfill{\scriptsize (operation)}}\\


 Arguments: \mbox{\texttt{\mdseries\slshape P}} \texttt{\symbol{45}}\texttt{\symbol{45}} a bounded below projective complex.\\
 \textbf{\indent Returns:\ }
An injective complex $I = D\mathrm{Hom}_{A}(P,A)$. 



$P$ and $I$ will always have the same length. Especially, if $P$ is unbounded above, then so is $I$. If $P$ is a finite complex (that is; \texttt{LengthOfComplex(P)} is an integer) then the simpler method \texttt{ProjectiveToInjectiveFiniteComplex} is used. }

 

\subsection{\textcolor{Chapter }{TauOfComplex}}
\logpage{[ 11, 2, 3 ]}\nobreak
\hyperdef{L}{X82CB35E07DD2D96D}{}
{\noindent\textcolor{FuncColor}{$\triangleright$\enspace\texttt{TauOfComplex({\mdseries\slshape C})\index{TauOfComplex@\texttt{TauOfComplex}}
\label{TauOfComplex}
}\hfill{\scriptsize (operation)}}\\


 Arguments: \mbox{\texttt{\mdseries\slshape C}} \texttt{\symbol{45}}\texttt{\symbol{45}} a finite complex over an algebra of
finite global dimension.\\
 \textbf{\indent Returns:\ }
A projective complex $P$ which is the AR\texttt{\symbol{45}}translation of \mbox{\texttt{\mdseries\slshape C}}. 



This function only works when the algebra has finite global dimension. It will
always assume that both the projective resolutions computed are finite. }

 
\subsection{\textcolor{Chapter }{Example}}\logpage{[ 11, 2, 4 ]}
\hyperdef{L}{X85861B017AEEC50B}{}
{
  The following example illustrates the above mentioned functions and
properties. Note that both \texttt{ProjectiveResolutionOfComplex} and \texttt{ProjectiveToInjectiveComplex} return complexes with a nonzero \emph{positive} part, whereas \texttt{TauOfComplex} always returns a complex for which \texttt{IsFiniteComplex} returns true. Also note that after the complex \texttt{C} in the example is found to have the \texttt{IsInjectiveComplex} property, the printing of the complex changes. 

 The algebra in the example is $kQ/I$, where $Q$ is the quiver $1 \longrightarrow 2 \longrightarrow 3$ and $I$ is generated by the composition of the arrows. We construct $C$ as the stalk complex with the injective $I_1$ in degree 0. 
\begin{Verbatim}[commandchars=!@|,fontsize=\small,frame=single,label=Example]
  !gapprompt@gap>| !gapinput@alg;|
  <Rationals[<quiver with 3 vertices and 2 arrows>]/
  <two-sided ideal in <Rationals[<quiver with 3 vertices and 2 arrows>]>, 
    (1 generators)>>
  !gapprompt@gap>| !gapinput@cat := CatOfRightAlgebraModules(alg);|
  <cat: right modules over algebra>
  !gapprompt@gap>| !gapinput@C := StalkComplex(cat, IndecInjectiveModules(alg)[1], 0);|
  0 -> 0:(1,0,0) -> 0
  !gapprompt@gap>| !gapinput@ProjC := ProjectiveResolutionOfComplex(C);|
  --- -> 0: P1 -> 0
  !gapprompt@gap>| !gapinput@InjC := ProjectiveToInjectiveComplex(ProjC);|
  --- -> 1: I2 -> 0: I1 -> 0
  !gapprompt@gap>| !gapinput@TauC := TauOfComplex(C);|
  0 -> 1: P3 -> 0
  !gapprompt@gap>| !gapinput@IsProjectiveComplex(C);|
  false
  !gapprompt@gap>| !gapinput@IsInjectiveComplex(C);|
  true
  !gapprompt@gap>| !gapinput@C;|
  0 -> 0: I1 -> 0 
\end{Verbatim}
 }

 

\subsection{\textcolor{Chapter }{StarOfMapBetweenProjectives}}
\logpage{[ 11, 2, 5 ]}\nobreak
\hyperdef{L}{X84D4D700876395AF}{}
{\noindent\textcolor{FuncColor}{$\triangleright$\enspace\texttt{StarOfMapBetweenProjectives({\mdseries\slshape f, list{\textunderscore}i, list{\textunderscore}j})\index{StarOfMapBetweenProjectives@\texttt{StarOfMapBetweenProjectives}}
\label{StarOfMapBetweenProjectives}
}\hfill{\scriptsize (operation)}}\\
\noindent\textcolor{FuncColor}{$\triangleright$\enspace\texttt{StarOfMapBetweenIndecProjectives({\mdseries\slshape f, i, list{\textunderscore}j})\index{StarOfMapBetweenIndecProjectives@\texttt{StarOfMapBetweenIndecProjectives}}
\label{StarOfMapBetweenIndecProjectives}
}\hfill{\scriptsize (operation)}}\\
\noindent\textcolor{FuncColor}{$\triangleright$\enspace\texttt{StarOfMapBetweenDecompProjectives({\mdseries\slshape f, list{\textunderscore}i, list{\textunderscore}j})\index{StarOfMapBetweenDecompProjectives@\texttt{StarOfMapBetweenDecompProjectives}}
\label{StarOfMapBetweenDecompProjectives}
}\hfill{\scriptsize (operation)}}\\


 Arguments: \mbox{\texttt{\mdseries\slshape f}} \texttt{\symbol{45}}\texttt{\symbol{45}} a map between to projective modules $P = \bigoplus P_i$ and $Q = \bigoplus Q_j$, each of which were constructed as direct sums of indecomposable projective
modules; \mbox{\texttt{\mdseries\slshape list{\textunderscore}i}} \texttt{\symbol{45}}\texttt{\symbol{45}} describes the summands of $P$; \mbox{\texttt{\mdseries\slshape list{\textunderscore}j}} \texttt{\symbol{45}}\texttt{\symbol{45}} describes the summands of $Q$. If $P = P_1 \oplus P_3 \oplus P_3$ (where $P_i$ is the indecomposable projective representation in vertex $i$), then \mbox{\texttt{\mdseries\slshape list{\textunderscore}i}} is [1,3,3].\\
 \textbf{\indent Returns:\ }
The map $f^* = \Hom_A(f,A): \Hom_A(Q,A) \rightarrow \Hom_A(P,A)$ in $A^{\mathrm{op}}$ (where $A$ is the original algebra). 



The function \texttt{StarOfMapBetweenProjectives} is supposed to be called from within the \texttt{ProjectiveToInjectiveComplex} method, and might not do as expected when called from somewhere else. 

 The other similarly named functions are called from within the first. }

 }

 }

 
\chapter{\textcolor{Chapter }{Combinatorial representation theory}}\label{Combinatorial_Rep_Theory}
\logpage{[ 12, 0, 0 ]}
\hyperdef{L}{X7F34F6A77A24AF1C}{}
{
  
\section{\textcolor{Chapter }{Introduction}}\logpage{[ 12, 1, 0 ]}
\hyperdef{L}{X7DFB63A97E67C0A1}{}
{
  Here we introduce the implementation of the software package CREP initially
designed for MAPLE. }

 
\section{\textcolor{Chapter }{Different unit forms}}\label{UnitForms}
\logpage{[ 12, 2, 0 ]}
\hyperdef{L}{X81C656897FC2CE5A}{}
{
  

\subsection{\textcolor{Chapter }{IsUnitForm}}
\logpage{[ 12, 2, 1 ]}\nobreak
\hyperdef{L}{X79E906FD7B3DF010}{}
{\noindent\textcolor{FuncColor}{$\triangleright$\enspace\texttt{IsUnitForm\index{IsUnitForm@\texttt{IsUnitForm}}
\label{IsUnitForm}
}\hfill{\scriptsize (Category)}}\\


 The category for unit forms, which we identify with symmetric integral
matrices with 2 along the diagonal. }

 

\subsection{\textcolor{Chapter }{BilinearFormOfUnitForm}}
\logpage{[ 12, 2, 2 ]}\nobreak
\hyperdef{L}{X81A3BD6C8027B193}{}
{\noindent\textcolor{FuncColor}{$\triangleright$\enspace\texttt{BilinearFormOfUnitForm({\mdseries\slshape B})\index{BilinearFormOfUnitForm@\texttt{BilinearFormOfUnitForm}}
\label{BilinearFormOfUnitForm}
}\hfill{\scriptsize (attribute)}}\\


 Arguments: \mbox{\texttt{\mdseries\slshape B}} \texttt{\symbol{45}}\texttt{\symbol{45}} a unit form.\\
 \textbf{\indent Returns:\ }
 the bilinear form associated to a unit form \mbox{\texttt{\mdseries\slshape B}}. 



The bilinear form associated to the unitform \mbox{\texttt{\mdseries\slshape B}} given by a matrix \texttt{B} is defined for two vectors \texttt{x} and \texttt{y} as: $x*B*y^T$. }

 

\subsection{\textcolor{Chapter }{IsWeaklyNonnegativeUnitForm}}
\logpage{[ 12, 2, 3 ]}\nobreak
\hyperdef{L}{X857561F27D797EEC}{}
{\noindent\textcolor{FuncColor}{$\triangleright$\enspace\texttt{IsWeaklyNonnegativeUnitForm({\mdseries\slshape B})\index{IsWeaklyNonnegativeUnitForm@\texttt{IsWeaklyNonnegativeUnitForm}}
\label{IsWeaklyNonnegativeUnitForm}
}\hfill{\scriptsize (property)}}\\


 Arguments: \mbox{\texttt{\mdseries\slshape B}} \texttt{\symbol{45}}\texttt{\symbol{45}} a unit form.\\
 \textbf{\indent Returns:\ }
 true is the unitform \mbox{\texttt{\mdseries\slshape B}} is weakly non\texttt{\symbol{45}}negative, otherwise false. 



The unit form \mbox{\texttt{\mdseries\slshape B}} is weakly non\texttt{\symbol{45}}negative is $B(x,y) \geq 0$ for all $x\neq 0$ in $\mathbb{Z}^n$, where $n$ is the dimension of the square matrix associated to \mbox{\texttt{\mdseries\slshape B}}. }

 

\subsection{\textcolor{Chapter }{IsWeaklyPositiveUnitForm}}
\logpage{[ 12, 2, 4 ]}\nobreak
\hyperdef{L}{X7F822DCA87F05BBC}{}
{\noindent\textcolor{FuncColor}{$\triangleright$\enspace\texttt{IsWeaklyPositiveUnitForm({\mdseries\slshape B})\index{IsWeaklyPositiveUnitForm@\texttt{IsWeaklyPositiveUnitForm}}
\label{IsWeaklyPositiveUnitForm}
}\hfill{\scriptsize (property)}}\\


 Arguments: \mbox{\texttt{\mdseries\slshape B}} \texttt{\symbol{45}}\texttt{\symbol{45}} a unit form.\\
 \textbf{\indent Returns:\ }
 true is the unitform \mbox{\texttt{\mdseries\slshape B}} is weakly positive, otherwise false. 



The unit form \mbox{\texttt{\mdseries\slshape B}} is weakly positive if $B(x,y) > 0$ for all $x\neq 0$ in $\mathbb{Z}^n$, where $n$ is the dimension of the square matrix associated to \mbox{\texttt{\mdseries\slshape B}}. }

 

\subsection{\textcolor{Chapter }{PositiveRootsOfUnitForm}}
\logpage{[ 12, 2, 5 ]}\nobreak
\hyperdef{L}{X79CAF63A806BCF80}{}
{\noindent\textcolor{FuncColor}{$\triangleright$\enspace\texttt{PositiveRootsOfUnitForm({\mdseries\slshape B})\index{PositiveRootsOfUnitForm@\texttt{PositiveRootsOfUnitForm}}
\label{PositiveRootsOfUnitForm}
}\hfill{\scriptsize (attribute)}}\\


 Arguments: \mbox{\texttt{\mdseries\slshape B}} \texttt{\symbol{45}}\texttt{\symbol{45}} a unit form.\\
 \textbf{\indent Returns:\ }
 the positive roots of a unit form, if the unit form is weakly positive. If
they have not been computed, an error message will be returned saying "no
method found!". 



 This attribute will be attached to \mbox{\texttt{\mdseries\slshape B}} when \texttt{IsWeaklyPositiveUnitForm} is applied to \mbox{\texttt{\mdseries\slshape B}} and it is weakly positive. }

 

\subsection{\textcolor{Chapter }{QuadraticFormOfUnitForm}}
\logpage{[ 12, 2, 6 ]}\nobreak
\hyperdef{L}{X799A82E58690F9C5}{}
{\noindent\textcolor{FuncColor}{$\triangleright$\enspace\texttt{QuadraticFormOfUnitForm({\mdseries\slshape B})\index{QuadraticFormOfUnitForm@\texttt{QuadraticFormOfUnitForm}}
\label{QuadraticFormOfUnitForm}
}\hfill{\scriptsize (attribute)}}\\


 Arguments: \mbox{\texttt{\mdseries\slshape B}} \texttt{\symbol{45}}\texttt{\symbol{45}} a unit form.\\
 \textbf{\indent Returns:\ }
 the quadratic form associated to a unit form \mbox{\texttt{\mdseries\slshape B}}. 



The quadratic form associated to the unitform \mbox{\texttt{\mdseries\slshape B}} given by a matrix \texttt{B} is defined for a vector \texttt{x} as: $\frac{1}{2}x*B*x^T$. }

 

\subsection{\textcolor{Chapter }{SymmetricMatrixOfUnitForm}}
\logpage{[ 12, 2, 7 ]}\nobreak
\hyperdef{L}{X7837555780590360}{}
{\noindent\textcolor{FuncColor}{$\triangleright$\enspace\texttt{SymmetricMatrixOfUnitForm({\mdseries\slshape B})\index{SymmetricMatrixOfUnitForm@\texttt{SymmetricMatrixOfUnitForm}}
\label{SymmetricMatrixOfUnitForm}
}\hfill{\scriptsize (attribute)}}\\


 Arguments: \mbox{\texttt{\mdseries\slshape B}} \texttt{\symbol{45}}\texttt{\symbol{45}} a unit form.\\
 \textbf{\indent Returns:\ }
 the symmetric integral matrix which defines the unit form \mbox{\texttt{\mdseries\slshape B}}. 

}

 

\subsection{\textcolor{Chapter }{TitsUnitFormOfAlgebra}}
\logpage{[ 12, 2, 8 ]}\nobreak
\hyperdef{L}{X8404227F7C6BB5D4}{}
{\noindent\textcolor{FuncColor}{$\triangleright$\enspace\texttt{TitsUnitFormOfAlgebra({\mdseries\slshape A})\index{TitsUnitFormOfAlgebra@\texttt{TitsUnitFormOfAlgebra}}
\label{TitsUnitFormOfAlgebra}
}\hfill{\scriptsize (operation)}}\\


 Arguments: \mbox{\texttt{\mdseries\slshape A}} \texttt{\symbol{45}}\texttt{\symbol{45}} a finite dimensional (quotient of a)
path algebra (by an admissible ideal).\\
 \textbf{\indent Returns:\ }
 the Tits unit form associated to the algebra \mbox{\texttt{\mdseries\slshape A}}. 



This function returns the Tits unitform associated to a finite dimensional
quotient of a path algebra by an admissible ideal or path algebra, given that
the underlying quiver has no loops or minimal relations that starts and ends
in the same vertex. That is, then it returns a symmetric matrix $B$ such that for $x = (x_1,...,x_n) (1/2)*(x_1,...,x_n)B(x_1,...,x_n)^T = \sum_{i=1}^n x_i^2 -
\sum_{i,j} \dim_k \Ext^1(S_i,S_j)x_ix_j + \sum_{i,j} \dim_k
\Ext^2(S_i,S_j)x_ix_j$, where $n$ is the number of vertices in $Q$. }

 

\subsection{\textcolor{Chapter }{EulerBilinearFormOfAlgebra}}
\logpage{[ 12, 2, 9 ]}\nobreak
\hyperdef{L}{X836F702A845D1A84}{}
{\noindent\textcolor{FuncColor}{$\triangleright$\enspace\texttt{EulerBilinearFormOfAlgebra({\mdseries\slshape A})\index{EulerBilinearFormOfAlgebra@\texttt{EulerBilinearFormOfAlgebra}}
\label{EulerBilinearFormOfAlgebra}
}\hfill{\scriptsize (operation)}}\\


 Arguments: \mbox{\texttt{\mdseries\slshape A}} \texttt{\symbol{45}}\texttt{\symbol{45}} a finite dimensional (quotient of a)
path algebra (by an admissible ideal).\\
 \textbf{\indent Returns:\ }
 the Euler (non\texttt{\symbol{45}}symmetric) bilinear form associated to the
algebra \mbox{\texttt{\mdseries\slshape A}}. 



This function returns the Euler (non\texttt{\symbol{45}}symmetric) bilinear
form associated to a finite dimensional (basic) quotient of a path algebra \mbox{\texttt{\mdseries\slshape A}}. That is, it returns a bilinear form (function) defined by\\
 $f(x,y) = x*\textrm{CartanMatrix}(A)^{(-1)}*y$\\
 It makes sense only in case \mbox{\texttt{\mdseries\slshape A}} is of finite global dimension. }

 

\subsection{\textcolor{Chapter }{UnitForm}}
\logpage{[ 12, 2, 10 ]}\nobreak
\hyperdef{L}{X7D52745781DF8980}{}
{\noindent\textcolor{FuncColor}{$\triangleright$\enspace\texttt{UnitForm({\mdseries\slshape B})\index{UnitForm@\texttt{UnitForm}}
\label{UnitForm}
}\hfill{\scriptsize (operation)}}\\


 Arguments: \mbox{\texttt{\mdseries\slshape B}} \texttt{\symbol{45}}\texttt{\symbol{45}} an integral matrix.\\
 \textbf{\indent Returns:\ }
 the unit form in the category \texttt{IsUnitForm} (\ref{IsUnitForm}) associated to the matrix \mbox{\texttt{\mdseries\slshape B}}. 



The function checks if \mbox{\texttt{\mdseries\slshape B}} is a symmetric integral matrix with 2 along the diagonal, and returns an error
message otherwise. In addition it sets the attributes, \texttt{BilinearFormOfUnitForm} (\ref{BilinearFormOfUnitForm}), \texttt{QuadraticFormOfUnitForm} (\ref{QuadraticFormOfUnitForm}) and \texttt{SymmetricMatrixOfUnitForm} (\ref{SymmetricMatrixOfUnitForm}). }

 }

 
\section{\textcolor{Chapter }{Combinatorial Maps}}\label{Combinatorial Maps}
\logpage{[ 12, 3, 0 ]}
\hyperdef{L}{X8779DC4586048F1F}{}
{
  

\subsection{\textcolor{Chapter }{IsCombinatorialMap}}
\logpage{[ 12, 3, 1 ]}\nobreak
\hyperdef{L}{X8413C62083116C92}{}
{\noindent\textcolor{FuncColor}{$\triangleright$\enspace\texttt{IsCombinatorialMap({\mdseries\slshape })\index{IsCombinatorialMap@\texttt{IsCombinatorialMap}}
\label{IsCombinatorialMap}
}\hfill{\scriptsize (filter)}}\\


 This defines the category \texttt{IsCombinatorialMap}. }

 

\subsection{\textcolor{Chapter }{CombinatorialMap (combinatorialmap)}}
\logpage{[ 12, 3, 2 ]}\nobreak
\hyperdef{L}{X814C0E20858B22D6}{}
{\noindent\textcolor{FuncColor}{$\triangleright$\enspace\texttt{CombinatorialMap({\mdseries\slshape n, sigma, iota, m})\index{CombinatorialMap@\texttt{CombinatorialMap}!combinatorialmap}
\label{CombinatorialMap:combinatorialmap}
}\hfill{\scriptsize (operation)}}\\


 Arguments : \mbox{\texttt{\mdseries\slshape n}} \texttt{\symbol{45}}\texttt{\symbol{45}} number of
half\texttt{\symbol{45}}edges, \mbox{\texttt{\mdseries\slshape sigma}} \texttt{\symbol{45}}\texttt{\symbol{45}} a permutation specifying the ordering
of half\texttt{\symbol{45}}edges around vertices, \mbox{\texttt{\mdseries\slshape iota}} \texttt{\symbol{45}}\texttt{\symbol{45}} an involution pairing the
half\texttt{\symbol{45}}edges, \mbox{\texttt{\mdseries\slshape m}} \texttt{\symbol{45}}\texttt{\symbol{45}} a list of marked
half\texttt{\symbol{45}}edges \textbf{\indent Returns:\ }
 a combinatorial map, which is an object from the category \texttt{IsCombinatorialMap} (\ref{IsCombinatorialMap}).



 The size of the underlying set given by \mbox{\texttt{\mdseries\slshape size}} is given a positive integer, the ordering and the pairing are given by two
permutations \mbox{\texttt{\mdseries\slshape ordering}} and \mbox{\texttt{\mdseries\slshape paring}} of a set of \mbox{\texttt{\mdseries\slshape size}} elements (\texttt{\symbol{123}}1, 2,..., \mbox{\texttt{\mdseries\slshape size}}\texttt{\symbol{125}}) and the marked half\texttt{\symbol{45}}edges are given
as a list of integers in \texttt{\symbol{123}}1, 2,..., \mbox{\texttt{\mdseries\slshape size}}\texttt{\symbol{125}}. }

 

\subsection{\textcolor{Chapter }{FacesOfCombinatorialMap (map)}}
\logpage{[ 12, 3, 3 ]}\nobreak
\hyperdef{L}{X82713DF87D1FA0D0}{}
{\noindent\textcolor{FuncColor}{$\triangleright$\enspace\texttt{FacesOfCombinatorialMap({\mdseries\slshape map})\index{FacesOfCombinatorialMap@\texttt{FacesOfCombinatorialMap}!map}
\label{FacesOfCombinatorialMap:map}
}\hfill{\scriptsize (attribute)}}\\


 Argument : \mbox{\texttt{\mdseries\slshape map}} \texttt{\symbol{45}}\texttt{\symbol{45}} a combinatorial map. \textbf{\indent Returns:\ }
 the list of faces of the combinatorial map \mbox{\texttt{\mdseries\slshape map}}. 



 The faces are given by a list of half\texttt{\symbol{45}}edges of \mbox{\texttt{\mdseries\slshape map}}. }

 

\subsection{\textcolor{Chapter }{GenusOfCombinatorialMap (map)}}
\logpage{[ 12, 3, 4 ]}\nobreak
\hyperdef{L}{X86E5D64679A503C5}{}
{\noindent\textcolor{FuncColor}{$\triangleright$\enspace\texttt{GenusOfCombinatorialMap({\mdseries\slshape map})\index{GenusOfCombinatorialMap@\texttt{GenusOfCombinatorialMap}!map}
\label{GenusOfCombinatorialMap:map}
}\hfill{\scriptsize (attribute)}}\\


 Argument : \mbox{\texttt{\mdseries\slshape map}} \texttt{\symbol{45}}\texttt{\symbol{45}} a combinatorial map. \textbf{\indent Returns:\ }
 the genus of the surface represented by the combinatorial map \mbox{\texttt{\mdseries\slshape map}}.

}

 

\subsection{\textcolor{Chapter }{DualOfCombinatorialMap (map)}}
\logpage{[ 12, 3, 5 ]}\nobreak
\hyperdef{L}{X83BC50DC822EB5F8}{}
{\noindent\textcolor{FuncColor}{$\triangleright$\enspace\texttt{DualOfCombinatorialMap({\mdseries\slshape map})\index{DualOfCombinatorialMap@\texttt{DualOfCombinatorialMap}!map}
\label{DualOfCombinatorialMap:map}
}\hfill{\scriptsize (attribute)}}\\


 Argument : \mbox{\texttt{\mdseries\slshape map}} \texttt{\symbol{45}}\texttt{\symbol{45}} a combinatorial map. \textbf{\indent Returns:\ }
 the dual combinatorial map of \mbox{\texttt{\mdseries\slshape map}}. 

}

 

\subsection{\textcolor{Chapter }{MaximalPathsOfGentleAlgebra (algebra)}}
\logpage{[ 12, 3, 6 ]}\nobreak
\hyperdef{L}{X86AFE052783F5542}{}
{\noindent\textcolor{FuncColor}{$\triangleright$\enspace\texttt{MaximalPathsOfGentleAlgebra({\mdseries\slshape A})\index{MaximalPathsOfGentleAlgebra@\texttt{MaximalPathsOfGentleAlgebra}!algebra}
\label{MaximalPathsOfGentleAlgebra:algebra}
}\hfill{\scriptsize (attribute)}}\\


 Argument : \mbox{\texttt{\mdseries\slshape A}} \texttt{\symbol{45}}\texttt{\symbol{45}} a gentle algebra. \textbf{\indent Returns:\ }
 the list of maximal paths of the quiver with relations defining \mbox{\texttt{\mdseries\slshape A}}. 



 Only returns paths of non\texttt{\symbol{45}}zero length. }

 

\subsection{\textcolor{Chapter }{RemoveEdgeOfCombinatorialMap (map, edge)}}
\logpage{[ 12, 3, 7 ]}\nobreak
\hyperdef{L}{X84E2EE177F234189}{}
{\noindent\textcolor{FuncColor}{$\triangleright$\enspace\texttt{RemoveEdgeOfCombinatorialMap({\mdseries\slshape combmap, E})\index{RemoveEdgeOfCombinatorialMap@\texttt{RemoveEdgeOfCombinatorialMap}!map, edge}
\label{RemoveEdgeOfCombinatorialMap:map, edge}
}\hfill{\scriptsize (operation)}}\\


 Argument : \mbox{\texttt{\mdseries\slshape map}} \texttt{\symbol{45}}\texttt{\symbol{45}} a combinatorial map, \mbox{\texttt{\mdseries\slshape edge}} \texttt{\symbol{45}}\texttt{\symbol{45}} a list of two paired
half\texttt{\symbol{45}}edges. \textbf{\indent Returns:\ }
 the combinatorial map obtained by removing \mbox{\texttt{\mdseries\slshape edge}} from the combinatorial map \mbox{\texttt{\mdseries\slshape map}}. 



 The argument \mbox{\texttt{\mdseries\slshape edge}} has to be a list consisting of two half\texttt{\symbol{45}}edges of \mbox{\texttt{\mdseries\slshape map}} that are paired. The returned combinatorial map does not change size, the two
removed half\texttt{\symbol{45}}edges are now fixed by the pairing and
ordering so do not appear in non\texttt{\symbol{45}}trivial orbits. }

 

\subsection{\textcolor{Chapter }{CombinatorialMapOfGentleAlgebra (algebra)}}
\logpage{[ 12, 3, 8 ]}\nobreak
\hyperdef{L}{X7A5333CC814054FF}{}
{\noindent\textcolor{FuncColor}{$\triangleright$\enspace\texttt{CombinatorialMapOfGentleAlgebra({\mdseries\slshape A})\index{CombinatorialMapOfGentleAlgebra@\texttt{CombinatorialMapOfGentleAlgebra}!algebra}
\label{CombinatorialMapOfGentleAlgebra:algebra}
}\hfill{\scriptsize (attribute)}}\\


 Argument : \mbox{\texttt{\mdseries\slshape A}} \texttt{\symbol{45}}\texttt{\symbol{45}} a gentle algebra. \textbf{\indent Returns:\ }
 the combinatorial map corresponding to the Brauer Graph of the trivial
extension of \mbox{\texttt{\mdseries\slshape algebra}} with regard to its dual. 



 The function assumes that \mbox{\texttt{\mdseries\slshape algebra}} is gentle. }

 

\subsection{\textcolor{Chapter }{MarkedBoundariesOfCombinatorialMap (map)}}
\logpage{[ 12, 3, 9 ]}\nobreak
\hyperdef{L}{X7E23962982506CB4}{}
{\noindent\textcolor{FuncColor}{$\triangleright$\enspace\texttt{MarkedBoundariesOfCombinatorialMap({\mdseries\slshape map})\index{MarkedBoundariesOfCombinatorialMap@\texttt{MarkedBoundariesOfCombinatorialMap}!map}
\label{MarkedBoundariesOfCombinatorialMap:map}
}\hfill{\scriptsize (attribute)}}\\


 Argument : \mbox{\texttt{\mdseries\slshape map}} \texttt{\symbol{45}}\texttt{\symbol{45}} a combinatorial map. \textbf{\indent Returns:\ }
 a list consisting of pairs [Face,numberofmarkedpoints] where Face is a face of
the combinatorial map \mbox{\texttt{\mdseries\slshape map}} and numberofmarkedpoints is the number of marked half\texttt{\symbol{45}}edges
on the corresponding boundary component. 

}

 

\subsection{\textcolor{Chapter }{WindingNumber (map, curve)}}
\logpage{[ 12, 3, 10 ]}\nobreak
\hyperdef{L}{X84D1EB9379681C62}{}
{\noindent\textcolor{FuncColor}{$\triangleright$\enspace\texttt{WindingNumber({\mdseries\slshape map, gamma})\index{WindingNumber@\texttt{WindingNumber}!map, curve}
\label{WindingNumber:map, curve}
}\hfill{\scriptsize (operation)}}\\


 Argument : \mbox{\texttt{\mdseries\slshape map}} \texttt{\symbol{45}}\texttt{\symbol{45}} a combinatorial map, \mbox{\texttt{\mdseries\slshape gamma}} \texttt{\symbol{45}}\texttt{\symbol{45}} a list of
half\texttt{\symbol{45}}edges forming a closed curve on the surface. \textbf{\indent Returns:\ }
 the combinatorial winding number of \mbox{\texttt{\mdseries\slshape gamma}} on the dissected surface given by \mbox{\texttt{\mdseries\slshape map}}. 

}

 

\subsection{\textcolor{Chapter }{BoundaryCurvesOfCombinatorialMap (map)}}
\logpage{[ 12, 3, 11 ]}\nobreak
\hyperdef{L}{X8599C842827591B3}{}
{\noindent\textcolor{FuncColor}{$\triangleright$\enspace\texttt{BoundaryCurvesOfCombinatorialMap({\mdseries\slshape map})\index{BoundaryCurvesOfCombinatorialMap@\texttt{BoundaryCurvesOfCombinatorialMap}!map}
\label{BoundaryCurvesOfCombinatorialMap:map}
}\hfill{\scriptsize (attribute)}}\\


 Argument : \mbox{\texttt{\mdseries\slshape map}} \texttt{\symbol{45}}\texttt{\symbol{45}} a combinatorial map. \textbf{\indent Returns:\ }
 a list whose elements are lists of half\texttt{\symbol{45}}edges of the
combinatorial map \mbox{\texttt{\mdseries\slshape map}}. Each of these lists corresponds to a closed curve homotopic to a boundary
component of the surface represented by \mbox{\texttt{\mdseries\slshape map}}. 



 The closed curves are represented by a list of adjacent
half\texttt{\symbol{45}}edges. The orientation of these curves is chosen so
that the corresponding boundary component is to the right. The returned curves
correspond to the curves in minimal position in regards to the dual dissection
of the surface. In the case of the disk, the boundary curve is trivial and the
empty list is returned. }

 

\subsection{\textcolor{Chapter }{DepthSearchCombinatorialMap (map, halfedge)}}
\logpage{[ 12, 3, 12 ]}\nobreak
\hyperdef{L}{X7C1850CF83BDBB05}{}
{\noindent\textcolor{FuncColor}{$\triangleright$\enspace\texttt{DepthSearchCombinatorialMap({\mdseries\slshape combmap, x})\index{DepthSearchCombinatorialMap@\texttt{DepthSearchCombinatorialMap}!map, halfedge}
\label{DepthSearchCombinatorialMap:map, halfedge}
}\hfill{\scriptsize (operation)}}\\


 Argument : \mbox{\texttt{\mdseries\slshape map}} \texttt{\symbol{45}}\texttt{\symbol{45}} a combinatorial map, \mbox{\texttt{\mdseries\slshape x}} \texttt{\symbol{45}}\texttt{\symbol{45}} a half\texttt{\symbol{45}}edge of \mbox{\texttt{\mdseries\slshape map}} \textbf{\indent Returns:\ }
 a list of pairs of paired half\texttt{\symbol{45}}edges of the combinatorial
map \mbox{\texttt{\mdseries\slshape map}} corresponding to the cover tree of the underlying graph of \mbox{\texttt{\mdseries\slshape map}} obtained by a depth first search with root the vertex to which is attached \mbox{\texttt{\mdseries\slshape x}}. 



 The paired half\texttt{\symbol{45}}edges of the cover tree are oriented
towards the root of the cover tree. }

 

\subsection{\textcolor{Chapter }{WidthSearchCombinatorialMap (map, halfedge)}}
\logpage{[ 12, 3, 13 ]}\nobreak
\hyperdef{L}{X80952CEC808FFB52}{}
{\noindent\textcolor{FuncColor}{$\triangleright$\enspace\texttt{WidthSearchCombinatorialMap({\mdseries\slshape combmap, x})\index{WidthSearchCombinatorialMap@\texttt{WidthSearchCombinatorialMap}!map, halfedge}
\label{WidthSearchCombinatorialMap:map, halfedge}
}\hfill{\scriptsize (operation)}}\\


 Argument : \mbox{\texttt{\mdseries\slshape map}} \texttt{\symbol{45}}\texttt{\symbol{45}} a combinatorial map, \mbox{\texttt{\mdseries\slshape x}} \texttt{\symbol{45}}\texttt{\symbol{45}} a half\texttt{\symbol{45}}edge of \mbox{\texttt{\mdseries\slshape map}} \textbf{\indent Returns:\ }
 a list of pairs of paired half\texttt{\symbol{45}}edges of the combinatorial
map \mbox{\texttt{\mdseries\slshape map}} corresponding to the cover tree of the underlying graph of \mbox{\texttt{\mdseries\slshape map}} obtained by a width first search with root the vertex to which is attached \mbox{\texttt{\mdseries\slshape x}}. 



 The paired half\texttt{\symbol{45}}edges of the cover tree are oriented
towards the root of the cover tree. }

 

\subsection{\textcolor{Chapter }{NonSeperatingCurve (map)}}
\logpage{[ 12, 3, 14 ]}\nobreak
\hyperdef{L}{X85C23E5180E0F1BB}{}
{\noindent\textcolor{FuncColor}{$\triangleright$\enspace\texttt{NonSeperatingCurve({\mdseries\slshape map})\index{NonSeperatingCurve@\texttt{NonSeperatingCurve}!map}
\label{NonSeperatingCurve:map}
}\hfill{\scriptsize (attribute)}}\\


 Argument : \mbox{\texttt{\mdseries\slshape map}} \texttt{\symbol{45}}\texttt{\symbol{45}} a combinatorial map. \textbf{\indent Returns:\ }
 a list of half\texttt{\symbol{45}}edges of the combinatorial map \mbox{\texttt{\mdseries\slshape map}} corresponding to a non\texttt{\symbol{45}}seperating closed curve of the
surface. 



 This function must be used on a combinatorial map of genus at least one. }

 

\subsection{\textcolor{Chapter }{CutNonSepCombinatorialMap (map, curve, index)}}
\logpage{[ 12, 3, 15 ]}\nobreak
\hyperdef{L}{X8233249683293DCD}{}
{\noindent\textcolor{FuncColor}{$\triangleright$\enspace\texttt{CutNonSepCombinatorialMap({\mdseries\slshape map, alpha, index})\index{CutNonSepCombinatorialMap@\texttt{CutNonSepCombinatorialMap}!map, curve, index}
\label{CutNonSepCombinatorialMap:map, curve, index}
}\hfill{\scriptsize (operation)}}\\


 Arguments : \mbox{\texttt{\mdseries\slshape map}} \texttt{\symbol{45}}\texttt{\symbol{45}} a combinatorial map, \mbox{\texttt{\mdseries\slshape alpha}} \texttt{\symbol{45}}\texttt{\symbol{45}} a list of
half\texttt{\symbol{45}}edges, \mbox{\texttt{\mdseries\slshape index}} \texttt{\symbol{45}}\texttt{\symbol{45}} a list. \textbf{\indent Returns:\ }
 a triplet [newmap, [boundary1,boundary2], newindex] where newmap is the
combinatorial map obtained by cutting the surface represented by \mbox{\texttt{\mdseries\slshape map}} along the closed curve \mbox{\texttt{\mdseries\slshape curve}}. The lists boundary1 and boundary2 are the boundaries of newmap that were
created by the cut. 



 The argument \mbox{\texttt{\mdseries\slshape index}} is a list whose i\texttt{\symbol{45}}th element is the
half\texttt{\symbol{45}}edge to which i corresponds in some original
combinatorial map. The result newindex is the updated index for newmap where
all created half\texttt{\symbol{45}}edges are added. }

 

\subsection{\textcolor{Chapter }{JoinCurveCombinatorialMap (map, boundary1, boundary2, index)}}
\logpage{[ 12, 3, 16 ]}\nobreak
\hyperdef{L}{X7ECDA130870127F0}{}
{\noindent\textcolor{FuncColor}{$\triangleright$\enspace\texttt{JoinCurveCombinatorialMap({\mdseries\slshape map, bound1, bound2, index})\index{JoinCurveCombinatorialMap@\texttt{JoinCurveCombinatorialMap}!map, boundary1, boundary2, index}
\label{JoinCurveCombinatorialMap:map, boundary1, boundary2, index}
}\hfill{\scriptsize (operation)}}\\


 Arguments : \mbox{\texttt{\mdseries\slshape map}} \texttt{\symbol{45}}\texttt{\symbol{45}} a combinatorial map, \mbox{\texttt{\mdseries\slshape boundary1}} \texttt{\symbol{45}}\texttt{\symbol{45}} a list of
half\texttt{\symbol{45}}edges, \mbox{\texttt{\mdseries\slshape boundary2}} \texttt{\symbol{45}}\texttt{\symbol{45}} a list of
half\texttt{\symbol{45}}edges, \mbox{\texttt{\mdseries\slshape index}} \texttt{\symbol{45}}\texttt{\symbol{45}} a list. \textbf{\indent Returns:\ }
 a list corresponding to a simple curve of \mbox{\texttt{\mdseries\slshape map}} joining the two boundaries \mbox{\texttt{\mdseries\slshape boundary1}} and \mbox{\texttt{\mdseries\slshape boundary2}} in such a way that using \mbox{\texttt{\mdseries\slshape index}} it corresponds to a closed curve on the original combinatorial map. 



 The list consists of adjacent half\texttt{\symbol{45}}edges of \mbox{\texttt{\mdseries\slshape map}} which form the closed curve. }

 

\subsection{\textcolor{Chapter }{CutJoinCurveCombinatorialMap (map, boundary1, boundary2, curve, index)}}
\logpage{[ 12, 3, 17 ]}\nobreak
\hyperdef{L}{X78D7425B7DB40019}{}
{\noindent\textcolor{FuncColor}{$\triangleright$\enspace\texttt{CutJoinCurveCombinatorialMap({\mdseries\slshape map, bound1, bound2, beta, index})\index{CutJoinCurveCombinatorialMap@\texttt{CutJoinCurveCombinatorialMap}!map, boundary1, boundary2, curve, index}
\label{CutJoinCurveCombinatorialMap:map, boundary1, boundary2, curve, index}
}\hfill{\scriptsize (operation)}}\\


 Arguments : \mbox{\texttt{\mdseries\slshape map}} \texttt{\symbol{45}}\texttt{\symbol{45}} a combinatorial map, \mbox{\texttt{\mdseries\slshape boundary1}} \texttt{\symbol{45}}\texttt{\symbol{45}} a list of
half\texttt{\symbol{45}}edges, \mbox{\texttt{\mdseries\slshape boundary2}} \texttt{\symbol{45}}\texttt{\symbol{45}} a list of
half\texttt{\symbol{45}}edges, \mbox{\texttt{\mdseries\slshape beta}} \texttt{\symbol{45}}\texttt{\symbol{45}} a list of
half\texttt{\symbol{45}}edges corresponding to a curve, \mbox{\texttt{\mdseries\slshape index}} \texttt{\symbol{45}}\texttt{\symbol{45}} a list. \textbf{\indent Returns:\ }
 the pair [newmap, newindex] where newmap is the combinatorial map obtained by
cutting the combinatorial map \mbox{\texttt{\mdseries\slshape map}} along \mbox{\texttt{\mdseries\slshape curve}} which has to be a simple curve joining \mbox{\texttt{\mdseries\slshape boundary1}} and \mbox{\texttt{\mdseries\slshape boundary2}}

}

 

\subsection{\textcolor{Chapter }{HomologyBasisOfCombinatorialMap (map)}}
\logpage{[ 12, 3, 18 ]}\nobreak
\hyperdef{L}{X874A2B0B84D732B2}{}
{\noindent\textcolor{FuncColor}{$\triangleright$\enspace\texttt{HomologyBasisOfCombinatorialMap({\mdseries\slshape map})\index{HomologyBasisOfCombinatorialMap@\texttt{HomologyBasisOfCombinatorialMap}!map}
\label{HomologyBasisOfCombinatorialMap:map}
}\hfill{\scriptsize (attribute)}}\\


 Argument : \mbox{\texttt{\mdseries\slshape map}} \texttt{\symbol{45}}\texttt{\symbol{45}} a combinatorial map. \textbf{\indent Returns:\ }
 a list of pairs [a{\textunderscore}i,b{\textunderscore}i] where
a{\textunderscore}i and b{\textunderscore}i are lists of
half\texttt{\symbol{45}}edges of the combinatorial map \mbox{\texttt{\mdseries\slshape map}} forming a closed curve such that the a{\textunderscore}i and
b{\textunderscore}i form a simplectic basis of the first homology group of the
surface represented by \mbox{\texttt{\mdseries\slshape map}} for the intersection form. 

}

 

\subsection{\textcolor{Chapter }{AreDerivedEquivalent (A,B)}}
\logpage{[ 12, 3, 19 ]}\nobreak
\hyperdef{L}{X82C2709B7CC86B07}{}
{\noindent\textcolor{FuncColor}{$\triangleright$\enspace\texttt{AreDerivedEquivalent({\mdseries\slshape A, B})\index{AreDerivedEquivalent@\texttt{AreDerivedEquivalent}!A,B}
\label{AreDerivedEquivalent:A,B}
}\hfill{\scriptsize (operation)}}\\


 Arguments : \mbox{\texttt{\mdseries\slshape A}} \texttt{\symbol{45}}\texttt{\symbol{45}} a gentle algebra, \mbox{\texttt{\mdseries\slshape B}} \texttt{\symbol{45}}\texttt{\symbol{45}} a gentle algebra. \textbf{\indent Returns:\ }
 true if \mbox{\texttt{\mdseries\slshape A}} and \mbox{\texttt{\mdseries\slshape B}} are derived equivalent and false otherwise. The arguments must be gentle
algebras over the same field. 

}

 }

 }

 
\chapter{\textcolor{Chapter }{Degeneration order for modules in finite type}}\label{DegOrderFiniteType}
\logpage{[ 13, 0, 0 ]}
\hyperdef{L}{X82CC1A63854C04F1}{}
{
  
\section{\textcolor{Chapter }{Introduction}}\label{introDO}
\logpage{[ 13, 1, 0 ]}
\hyperdef{L}{X7DFB63A97E67C0A1}{}
{
  This is an implementation of several tools for computing degeneration order
for modules over algebras of finite type. It can be treated as a "subpackage"
of QPA and used separately since the functions do not use any of QPA routines
so far.\\
 This subpackage has a little bit different philosophy than QPA in general.
Namely, the "starting point" is not an algebra A defined by a Gabriel quiver
with relations but an Auslander\texttt{\symbol{45}}Reiten
(A\texttt{\symbol{45}}R) quiver of the category mod A, defined by numerical
data (see \texttt{ARQuiverNumerical} (\ref{ARQuiverNumerical})). All the indecomposables (actually their isoclasses) have unique natural
numbers established at the beginning, by invoking \texttt{ARQuiverNumerical} (\ref{ARQuiverNumerical}). This function should be used before all further computations. An arbitrary
module M is identified by its multiplicity vector (the sequence of
multiplicities of all the indecomposables appearing in a direct sum
decomposition of M). \\
 Here we always assume that A is an algebra of finite representation type. Note
that in this case deg\texttt{\symbol{45}}order coincide with
Hom\texttt{\symbol{45}}order, and this fact is used in the algorithms of this
subpackage. The main goal of this subpackage is to give tools for testing a
deg\texttt{\symbol{45}}order relation between two A\texttt{\symbol{45}}modules
and determining (direct) deg\texttt{\symbol{45}}order predecessors and
successors (see \ref{bdef} for basic definitions from this theory). As a side effect one can also obtain
the dimensions of Hom\texttt{\symbol{45}}spaces between arbitrary modules (and
in particular the dimension vectors of indecomposable modules). }

 
\section{\textcolor{Chapter }{Basic definitions}}\label{bdef}
\logpage{[ 13, 2, 0 ]}
\hyperdef{L}{X78ED07E37FC2BD46}{}
{
  Here we briefly recall the basic notions we use in all the functions from this
chapter. 

 Let A be an algebra. We say that for two A\texttt{\symbol{45}}modules M and N
of the same dimension vector d, M degenerates to N (N is a degeneration of M)
iff N belongs to a Zariski closure of the orbit of M in a variety $mod_A(d)$ of A\texttt{\symbol{45}}modules of dimension vector d. If it is the case, we
write $M <= N$. It is well known that 

 (1) The relation $<=$ is a partial order on the set of isomorphism classes of
A\texttt{\symbol{45}}modules of dimension vector d.\\
\\
 (2) If A is an algebra of finite representation type, $<=$ coincides with so\texttt{\symbol{45}}called Hom\texttt{\symbol{45}}order $<=_{Hom}$, defined as follows: $M <=_{Hom} N$ iff $[X,M] <= [X,N]$ for all indecomposable A\texttt{\symbol{45}}modules X, where by $[Y,Z]$ we denote always the dimension of a Hom\texttt{\symbol{45}}space between Y and
Z.

 Further, if $M < N$ (i.e. $M <= N$ and M is not isomorphic to N), we say that M is a deg\texttt{\symbol{45}}order
predecessor of N (resp. N is a deg\texttt{\symbol{45}}order successor of M).
Moreover, we say that M is a direct deg\texttt{\symbol{45}}order predecessor
of N if $M < N$ and there is no M' such that $M < M' < N$ (similarly for successors). }

 
\section{\textcolor{Chapter }{Defining Auslander\texttt{\symbol{45}}Reiten quiver in finite type}}\label{DefARQuiver}
\logpage{[ 13, 3, 0 ]}
\hyperdef{L}{X7B3AC0B87DF8A219}{}
{
  

\subsection{\textcolor{Chapter }{ARQuiverNumerical}}
\logpage{[ 13, 3, 1 ]}\nobreak
\hyperdef{L}{X841AD2EF7FD636B9}{}
{\noindent\textcolor{FuncColor}{$\triangleright$\enspace\texttt{ARQuiverNumerical({\mdseries\slshape ind, proj, list, or, name, or, name, parm1, or, name, param1, param2})\index{ARQuiverNumerical@\texttt{ARQuiverNumerical}}
\label{ARQuiverNumerical}
}\hfill{\scriptsize (function)}}\\


 Arguments: \mbox{\texttt{\mdseries\slshape ind}} \texttt{\symbol{45}} number of indecomposable modules in our category;\\
 \mbox{\texttt{\mdseries\slshape proj}} \texttt{\symbol{45}} number of indecomposable projective modules in our
category;\\
 \mbox{\texttt{\mdseries\slshape list}} \texttt{\symbol{45}} list of lists containing description of meshes in
A\texttt{\symbol{45}}R quiver defined as follows:\\
 \mbox{\texttt{\mdseries\slshape list}}[i] = description of mesh ending in vertex (indec. mod.) number i having the
shape [a1,...,an,t] where\\
 a1,...,an = numbers of direct predecessors of i in A\texttt{\symbol{45}}R
quiver;\\
 t = number of tau(i), or 0 if tau i does not exist (iff i is projective).\\
 In particular if i is projective \mbox{\texttt{\mdseries\slshape list}}[i]=[a1,...,an,0] where a1,...,an are indec. summands of rad(i).

 OR:\\
 \mbox{\texttt{\mdseries\slshape list}} second version \texttt{\symbol{45}} if the first element of \mbox{\texttt{\mdseries\slshape list}} is a string "orbits" then the remaining elements should provide an alternative
(shorter than above) description of A\texttt{\symbol{45}}R quiver as follows. \\
 \mbox{\texttt{\mdseries\slshape list}}[2] is a list of descriptions of orbits identified by chosen representatives.
We assume that in case an orbit is non\texttt{\symbol{45}}periodic, then a
projective module is its representative. Each element of list \mbox{\texttt{\mdseries\slshape list}}[2] is a description of i\texttt{\symbol{45}}th orbit and has the shape:\\
 \texttt{[l, [i1,t1], ... , [is,ts]]} where\\
 l = length of orbit \texttt{\symbol{45}} 1\\
 \texttt{[i1,t1], ... , [is,ts]} all the direct predecessors of a representative of this orbit, of the shape
tau\texttt{\symbol{94}}\texttt{\symbol{123}}\texttt{\symbol{45}}t1\texttt{\symbol{125}}(i1),
and i1 denotes the representative of orbit no. i1, and so on.\\
 We assume first p elements of \mbox{\texttt{\mdseries\slshape list}}[2] are the orbits of projectives.

 REMARK: we ALWAYS assume that indecomposables with numbers 1..\mbox{\texttt{\mdseries\slshape proj}} are projectives and the only projectives (further dimension vectors are
interpreted according to this order of projectives!).

 Alternative arguments:\\
 \mbox{\texttt{\mdseries\slshape name}} = string with the name of predefined A\texttt{\symbol{45}}R quiver;\\
 \mbox{\texttt{\mdseries\slshape param1}} = (optional) parameter for \mbox{\texttt{\mdseries\slshape name}};\\
 \mbox{\texttt{\mdseries\slshape param2}} = (optional) second parameter for \mbox{\texttt{\mdseries\slshape name}}.\\
 \\
 Call ARQuiverNumerical("what") to get a description of all the names and
parameters for currently available predefined A\texttt{\symbol{45}}R quivers. \textbf{\indent Returns:\ }
an object from the category \texttt{IsARQuiverNumerical} (\ref{IsARQuiverNumerical}).



 This function "initializes" Auslander\texttt{\symbol{45}}Reiten quiver and
performs all necessary preliminary computations concerning mainly determining
the matrix of dimensions of all Hom\texttt{\symbol{45}}spaces between
indecomposables. }

 Examples.\\
 Below we define an A\texttt{\symbol{45}}R quiver of a path algebra of the
Dynkin quiver D4 with subspace orientation of arrows. 
\begin{Verbatim}[commandchars=!@|,fontsize=\small,frame=single,label=Example]
  !gapprompt@gap>| !gapinput@a := ARQuiverNumerical(12, 4, [ [0],[1,0],[1,0],[1,0],[2,3,4,1],[5,2],[5,3],[5,4],[6,7,8,5],[9,6],[9,7],[9,8] ]);|
  <ARQuiverNumerical with 12 indecomposables and 4 projectives>
   
\end{Verbatim}
 The same A\texttt{\symbol{45}}R quiver (with possibly little bit different
enumeration of indecomposables) can be obtained by invoking: 
\begin{Verbatim}[commandchars=!@|,fontsize=\small,frame=single,label=Example]
  !gapprompt@gap>| !gapinput@b := ARQuiverNumerical(12, 4, ["orbits", [ [2], [2,[1,0]], [2,[1,0]], [2,[1,0]] ] ]);|
  <ARQuiverNumerical with 12 indecomposables and 4 projectives>
   
\end{Verbatim}
 This A\texttt{\symbol{45}}R quiver can be also obtained by: 
\begin{Verbatim}[commandchars=!@|,fontsize=\small,frame=single,label=Example]
  !gapprompt@gap>| !gapinput@a := ARQuiverNumerical("D4 subspace");|
  <ARQuiverNumerical with 12 indecomposables and 4 projectives>
   
\end{Verbatim}
 since this is one of the predefined A\texttt{\symbol{45}}R quivers. \\
 Another example of predefined A\texttt{\symbol{45}}R quiver: for an algebra
from Bongartz\texttt{\symbol{45}}Gabriel list of maximal finite type algebras
with two simple modules. This is an algebra with number 5 on this list. 
\begin{Verbatim}[commandchars=!@|,fontsize=\small,frame=single,label=Example]
  !gapprompt@gap>| !gapinput@a := ARQuiverNumerical("BG", 5);|
  <ARQuiverNumerical with 72 indecomposables and 2 projectives>
   
\end{Verbatim}
 

\subsection{\textcolor{Chapter }{IsARQuiverNumerical}}
\logpage{[ 13, 3, 2 ]}\nobreak
\hyperdef{L}{X7CFA7C7D861B2D4E}{}
{\noindent\textcolor{FuncColor}{$\triangleright$\enspace\texttt{IsARQuiverNumerical\index{IsARQuiverNumerical@\texttt{IsARQuiverNumerical}}
\label{IsARQuiverNumerical}
}\hfill{\scriptsize (Category)}}\\


 Objects from this category represent Auslander\texttt{\symbol{45}}Reiten
(finite) quivers and additionally contain all data necessary for further
computations (as components accessed as usual by
!.name\texttt{\symbol{45}}of\texttt{\symbol{45}}component):\\
 ARdesc = numerical description of AR quiver (as \mbox{\texttt{\mdseries\slshape list}} in \texttt{ARQuiverNumerical} (\ref{ARQuiverNumerical})),\\
 DimHomMat = matrix [dim Hom (i,j)] (={\textgreater} rows 1..p contain dim.
vectors of all indecomposables),\\
 Simples = list of numbers of simple modules. }

 

\subsection{\textcolor{Chapter }{NumberOfIndecomposables}}
\logpage{[ 13, 3, 3 ]}\nobreak
\hyperdef{L}{X861D4F9E8238AA6E}{}
{\noindent\textcolor{FuncColor}{$\triangleright$\enspace\texttt{NumberOfIndecomposables({\mdseries\slshape AR})\index{NumberOfIndecomposables@\texttt{NumberOfIndecomposables}}
\label{NumberOfIndecomposables}
}\hfill{\scriptsize (attribute)}}\\


 Argument: \mbox{\texttt{\mdseries\slshape AR}} \texttt{\symbol{45}} an object from the category \texttt{IsARQuiverNumerical} (\ref{IsARQuiverNumerical}).\\
 \textbf{\indent Returns:\ }
 the number of indecomposable modules in \mbox{\texttt{\mdseries\slshape AR}}. 

}

 

\subsection{\textcolor{Chapter }{NumberOfProjectives}}
\logpage{[ 13, 3, 4 ]}\nobreak
\hyperdef{L}{X80BCA6AE80665A4B}{}
{\noindent\textcolor{FuncColor}{$\triangleright$\enspace\texttt{NumberOfProjectives({\mdseries\slshape AR})\index{NumberOfProjectives@\texttt{NumberOfProjectives}}
\label{NumberOfProjectives}
}\hfill{\scriptsize (attribute)}}\\


 Argument: \mbox{\texttt{\mdseries\slshape AR}} \texttt{\symbol{45}} an object from the category \texttt{IsARQuiverNumerical} (\ref{IsARQuiverNumerical}).\\
 \textbf{\indent Returns:\ }
 the number of indecomposable projective modules in \mbox{\texttt{\mdseries\slshape AR}}. 

}

 }

 
\section{\textcolor{Chapter }{Elementary operations}}\label{elopers}
\logpage{[ 13, 4, 0 ]}
\hyperdef{L}{X7D5967788791362B}{}
{
  

\subsection{\textcolor{Chapter }{DimensionVector (DimVectFT)}}
\logpage{[ 13, 4, 1 ]}\nobreak
\hyperdef{L}{X7C89E06784FB86C4}{}
{\noindent\textcolor{FuncColor}{$\triangleright$\enspace\texttt{DimensionVector({\mdseries\slshape AR, M})\index{DimensionVector@\texttt{DimensionVector}!DimVectFT}
\label{DimensionVector:DimVectFT}
}\hfill{\scriptsize (operation)}}\\


 Arguments: \mbox{\texttt{\mdseries\slshape AR}} \texttt{\symbol{45}} an object from the category \texttt{IsARQuiverNumerical} (\ref{IsARQuiverNumerical});\\
 \mbox{\texttt{\mdseries\slshape M}} \texttt{\symbol{45}} a number of an indecomposable module in \mbox{\texttt{\mdseries\slshape AR}} or a multiplicity vector (cf. \ref{introDO}). \textbf{\indent Returns:\ }
 a dimension vector of a module \mbox{\texttt{\mdseries\slshape M}} in the form of a list. The order of dimensions in this list corresponds to an
order of projectives defined in \mbox{\texttt{\mdseries\slshape AR}} (cf. \texttt{ARQuiverNumerical} (\ref{ARQuiverNumerical})). 

}

 
\begin{Verbatim}[commandchars=!@|,fontsize=\small,frame=single,label=Example]
  !gapprompt@gap>| !gapinput@a := ARQuiverNumerical("D4 subspace");|
  <ARQuiverNumerical with 12 indecomposables and 4 projectives>
  !gapprompt@gap>| !gapinput@DimensionVector(a, 7);|
  [ 1, 1, 0, 1 ]
  !gapprompt@gap>| !gapinput@DimensionVector(a, [0,1,0,0,0,0,2,0,0,0,0,0]);|
  [ 3, 3, 0, 2 ]
   
\end{Verbatim}
 

\subsection{\textcolor{Chapter }{DimHom}}
\logpage{[ 13, 4, 2 ]}\nobreak
\hyperdef{L}{X869602A97DFBCA30}{}
{\noindent\textcolor{FuncColor}{$\triangleright$\enspace\texttt{DimHom({\mdseries\slshape AR, M, N})\index{DimHom@\texttt{DimHom}}
\label{DimHom}
}\hfill{\scriptsize (operation)}}\\


 Arguments: \mbox{\texttt{\mdseries\slshape AR}} \texttt{\symbol{45}} an object from the category \texttt{IsARQuiverNumerical} (\ref{IsARQuiverNumerical});\\
 \mbox{\texttt{\mdseries\slshape M}} \texttt{\symbol{45}} a number of indecomposable module in \mbox{\texttt{\mdseries\slshape AR}} or a multiplicity vector;\\
 \mbox{\texttt{\mdseries\slshape N}} \texttt{\symbol{45}} a number of indecomposable module in \mbox{\texttt{\mdseries\slshape AR}} or a multiplicity vector (cf. \ref{introDO}).\\
 \textbf{\indent Returns:\ }
 the dimension of the homomorphism space between modules \mbox{\texttt{\mdseries\slshape M}} and \mbox{\texttt{\mdseries\slshape N}}. 

}

 

\subsection{\textcolor{Chapter }{DimEnd}}
\logpage{[ 13, 4, 3 ]}\nobreak
\hyperdef{L}{X8017F302837B903A}{}
{\noindent\textcolor{FuncColor}{$\triangleright$\enspace\texttt{DimEnd({\mdseries\slshape AR, M})\index{DimEnd@\texttt{DimEnd}}
\label{DimEnd}
}\hfill{\scriptsize (operation)}}\\


 Arguments: \mbox{\texttt{\mdseries\slshape AR}} \texttt{\symbol{45}} an object from the category \texttt{IsARQuiverNumerical} (\ref{IsARQuiverNumerical});\\
 \mbox{\texttt{\mdseries\slshape M}} \texttt{\symbol{45}} a number of indecomposable module in \mbox{\texttt{\mdseries\slshape AR}} or a multiplicity vector (cf. \ref{introDO}).\\
 \textbf{\indent Returns:\ }
 the dimension of the endomorphism algebra of a module \mbox{\texttt{\mdseries\slshape M}}. 

}

 

\subsection{\textcolor{Chapter }{OrbitDim}}
\logpage{[ 13, 4, 4 ]}\nobreak
\hyperdef{L}{X835E686B7BDF5A09}{}
{\noindent\textcolor{FuncColor}{$\triangleright$\enspace\texttt{OrbitDim({\mdseries\slshape AR, M})\index{OrbitDim@\texttt{OrbitDim}}
\label{OrbitDim}
}\hfill{\scriptsize (operation)}}\\


 Arguments: \mbox{\texttt{\mdseries\slshape AR}} \texttt{\symbol{45}} an object from the category \texttt{IsARQuiverNumerical} (\ref{IsARQuiverNumerical});\\
 \mbox{\texttt{\mdseries\slshape M}} \texttt{\symbol{45}} a number of indecomposable module in \mbox{\texttt{\mdseries\slshape AR}} or a multiplicity vector (cf. \ref{introDO}).\\
 \textbf{\indent Returns:\ }
 the dimension of the orbit of module \mbox{\texttt{\mdseries\slshape M}} (in the variety of representations of quiver with relations). 



 OrbitDim(\mbox{\texttt{\mdseries\slshape M}}) =
d{\textunderscore}1\texttt{\symbol{94}}2+...+d{\textunderscore}p\texttt{\symbol{94}}2
\texttt{\symbol{45}} dim End(\mbox{\texttt{\mdseries\slshape M}}), where (d{\textunderscore}i){\textunderscore}i = DimensionVector(\mbox{\texttt{\mdseries\slshape M}}). }

 

\subsection{\textcolor{Chapter }{OrbitCodim}}
\logpage{[ 13, 4, 5 ]}\nobreak
\hyperdef{L}{X782CF3C47A7DDFC0}{}
{\noindent\textcolor{FuncColor}{$\triangleright$\enspace\texttt{OrbitCodim({\mdseries\slshape AR, M, N})\index{OrbitCodim@\texttt{OrbitCodim}}
\label{OrbitCodim}
}\hfill{\scriptsize (operation)}}\\


 Arguments: \mbox{\texttt{\mdseries\slshape AR}} \texttt{\symbol{45}} an object from the category \texttt{IsARQuiverNumerical} (\ref{IsARQuiverNumerical});\\
 \mbox{\texttt{\mdseries\slshape M}} \texttt{\symbol{45}} a number of indecomposable module in \mbox{\texttt{\mdseries\slshape AR}} or a multiplicity vector;\\
 \mbox{\texttt{\mdseries\slshape N}} \texttt{\symbol{45}} a number of indecomposable module in \mbox{\texttt{\mdseries\slshape AR}} or a multiplicity vector (cf. \ref{introDO}).\\
 \textbf{\indent Returns:\ }
 the codimension of orbits of modules \mbox{\texttt{\mdseries\slshape M}} and \mbox{\texttt{\mdseries\slshape N}} (= dim End(\mbox{\texttt{\mdseries\slshape N}}) \texttt{\symbol{45}} dim End(\mbox{\texttt{\mdseries\slshape M}})). [explain more???] 



 NOTE: The function does not check if it makes sense, i.e. if \mbox{\texttt{\mdseries\slshape M}} and \mbox{\texttt{\mdseries\slshape N}} are in the same variety ( = dimension vectors coincide)! }

 

\subsection{\textcolor{Chapter }{DegOrderLEQ}}
\logpage{[ 13, 4, 6 ]}\nobreak
\hyperdef{L}{X85D9166A85499D99}{}
{\noindent\textcolor{FuncColor}{$\triangleright$\enspace\texttt{DegOrderLEQ({\mdseries\slshape AR, M, N})\index{DegOrderLEQ@\texttt{DegOrderLEQ}}
\label{DegOrderLEQ}
}\hfill{\scriptsize (operation)}}\\


 Arguments: \mbox{\texttt{\mdseries\slshape AR}} \texttt{\symbol{45}} an object from the category \texttt{IsARQuiverNumerical} (\ref{IsARQuiverNumerical});\\
 \mbox{\texttt{\mdseries\slshape M}} \texttt{\symbol{45}} a number of indecomposable module in \mbox{\texttt{\mdseries\slshape AR}} or a multiplicity vector;\\
 \mbox{\texttt{\mdseries\slshape N}} \texttt{\symbol{45}} a number of indecomposable module in \mbox{\texttt{\mdseries\slshape AR}} or a multiplicity vector (cf. \ref{introDO}).\\
 \textbf{\indent Returns:\ }
 true if \mbox{\texttt{\mdseries\slshape M}}$<=$\mbox{\texttt{\mdseries\slshape N}} in a degeneration order i.e. if \mbox{\texttt{\mdseries\slshape N}} is a degeneration of \mbox{\texttt{\mdseries\slshape M}} (see \ref{bdef}), and false otherwise. 



 NOTE: Function checks if it makes sense, i.e. if \mbox{\texttt{\mdseries\slshape M}} and \mbox{\texttt{\mdseries\slshape N}} are in the same variety ( = dimension vectors coincide). If not, it returns
false and additionally prints warning. }

 
\begin{Verbatim}[commandchars=!@|,fontsize=\small,frame=single,label=Example]
  !gapprompt@gap>| !gapinput@a := ARQuiverNumerical("R nilp");|
  <ARQuiverNumerical with 7 indecomposables and 2 projectives>
  !gapprompt@gap>| !gapinput@DimensionVector(a, 2);  DimensionVector(a, 3);|
  [ 2, 1 ]
  [ 2, 1 ]
  !gapprompt@gap>| !gapinput@DegOrderLEQ(a, 2, 3);|
  true
  !gapprompt@gap>| !gapinput@DegOrderLEQ(a, 3, 2);|
  false
   
\end{Verbatim}
 

\subsection{\textcolor{Chapter }{DegOrderLEQNC}}
\logpage{[ 13, 4, 7 ]}\nobreak
\hyperdef{L}{X790BD2D9813320B6}{}
{\noindent\textcolor{FuncColor}{$\triangleright$\enspace\texttt{DegOrderLEQNC({\mdseries\slshape AR, M, N})\index{DegOrderLEQNC@\texttt{DegOrderLEQNC}}
\label{DegOrderLEQNC}
}\hfill{\scriptsize (operation)}}\\


 Arguments: \mbox{\texttt{\mdseries\slshape AR}} \texttt{\symbol{45}} an object from the category \texttt{IsARQuiverNumerical} (\ref{IsARQuiverNumerical});\\
 \mbox{\texttt{\mdseries\slshape M}} \texttt{\symbol{45}} a number of indecomposable module in \mbox{\texttt{\mdseries\slshape AR}} or a multiplicity vector;\\
 \mbox{\texttt{\mdseries\slshape N}} \texttt{\symbol{45}} a number of indecomposable module in \mbox{\texttt{\mdseries\slshape AR}} or a multiplicity vector (cf. \ref{introDO}).\\
 \textbf{\indent Returns:\ }
 true if \mbox{\texttt{\mdseries\slshape M}}$<=$\mbox{\texttt{\mdseries\slshape N}} in a degeneration order i.e. if \mbox{\texttt{\mdseries\slshape N}} is a degeneration of \mbox{\texttt{\mdseries\slshape M}} (see \ref{bdef}), and false otherwise. 



 NOTE: Function does Not Check ("NC") if it makes sense, i.e. if \mbox{\texttt{\mdseries\slshape M}} and \mbox{\texttt{\mdseries\slshape N}} are in the same variety ( = dimension vectors coincide). If not, the result
doesn't make sense!\\
 It is useful when one wants to speed up computations (does not need to check
the dimension vectors). }

 

\subsection{\textcolor{Chapter }{PrintMultiplicityVector}}
\logpage{[ 13, 4, 8 ]}\nobreak
\hyperdef{L}{X877D07357EEA1418}{}
{\noindent\textcolor{FuncColor}{$\triangleright$\enspace\texttt{PrintMultiplicityVector({\mdseries\slshape M})\index{PrintMultiplicityVector@\texttt{PrintMultiplicityVector}}
\label{PrintMultiplicityVector}
}\hfill{\scriptsize (operation)}}\\


 \mbox{\texttt{\mdseries\slshape M}} \texttt{\symbol{45}} a list = multiplicity vector (cf. \ref{introDO}).\\
 

 This function prints the multiplicity vector \mbox{\texttt{\mdseries\slshape M}} in a more "readable" way (especially useful if \mbox{\texttt{\mdseries\slshape M}} is long and sparse). It prints a "sum" of non\texttt{\symbol{45}}zero
multiplicities in the form "multiplicity *
(number\texttt{\symbol{45}}of\texttt{\symbol{45}}indecomposable)". }

 

\subsection{\textcolor{Chapter }{PrintMultiplicityVectors}}
\logpage{[ 13, 4, 9 ]}\nobreak
\hyperdef{L}{X83C7434A7AC586CE}{}
{\noindent\textcolor{FuncColor}{$\triangleright$\enspace\texttt{PrintMultiplicityVectors({\mdseries\slshape list})\index{PrintMultiplicityVectors@\texttt{PrintMultiplicityVectors}}
\label{PrintMultiplicityVectors}
}\hfill{\scriptsize (operation)}}\\


 \mbox{\texttt{\mdseries\slshape list}} \texttt{\symbol{45}} a list of multiplicity vectors (cf. \ref{introDO}).\\
 

 This function prints all the multiplicity vectors from the \mbox{\texttt{\mdseries\slshape list}} in a more "readable" way, as \texttt{PrintMultiplicityVector} (\ref{PrintMultiplicityVector}). }

 }

 
\section{\textcolor{Chapter }{Operations returning families of modules}}\label{opers}
\logpage{[ 13, 5, 0 ]}
\hyperdef{L}{X7B0F730F82B4FACA}{}
{
  The functions from this section use quite advanced algorithms on (potentially)
big amount of data, so their runtimes can be long for "big"
A\texttt{\symbol{45}}R quivers! 

\subsection{\textcolor{Chapter }{ModulesOfDimVect}}
\logpage{[ 13, 5, 1 ]}\nobreak
\hyperdef{L}{X7DE8B7ED854B4EA1}{}
{\noindent\textcolor{FuncColor}{$\triangleright$\enspace\texttt{ModulesOfDimVect({\mdseries\slshape AR, which})\index{ModulesOfDimVect@\texttt{ModulesOfDimVect}}
\label{ModulesOfDimVect}
}\hfill{\scriptsize (operation)}}\\


 Arguments: \mbox{\texttt{\mdseries\slshape AR}} \texttt{\symbol{45}} an object from the category \texttt{IsARQuiverNumerical} (\ref{IsARQuiverNumerical});\\
 \mbox{\texttt{\mdseries\slshape which}} \texttt{\symbol{45}} a number of an indecomposable module in \mbox{\texttt{\mdseries\slshape AR}} or a dimension vector (see \texttt{DimensionVector} (\ref{DimensionVector:DimVectFT})). \textbf{\indent Returns:\ }
 a list of all modules (= multiplicity vectors, see \ref{introDO}) with dimension vector equal to \mbox{\texttt{\mdseries\slshape which}}. 

}

 

\subsection{\textcolor{Chapter }{DegOrderPredecessors}}
\logpage{[ 13, 5, 2 ]}\nobreak
\hyperdef{L}{X871CF6EE8579FEC8}{}
{\noindent\textcolor{FuncColor}{$\triangleright$\enspace\texttt{DegOrderPredecessors({\mdseries\slshape AR, M})\index{DegOrderPredecessors@\texttt{DegOrderPredecessors}}
\label{DegOrderPredecessors}
}\hfill{\scriptsize (operation)}}\\


 Arguments: \mbox{\texttt{\mdseries\slshape AR}} \texttt{\symbol{45}} an object from the category \texttt{IsARQuiverNumerical} (\ref{IsARQuiverNumerical});\\
 \mbox{\texttt{\mdseries\slshape M}} \texttt{\symbol{45}} a number of indecomposable module in \mbox{\texttt{\mdseries\slshape AR}} or a multiplicity vector (cf. \ref{introDO}).\\
 \textbf{\indent Returns:\ }
 a list of all modules (= multiplicity vectors) which are the predecessors of
module \mbox{\texttt{\mdseries\slshape M}} in a degeneration order (see \ref{bdef}). 

}

 
\begin{Verbatim}[commandchars=!@|,fontsize=\small,frame=single,label=Example]
  !gapprompt@gap>| !gapinput@a := ARQuiverNumerical("BG", 5);|
  <ARQuiverNumerical with 72 indecomposables and 2 projectives>
  !gapprompt@gap>| !gapinput@preds := DegOrderPredecessors(a, 60);; Length(preds);|
  18
  !gapprompt@gap>| !gapinput@DegOrderLEQ(a, preds[7], 60);|
  true
  !gapprompt@gap>| !gapinput@dpreds := DegOrderDirectPredecessors(a, 60);; Length(dpreds);|
  5
  !gapprompt@gap>| !gapinput@PrintMultiplicityVectors(dpreds);|
  1*(14) + 1*(64)
  1*(10) + 1*(71)
  1*(9) + 1*(67)
  1*(5) + 1*(17) + 1*(72)
  1*(1) + 1*(5) + 1*(20)
   
\end{Verbatim}
 

\subsection{\textcolor{Chapter }{DegOrderDirectPredecessors}}
\logpage{[ 13, 5, 3 ]}\nobreak
\hyperdef{L}{X7EBEA7DA85E645EB}{}
{\noindent\textcolor{FuncColor}{$\triangleright$\enspace\texttt{DegOrderDirectPredecessors({\mdseries\slshape AR, M})\index{DegOrderDirectPredecessors@\texttt{DegOrderDirectPredecessors}}
\label{DegOrderDirectPredecessors}
}\hfill{\scriptsize (operation)}}\\


 Arguments: \mbox{\texttt{\mdseries\slshape AR}} \texttt{\symbol{45}} an object from the category \texttt{IsARQuiverNumerical} (\ref{IsARQuiverNumerical});\\
 \mbox{\texttt{\mdseries\slshape M}} \texttt{\symbol{45}} a number of indecomposable module in \mbox{\texttt{\mdseries\slshape AR}} or a multiplicity vector (cf. \ref{introDO}).\\
 \textbf{\indent Returns:\ }
 a list of all modules (= multiplicity vectors) which are the direct
predecessors of module \mbox{\texttt{\mdseries\slshape M}} in a degeneration order (see \ref{bdef}). 

}

 

\subsection{\textcolor{Chapter }{DegOrderPredecessorsWithDirect}}
\logpage{[ 13, 5, 4 ]}\nobreak
\hyperdef{L}{X84930921811486F3}{}
{\noindent\textcolor{FuncColor}{$\triangleright$\enspace\texttt{DegOrderPredecessorsWithDirect({\mdseries\slshape AR, M})\index{DegOrderPredecessorsWithDirect@\texttt{DegOrderPredecessorsWithDirect}}
\label{DegOrderPredecessorsWithDirect}
}\hfill{\scriptsize (operation)}}\\


 Arguments: \mbox{\texttt{\mdseries\slshape AR}} \texttt{\symbol{45}} an object from the category \texttt{IsARQuiverNumerical} (\ref{IsARQuiverNumerical});\\
 \mbox{\texttt{\mdseries\slshape M}} \texttt{\symbol{45}} a number of indecomposable module in \mbox{\texttt{\mdseries\slshape AR}} or a multiplicity vector (cf. \ref{introDO}).\\
 \textbf{\indent Returns:\ }
 a pair (2\texttt{\symbol{45}}element list) [\mbox{\texttt{\mdseries\slshape p}}, \mbox{\texttt{\mdseries\slshape dp}}] where\\
 \mbox{\texttt{\mdseries\slshape p}} = the same as a result of \texttt{DegOrderPredecessors} (\ref{DegOrderPredecessors});\\
 \mbox{\texttt{\mdseries\slshape dp}} = the same as a result of \texttt{DegOrderDirectPredecessors} (\ref{DegOrderDirectPredecessors});\\
 



 The function generates predecessors only once, so the runtime is exactly the
same as DegOrderDirectPredecessors. }

 

\subsection{\textcolor{Chapter }{DegOrderSuccessors}}
\logpage{[ 13, 5, 5 ]}\nobreak
\hyperdef{L}{X7DB3AC3685859104}{}
{\noindent\textcolor{FuncColor}{$\triangleright$\enspace\texttt{DegOrderSuccessors({\mdseries\slshape AR, M})\index{DegOrderSuccessors@\texttt{DegOrderSuccessors}}
\label{DegOrderSuccessors}
}\hfill{\scriptsize (operation)}}\\


 Arguments: \mbox{\texttt{\mdseries\slshape AR}} \texttt{\symbol{45}} an object from the category \texttt{IsARQuiverNumerical} (\ref{IsARQuiverNumerical});\\
 \mbox{\texttt{\mdseries\slshape M}} \texttt{\symbol{45}} a number of indecomposable module in \mbox{\texttt{\mdseries\slshape AR}} or a multiplicity vector (cf. \ref{introDO}).\\
 \textbf{\indent Returns:\ }
 a list of all modules (= multiplicity vectors) which are the successors of
module \mbox{\texttt{\mdseries\slshape M}} in a degeneration order (see \ref{bdef}). 

}

 

\subsection{\textcolor{Chapter }{DegOrderDirectSuccessors}}
\logpage{[ 13, 5, 6 ]}\nobreak
\hyperdef{L}{X7C477F4E81CA6791}{}
{\noindent\textcolor{FuncColor}{$\triangleright$\enspace\texttt{DegOrderDirectSuccessors({\mdseries\slshape AR, M})\index{DegOrderDirectSuccessors@\texttt{DegOrderDirectSuccessors}}
\label{DegOrderDirectSuccessors}
}\hfill{\scriptsize (operation)}}\\


 Arguments: \mbox{\texttt{\mdseries\slshape AR}} \texttt{\symbol{45}} an object from the category \texttt{IsARQuiverNumerical} (\ref{IsARQuiverNumerical});\\
 \mbox{\texttt{\mdseries\slshape M}} \texttt{\symbol{45}} a number of indecomposable module in \mbox{\texttt{\mdseries\slshape AR}} or a multiplicity vector (cf. \ref{introDO}).\\
 \textbf{\indent Returns:\ }
 a list of all modules (= multiplicity vectors) which are the direct successors
of module \mbox{\texttt{\mdseries\slshape M}} in a degeneration order (see \ref{bdef}). 

}

 

\subsection{\textcolor{Chapter }{DegOrderSuccessorsWithDirect}}
\logpage{[ 13, 5, 7 ]}\nobreak
\hyperdef{L}{X7EA7BB8986E8FAF3}{}
{\noindent\textcolor{FuncColor}{$\triangleright$\enspace\texttt{DegOrderSuccessorsWithDirect({\mdseries\slshape AR, M})\index{DegOrderSuccessorsWithDirect@\texttt{DegOrderSuccessorsWithDirect}}
\label{DegOrderSuccessorsWithDirect}
}\hfill{\scriptsize (operation)}}\\


 Arguments: \mbox{\texttt{\mdseries\slshape AR}} \texttt{\symbol{45}} an object from the category \texttt{IsARQuiverNumerical} (\ref{IsARQuiverNumerical});\\
 \mbox{\texttt{\mdseries\slshape M}} \texttt{\symbol{45}} a number of indecomposable module in \mbox{\texttt{\mdseries\slshape AR}} or a multiplicity vector (cf. \ref{introDO}).\\
 \textbf{\indent Returns:\ }
 a pair (2\texttt{\symbol{45}}element list) [\mbox{\texttt{\mdseries\slshape s}}, \mbox{\texttt{\mdseries\slshape ds}}] where\\
 \mbox{\texttt{\mdseries\slshape s}} = the same as a result of \texttt{DegOrderSuccessors} (\ref{DegOrderSuccessors});\\
 \mbox{\texttt{\mdseries\slshape ds}} = the same as a result of \texttt{DegOrderDirectSuccessors} (\ref{DegOrderDirectSuccessors});\\
 



 The function generates successors only once, so the runtime is exactly the
same as DegOrderDirectSuccessors. }

 }

 }

 \def\bibname{References\logpage{[ "Bib", 0, 0 ]}
\hyperdef{L}{X7A6F98FD85F02BFE}{}
}

\bibliographystyle{alpha}
\bibliography{qpadocumentation.bib}

\addcontentsline{toc}{chapter}{References}

\def\indexname{Index\logpage{[ "Ind", 0, 0 ]}
\hyperdef{L}{X83A0356F839C696F}{}
}

\cleardoublepage
\phantomsection
\addcontentsline{toc}{chapter}{Index}


\printindex

\newpage
\immediate\write\pagenrlog{["End"], \arabic{page}];}
\immediate\closeout\pagenrlog
\end{document}
